\documentclass[UTF8,11pt,twoside,openany]{ctexbook}
\usepackage[paperwidth=176mm,paperheight=250mm,inner=21mm,outer=18mm,top=20mm,bottom=20mm,headheight=16pt,headsep=6mm,footskip=10mm]{geometry}
\usepackage{amsmath,amssymb,amsthm,bm,mathtools,booktabs,array,enumitem,float}
\usepackage{tikz,fancyhdr,hyperref}
\usetikzlibrary{arrows.meta,calc,patterns,angles,quotes,decorations.pathmorphing}
\hypersetup{hidelinks,pdftitle={热学进阶：物理竞赛二轮讲义},pdfauthor={}}
\setlength{\parindent}{2em}\setlength{\parskip}{3pt}\setlength{\emergencystretch}{2em}
\linespread{1.18}\raggedbottom\let\cleardoublepage\clearpage\setcounter{tocdepth}{1}
\ctexset{chapter={name={第,讲},number=\chinese{chapter},format=\LARGE\bfseries\centering,beforeskip=16pt,afterskip=22pt},section={format=\Large\bfseries},subsection={format=\large\bfseries}}
\pagestyle{fancy}\fancyhf{}\fancyhead[LE,RO]{\small\thepage}\fancyhead[LO]{\small 热学进阶}\fancyhead[RE]{\small\leftmark}
\renewcommand{\headrulewidth}{0pt}\fancypagestyle{plain}{\fancyhf{}\fancyfoot[C]{\small\thepage}\renewcommand{\headrulewidth}{0pt}}
\theoremstyle{definition}\newtheorem{example}{例}[chapter]\newtheorem{exercise}{习题}[chapter]
\newenvironment{solution}{\par\medskip\noindent\textbf{解}\quad\ignorespaces}{\par\medskip}
\numberwithin{equation}{chapter}
\newcommand{\dd}{\mathop{}\!\mathrm d}\newcommand{\ee}{\mathrm e}\newcommand{\kb}{k_{\mathrm B}}\newcommand{\vct}[1]{\bm{#1}}
\DeclareMathOperator{\erf}{erf}\DeclareMathOperator{\erfc}{erfc}
\tikzset{axis/.style={->,>=Stealth,line width=.45pt},curve/.style={line width=.7pt},guide/.style={dashed,line width=.35pt},every node/.style={font=\small},particle/.style={circle,fill,inner sep=1.3pt}}
\setlist[enumerate]{itemsep=3pt,topsep=4pt,leftmargin=2em}
\let\LectureContents\tableofcontents
\renewcommand{\tableofcontents}{\begingroup\small\linespread{1}\selectfont\LectureContents\endgroup}
\newcommand{\Diagram}[3]{\begin{figure}[H]\centering\begin{tikzpicture}#3\end{tikzpicture}\caption{#2}\label{#1}\end{figure}}
\newcommand{\Chamber}[4]{\draw[curve](#1,#2)rectangle++(#3,#4);}
\newcommand{\Bath}[4]{\draw[curve](#1,#2)rectangle++(2.3,.6);\node at(#1+1.15,#2+.3){$#3$};\node[anchor=west]at(#1+2.45,#2+.3){#4};}
\providecommand{\ExpansionDiagram}[3]{\begin{figure}[H]\centering\begin{tikzpicture}#3\end{tikzpicture}\caption{#2}\label{#1}\end{figure}}


% 分册与十编合订本共用的定位页码；局部章、图、题号保持原有对应。
\newcommand{\CoursePartNumber}{8}
\geometry{inner=19.5mm,outer=19.5mm,headheight=16pt}
\AtBeginDocument{%
  \fancyhead{}\fancyhead[C]{\small 热学进阶}%
  \fancyfoot[C]{\small\CoursePartNumber-\thepage}%
  \fancypagestyle{plain}{\fancyhf{}\fancyfoot[C]{\small\CoursePartNumber-\thepage}\renewcommand{\headrulewidth}{0pt}}%
}

\begin{document}
\frontmatter
\begin{titlepage}\centering\vspace*{28mm}
{\Large 物理竞赛二轮讲义\quad 第八编\par}\vspace{8mm}
{\Huge\bfseries 热学进阶\par}\vspace{12mm}{\Large 物理竞赛二轮讲义\par}\vspace{20mm}{\large 分子分布\quad 热力学关系\quad 熵与热机\quad 相变与外场\par}\vfill\end{titlepage}
\chapter*{编写说明}
本编围绕分子运动、热力学状态函数、熵与热机、相变和外场平衡展开。各讲先建立模型与适用条件，再推导关系式，随后用例题讨论计算方法与模型变式。习题均附解答，图中的方向、边界和状态标记贯穿题设与计算。

概率密度与积分测度共同决定区间概率，偏导数的意义则取决于保持不变的变量。开始计算前，还需确定系统边界：粒子能否出入，热量与功怎样交换，外场势能是否计入总能量。这些约定决定了状态关系与过程方程各自的适用范围。

本编用 $N$ 表示粒子数，$\nu$ 表示物质的量，$n=N/V$ 表示数密度；$m$ 为单粒子质量，$M$ 为摩尔质量。热量以传入系统为正，功以系统对外做功为正，故第一定律写为 $\dd U=\delta Q-\delta W$。除非另有说明，温度均为热力学温度。
\tableofcontents
\mainmatter
\chapter{分子运动的统计描述}

\section{从状态概率到速度分布}
\label{th:mb}
宏观气体处于平衡时，单个分子的速度仍不断变化。统计描述用概率表示相同宏观条件下各微观状态的占据情况。若有 $N_i$ 个粒子占据第 $i$ 个单粒子状态，则 $p_i=N_i/N$，且 $\sum_i p_i=1$。在经典稀薄气体近似下，温度为 $T$ 的热平衡分布为
\begin{equation}
 p_i=A\ee^{-E_i/(\kb T)},\qquad A^{-1}=\sum_i\ee^{-E_i/(\kb T)}.
\end{equation}
若能量 $E_i$ 对应 $r_i$ 个不同状态，整个能级的概率应为 $r_iA\ee^{-E_i/(\kb T)}$。简并度计数与指数因子承担不同作用：前者计算可占据状态的数目，后者比较每个状态的权重。

这里的单粒子状态必须区分清楚。能量相同并不意味着状态相同；速度大小相同而方向不同，就是不同状态。经典分布还要求量子简并效应可忽略，粒子间相互作用对能量和计数的影响足够小。平动服从本讲分布，并不意味着转动、振动等内部自由度在任意温度都充分激发。

\subsection{速度分量的归一化}
\label{th:gaussian}
没有宏观流动与外场方向偏置时，平动能
\[
 E_k=\frac m2(v_x^2+v_y^2+v_z^2)
\]
使三维概率密度可分离为三个相同的一维因子。记 $\alpha=m/(2\kb T)$，一维密度为 $f_x(v_x)=A\ee^{-\alpha v_x^2}$。概率是密度乘速度区间，$\dd P=f_x(v_x)\dd v_x$，故 $f_x$ 的量纲是速率的倒数，而不是无量纲概率。

归一化需要 $\int_{-\infty}^{\infty}f_x\dd v_x=1$。先令 $I=\int_{-\infty}^{\infty}\ee^{-\alpha x^2}\dd x$，利用两个独立积分相乘，有
\begin{align}
 I^2&=\iint_{\mathbb R^2}\ee^{-\alpha(x^2+y^2)}\dd x\dd y\\
 &=\int_0^{2\pi}\dd\theta\int_0^\infty \ee^{-\alpha r^2}r\dd r\\
 &=2\pi\lim_{L\to\infty}
 \left[-\frac{\ee^{-\alpha r^2}}{2\alpha}\right]_0^L
 =\frac\pi\alpha.
\end{align}
因 $I>0$，得到 $I=\sqrt{\pi/\alpha}$，从而
\begin{equation}\label{eq:maxwellvector}
 f_x(v_x)=\sqrt{\frac{m}{2\pi\kb T}}\ee^{-mv_x^2/(2\kb T)},\quad
 f(\vct v)=\left(\frac{m}{2\pi\kb T}\right)^{3/2}\ee^{-mv^2/(2\kb T)}.
\end{equation}
三维归一化测度为 $\dd^3v=\dd v_x\dd v_y\dd v_z$。速度分量的独立性来自这里的平衡分布。非平衡气流若存在方向关联，联合密度就需要另行确定。

\Diagram{fig:gaussianarea}{高斯积分中直角坐标面积与极坐标面积的对应}{
 \draw[axis](-.4,0)--(4.3,0)node[right]{$x$};\draw[axis](0,-.4)--(0,3.3)node[above]{$y$};
 \draw[curve](0,0)circle(2.4);\draw[guide](0,0)circle(2.7);
 \fill[gray!20](0,0)--(35:2.7)arc(35:47:2.7)--(47:2.4)arc(47:35:2.4)--cycle;
 \draw(0,0)--(41:2.7);\node at(1.2,1.35){$r$};\node[anchor=west]at(2.9,1.9){$r\dd r\dd\theta$};
 \draw[->](.8,0)arc(0:41:.8);\node at(1.1,.35){$\theta$};
}
图\ref{fig:gaussianarea}中小扇环的径向宽度为 $\dd r$，切向长度为 $r\dd\theta$，两者相乘得到面积元。因子 $r$ 就来自切向长度随半径的增长。

\subsection{计算矩所需的积分}
\label{th:integrals}
定义半轴高斯矩
\[
 I_s(\alpha)=\int_0^\infty x^s\ee^{-\alpha x^2}\dd x,\qquad s>-1.
\]
偶函数的半轴积分为整轴的一半，故$I_0=\sqrt\pi/(2\sqrt\alpha)$。
对$I_1$取$u=\alpha x^2$，则$\dd u=2\alpha x\dd x$，端点仍为$0,\infty$，所以
\[
 I_1=\frac1{2\alpha}\int_0^\infty\ee^{-u}\dd u
 =\frac1{2\alpha}\lim_{U\to\infty}[-\ee^{-u}]_0^U
 =\frac1{2\alpha}.
\]
在积分号内对参数求导，$\partial_\alpha\ee^{-\alpha x^2}=-x^2\ee^{-\alpha x^2}$，因此
\[
 \frac{\dd I_s}{\dd\alpha}
 =-\int_0^\infty x^{s+2}\ee^{-\alpha x^2}\dd x=-I_{s+2}.
\]
递推式用分部积分建立。将$I_{s+2}$中的$x\ee^{-\alpha x^2}\dd x$
写成$-\dd(\ee^{-\alpha x^2})/(2\alpha)$，便有
\[
 \begin{aligned}
 I_{s+2}&=-\frac1{2\alpha}
 [x^{s+1}\ee^{-\alpha x^2}]_0^\infty
 +\frac{s+1}{2\alpha}\int_0^\infty x^s\ee^{-\alpha x^2}\dd x\\
 &=\frac{s+1}{2\alpha}I_s.
 \end{aligned}
\]
$s>-1$保证零端点的$x^{s+1}$趋零，无穷远端则由指数衰减消去边界项。
例如$I_4=3I_2/(2\alpha)=3I_0/(4\alpha^2)$、
$I_5=2I_3/\alpha=2I_1/\alpha^2$。需要频繁使用的结果为
\[
 \begin{array}{c|cccccc}
 s&0&1&2&3&4&5\\ \hline
 I_s&\dfrac{\sqrt\pi}{2\sqrt\alpha}&\dfrac1{2\alpha}&
 \dfrac{\sqrt\pi}{4\alpha^{3/2}}&\dfrac1{2\alpha^2}&
 \dfrac{3\sqrt\pi}{8\alpha^{5/2}}&\dfrac1{\alpha^3}
 \end{array}
\]
积分区间是半轴还是整轴，要在代入前确定。偶函数的整轴积分是半轴的两倍；奇函数的整轴积分为零，不能用半轴积分代替。

\subsection{平衡权重的来源与能量零点}
\label{th:reservoirweight}
Boltzmann 因子可由粒子与热库的能量交换推导。把两者组成孤立整体，忽略耦合所占的能量，总能量记为 $E_{\rm tot}$。粒子处于能量为 $E_i$ 的一个微观状态时，热库能量为 $E_{\rm tot}-E_i$；该状态的概率与热库在此能量下的可占据状态数 $\Omega_b(E_{\rm tot}-E_i)$ 成正比。用 $S_b=\kb\ln\Omega_b$ 表示热库熵，在参考能量附近展开，
\[
 S_b(E_b-E_i)=S_b(E_b)-\frac{E_i}{T}
 -\frac{E_i^2}{2C_bT^2}+\cdots.
\]
这里用到$\partial S_b/\partial E_b=1/T$。对它再用链式法则，并用$C_b=\partial E_b/\partial T$，
\[\frac{\partial^2S_b}{\partial E_b^2}
=\frac{\dd(T^{-1})}{\dd T}\frac{\partial T}{\partial E_b}
=-\frac1{T^2}\frac1{C_b}.\]
热库的热容$C_b$足够大时，粒子交换的有限能量对热库温度的改变可忽略，指数中保留一次项便得到$\ee^{-E_i/(\kb T)}$。这里的温度是共同平衡温度。若热库很小，二次项不可忽略，单一的正则指数权重便只是近似。

\Diagram{fig:reservoirweight}{粒子的状态权重由热库可占据状态数决定}{
 \node[draw,minimum width=2cm,minimum height=1.2cm](a)at(0,0){粒子 $E_i$};
 \node[draw,minimum width=3.3cm,minimum height=1.5cm](b)at(5,0){热库 $E_b-E_i$};
 \draw[<->](a)--(b)node[midway,above]{能量交换};
 \node at(2.5,-1.35){整体能量固定；热库温度近似不变};
}
图\ref{fig:reservoirweight}比较的是单个微观状态的权重。同一能级有若干状态时，将各状态概率相加便出现简并度因子。连续能量情形相应地乘以态密度，因此能量分布的峰值由态密度与指数权重共同决定。

设单粒子配分和为 $Z_1=\sum_i r_i\ee^{-\beta E_i}$，其中 $\beta=1/(\kb T)$。若所有能级同时平移常数 $E_\ast$，则
\[
 Z_1'= \ee^{-\beta E_\ast}Z_1,\qquad
 \frac{r_i\ee^{-\beta(E_i+E_\ast)}}{Z_1'}
 =\frac{r_i\ee^{-\beta E_i}}{Z_1}.
\]
占有概率不变，平均能量只增加 $E_\ast$。粒子数固定时，能量差、做功和热容都不受这一平移影响；粒子数变化的开口系统则须在储能和粒子携带能量两处采用同一零点。

\subsection{由各向同性与分量独立性确定高斯形状}
各向同性规定分布只能依赖 $v_x^2+v_y^2+v_z^2$；分量独立性又规定联合密度可写成三个一维因子的乘积。这两个条件共同限制了函数形状。写成
\[
 f(\vct v)=\phi(v_x)\phi(v_y)\phi(v_z)
 =\Psi(v_x^2+v_y^2+v_z^2).
\]
在密度为正、二次连续可微的区域取对数，记 $L(s)=\ln\Psi(s)$。对 $v_x$ 求导，再对 $v_y$ 求导，
\[
 \frac{\partial^2\ln f}{\partial v_y\partial v_x}
 =4v_xv_yL''(v^2)=0.
\]
左端为零是因为乘积的对数中没有同时含$v_x,v_y$的项。任取非零的$v_x,v_y$可知$L''(s)=0$。两次积分分别给$L'(s)=-\alpha$、$L(s)=a-\alpha s$，其中$a,\alpha$为常数。归一化要求$\alpha>0$，于是得到高斯形状。由其二阶矩可得$m\langle v_x^2\rangle=m/(2\alpha)$；与平衡态的热能尺度$\kb T$对应，便有$\alpha=m/(2\kb T)$。

各向同性单独并不保证高斯分布。例如，速率固定而方向均匀的粒子群也是各向同性的，但其三个速度分量受 $v_x^2+v_y^2+v_z^2=v_0^2$ 约束，彼此不独立。后文固定速率背景的碰撞题正属于这种分布，不能把它的某一分量再当成 Maxwell 高斯分量。

\subsection{高斯积分的缩放、递推与边界}
已有的 $I_s$ 表可以由一个无量纲积分统一使用。令 $y=\sqrt\alpha\,x$，则
\[
 I_s(\alpha)=\alpha^{-(s+1)/2}
 \int_0^\infty y^s\ee^{-y^2}\dd y.
\]
幂次 $\alpha^{-(s+1)/2}$ 同时包含变量幂的缩放和微分 $\dd x$ 的缩放。它可先用来判断答案的温度、质量依赖，再计算常数系数。对连续参数求导时，积分区间须与该参数无关，并要求导数的积分收敛；本讲的 $\alpha>0$ 高斯积分满足这些条件。

若下限为非零阈值 $a$，则定义
$J_s(a)=\int_a^\infty x^s\ee^{-\alpha x^2}\dd x$。分部积分保留下限项，得到
\[
\begin{aligned}
 J_{s+2}(a)&=-\frac1{2\alpha}[x^{s+1}e^{-\alpha x^2}]_a^\infty
+\frac{s+1}{2\alpha}\int_a^\infty x^se^{-\alpha x^2}\dd x\\
&=\frac{a^{s+1}\ee^{-\alpha a^2}}{2\alpha}
 +\frac{s+1}{2\alpha}J_s(a).
\end{aligned}
\]
取 $a=0$ 才回到 $I_s$ 的递推。阈值问题出现的“指数乘多项式”，正是反复保留下限项的结果；半轴归一化积分中的边界项为零，并不意味着截断积分的边界项也为零。


\begin{example}[两个能级的占有与热容]\label{ex:twolevels}
独立粒子只有能量 $0$ 和 $\varepsilon$ 两个能级，简并度分别为 $1$、$r$。粒子数 $N$ 固定，求激发态占比、内能和定容热容，并讨论高、低温极限。
\Diagram{fig:twolevelq}{两个能级及其简并度}{
 \draw[axis](0,0)--(0,2.4)node[above]{$E$};\draw[curve](.5,.3)--(3.5,.3)node[right]{$0\;(1)$};
 \foreach \x in{.5,1.55,2.6}{\draw[curve](\x,1.65)--({\x+.8},1.65);}\node[right]at(3.5,1.65){$\varepsilon\;(r)$};
 \draw[->](2,.5)--(2,1.35)node[midway,right]{$T$};
}
\end{example}
\begin{solution}
图\ref{fig:twolevelq}中的上能级包含 $r$ 个状态。把各状态权重相加，单粒子归一化因子为 $Z=1+r\ee^{-x}$，其中 $x=\varepsilon/(\kb T)$。于是
\[
 P_1=\frac{r\ee^{-x}}{1+r\ee^{-x}},\qquad U=N\varepsilon P_1.
\]
由于$x=\varepsilon T^{-1}/\kb$，有$\dd x/\dd T=-\varepsilon/(\kb T^2)=-x/T$。
占有率的商式求导逐项给出
\[
 \frac{\dd P_1}{\dd x}
 =\frac{(-r\ee^{-x})(1+r\ee^{-x})-(r\ee^{-x})(-r\ee^{-x})}
 {(1+r\ee^{-x})^2}
 =-\frac{r\ee^{-x}}{(1+r\ee^{-x})^2}.
\]
再对$U=N\varepsilon P_1(x(T))$使用链式法则，
\[
 C_V=N\varepsilon\frac{\dd P_1}{\dd x}\frac{\dd x}{\dd T}
 =N\kb x^2\frac{r\ee^{-x}}{(1+r\ee^{-x})^2}.
\]
热容的温度导数是在粒子数不变的条件下取的。还可把结果中的分式写成 $P_1(1-P_1)$：占有率接近0或1时，粒子很少改变能级，内能随温度的变化便很小。对单粒子有 $\langle E^2\rangle=\varepsilon^2P_1$，所以
\[
 C_V=\frac{N}{\kb T^2}
 \bigl(\langle E^2\rangle-\langle E\rangle^2\bigr).
\]
这里的方差指一个粒子的能量涨落。$N$ 个独立粒子的总方差为它的 $N$ 倍，与上面直接求导相符。高温极限中各微观状态等概率，上能级占比仍由其简并度决定。

\Diagram{fig:twolevela}{能量与温度的两个极限}{
 \draw[axis](0,0)--(5,0)node[right]{$T$};\draw[axis](0,0)--(0,2.3)node[above]{$U/N$};
 \draw[curve]plot[domain=.12:4.5,samples=60](\x,{1.8*exp(-1/\x)/(1+exp(-1/\x))});
 \draw[guide](0,.9)--(4.6,.9)node[right]{$r\varepsilon/(1+r)$};\node at(3.5,1.7){示意取 $r=1$};
}
图\ref{fig:twolevela}显示有限能级系统的内能有上界。低温 $x\gg1$ 时激发被冻结，$C_V\sim N\kb r x^2\ee^{-x}$；高温 $x\ll1$ 时各微观状态近乎等概率，$P_1\to r/(1+r)$。此时内能接近饱和值，继续升温引起的内能增量趋于零，热容也趋于零。这与具有无界二次型能量的经典自由度不同。
\end{solution}

\section{速率、方向与抽样方式}
\label{th:speed}
速度矢量落入速度空间体积元的概率为 $f(\vct v)\dd^3v$。若只问速率在 $v$ 与 $v+\dd v$ 之间，就要把所有方向相加。各向同性时，球壳体积为 $4\pi v^2\dd v$，故
\begin{equation}\label{eq:speeddensity}
 g(v)=4\pi\left(\frac{m}{2\pi\kb T}\right)^{3/2}v^2\ee^{-mv^2/(2\kb T)},\qquad v\ge0.
\end{equation}
\Diagram{fig:velocityshell}{速度空间中厚度为\(\dd v\)的球壳}{
 \draw[axis](-2.1,0)--(2.4,0)node[right]{$v_x$};\draw[axis](0,-1.8)--(0,2.1)node[above]{$v_z$};
 \draw[axis](-1.7,-1.2)--(1.7,1.2)node[right]{$v_y$};
 \draw[curve](0,0)circle(1.45);\draw[guide](0,0)circle(1.62);
 \draw[curve](0,0)ellipse(1.45 and .4);\draw[guide](0,0)ellipse(1.62 and .45);
 \draw[->](0,0)--(35:1.45)node[midway,above]{$v$};
 \draw[<->](315:1.45)--(315:1.62);\node[right]at(1.5,-1.2){$\dd v$};
}
$f(\vct v)$ 在原点最大，而 $g(v)$ 在 $v=0$ 为零。这两个结论并不矛盾：球壳因子 $v^2$ 把不同方向的状态数计入速率概率。图\ref{fig:velocityshell}中的球壳是概率求和的几何区域，并非分子在真实空间的轨道。

\subsection{三个特征速率与一般矩}
\label{th:moments}
令 $B=4\pi[m/(2\pi\kb T)]^{3/2}$，则
\[
 \langle v^s\rangle=B I_{s+2}(\alpha)
 =\frac2{\sqrt\pi}\Gamma\!\left(\frac{s+3}{2}\right)
 \left(\frac{2\kb T}{m}\right)^{s/2}.
\]
这里的一般矩由$u=\alpha v^2$得到：$\dd v=\dd u/(2\sqrt{\alpha u})$，因而
\[
 I_{s+2}=\frac1{2\alpha^{(s+3)/2}}
 \int_0^\infty u^{(s+1)/2}\ee^{-u}\dd u
 =\frac{\Gamma((s+3)/2)}{2\alpha^{(s+3)/2}},\qquad s>-3.
\]
最后的积分定义了$\Gamma$函数。整数矩也可直接使用前节递推。平均速率由$I_3$给出：
\[
 \bar v=B\frac1{2\alpha^2}=\sqrt{\frac{8\kb T}{\pi m}}.
\]
均方速率由 $I_4$ 给出，$\langle v^2\rangle=3\kb T/m$，故 $v_{\rm rms}=\sqrt{3\kb T/m}$。最概然速率需对密度 $g(v)$ 求极值：
\[
 g'(v)=B[2v\ee^{-\alpha v^2}+v^2(-2\alpha v)\ee^{-\alpha v^2}]
 =2Bv\ee^{-\alpha v^2}(1-\alpha v^2)=0,
 \qquad v_p=\sqrt{\frac{2\kb T}{m}}.
\]
除 $v=0$ 端点外，只有一个峰值。因此 $v_p<\bar v<v_{\rm rms}$。此外
\[
 \langle v^3\rangle=\sqrt{\frac{128}\pi}\left(\frac{\kb T}{m}\right)^{3/2},\quad
 \langle v^4\rangle=15\left(\frac{\kb T}{m}\right)^2.
\]
等分定理在这里可直接验证：$\langle v_x^2\rangle=\kb T/m$，每个平动二次项的平均能量为 $\kb T/2$。能量由二阶矩决定，碰撞通量由一阶矩决定，不能用同一个“平均速度”处理所有问题。

\begin{example}[能量密度的最概然值]
求平动能 $E=mv^2/2$ 的概率密度 $h(E)$、平均值与最概然值。为什么把 $v_p$ 代入 $mv^2/2$，不等于能量密度的峰值位置？
\Diagram{fig:energydensityq}{速率区间映射为不等宽的能量区间}{
 \draw[axis](0,1.7)--(5,1.7)node[right]{$v$};\draw[axis](0,0)--(5,0)node[right]{$E$};
 \draw[line width=2pt](1,1.7)--(1.5,1.7);\draw[line width=2pt](1,0)--(1.8,0);
 \draw[->](1.25,1.4)--(1.4,.3)node[midway,right]{$\dd E=mv\dd v$};
}
\end{example}
\begin{solution}
同一批粒子的概率不因变量替换而变，图\ref{fig:energydensityq}中两段对应 $h(E)\dd E=g(v)\dd v$。由 $v=\sqrt{2E/m}$、$\dd v/\dd E=1/\sqrt{2mE}$，得到
\[
 h(E)=\frac2{\sqrt\pi}\frac{\sqrt E}{(\kb T)^{3/2}}\ee^{-E/(\kb T)},\quad E\ge0.
\]
平均值可直接由二阶速率矩得到 $\langle E\rangle=3\kb T/2$。对密度取对数后求导，
\[
 \frac{\dd\ln h}{\dd E}=\frac1{2E}-\frac1{\kb T}=0
 \quad\Rightarrow\quad E_p=\frac12\kb T.
\]
变量替换中的幂次可以逐项追踪：原密度里的 $v^2$ 变为与 $E$ 成正比，雅可比里的 $1/v$ 又给出 $E^{-1/2}$，相乘才得到 $E^{1/2}$。密度的量纲也从 $1/v$ 变为 $1/E$。例如一小段能量区间的概率近似为 $h(E)\Delta E$；比较两个区间的粒子数时，既要比较密度，又要比较区间宽度。以速率等宽和以能量等宽划分样本，所比较的本来就是不同的区间。

\Diagram{fig:energydensitya}{能量概率密度的峰值与平均值}{
 \begin{scope}[yscale=3]
 \draw[axis](0,0)--(5,0)node[right]{$E/(\kb T)$};\draw[axis](0,0)--(0,.85)node[above]{$\kb T h$};
 \draw[curve]plot[domain=0:4.5,samples=90](\x,{2/sqrt(pi)*sqrt(\x)*exp(-\x)});
 \draw[guide](.5,0)--(.5,.484)node[above left]{$E_p$};\draw[guide](1.5,0)--(1.5,.308)node[above right]{$\langle E\rangle$};
 \end{scope}
}
图\ref{fig:energydensitya}的横轴为能量。用 $mv_p^2/2$ 得到 $\kb T$，它只表示速率密度峰值所对应的能量，不是 $h(E)$ 的峰值。概率密度的极大位置会受变量变换中的雅可比影响；具有直接不变意义的是某一区间的积分概率。
\end{solution}

\begin{example}[速度分量的阈值与尾概率]
平衡气体温度为 $T$。求 $v_x>u$ 的粒子比例及其中的条件平均 $\langle v_x\mid v_x>u\rangle$，其中 $u\ge0$。
\Diagram{fig:tailq}{一维速度分布中的阈值区域}{
 \draw[axis](-3,0)--(3,0)node[right]{$v_x$};\draw[axis](0,0)--(0,2.3)node[above]{$f_x$};
 \fill[gray!20](1,0)--plot[domain=1:2.7](\x,{1.8*exp(-\x*\x)})--(2.7,0)--cycle;
 \draw[curve]plot[domain=-2.7:2.7,samples=70](\x,{1.8*exp(-\x*\x)});\node[below]at(1,0){$u$};
}
\end{example}
\begin{solution}
按图\ref{fig:tailq}只积分右侧阴影，不再乘二。设 $z=\sqrt\alpha u$，定义
\[
 \erfc z=(2/\sqrt\pi)\int_z^\infty\ee^{-t^2}\dd t,
\]
则
\[
 P(v_x>u)=\int_u^\infty\sqrt{\frac\alpha\pi}\ee^{-\alpha v_x^2}\dd v_x
=\frac1{\sqrt\pi}\int_z^\infty\ee^{-t^2}\dd t
=\frac12\erfc z.
\]
其中$t=\sqrt\alpha v_x$，$\dd v_x=\dd t/\sqrt\alpha$，下限$u$变为$z$，上限仍为无穷大。
条件平均的分子为 $\int_u^\infty v_x\sqrt{\alpha/\pi}\ee^{-\alpha v_x^2}\dd v_x$。利用 $\dd(\ee^{-\alpha v_x^2})=-2\alpha v_x\ee^{-\alpha v_x^2}\dd v_x$，有
\[
 \int_u^\infty v_x f_x\dd v_x
 =\sqrt{\frac\alpha\pi}\lim_{L\to\infty}
 \left[-\frac{\ee^{-\alpha v_x^2}}{2\alpha}\right]_u^L
 =\frac{\ee^{-z^2}}{2\sqrt{\pi\alpha}},\qquad
 \langle v_x\mid v_x>u\rangle=
 \frac{\ee^{-z^2}}{\sqrt{\pi\alpha}\,\erfc z}.
\]
条件抽样可直接写成归一化密度
\[
 q(v_x)=\frac{f_x(v_x)}{P(v_x>u)},\qquad v_x>u,
\]
阈值之外密度为零。分子积分只给出原总体中这部分粒子贡献的速度一阶矩，除以尾概率才是子样本的平均值。对 $u=0$，原总体一半粒子被选中；对 $u>0$，所选比例进一步减小。由于保留的每个速度都大于 $u$，所得条件平均也应大于 $u$，这与把一个未经归一化的一阶矩误当平均值的结果有明显区别。

\Diagram{fig:taila}{阈值右侧重新归一化后的条件密度}{
 \draw[axis](-.3,0)--(4.5,0)node[right]{$v_x$};\draw[axis](0,0)--(0,3.3)node[above]{$q$};
 \fill[gray!15](1,0)--plot[domain=1:3.5,samples=70](\x,{7.1739*exp(-\x*\x)})--(3.5,0)--cycle;
 \draw[curve](1,0)--(1,{7.1739*exp(-1)});\draw[curve]plot[domain=1:3.5,samples=70](\x,{7.1739*exp(-\x*\x)});
 \node[below]at(1,0){$u$};\node at(2.8,1.8){面积为1};\node at(2.3,-.7){$q=f_x/P(v_x>u)$};
}
$u=0$ 时结果为 $\sqrt{2\kb T/(\pi m)}$。图\ref{fig:taila}强调，这是从某一时刻的气体中筛选速度分量；若按单位时间撞壁的粒子抽样，还要额外乘碰撞速率，不能沿用这个条件密度。
\end{solution}

\subsection{同一粒子群的不同概率测度}
\label{th:samplingmeasure}
概率密度总是相对于指定的积分变量定义。若一个单调变量变换为 $y=F(x)$，相同小区间中的概率应相等，
\[
 q(y)\dd y=p(x)\dd x,\qquad
 q(y)=p(x(y))\left|\frac{\dd x}{\dd y}\right|.
\]
密度的峰值会随变量变换而改变，区间概率和由同一事件定义的期望值则保持一致。例如直接计算 $\langle mv^2/2\rangle$，与先变成能量密度再计算 $\langle E\rangle$，必得相同结果。两种算法若不一致，应检查微分因子和积分范围。

\Diagram{fig:measuremap}{变量变换保持概率而改变区间宽度}{
 \draw[axis](0,1.4)--(6,1.4)node[right]{$x$};
 \draw[axis](0,0)--(6,0)node[right]{$y$};
 \draw[line width=2pt](1,1.4)--(1.8,1.4);
 \draw[line width=2pt](2.3,0)--(4.3,0);
 \draw[->](1.4,1.15)--(3.3,.25);
 \node[above]at(1.4,1.55){$\dd x$};
 \node[below]at(3.3,-.15){$\dd y$};
 \node[right]at(4,.8){$p(x)\dd x=q(y)\dd y$};
}
图\ref{fig:measuremap}中，两条粗线代表同一批事件在两个变量下占据的区间。速度矢量变为速率时，还须合并所有方向上的事件，球壳面积便来自这一步求和。这与单纯将一个一维变量改写成另一个变量有所不同。

另一种变化来自抽样方式。若一个粒子处于状态 $x$ 时，每单位时间被记录的机会正比于 $w(x)\ge0$，原有归一化密度为 $p(x)$，则记录事件的密度为
\[
 p_{\rm event}(x)=\frac{w(x)p(x)}{\langle w\rangle},
 \qquad
 \langle A\rangle_{\rm event}
 =\frac{\langle Aw\rangle}{\langle w\rangle}.
\]
分母是总记录率，用来把加权结果重新归一化。两次平均一般不同；差值可写成
\[
 \langle A\rangle_{\rm event}-\langle A\rangle
 =\frac{\langle Aw\rangle-\langle A\rangle\langle w\rangle}
 {\langle w\rangle}.
\]
因此事件样本中某个量偏大还是偏小，取决于它与记录权重的关联。快分子更常穿孔，所以穿孔样本的能量偏大；在一段束流中作瞬时空间计数，慢分子停留更久，权重又不同。这两种偏置都可由同一归一化式推导。

\subsection{条件筛选与通量加权的次序}
探测器若只接受某一区域 $\mathcal A$，先用指示函数 $\mathbf1_{\mathcal A}$ 表示筛选，则被记录事件的概率为
\[
 \dd P_{\rm det}=
 \frac{\mathbf1_{\mathcal A}(x)w(x)p(x)\dd x}
 {\int_{\mathcal A}w(x)p(x)\dd x}.
\]
分子中的筛选因子和权重可交换乘法次序，分母则始终取最终被接受事件的总记录率。体内满足阈值的粒子比例按粒子数计算，通量比例还包含各粒子的穿面机会，两者应分别积分。

各向同性小孔泄流中，速率因子 $vg(v)$ 与角度因子 $\cos\theta\,\dd\Omega$ 分离。单独限制接受圆锥不会改变速率分布；单独限制速率也不会改变归一化的角分布。这一结论以孔的传输几何与速率无关为条件。带电粒子经偏转场、长孔中的碰壁损失或随速率变化的探测效率，都可能使速率和方向在记录样本中重新关联。

\subsection{累积分布、典型宽度与尾部估计}
分量高斯密度的自然宽度是 $s=\sqrt{\kb T/m}$，其分布写成
\[
 f_x(v_x)=\frac1{\sqrt{2\pi}s}
 \exp\!\left(-\frac{v_x^2}{2s^2}\right).
\]
两个对称阈值之间的概率是
$P(|v_x|<a)=\erf(a/(\sqrt2s))$，阈值之外则为相应的 $\erfc$。所谓“宽度约为热速率”，只表示密度发生显著变化的尺度；精确的区间概率仍须积分。

对于远尾 $z\gg1$，令 $K(z)=\int_z^\infty\ee^{-t^2}\dd t$，用 $\ee^{-t^2}\dd t=-\dd(\ee^{-t^2})/(2t)$ 分部积分，
\[
 K(z)=\frac{\ee^{-z^2}}{2z}
 -\frac12\int_z^\infty\frac{\ee^{-t^2}}{t^2}\dd t.
\]
被减去的积分为正，因此首项给上界。再利用 $t^{-2}\le z^{-2}$，
\[
 \frac{\ee^{-z^2}}{2z+z^{-1}}
 \le K(z)\le\frac{\ee^{-z^2}}{2z}.
\]
两界之比趋于1，故
$\erfc z\sim\ee^{-z^2}/(\sqrt\pi z)$。指数给主要压低作用，前面的 $1/z$ 来自尾区有效宽度。对较高阈值估算泄流、反应或探测计数时，这一前因子也应保留。

\subsection{速率谱、能量谱与仪器记录的转换}
实际记录的横坐标可以是速度、能量或飞行时间。先把记录事件的密度写成 $p_v(v)$，再换横坐标，才能保持同一批样本。设测量路程固定为 $\ell$，到达时间 $t=\ell/v$，则
\[
 p_t(t)=p_v(\ell/t)\frac{\ell}{t^2},\qquad
 p_E(E)=p_v\!\left(\sqrt{\frac{2E}{m}}\right)
 \frac1{\sqrt{2mE}}.
\]
两个雅可比分别来自 $\left|\dd v/\dd t\right|$ 与
$\left|\dd v/\dd E\right|$。例如，等宽时间窗口并不对应等宽速度窗口：早到粒子的相邻时间窗口覆盖更大的速度区间。图上计数峰的位置因此随记录变量变化，测温时应拟合整个已归一化的记录模型，而不宜只把某个峰当作最概然速率。

若仪器记录效率为 $\eta(v)$，实际事件密度还需乘此因子，
$p_{\rm rec}(v)=\eta(v)p_v(v)/\langle\eta\rangle$。
反推出原谱时，可以在效率非零的区域除去该权重并重新归一化；效率为零的区间没有记录信息，单凭已测区间无法恢复其概率。有限飞行路径分散和有限时间分辨率则引入不同初始条件的叠加，属于在变量变换之外的另一层平均。


\section{相对运动、碰撞与自由程}
\label{th:relative}
取同温、互不相关的两个粒子，质量为 $m_1,m_2$，令 $M_t=m_1+m_2$、$\mu=m_1m_2/M_t$。相对速度与质心速度定义为
\[
 \vct u=\vct v_1-\vct v_2,\qquad
 \vct V_c=\frac{m_1\vct v_1+m_2\vct v_2}{M_t}.
\]
反解为 $\vct v_1=\vct V_c+(m_2/M_t)\vct u$、$\vct v_2=\vct V_c-(m_1/M_t)\vct u$。沿任一坐标方向的变换行列式为
\[
 \det\frac{\partial(v_1,v_2)}{\partial(u,V_c)}
 =\begin{vmatrix}m_2/M_t&1\\-m_1/M_t&1\end{vmatrix}=1.
\]
因而六维测度也满足 $\dd^3v_1\dd^3v_2=\dd^3u\dd^3V_c$。概率变量变换中使用的是行列式绝对值。一般非正交变换把小矩形变成平行四边形，面积不是随意把两个微分相乘所得的结果。

把反解代入动能，交叉项的系数为$2[m_1m_2/M_t-m_2m_1/M_t]=0$，相对速度平方的系数为$[m_1m_2^2+m_2m_1^2]/M_t^2=\mu$，所以
\[
 m_1v_1^2+m_2v_2^2=M_tV_c^2+\mu u^2.
\]
联合密度的指数因而分成$\exp[-M_tV_c^2/(2\kb T)]$与$\exp[-\mu u^2/(2\kb T)]$两因子。对质心速度的积分使用式\eqref{eq:maxwellvector}已证的三维高斯归一化，其值为$(2\pi\kb T/M_t)^{3/2}$。与原联合密度的系数相乘后，利用$m_1m_2/M_t=\mu$，得到
\begin{equation}\label{eq:relative}
 f_u(\vct u)=\left(\frac\mu{2\pi\kb T}\right)^{3/2}
 \ee^{-\mu u^2/(2\kb T)},\qquad
 \bar u=\sqrt{\frac{8\kb T}{\pi\mu}},\quad
 \langle u^2\rangle=\frac{3\kb T}\mu.
\end{equation}
这里 $\langle\vct u\rangle=\vct0$，而 $\bar u=\langle|\vct u|\rangle>0$。矢量平均与速率平均不能混写。

\begin{example}[不同分子之间的碰撞速率]\label{ex:relative}
温度 $T$ 下，气体组分1的质量为 $m_1$、数密度为 $n_1$，组分2为 $m_2,n_2$。硬球半径分别为 $r_1,r_2$。求每个组分1粒子与组分2粒子的平均碰撞频率，以及单位体积、单位时间的异种碰撞次数。
\Diagram{fig:relativeq}{实验室系的两速度与粒子2参照系中的相对运动}{
 \node at(1.1,1.7){实验室系};
 \node[particle,label=below:$1$](a)at(0,0){};\node[particle,label=below:$2$](b)at(2.5,0){};
 \draw[axis](a)--(1.1,.5)node[above]{$\vct v_1$};\draw[axis](b)--(3.3,.5)node[right]{$\vct v_2$};
 \begin{scope}[xshift=5.4cm]
 \node at(1.2,1.7){以粒子2为参照};\node[particle,label=below:$2$]at(2.2,0){};
 \draw[guide](2.2,0)circle(.7);\node[particle,label=below:$1$]at(0,.4){};
 \draw[axis](0,.4)--(1.25,.4)node[midway,above]{$\vct u$};
 \draw[<->](2.2,0)--(2.2,.7);\node[right]at(2.3,.8){$r_1+r_2$};
 \end{scope}
}
\end{example}
\begin{solution}
图\ref{fig:relativeq}中两球接触时中心距为 $r_1+r_2$。把粒子2的中心固定，另一中心以相对速度 $\vct u$ 运动，对应碰撞截面
\[
 \sigma_{12}=\pi(r_1+r_2)^2.
\]
在一小段时间 $\dd t$ 内，相对运动扫过体积为 $\sigma_{12}u\dd t$。先对速度平均，再乘靶粒子数密度，得到单个粒子的平均频率
\[
 z_{12}=n_2\sigma_{12}\bar u
 =n_2\pi(r_1+r_2)^2\sqrt{\frac{8\kb T}{\pi\mu}}.
\]
式中需要的是两粒子中心之间的相对扫掠长度。若直接把组分1的平均速率代入，会把组分2也在运动这一事实漏掉。对独立速度分量有
\[
 \langle u_x^2\rangle=\langle v_{1x}^2\rangle+
 \langle v_{2x}^2\rangle
 =\kb T\left(\frac1{m_1}+\frac1{m_2}\right).
\]
括号正好是 $1/\mu$，与前文相对速度分布的宽度一致。只有求出这一相对分布的一阶速率矩后，才能与截面、靶密度相乘。截面按中心接触条件定义，半径已经包含两球半径之和。

\Diagram{fig:relativea}{相对运动扫掠的碰撞圆柱}{
 \draw[curve](0,-.6)--(4,-.6)arc(-90:90:.2 and .6)--(0,.6);
 \draw[curve](0,0)ellipse(.2 and .6);\draw[axis](.2,0)--(3.8,0)node[midway,above]{$u\dd t$};
 \node at(2,-1.1){横截面 $\sigma_{12}$};
}
图\ref{fig:relativea}的每个圆柱以一个组分1粒子为起点，故体积碰撞率为 $R_{12}=n_1z_{12}$。每一对异种粒子按“1碰2”计数一次，无须再除以二。若统计同一种气体的所有碰撞，同一事件会从两个参与者各计一次，结果应除以二。
\end{solution}

\subsection{平均自由程与无碰撞概率}
\label{th:collisions}
同种硬球半径为 $r$ 时，$\sigma=4\pi r^2$，$\bar u=\sqrt2\bar v$。用平均率 $z=n\sigma\bar u$ 近似碰撞过程，若 $P(t)$ 表示从一次碰撞后持续到 $t$ 仍未再次碰撞的概率，则
\[
 P(t+\dd t)=P(t)(1-z\dd t),\qquad
 \dot P=-zP,\qquad P(0)=1.
\]
分离变量并用初始概率归一化，
\[\int_1^{P(t)}\frac{\dd p}{p}=-z\int_0^t\dd s,
\qquad[\ln p]_1^{P(t)}=-zt,\qquad P(t)=\ee^{-zt}.\]
等待时间密度为$w(t)=-\dot P=z\ee^{-zt}$。其归一化与平均时间分别为
\[\begin{aligned}
\int_0^\infty w(t)\dd t&=[-e^{-zt}]_0^\infty=1,\\
\tau=\int_0^\infty tw(t)\dd t
&=\left[-\left(t+\frac1z\right)e^{-zt}\right]_0^\infty=\frac1z.
\end{aligned}\]
第二个原函数可由$\int t e^{-zt}\dd t$分部积分得到。平均自由程采用分子运动论的平均定义
\begin{equation}
 \lambda=\bar v\tau=\frac{\bar v}{n\sigma\bar u}
=\frac{\bar v}{n\sigma\sqrt2\bar v}=\frac1{\sqrt2n\sigma}.
\end{equation}
严格说，给定速度的分子碰撞率依赖该速度，混合不同速度后的生存概率不一定是单一指数。上面的泊松模型用于平均率估算；把固定靶粒子的扫掠模型与全部粒子均运动的模型混用，会漏掉 $\sqrt2$。

\begin{example}[自由程与稀薄程度]\label{ex:freepath}
硬球气体的有效直径为 $d$，压强 $p$、温度 $T$。求平均自由程；在长为 $L$ 的管中，用平均碰撞模型估计分子穿越全管而不碰撞的概率。说明 $\lambda\gg L$ 与 $\lambda\ll L$ 各自意味着什么。
\Diagram{fig:freepathq}{管长与平均自由程的比较}{
 \draw[curve](0,0)rectangle(5,1);\draw[axis](.3,.5)--(4.6,.5);\draw[<->](0,-.4)--(5,-.4)node[midway,below]{$L$};
 \draw[guide](1,0)--(1,1);\draw[guide](2,0)--(2,1);\node at(1.5,1.4){典型长度 $\lambda$};
}
\end{example}
\begin{solution}
碰撞中心允许的最小距离为 $d$，截面 $\sigma=\pi d^2$。由 $n=p/(\kb T)$ 得
\[
 \lambda=\frac{\kb T}{\sqrt2\pi d^2p}.
\]
若用图\ref{fig:freepathq}中沿管轴的平均运动近似，取穿越时间 $t\simeq L/\bar v$，则 $P\simeq\ee^{-t/\tau}=\ee^{-L/\lambda}$。这不是对所有入射角、所有速度严格平均后的透射率；它衡量的是体碰撞是否重要。
定义 $\mathrm{Kn}=\lambda/L$，则这里的估计是 $P\simeq\ee^{-1/\mathrm{Kn}}$。当 $\mathrm{Kn}\gg1$，展开指数可得 $P\simeq1-L/\lambda$，一次体碰撞都不发生的粒子占绝大多数；当 $\mathrm{Kn}\ll1$，这类粒子的比例指数小。压强在定温下增大一倍，会使数密度和平均碰撞频率增大一倍，自由程减半。上述判断只涉及分子间碰撞；即使体碰撞很少，管壁的反射、吸附和孔道几何仍会影响传输。

\Diagram{fig:freepatha}{两种尺度下的典型轨迹}{
 \draw(0,0)rectangle(2.5,1);\draw[->](.15,.3)--(2.3,.7);\node at(1.25,-.4){$L/\lambda\ll1$};
 \draw(3.6,0)rectangle(6.1,1);\draw[->](3.75,.3)--(4.1,.7)--(4.4,.2)--(4.8,.8)--(5.1,.5)--(5.8,.7);\node at(4.85,-.4){$L/\lambda\gg1$};
}
图\ref{fig:freepatha}左图主要由壁面碰撞和几何边界决定输运；右图体碰撞多，局部平衡较容易建立。但“碰撞多”本身仍不能保证过程可逆，还需宏观变化足够慢、压强和温度梯度足够小。
\end{solution}

\section{活塞、分子流与壁面通量}
\begin{example}[自由管与活塞的相对速度]\label{ex:piston}
水平光滑面上有长管，管质量为 $m_2$，可无摩擦滑动的活塞质量为 $m_1$。管的右端开口，左端封闭，活塞最初位于中点。初始气体体积为 $V_0/2$、温度为 $T_0$、物质的量为 $\nu$。外界真空，装置起初静止。忽略气体质量相对于 $m_1,m_2$ 的贡献，气体可视为准静态绝热膨胀，$C_{V,m}$ 恒定。求活塞刚离开管口时相对于管的速度。
\Diagram{fig:pistonq}{初始半管气体与可反冲的管体}{
 \draw[curve](0,1.3)--(0,0)--(5,0);\draw[curve](0,1.3)--(5,1.3);\draw[line width=2pt](2.5,0)--(2.5,1.3);
 \node at(1.25,.65){$T_0,V_0/2$};\node[above]at(2.5,1.3){$m_1$};\node[below]at(1,-.15){管 $m_2$};\node at(4,.65){真空};
}
\end{example}
\begin{solution}
图\ref{fig:pistonq}的管体也可运动，不能把它当作固定边界。水平总动量为零；设两物体速度为 $v_1,v_2$，有 $m_1v_1+m_2v_2=0$。令 $u=v_1-v_2$，解得
\[
 v_1=\frac{m_2}{m_1+m_2}u,\qquad v_2=-\frac{m_1}{m_1+m_2}u.
\]
将两速度代入动能，得到
\[
 \frac12m_1v_1^2+\frac12m_2v_2^2
 =\frac12\frac{m_1m_2(m_1+m_2)}{(m_1+m_2)^2}u^2
 =\frac12\mu_pu^2,
 \quad \mu_p=\frac{m_1m_2}{m_1+m_2}.
\]
到管口时气体体积从半管增至全管。准静态绝热过程中，第一定律为
\[
 \nu C_{V,m}\dd T=-p\dd V
 =-\frac{\nu RT}{V}\dd V.
\]
除以$\nu C_{V,m}T$，把温度和体积分离后按实际端点积分：
\[
 \int_{T_0}^{T_f}\frac{\dd T}{T}
 =-\frac R{C_{V,m}}\int_{V_0/2}^{V_0}\frac{\dd V}{V},
 \quad [\ln T]_{T_0}^{T_f}
 =-\frac R{C_{V,m}}[\ln V]_{V_0/2}^{V_0}.
\]
于是$\ln(T_f/T_0)=-(R/C_{V,m})\ln2$，指数化得
\[
 T_f=T_0\,2^{1-\gamma},\qquad\gamma=1+\frac R{C_{V,m}},
 \quad T_0(V_0/2)^{\gamma-1}=T_fV_0^{\gamma-1}.
\]
这里的体积是管与活塞之间的相对几何体积，管体反冲不改变它由半管增至全管的事实。对整体使用能量守恒时，管与活塞所受气体力做的功都是内部能量交换，不再另加一次$\int p\dd V$。

\Diagram{fig:pistona}{离开管口前的能量和动量关系}{
 \draw[curve](0,0)rectangle(4.5,1);\draw[line width=2pt](4.5,0)--(4.5,1);\node at(2.1,.5){$T_f,V_0$};
 \draw[axis](4.5,1.4)--(5.7,1.4)node[right]{$v_1$};\draw[axis](1,1.4)--(-.2,1.4)node[left]{$v_2$};
 \node at(2.8,-.6){$\Delta U+\tfrac12\mu_pu^2=0$};
}
图\ref{fig:pistona}中内能减少转化为管和活塞的动能，故
\[
 \frac12\mu_pu^2=\nu C_{V,m}(T_0-T_f),\qquad
 u=\sqrt{\frac{2\nu C_{V,m}T_0}{\mu_p}(1-2^{1-\gamma})}.
\]
该结果以气体能及时重新平衡为前提。估计 $u^2/v_{\rm th}^2\sim\nu M/\mu_p$，所以 $\nu M\ll\mu_p$ 有利于活塞运动远慢于热运动；分子碰撞和声学松弛的时间也须远短于体积显著改变的时间。在这些条件下，气体才可沿一列平衡态演化并采用可逆绝热关系 $TV^{\gamma-1}=\text{常量}$。
\end{solution}

\subsection{运动壁面的碰撞积分}
\label{th:movingwall}
当 $L\ll\lambda$ 时，分子主要在壁面间自由飞行，瞬时速度分布通常带有历史信息。此时要计算碰撞次数及每次交换的动量，不能预先给气体套用平衡绝热线。先处理局部已知入射分布的壁面模型：平面壁沿 $+x$ 方向以速度 $U$ 运动，气体在壁左侧，入射分子必须满足 $v_x>U$。单位面积、单位时间入射数为 $n(v_x-U)f_x(v_x)\dd v_x$，弹性反射给予壁面的冲量为 $2m(v_x-U)$。因此
\begin{equation}\label{eq:movingwall}
 p_w(U)=2mn\int_U^\infty(v_x-U)^2f_x(v_x)\dd v_x.
\end{equation}
这是给定壁前入射分布时的瞬时力，不表示整个有限气体在后续时刻始终保持同一 Maxwell 分布。反射后的分子可能再次撞壁，须随边界更新分布。

\begin{example}[缓慢退让壁面的压强]\label{ex:movingwall}
气体入射分布在实验室系中保持温度 $T$、数密度 $n$ 的 Maxwell 分布。求运动壁面压强的闭式表达，并求 $|U|\ll\sqrt{\kb T/m}$ 时的一阶修正。
\Diagram{fig:movingwallq}{退让壁面只接收追得上它的分子}{
 \draw[line width=1.6pt](3,-.5)--(3,1.5);\node[particle]at(0,.4){};\draw[axis](0,.4)--(1.6,.4)node[midway,above]{$v_x$};
 \draw[axis](3,1.2)--(4.2,1.2)node[right]{$U$};\node at(1.4,-.5){$v_x>U$};
}
\end{example}
\begin{solution}
图\ref{fig:movingwallq}的碰撞速度是 $v_x-U$，它同时出现在通量和冲量中。设 $z=\sqrt\alpha U$，代入式\eqref{eq:movingwall}并令 $t=\sqrt\alpha v_x$，得
\[
 p_w=\frac{2mn}{\alpha\sqrt\pi}\int_z^\infty(t-z)^2\ee^{-t^2}\dd t.
\]
记$A_j=\int_z^\infty t^j\ee^{-t^2}\dd t$。零阶项按$\erfc$的定义为
$A_0=\sqrt\pi\erfc z/2$，一阶项直接使用指数原函数：
\[
 A_1=\lim_{L\to\infty}\left[-\frac12\ee^{-t^2}\right]_z^L
 =\frac12\ee^{-z^2}.
\]
二阶项分部积分，得到
\[
 \begin{aligned}
 A_2&=\left[-\frac t2\ee^{-t^2}\right]_z^\infty
 +\frac12\int_z^\infty\ee^{-t^2}\dd t\\
 &=\frac z2\ee^{-z^2}+\frac{\sqrt\pi}4\erfc z.
 \end{aligned}
\]
展开$(t-z)^2$，原积分为$A_2-2zA_1+z^2A_0$。代入三项并用$1/\alpha=2\kb T/m$，得到
\[
 p_w=n\kb T\left[(1+2z^2)\erfc z-\frac{2z}{\sqrt\pi}\ee^{-z^2}\right].
\]
分部积分的边界项在$z$处保留，在无穷远处为零。压强来自非负被积函数的积分，始终非负；一阶速度修正同时包含碰撞频率和单次冲量的变化。

\Diagram{fig:movingwalla}{壁面功率来自相对速度决定的力}{
 \draw[line width=1.5pt](2,0)--(2,1.5);\draw[axis](2,.4)--(3.4,.4)node[right]{$F=Ap_w$};
 \draw[axis](2,1.2)--(3.1,1.2)node[right]{$U$};\node at(.4,.7){气体};\node at(2.1,-.6){$\dot W=Ap_wU$};
}
在 $z=0$ 邻域，$\erfc z=1-2z/\sqrt\pi+O(z^3)$，所以
\[
 p_w=n\kb T\left[1-\frac4{\sqrt\pi}\sqrt{\frac m{2\kb T}}U
 +O\!\left(\frac{mU^2}{\kb T}\right)\right].
\]
退让 $U>0$ 时压强减小；迎面运动时增大。图\ref{fig:movingwalla}中壁所获功率为 $Ap_wU$，正负随运动方向改变。$U=0$ 时恢复 $n\kb T$，也是积分系数的直接校核。
\end{solution}

\subsection{壁面通量及余弦定律}
\label{th:wallflux}
取面元 $\dd S$，法线与速度夹角为 $\theta$。能够在 $\dd t$ 内撞上该面元的粒子，来自体积 $v\cos\theta\,\dd t\dd S$。各向同性分布的方向概率为 $\dd\Omega/(4\pi)$，且 $\dd\Omega=2\pi\sin\theta\dd\theta$。因此
\[
 \dd N=n g(v)\dd v\frac{\dd\Omega}{4\pi}
 v\cos\theta\dd t\dd S.
\]
只积入射半球 $0\le\theta\le\pi/2$，有
\[
 j=\frac{\dd N}{\dd S\dd t}=n\bar v\int_0^{\pi/2}\frac12\sin\theta\cos\theta\dd\theta
=\frac{n\bar v}{4}[\sin^2\theta]_0^{\pi/2}
 =\frac14n\bar v.
\]
单独给出方向时，撞壁样本的立体角密度为 $\cos\theta/\pi$；对方位角积分后的极角密度为 $2\sin\theta\cos\theta$。它在 $45^\circ$ 处最大，并不意味着单位立体角的通量在斜向最大。

\begin{example}[由撞壁冲量求气体压强]\label{ex:gaspressure}
用各向同性速率分布 $g(v)$ 计算理想弹性壁面的压强，再代入 Maxwell 分布，推出理想气体状态方程。
\Diagram{fig:gaspressureq}{入射与反射的法向动量变化}{
 \draw[line width=1.4pt](3,-1)--(3,1.6);\draw[->](.5,1.4)--(3,0);\draw[->](3,0)--(.5,-1.4);
 \draw[guide](1,0)--(4,0)node[right]{法线};\node at(1.8,.9){$\theta$};\node at(1.7,-1){$\theta$};
}
\end{example}
\begin{solution}
图\ref{fig:gaspressureq}中粒子法向动量从 $mv\cos\theta$ 变为 $-mv\cos\theta$，壁面得到 $2mv\cos\theta$ 的冲量。把该冲量乘入射数，并除以面积和时间，
\begin{align*}
 p&=n\int_0^\infty g(v)\dd v\int_0^{\pi/2}
 \frac12\sin\theta\dd\theta\,(v\cos\theta)(2mv\cos\theta)\\
 &=nm\langle v^2\rangle\int_0^{\pi/2}\sin\theta\cos^2\theta\dd\theta
 =\frac13nm\langle v^2\rangle.
\end{align*}
角积分中令$\mu=\cos\theta$，则$\dd\mu=-\sin\theta\dd\theta$，入射半球对应$\mu$从1降至0，所以
\[
 \int_0^{\pi/2}\sin\theta\cos^2\theta\dd\theta
 =-\int_1^0\mu^2\dd\mu
 =\left[\frac{\mu^3}3\right]_0^1=\frac13.
\]
径向积分留下的是二阶速率矩，而非平均速率的平方。对各向异性气体，沿$x$法向的壁面应使用$p_x=nm\langle v_x^2\rangle$；本题因三个方向等价，才有$\langle v_x^2\rangle=\langle v^2\rangle/3$，并可用一个标量压强描述各个壁面。

\Diagram{fig:gaspressurea}{通量中的法向投影与弹性反射冲量}{
 \draw[line width=1.3pt](2.5,-1.2)--(2.5,1.5);
 \draw[axis](.1,1.1)--(2.5,0)node[midway,above]{$\vct v$};\draw[axis](2.5,0)--(.1,-1.1);
 \draw[guide](.2,0)--(3.2,0);\draw[axis](.2,-1.5)--(2.2,-1.5)node[midway,below]{$v_n=v\cos\theta$};
 \draw(1.6,0)arc[start angle=180,end angle=155.4,radius=.9];\node at(1.36,.23){$\theta$};
 \node[right,align=left]at(3.6,.8){单位时间入射数\\含因子 $v_n$};
 \node[right,align=left]at(3.6,-.6){每次反射给壁面\\冲量 $2mv_n$};
}
图\ref{fig:gaspressurea}中的两个余弦分别来自扫掠体积和冲量；漏掉任何一个都会得到错误角因子。代入 $m\langle v^2\rangle=3\kb T$，得 $p=n\kb T$，再用 $N=\nu N_A$、$R=N_A\kb$，即 $pV=\nu RT$。
\end{solution}

\begin{example}[逸出粒子的速率与能量]\label{ex:effusivesample}
容器中气体处于温度 $T$ 的平衡态。壁上小孔尺寸远小于平均自由程，孔外真空，孔薄且不额外筛选方向。求穿孔粒子的归一化速率密度、平均速率和平均平动能。
\Diagram{fig:effusiveq}{小孔泄流与容器内的粒子抽样}{
 \draw[curve](0,0)--(0,2)--(3,2)--(3,1.15);\draw[curve](3,.85)--(3,0)--(0,0);
 \foreach \p in{(.4,.5),(.8,1.4),(1.6,.8),(2.3,1.5)}\node[particle]at\p{};
 \draw[axis](3,1)--(4.7,1)node[right]{真空};\node at(1.5,-.5){$T,n$};
}
\end{example}
\begin{solution}
图\ref{fig:effusiveq}中的小孔按“单位时间经过”抽样，同速率下快粒子有更多穿孔机会。角积分对各速率都给相同的 $1/4$，故归一化后的密度
\[
 g_e(v)=\frac{vg(v)}{\bar v}=2\alpha^2v^3\ee^{-\alpha v^2},\qquad
 \int_0^\infty g_e\dd v=2\alpha^2 I_3=1.
\]
于是
\[
 \bar v_e=\frac{\langle v^2\rangle}{\bar v}
 =\frac{3\pi}8\bar v,\qquad
 \langle v^2\rangle_e=2\alpha^2 I_5=\frac{4\kb T}{m}.
\]
加权后的密度仍按 $\dd v$ 归一化，平均值计算中的分母始终为同一个总通量。平均速率与平均能量分别涉及原分布的二阶、三阶矩，因而两者受快粒子的影响程度不同。逸出速率密度的峰值由 $3/v-2\alpha v=0$ 确定，对应 $v^2=3/(2\alpha)$；它比体内速率分布的峰向右移动。此处只改变了被观察的样本，容器在这一瞬时的温度仍是原来的 $T$，不能据束流平均能量另给它定义一个平衡温度。

\Diagram{fig:effusivea}{相同温度下体内分布与逸出分布的比较}{
 \draw[axis](0,0)--(5,0)node[right]{$v$};\draw[axis](0,0)--(0,2.3)node[above]{密度};
 \draw[curve]plot[domain=0:4.5,samples=70](\x,{4/sqrt(pi)*\x*\x*exp(-\x*\x)});
 \draw[curve,dashed]plot[domain=0:4.5,samples=70](\x,{2*\x*\x*\x*exp(-\x*\x)});
 \node at(3.3,1.9){实线 $g$};\node at(3.4,1.4){虚线 $g_e$（示意）};
}
逸出粒子的平均平动能为 $2\kb T$，高于体内的 $3\kb T/2$。图\ref{fig:effusivea}显示快粒子权重增加；容器若绝热且不补充气体，余气会冷却。冷却程度要进一步对变化的 $N,T$ 写能量守恒，不能只以“快粒子少了”代替计算。
\end{solution}

\subsection{相对变量独立性的温度条件}
\label{th:relativecorrelation}
质心与相对动能的代数分离对任何速度都成立，但它们在概率上独立，还依赖初始分布。考虑两个分别按温度 $T_1,T_2$ 独立制备的粒子群，在尚未通过碰撞建立共同温度的瞬间，任一方向均有
$\langle v_{1x}^2\rangle=\kb T_1/m_1$、
$\langle v_{2x}^2\rangle=\kb T_2/m_2$。
用第一讲已有的线性变换，
\begin{align*}
 \langle u_x^2\rangle&=\kb\left(\frac{T_1}{m_1}+\frac{T_2}{m_2}\right),\\
 \langle V_{cx}^2\rangle&=
 \frac{\kb(m_1T_1+m_2T_2)}{(m_1+m_2)^2},\\
 \langle V_{cx}u_x\rangle&=
 \frac{\kb(T_1-T_2)}{m_1+m_2}.
\end{align*}
最后一式由展开乘积后使用 $\langle v_{1x}v_{2x}\rangle=0$ 得到。不同温度时交叉协方差一般不为零，联合高斯密度不能再分解为质心、相对两个因子的乘积。等温时交叉项消失，对联合高斯变量，这又足以保证独立，恢复前文的分离结果。

\Diagram{fig:relativecorrelation}{同温与不同温度下联合分布等密度线的示意}{
 \begin{scope}
  \draw[axis](-1.45,0)--(1.6,0)node[right]{$V_{cx}$};
  \draw[axis](0,-1.1)--(0,1.3)node[above]{$u_x$};
  \draw[curve](0,0)ellipse(1 and .55);
  \node at(0,-1.6){$T_1=T_2$};
 \end{scope}
 \begin{scope}[xshift=5.5cm]
  \draw[axis](-1.45,0)--(1.6,0)node[right]{$V_{cx}$};
  \draw[axis](0,-1.1)--(0,1.3)node[above]{$u_x$};
  \draw[curve,rotate=30](0,0)ellipse(1.2 and .4);
  \node at(0,-1.6){$T_1>T_2$，存在正关联};
 \end{scope}
}
图\ref{fig:relativecorrelation}中的坐标各自按合适的速度尺度绘制；倾斜表示条件均值随另一变量改变。即使有这种关联，单独对质心变量积分后的相对速度边缘分布仍是高斯，其宽度就是上式的 $\langle u_x^2\rangle$。因此不同温度的相对速率习题可以只用边缘分布求平均碰撞率；若求碰撞后实验室系的能量分配，则还应保留质心与相对变量之间的关联。

在每个方向上，联合高斯变量的条件均值满足
$\langle V_{cx}\mid u_x\rangle
=\langle V_{cx}u_x\rangle u_x/\langle u_x^2\rangle$。
因此在非同温情况下，选择某个有符号的相对速度分量，也会改变样本的平均质心速度。这里 $T_1,T_2$ 描述制备或瞬时分布；持续碰撞会使两组分交换能量，最终建立共同温度。

\subsection{宏观漂移、热运动与动量通量}
\label{th:pressuretensor}
有整体流动时，实验室速度可分成平均流速与相对流体的热运动：
$\vct v=\vct V_0+\vct c$，其中 $\langle\vct c\rangle=0$。
局域平衡的 Maxwell 分布作用于 $\vct c$。温度由无规则运动的二阶矩确定，实验室系的总平均动能则为
\[
 \frac m2\langle v^2\rangle
 =\frac m2V_0^2+\frac m2\langle c^2\rangle.
\]
第一项是宏观流动能，不能计入温度。不同惯性系中的 $\vct V_0$ 改变，而 $\vct c$ 的统计宽度不变。

穿过法向为 $j$ 方向面元的粒子携带 $i$ 方向动量。每粒子动量为 $mv_i$，穿面通量含 $v_j$，故总动量通量张量为
\[
 \Pi_{ij}=nm\langle v_iv_j\rangle
 =\rho V_{0i}V_{0j}+P_{ij},\qquad
 P_{ij}=nm\langle c_ic_j\rangle,\quad \rho=nm.
\]
交叉项因 $\langle c_i\rangle=0$ 消失。各向同性平衡时 $P_{ij}=n\kb T\,\delta_{ij}$，才可用一个标量压强 $p$ 表示。若分布各向异性，三块不同朝向的壁可能承受不同的法向应力。

\begin{example}[各向异性粒子群的壁面压力]
\label{ex:anisotropicpressure}
长方体容器内有 $N$ 个质量为 $m$ 的独立粒子，速度分布关于各坐标反演对称，分量互不相关，且 $\langle v_i^2\rangle=s_i^2$。壁面均作镜面弹性反射。求三组壁面的压力及总平动能，并说明能否只用平均动能给任意一块壁确定压力。
\Diagram{fig:anisotropicpressureq}{三个方向的速度宽度与相应壁面}{
 \draw[curve](0,0)rectangle(3,2);
 \draw[curve](0,2)--(.9,2.8)--(3.9,2.8)--(3,2);
 \draw[curve](3,0)--(3.9,.8)--(3.9,2.8);
 \draw[axis](1.3,1)--(2.6,1)node[right]{$s_x$};
 \draw[axis](1.3,1)--(1.3,2.2)node[above]{$s_z$};
 \draw[axis](1.3,1)--(2,1.6)node[above left]{$s_y$};
 \node[right]at(4.4,1.3){$n=N/V$};
}
\end{example}
\begin{solution}
图\ref{fig:anisotropicpressureq}中，垂直 $x$ 轴的壁面只反转 $v_x$。对入射半空间求撞壁冲量，沿用前文的“通量乘单次冲量”，有
\[
 p_x=2nm\int_{v_x>0}v_x^2f_x(v_x)\dd v_x
 =nm s_x^2.
\]
第二步利用分量分布的反演对称性，把半轴二阶矩化成整轴的一半。相同方法给 $p_y=nm s_y^2$、$p_z=nm s_z^2$。总能量为
\[
 U_{\rm tr}=\frac{Nm}{2}(s_x^2+s_y^2+s_z^2)
 =\frac V2(p_x+p_y+p_z).
\]
\Diagram{fig:anisotropicpressurea}{法向二阶矩分别决定三组壁面压力}{
 \draw[axis](0,0)--(5,0)node[right]{壁面法向};
 \foreach \x/\h/\t in{1/1.1/x,2.5/1.8/y,4/.7/z}{
  \draw[fill=gray!15](\x-.3,0)rectangle(\x+.3,\h);
  \node[below]at(\x,0){$\t$};
  \node[above]at(\x,\h){$p_{\t}$};
 }
}
图\ref{fig:anisotropicpressurea}示意三方向不等的压力。给定 $U_{\rm tr}$ 只固定三个压力之和，仍可有不同的分配。只有各向同性时 $p_x=p_y=p_z=p$，才能进一步写成 $pV=2U_{\rm tr}/3$。若分子间碰撞逐渐使三方向能量均分，在体积固定且不向外散热的条件下，总能量不变，最终共同压力等于原先三个压力的算术平均。
\end{solution}

\subsection{碰撞事件的速率权重}
两组分在同温下的相对速率密度已由式\eqref{eq:relative}确定。某个相对速率 $u$ 对碰撞率的贡献，除概率密度外还乘扫掠速率 $u$ 和截面 $\sigma(u)$，故
\[
 K(T)=\langle\sigma u\rangle
 =\int_0^\infty\sigma(u)u\,g_u(u)\dd u.
\]
$K$ 的量纲为体积除时间，乘 $n_1n_2$ 后才成为单位体积、单位时间的异种碰撞数。按碰撞事件抽样的密度是 $\sigma(u)u g_u(u)/K$。截面随能量变化时，即使热平衡的 $g_u$ 未变，碰撞样本的分布也会改变。

\begin{example}[速率相关截面的碰撞平均]
\label{ex:powercrosssection}
两种粒子处于共同温度 $T$，约化质量为 $\mu$。在所讨论速率范围内，碰撞截面可近似为 $\sigma(u)=\sigma_\ast(u/u_\ast)^q$，其中 $u_\ast>0$、$q>-4$。求速率系数 $K(T)$，以及按碰撞事件抽样的平均相对动能。
\Diagram{fig:powercrosssectionq}{相对速率分布经扫掠与截面两次取权}{
 \node[draw](a)at(0,0){$g_u(u)\dd u$};
 \node[draw](b)at(3.1,0){$u\,\sigma(u)$};
 \draw[axis](a)--(b);
 \node at(1.55,.6){碰撞机会};
 \node[draw](c)at(6.2,0){事件样本};
 \draw[axis](b)--(c);
}
\end{example}
\begin{solution}
图\ref{fig:powercrosssectionq}中的权重为
$\sigma_\ast u^{q+1}/u_\ast^q$。用前文的一般 Maxwell 矩，把其中的质量换成 $\mu$，
\[
 K=\frac{\sigma_\ast}{u_\ast^q}
 \frac2{\sqrt\pi}\Gamma\!\left(\frac{q+4}{2}\right)
 \left(\frac{2\kb T}{\mu}\right)^{(q+1)/2}.
\]
积分原点附近的被积函数正比于 $u^{q+3}$，收敛要求 $q>-4$；无穷远端由高斯因子保证收敛。这个条件对应题设的指数范围。若真实截面只在有限速率区间服从幂律，则上式只适用于主要碰撞事件集中在该区间的温度。

事件平均必须用同一总碰撞率作分母，因此
\[
\begin{aligned}
 \langle E_r\rangle_{\rm coll}
 &=\frac{\mu}{2}\frac{\langle u^{q+3}\rangle}
 {\langle u^{q+1}\rangle}
 =\frac{\mu}{2}\frac{I_{q+5}(\alpha_u)}{I_{q+3}(\alpha_u)}\\
 &=\frac{\mu}{2}\frac{q+4}{2\alpha_u}
 =\frac{q+4}{2}\kb T,\qquad \alpha_u=\frac\mu{2\kb T}.
\end{aligned}
\]
第二行使用第\ref{th:integrals}节已由分部积分证明的递推式，取$s=q+3>-1$；两个速率矩具有相同归一化系数，作比后已消去。
\Diagram{fig:powercrosssectiona}{截面指数改变碰撞速率与事件能量}{
 \draw[axis](0,0)--(5,0)node[right]{$q$};
 \draw[axis](0,0)--(0,3)node[above]{$\langle E_r\rangle_{\rm coll}/(\kb T)$};
 \draw[curve](.6,.3)--(4.5,2.25);
 \fill(2,1)circle(1.2pt)node[below right]{$(-2,1)$};
 \fill(4,2)circle(1.2pt)node[above left]{$(0,2)$};
 \node[below]at(0,0){$-4$};
}
图\ref{fig:powercrosssectiona}中的直线只在 $q>-4$ 适用。常截面 $q=0$ 给出 $2\kb T$，比随机粒子对的 $3\kb T/2$ 大；$q=-1$ 时 $\sigma u$ 为常数，事件取样与原相对速率取样一致，平均能量恢复为 $3\kb T/2$。两个极限分别检验了碰撞加权和归一化。
\end{solution}

\subsection{局域平衡与边界过程的时间尺度}
一个容器中的总粒子数和总能量不足以确定随时间演化的分布。采用瞬时平衡的 Maxwell 形式，还要求碰撞与宏观变化之间有时间尺度分离。典型碰撞时间为 $\tau_c\sim\lambda/\bar v$，声学信息穿越尺度 $L$ 的时间约为 $\tau_s\sim L/c_s$。若体积或温度显著改变的时间为 $\tau_{\rm mac}$，准静态处理至少要求相关的松弛时间远小于 $\tau_{\rm mac}$。

小孔泄流还要求“孔附近的传输”与“容器内部的重新平衡”分别满足条件。孔径远小于平均自由程，可忽略孔口内的分子间碰撞；容器内却可通过足够多的体碰撞或壁面热化保持平衡。因而分子流小孔与平衡余气可以同时存在。若整个容器都处于无碰撞极限且壁面只作镜面反射，最初分布随边界运动会发生形变，须跟踪速度方向及粒子飞行历史。

宏观运动缓慢时，还应判断各处是否接近平衡。有限压差推动活塞、有限温差传热及流动耗散都会产生熵；降低平均速度未必能同时消除这些因素。后两讲将按具体能量交换方式和过程约束计算相应的熵产生。


\section{碰撞抽样与运动边界的进一步应用}
同一气体的体内样本、穿孔样本和碰撞样本具有不同的速率权重。体内样本按粒子数计数；穿孔和碰撞则按单位时间发生的事件计数，要在体内密度上乘相应的相对通量，再归一化。

当微观速度与宏观边界速度相差很大时，平衡压强可作为边界受力的首项；当二者可比时，要保留相对速度、入射阈值和反射规律。下面将这些条件分别展开。

\begin{example}[由球坐标变换核算方向概率]
三维速度分布各向同性，矢量密度为 $f(v)$。从球坐标的雅可比出发求 $(v,\theta,\phi)$ 的联合概率密度，证明速率与方向独立，并求 $v_z/v$ 的密度。
\Diagram{fig:angularjacq}{球坐标中的极角、方位角和速度分量}{
 \begin{scope}[x={(1cm,0)},y={(-.55cm,-.4cm)},z={(0,1cm)}]
 \draw[axis](0,0,0)--(2.5,0,0)node[right]{$v_x$};\draw[axis](0,0,0)--(0,2.3,0)node[below left]{$v_y$};\draw[axis](0,0,0)--(0,0,2.7)node[above]{$v_z$};
 \draw[axis](0,0,0)--(1.448,1.014,1.768)node[above right]{$\vct v$};
 \draw[guide](1.448,1.014,1.768)--(1.448,1.014,0);\draw[guide](0,0,0)--(1.448,1.014,0);
 \draw[curve]plot[domain=0:45]({.75*sin(\x)*cos(35)},{.75*sin(\x)*sin(35)},{.75*cos(\x)});
 \node at(.5,.15,.9){$\theta$};
 \draw[curve]plot[domain=0:35]({.8*cos(\x)},{.8*sin(\x)},0);\node at(.95,.35,0){$\phi$};
 \end{scope}
}
\end{example}
\begin{solution}
图\ref{fig:angularjacq}的坐标变换为 $v_x=v\sin\theta\cos\phi$、$v_y=v\sin\theta\sin\phi$、$v_z=v\cos\theta$。三个坐标微元方向彼此正交，其长度分别为 $\dd v$、$v\dd\theta$、$v\sin\theta\dd\phi$，故体积元为 $v^2\sin\theta\dd v\dd\theta\dd\phi$。于是
\[
 \dd P=f(v)v^2\dd v\sin\theta\dd\theta\dd\phi
 =g(v)\dd v\frac{\sin\theta\dd\theta\dd\phi}{4\pi}.
\]
右侧两个因子分别归一化，故速率与方向独立。方位角均匀分布；极角则按球带面积取权，密度为 $q(\theta)=\sin\theta/2$。
方向分布的归一化由$[\phi]_0^{2\pi}=2\pi$和$[-\cos\theta]_0^\pi=2$给出。令$\mu_z=\cos\theta=v_z/v$，则$|\dd\mu_z|=\sin\theta\dd\theta$，所以$\mu_z$在$[-1,1]$上有均匀密度$q_\mu=1/2$。其前两个矩为
\[
 \langle\mu_z\rangle=\frac12\left[\frac{\mu_z^2}{2}\right]_{-1}^1=0,
 \qquad \langle\mu_z^2\rangle=\frac12\left[\frac{\mu_z^3}{3}\right]_{-1}^1=\frac13.
\]
一阶矩为零说明没有优选方向；二阶矩$1/3$使速度平方的平均值均分到三个坐标方向。

\Diagram{fig:angularjaca}{相等极角宽度的球带面积不同}{
 \draw[curve](0,0)circle(1.6);\draw[curve](-1.6,0)--(1.6,0);\draw[guide](-1.38,.8)--(1.38,.8);\draw[guide](-.7,1.44)--(.7,1.44);\node[right]at(2,1){球带面积 $2\pi v^2\sin\theta\dd\theta$};
}
图\ref{fig:angularjaca}的球带面积随$\theta$变化。各向同性对应等立体角等概率；若把极角本身误设为均匀分布，便会得到错误的$\langle\cos^2\theta\rangle=1/2$。
\end{solution}

\begin{example}[固定速度粒子的相对碰撞频率]
背景气体分子速度大小均为 $w$，方向各向同性，数密度 $n$。一个试探粒子以固定速度大小 $v$ 运动，与背景粒子的硬球截面为常量 $\sigma$。求角度平均相对速率和试探粒子碰撞频率；分别讨论 $v>w$、$v<w$。
\Diagram{fig:fixedrelativeq}{固定试探速度与背景方向之间的夹角}{
 \draw[axis](0,0)--(3.2,0)node[right]{$\vct v$};\draw[axis](0,0)--(1.4,1.8)node[above]{$\vct w$};\draw[axis](1.4,1.8)--(3.2,0)node[midway,right]{$\vct v-\vct w$};\node at(.8,.35){$\theta$};
}
\end{example}
\begin{solution}
图\ref{fig:fixedrelativeq}中 $u=\sqrt{v^2+w^2-2vw\cos\theta}$。由于 $\cos\theta$ 在 $[-1,1]$ 上均匀，
\[
 \bar u(v,w)=\frac12\int_{-1}^1\sqrt{v^2+w^2-2vw\mu}\dd\mu.
\]
代换 $z=v^2+w^2-2vw\mu$，$\dd\mu=-\dd z/(2vw)$，积分端点为 $(v-w)^2$ 与 $(v+w)^2$，所以
\[
\begin{aligned}
\bar u(v,w)&=\frac1{4vw}\int_{(v-w)^2}^{(v+w)^2}\sqrt z\dd z
=\frac1{4vw}\left[\frac23z^{3/2}\right]_{(v-w)^2}^{(v+w)^2}\\
&=\frac{(v+w)^3-|v-w|^3}{6vw}.
\end{aligned}
\]
当$v>w>0$时，分子为$(v+w)^3-(v-w)^3=6v^2w+2w^3$；当$w>v>0$时，交换$v,w$即可。因此
\[\bar u(v,w)=\begin{cases}
v+w^2/(3v),&v\ge w>0,\\
w+v^2/(3w),&w\ge v>0.
\end{cases}\]
某一速率为零时取连续极限，分别为$\bar u(v,0)=v$和$\bar u(0,w)=w$；两者都为零时碰撞率为零。
同一变量变换也给出相对速率的允许区间：两速度平行时相对速率最小，反平行时最大，即 $|v-w|\le u\le v+w$。由 $2u\dd u=-2vw\dd\mu$，得到方向样本的密度 $q(u)=u/(2vw)$。积分区间的两端平方相差 $4vw$，因而该密度归一化。碰撞率还要再乘一个 $u$，所以即使背景速度大小固定，接近反向运动的粒子对也会在碰撞事件中占更高权重。

\Diagram{fig:fixedrelativea}{两种速度大小次序给出分段表达}{
 \node[draw,align=center]at(0,0){$v\ge w$\\$v+w^2/(3v)$};\node[draw,align=center]at(4.6,0){$w\ge v$\\$w+v^2/(3w)$};\node at(2.3,-1){$z(v)=n\sigma\bar u(v,w)$};
}
展开立方即得图\ref{fig:fixedrelativea}的分段式。在 $v=w$ 时两式均为 $4v/3$，连续衔接；$v\to0$ 时结果为 $w$，快速试探粒子 $v\gg w$ 时趋于 $v$。碰撞频率 $z(v)=n\sigma\bar u(v,w)$ 随试探粒子的速度而变。对 Maxwell 背景，还须用背景速率密度对 $w$ 再作一次平均。
\end{solution}

\begin{example}[按碰撞事件抽样的相对能量]
两个同温平衡组分相互独立，相对速度对应约化质量 $\mu$。碰撞截面为与速率无关的常量。求随机抽取一次碰撞时的相对速率密度、平均相对平动能，以及相对能量超过阈值 $E_a$ 的碰撞比例。
\Diagram{fig:collisionenergyq}{一次事件的权重含有相对扫掠速度}{
 \draw[curve](0,0)ellipse(.25 and .6);\draw[curve](0,.6)--(3,.6);\draw[curve](0,-.6)--(3,-.6);\draw[axis](.3,0)--(2.7,0)node[midway,above]{$u$};\node at(1.5,-1.1){事件率 $\propto u\sigma$};
}
\end{example}
\begin{solution}
相对速率的体样本密度为 $g_u(u)\propto u^2\ee^{-\alpha_u u^2}$，$\alpha_u=\mu/(2\kb T)$。图\ref{fig:collisionenergyq}中单位时间碰撞次数还乘 $u\sigma$，故事件样本密度为
\[
 g_{\rm coll}(u)=\frac{ug_u(u)}{\bar u}=2\alpha_u^2u^3\ee^{-\alpha_uu^2}.
\]
这与小孔泄流的速率加权结构相同，平均相对能量为 $\mu\langle u^2\rangle_{\rm coll}/2=2\kb T$。
把速率密度变为能量密度时，$\dd E=\mu u\dd u$。因此 $u^3\dd u=u^2\dd E/\mu$，速率三次因子在能量变量中变成一次因子。尾概率的积分可用
\[
 \int x\ee^{-x}\dd x=-(x+1)\ee^{-x}
\]
计算，下限贡献 $(x_a+1)\ee^{-x_a}$，上限为零。这里是碰撞发生以后再问其中多少超过阈值，分母已经是全部碰撞次数；若从瞬时任意粒子对中抽样，则应使用未经速率加权的相对能量分布。

\Diagram{fig:collisionenergya}{碰撞相对能量的归一化密度}{
 \draw[axis](0,0)--(5,0)node[right]{$E/(\kb T)$};\draw[axis](0,0)--(0,1.5)node[above]{$\kb T h_{\rm coll}$};\draw[curve]plot[domain=0:4.6](\x,{\x*exp(-\x)});\draw[guide](2.2,0)--(2.2,.8);\node[below]at(2.2,0){$E_a/(\kb T)$};
}
令 $E=\mu u^2/2$，变换密度得 $h_{\rm coll}(E)=E\ee^{-E/(\kb T)}/(\kb T)^2$。图\ref{fig:collisionenergya}阈值右侧积分，令 $x=E/(\kb T)$，
\[
 P(E>E_a)=\int_{x_a}^\infty x\ee^{-x}\dd x=(1+x_a)\ee^{-x_a},\quad x_a=E_a/(\kb T).
\]
所得比例对应“总相对动能超过阈值”的判据。若反应要求沿连心线的能量超过阈值，允许的碰撞参数范围也随能量变化，须进一步对反应截面取平均，下一例讨论这一情形。
\end{solution}

\begin{example}[带能量阈值的硬球反应截面]
两组分碰撞时，只有接触前沿两球连心线的相对动能超过 $E_a$ 才反应。硬球接触中心距为 $d$，相对速率为 $u$，碰撞参数为 $b$。求反应截面 $\sigma_r(u)$，再求同温 Maxwell 相对运动下的速率系数 $\langle\sigma_ru\rangle$。
\Diagram{fig:lineofcentersq}{接触点\(P\)处的法向分量与碰撞参数}{
 \draw[curve](3,0)circle(1);\draw[axis](0,.6)--(2.2,.6)node[midway,above]{$u$};
 \draw[guide](0,0)--(3,0);\draw[<->](.4,0)--(.4,.6)node[midway,left]{$b$};
 \draw[curve](3,0)--(2.2,.6)node[midway,below left]{$d$};\fill(2.2,.6)circle(1.1pt)node[above left]{$P$};
 \draw[guide](2.2,.6)--(3,.6);\draw[curve](2.75,.6)arc(0:-36.87:.55);\node at(2.94,.38){$\theta$};
}
\end{example}
\begin{solution}
图\ref{fig:lineofcentersq}接触时，入射方向与连心线夹角满足 $\sin\theta=b/d$，所以法向相对能量为 $E_n=\mu u^2\cos^2\theta/2=E(1-b^2/d^2)$。反应要求 $b^2<d^2(1-E_a/E)$，且先要有 $E\ge E_a$。因此
\[
 \sigma_r(E)=\pi d^2\left(1-\frac{E_a}E\right),\quad E\ge E_a;
 \qquad \sigma_r=0,\ E<E_a.
\]
入射中心的落点在截面平面上按面积计数，碰撞参数位于$b$到$b+\dd b$的环带面积为$2\pi b\dd b$。因此允许的截面为$\int_0^{b_{\max}}2\pi b\dd b=[\pi b^2]_0^{b_{\max}}=\pi b_{\max}^2$。阈值恰被达到时，只允许正碰这一零面积集合，截面从零连续开始；能量远大于阈值时，允许圆盘趋于整个硬球截面。下面将这个面积因子与各相对能量处的事件密度相乘。

\Diagram{fig:lineofcentersa}{允许反应的碰撞参数圆盘随能量增大}{
 \draw[curve](0,0)circle(1.4);\fill[gray!20](0,0)circle(.9);\draw[curve](0,0)circle(.9);\draw[axis](0,0)--(.9,0)node[midway,above]{$b_{\max}$};\node[right]at(1.7,0){全碰撞半径 $d$};
}
图\ref{fig:lineofcentersa}的有效截面已经包含碰撞几何筛选。将其乘 $u$ 后对体样本平均，也可利用上一例的事件能量密度：
\begin{align*}
 \langle\sigma_ru\rangle
 &=\pi d^2\bar u\int_{E_a}^\infty
 \left(1-\frac{E_a}E\right)\frac{E}{(\kb T)^2}\ee^{-E/(\kb T)}\dd E\\
 &=\pi d^2\bar u\int_{x_a}^\infty(x-x_a)\ee^{-x}\dd x\\
 &=\pi d^2\bar u[-(x+1-x_a)e^{-x}]_{x_a}^\infty
 =\pi d^2\bar u\ee^{-x_a}.
\end{align*}
这里指数因子与上一例的 $(1+x_a)\ee^{-x_a}$ 不同，差别来自允许反应的碰撞参数面积随能量而变化。反应体积率为 $n_1n_2\langle\sigma_ru\rangle$，同种反应的无序对计数还须除二。
\end{solution}

\begin{example}[撞壁样本的法向与切向能量]
壁面静止，气体处于温度 $T$ 的平衡。按撞上壁面的事件随机选一个分子，求其入射法向速度 $v_n>0$ 的密度，以及法向动能和两个切向动能的平均值。
\Diagram{fig:wallcomponentsq}{碰撞抽样只对法向速度加权}{
 \draw[line width=1.2pt](3,-1)--(3,1.8);\draw[axis](0,0)--(3,1.2)node[midway,above]{$\vct v$};\draw[axis](0,0)--(3,0)node[midway,below]{$v_n$};\draw[axis](3,0)--(3,1.2)node[right]{$v_t$};
}
\end{example}
\begin{solution}
图\ref{fig:wallcomponentsq}的壁面每秒截获厚度 $v_n\dd t$ 内的粒子，所以在原高斯因子上乘 $v_n$。归一化分母为 $\int_0^\infty v_nf_n(v_n)\dd v_n=1/(2\sqrt{\pi\alpha})$，得到
\[
 q_n(v_n)=2\alpha v_n\ee^{-\alpha v_n^2},\quad v_n\ge0.
\]
法向二阶矩为$\int_0^\infty v_n^2q_n(v_n)\dd v_n=2\alpha I_3=1/\alpha=2\kb T/m$，故平均法向动能为$\kb T$。联合事件密度仍分离为$q_n(v_n)f_y(v_y)f_z(v_z)$，归一化权重只含法向变量，两个切向高斯因子各自积分为1，平均动能仍各为$\kb T/2$。三方向相加为$2\kb T$；法向能量偏高来自事件筛选，未经过碰撞的体内分布并没有被壁面加热。

\Diagram{fig:wallcomponentsa}{逸出平均平动能按壁面方向分解}{
 \node[draw,align=center]at(0,0){法向\\$\kb T$};\node[draw,align=center]at(3,0){切向1\\$\kb T/2$};\node[draw,align=center]at(6,0){切向2\\$\kb T/2$};\node at(3,-1){总和 $2\kb T$};
}
图\ref{fig:wallcomponentsa}给出了逸出能量 $2\kb T$ 的另一路线。法向速度的平均值为 $2\alpha I_2=\sqrt\pi/(2\sqrt\alpha)$；它不同于从体内只筛选 $v_n>0$ 得到的条件平均，也不同于三维速率平均。
\end{solution}

\begin{example}[运动镜面的单次碰撞能量交换]
质量 $m$ 的粒子沿 $+x$ 方向以速度 $v$ 追上以速度 $U<v$ 退让的重壁，碰撞弹性，壁速度不变。求反射后的粒子速度、粒子能量减少与壁获得的功，并将结果与壁面受力积分相联系。
\Diagram{fig:singlewallq}{在实验室系与壁面系中比较碰撞}{
 \draw[line width=1.4pt](3,-.4)--(3,1.5);\draw[axis](0,.5)--(1.7,.5)node[midway,above]{$v$};\draw[axis](3,1.2)--(4.3,1.2)node[right]{$U$};\node at(1.5,-.6){壁面系入射速度 $v-U$};
}
\end{example}
\begin{solution}
先转到图\ref{fig:singlewallq}的壁面系，碰撞前相对速度 $v-U$，碰后变为 $-(v-U)$。变回实验室系，$v'=U-(v-U)=2U-v$。壁面获得的冲量为 $\Delta p_w=m(v-v')=2m(v-U)$，对应功 $\Delta W_w=U\Delta p_w$。
本题的壁速在碰撞前后保持不变，可由大质量壁面或外界驱动实现。若换成有限质量的自由活塞，就须联立两物体的动量、动能守恒方程，求碰后的粒子速度和活塞速度。

\Diagram{fig:singlewalla}{壁面系反射后变回实验室系}{
 \node[draw]at(0,0){$v-U$};\draw[axis](.8,0)--(2.2,0)node[midway,above]{反射};\node[draw]at(3.4,0){$-(v-U)$};\draw[axis](4.6,0)--(5.8,0)node[midway,above]{$+U$};\node[draw]at(6.5,0){$v'$};
}
直接算粒子能量减少，
\[
 \frac m2[v^2-(2U-v)^2]
 =\frac m2(4Uv-4U^2)=2mU(v-U)=\Delta W_w.
\]
该式保留$U^2$项，对满足追赶条件$v>U$的碰撞没有使用慢壁近似。$U>0$时分子失能，迎面运动的$U<0$情形则使分子增能。碰后$v'$可以仍为正：粒子在实验室系中继续向前，但速度小于壁速，因而与壁面分离。

将图\ref{fig:singlewalla}的单次能量交换乘事件率 $n(v-U)f_x(v)\dd v$，并在 $v>U$ 上积分，得到能量传递率 $Up_w(U)$。它与压强积分对应，反映了同一批碰撞中动量和能量的交换。
\end{solution}

\begin{example}[一维无碰撞气体的缓慢活塞膨胀]
粒子被限制在左端固定、右端缓慢退让的长为 $L(t)$ 的一维盒中。粒子间不碰撞，活塞速度 $U=\dot L$ 远小于粒子往返速度，左壁弹性反射。求单粒子速度大小随 $L$ 的变化，进而求均匀初始相位样本的轴向能量标度。与三维平衡气体绝热关系比较。
\Diagram{fig:ballisticpistonq}{粒子在两壁之间往返，仅运动壁改变能量}{
 \draw[line width=1.4pt](0,-.6)--(0,1);\draw[line width=1.4pt](4,-.6)--(4,1);\draw[<->](.3,.2)--(3.7,.2)node[midway,above]{$v$};\draw[axis](4,.8)--(5,.8)node[right]{$U$};\draw[<->](0,-1)--(4,-1)node[midway,below]{$L(t)$};
}
\end{example}
\begin{solution}
图\ref{fig:ballisticpistonq}中粒子每与退让活塞碰撞一次，速率由 $v$ 减为 $v-2U$，忽略 $U/v$ 的高阶修正。相邻两次撞活塞的间隔约为 $2L/v$，所以平均变化率
\[
 \frac{\dd v}{\dd t}\simeq\frac{-2U}{2L/v}=-\frac{\dot L}L v.
\]
设$L=L_0$时该粒子的速率为$v_0>0$，分离变量并积分，
\[\int_{v_0}^{v(L)}\frac{\dd s}{s}=-\int_{L_0}^{L}\frac{\dd\ell}{\ell},
\qquad[\ln s]_{v_0}^{v(L)}=-[\ln\ell]_{L_0}^{L}.\]
故$\ln[v(L)/v_0]=-\ln(L/L_0)$，即$v(L)=v_0L_0/L$。因此$vL$为缓慢变化下的不变量，单粒子轴向能量为$mv_0^2(L_0/L)^2/2$。
若初态为各向同性温度 $T_0$，而随后只有 $x$ 方向长度由 $L_0$ 改为 $L$，可把三方向二阶矩分别记为 $m\langle v_i^2\rangle=\kb T_i$。于是
\[
 T_x=T_0(L_0/L)^2,\qquad T_y=T_z=T_0,
\]
\[
 U_{\rm tr}=\frac{N\kb T_0}{2}\bigl[(L_0/L)^2+2\bigr].
\]
这些 $T_i$ 只是分方向能量参数。它们不同，表明分布已偏离单温度 Maxwell 平衡态；横向能量不会自动补给不断损失轴向能量的分子。

\Diagram{fig:ballisticpistona}{缓慢循环碰撞给出轴向能量的平方反比律}{
 \draw[axis](0,0)--(5,0)node[right]{$L/L_0$};\draw[axis](0,0)--(0,2.8)node[above]{$E_x/E_{x0}$};\draw[curve]plot[domain=1:4.5](\x,{2.4/\x^2});\node[below]at(1,0){$1$};\node[left]at(0,2.4){$1$};
}
图\ref{fig:ballisticpistona}的标度也适用于对初始相位平均的独立粒子样本，但要求变化慢于各粒子的往返时间。非常慢的尾部粒子可能违反这一条件，不能把近似无限外推。

三维无碰撞盒只沿 $x$ 膨胀时，$E_x\propto L^{-2}$，横向 $E_y,E_z$ 不变，分布因而具有各向异性，需要分别描述三个方向的能量尺度。若碰撞能及时重新分配能量，气体才恢复单一平衡温度，并满足单原子气体的 $TV^{2/3}=\text{常量}$。边界运动缓慢与内部达到热平衡须分别判断。
\end{solution}

\begin{example}[时间飞行测温中的两种探测权重]
一个小孔发出的短脉冲分子束沿法线传播，速度分布取逸出速率密度 $g_e(v)$。探测器在距离 $L$ 处，接受极窄角度，对到达的每个粒子等效计数，忽略碰撞。求到达时间密度及峰值；若探测器测量固定时刻某薄层中的粒子数，权重有何变化？另比较同一速率谱的稳态束流在薄层中的瞬时计数。
\Diagram{fig:tofthermq}{短脉冲穿过固定飞行距离后到达计数器}{
 \draw[line width=1.2pt](0,-.6)--(0,1);\draw[line width=1.2pt](5,-.6)--(5,1);\draw[axis](.2,.2)--(4.7,.2)node[midway,above]{$v=L/t$};\draw[<->](0,-1)--(5,-1)node[midway,below]{$L$};\node at(0,1.4){孔};\node at(5,1.4){计数器};
}
\end{example}
\begin{solution}
图\ref{fig:tofthermq}中速度与时间一一对应，$v=L/t$、$|\dd v/\dd t|=L/t^2$。窄角度筛选不改变速率分布，因此到达时间密度
\[
 q_t(t)=g_e(L/t)\frac L{t^2}
 =2\alpha^2L^4t^{-5}\exp(-\alpha L^2/t^2),\quad t>0.
\]
对数求导，$\dd\ln q_t/\dd t=-5/t+2\alpha L^2/t^3$，峰值满足 $t_p^2=2\alpha L^2/5$，所以 $t_p=L\sqrt{m/(5\kb T)}$。到达时间密度的峰值还受雅可比 $L/t^2$ 影响，因此应直接对 $q_t$ 求极值，而不能把速率峰值简单取倒数。
到达时间密度的归一化可直接用 $v=L/t$ 验证：当 $t$ 从0增至无穷大，$v$ 从无穷大降至0，雅可比的绝对值抵消积分方向的反转，故总概率仍为1。长时间尾部有 $q_t\propto t^{-5}$，来自束流中很慢的分子；极短时间则被指数因子强烈抑制。仪器若按粒子到达事件计数，每个粒子只记一次；这与在固定薄层内持续观察、因停留时间而多看见慢粒子的情形不同。

\Diagram{fig:toftherma}{时间分布的峰值由雅可比与指数共同决定}{
 \draw[axis](0,0)--(5.4,0)node[right]{$t$};\draw[axis](0,0)--(0,2.8)node[above]{$q_t$};\draw[curve]plot[domain=.25:4.8,samples=90](\x,{2*exp(-1/(\x*\x))/\x^5});\node at(3.7,2.4){事件计数};
}
图\ref{fig:toftherma}是脉冲的到达时间密度。固定时刻的空间计数须另用$x=vt$作变量变换。把短脉冲的发射历时近似为零，接受方向内共含$N_p$个粒子；固定$t>0$时，$\dd x=t\dd v$，因此
\[
 \dd N=N_pg_e(v)\dd v
 =\frac{N_p}{t}g_e(x/t)\dd x.
\]
对图\ref{fig:toftherm-layer-a}中以$L$为中心、厚度为$\Delta\ell$的薄层，瞬时粒子数为
\[
 N_{\rm layer}(t)
 =N_p\int_{(L-\Delta\ell/2)/t}^{(L+\Delta\ell/2)/t}g_e(v)\dd v
 \simeq\frac{N_p\Delta\ell}{t}g_e(L/t).
\]
末式要求该窄速度区间内$g_e$变化很小。固定$t$和$L$已把速率限制在$L/t$附近，不能把整个脉冲的速率密度直接乘$1/v$来表示这一快照。

\Diagram{fig:toftherm-layer-a}{脉冲快照按位置筛选；稳态薄层计数另含停留时间}{
 \node[anchor=west]at(0,2.9){脉冲：固定$t$};
 \draw[axis](0,1.7)--(6.3,1.7)node[right]{$x$};
 \fill[black!8](3.4,1.8)rectangle(4.4,2.5);
 \draw[dashed](3.4,1.6)--(3.4,2.55);\draw[dashed](4.4,1.6)--(4.4,2.55);
 \foreach \x in{1.1,2.2,3.9,5.2}\fill(\x,2.15)circle(1.5pt);
 \draw[<->](3.4,2.7)--(4.4,2.7)node[midway,above]{$\Delta\ell$};
 \node[below]at(3.9,1.7){$L$};\node[above]at(1.6,2.15){$x=vt$};
 \node[anchor=west]at(0,.65){稳态：连续发射};
 \draw[axis](0,-.15)--(6.3,-.15)node[right]{$x$};
 \fill[black!8](3.4,-.05)rectangle(4.4,.45);
 \draw[axis](2.3,.2)--(4.9,.2)node[right]{$v$};
 \node[below]at(3.9,-.4){停留时间$\Delta\ell/v$};
}
稳态束流连续补充各速率的粒子。若总发射率为$J$，速率区间内的到达率为$Jg_e(v)\dd v$，每个粒子在薄层中停留$\Delta\ell/v$，故
\[
 \dd N_{\rm steady}=Jg_e(v)\dd v\frac{\Delta\ell}{v},\qquad
 q_{\rm space}(v)=\frac{g_e(v)/v}{\int_0^\infty g_e(w)\dd w/w}=g(v).
\]
最后用了$g_e(v)=vg(v)/\bar v$及$\int_0^\infty g(v)\dd v=1$。这里的$1/v$来自稳态停留时间；若对一个脉冲在薄层中的占据数作全时间累计，也会得到相同的停留时间权重。实际测温还须按装置条件计入角接受范围、孔道传输和仪器响应。
\end{solution}

这些计算中的权重，都能追溯到图中的扫掠体积或变量变换。先按面元和相对速度写出事件数，再对实际允许的速度区域积分，便能确定相应的平均量。

\section{自由飞行、方向随机化与扩散尺度}
\subsection{一次碰撞的抽样怎样变成长期输运}
先用恒速随机飞行模型求扩散系数。粒子速率为$v$，碰撞时刻构成平均间隔$\tau$的泊松过程，每次碰撞后方向独立且各向同性，平均自由程为$\ell=v\tau$。这一模型保留了自由飞行和方向随机化，但省去了Maxwell气体碰撞中速率的变化。

初始方向取各向同性。经过时间$t$没有碰撞的概率为$\ee^{-t/\tau}$；只要发生过至少一次方向随机化，当前速度与初始速度的平均点积便为零。因而速度相关函数为
\[
 C_v(t)=\langle\vct v(0)\cdot\vct v(t)\rangle=v^2\ee^{-t/\tau}.
\]
位移为$\vct r(t)=\int_0^t\vct v(s)\dd s$，将平方展开并取平均，得到$\langle r^2(t)\rangle=\int_0^t\int_0^t C_v(|s_1-s_2|)\dd s_1\dd s_2$。积分正方形中两个三角形的贡献相等；在$s_1\ge s_2$的三角形内，固定时间差$s=s_1-s_2$后，$s_2$从$0$到$t-s$，因此
\begin{align*}
 \langle r^2(t)\rangle
 &=2\int_0^t(t-s)C_v(s)\dd s\\
 &=2v^2\int_0^t(t-s)\ee^{-s/\tau}\dd s\\
 &=2v^2\left[\tau(s+\tau-t)\ee^{-s/\tau}\right]_0^t\\
 &=2v^2[\tau t-\tau^2(1-\ee^{-t/\tau})].
\end{align*}
短时间$t\ll\tau$时展开指数，首个非零项为$v^2t^2$，对应弹道运动；长时间$t\gg\tau$时为$2v^2\tau t=6Dt$，所以
\[
 D=\frac{v^2\tau}{3}=\frac{v\ell}{3}.
\]
系数$1/3$来自三维各向同性，不应移到一维仅左右随机的模型中照用。若速率变化、散射保留方向记忆或不同速率具有不同碰撞时间，必须重新求相关函数，不能把任何平均速率与平均自由程任意相乘。
\ExpansionDiagram{fig:dep-th1-walk}{短时间近似直线飞行，长期由许多随机方向段累积}{
\draw[->,thick](0,0)--(1,.4)--(1.6,1.3)--(2.6,1)--(3,.2)--(4,.7)--(4.4,1.8)--(5.2,1.4);\fill(0,0)circle(1.5pt)node[below]{$O$};\draw[dashed,->](0,0)--(5.2,1.4)node[midway,below]{$\vct r$};\node[above]at(2.5,2.2){每次碰撞后方向重新抽样};}
图\ref{fig:dep-th1-walk}中的实际路程始终为$vt$，端点位移却因方向抵消而只有典型大小$\sqrt{6Dt}$。路程与位移不能混用来估计扩散所需时间。

\subsection{从粗粒化位移到扩散方程}
在比$\tau$长、但比宏观扩散时间短的区间$\Delta t$中，单个方向满足$\langle\Delta x\rangle=0$、$\langle\Delta x^2\rangle=2D\Delta t$。设位移$\xi$的归一化分布为$P(\xi)$，且介质均匀、没有平均漂移，则
\[
 n(x,t+\Delta t)=\int_{-\infty}^{\infty}P(\xi)n(x-\xi,t)\dd\xi.
\]
在密度变化尺度上展开两边，保留时间一阶项与位移二阶项，得
\begin{align*}
 n+\Delta t\,\partial_tn
 &=n-\langle\xi\rangle\partial_xn
   +\frac12\langle\xi^2\rangle\partial_x^2n\\
 &=n+D\Delta t\,\partial_x^2n.
\end{align*}
约去$n$并除以$\Delta t$，再与粒子数守恒式$\partial_tn=-\partial_xj_N$比较；取均匀密度时通量为零，得到
\[
 \partial_t n=D\partial_x^2n,\qquad j_N=-D\partial_xn.
\]
其中$j_N$是粒子数通量。该连续近似要求密度变化尺度远大于$\ell$；壁面附近一个自由程内常有非各向同性层，边界条件应与吸收、反射或穿透的具体模型相符。

\begin{example}[两吸收端之间的平均逃逸时间]
粒子在$0<x<L$中扩散，两端到达即吸收。扩散系数为$D$，初始位置为$x_0$且远离壁面的微观边界层。求先从右端逃逸的概率，以及到达任一端的平均时间。
\ExpansionDiagram{fig:dep-th1-exit-q}{两吸收端的首次到达问题，初始点为$x_0$}{
\draw[->](0,0)--(6,0)node[right]{$x$};\draw[very thick](.5,-.6)--(.5,.8)node[above]{吸收};\draw[very thick](5.5,-.6)--(5.5,.8)node[above]{吸收};\fill(2.3,0)circle(2pt)node[below]{$x_0$};\node[below]at(.5,-.7){0};\node[below]at(5.5,-.7){$L$};\draw[<->](1.5,.4)--(3.1,.4);}
\end{example}
\begin{solution}
用$p(x)$表示右端先到达的概率。图\ref{fig:dep-th1-exit-q}中从内点经过短粗粒化时间后，后续概率等于$p(x+\Delta x)$的平均。于是
\[
 p(x)=p(x)+p'(x)\langle\Delta x\rangle
       +\frac12p''(x)\langle\Delta x^2\rangle
     =p(x)+D\Delta t\,p''(x),
\]
故$p''=0$。逐次积分得$p'=A$、$p=Ax+B$。从左端出发已经被左端吸收，所以$p(0)=0$，给出$B=0$；从右端出发已经到达右端，所以$p(L)=1$，给出$A=1/L$。因此$p(x_0)=x_0/L$。

用$T(x)$表示剩余平均到达时间。时间已经经过$\Delta t$，所以$T(x)=\Delta t+\langle T(x+\Delta x)\rangle$。展开右边得$T=\Delta t+T+D\Delta t\,T''$，约去$T$并除以$\Delta t$，得到$DT''=-1$。逐次积分为
\[
 T'(x)=-\frac{x}{D}+C_1,\qquad
 T(x)=-\frac{x^2}{2D}+C_1x+C_2.
\]
两端一旦到达就不再等待，故$T(0)=T(L)=0$。前一条件给出$C_2=0$，后一条件给出$C_1=L/(2D)$，因此
\[
 T(x)=\frac{x(L-x)}{2D}.
\]
\ExpansionDiagram{fig:dep-th1-exit-a}{右端到达概率为直线，平均等待时间为抛物线}{
\draw[->](0,0)--(5,0)node[right]{$x/L$};\draw[->](0,0)--(0,2.7)node[above]{归一化量};\draw[dotted](0,2)--(4,2);\draw(-.06,2)--(.06,2)node[left,xshift=-2pt]{$1$};\draw[dashed](0,0)--(4,2);\draw[thick,domain=0:4,samples=70]plot(\x,{2*\x*(4-\x)/4});\node[above]at(2,2){$T/T_{\max}$};\node[right]at(4,2){$p$};\node[below]at(4,0){1};}
图\ref{fig:dep-th1-exit-a}中最长平均等待发生在中央，$T_{\max}=L^2/(8D)$，而非$L/v$。在本节的方向随机化模型下，这等于$3L^2/(8v\ell)$；$L\gg\ell$时远大于一次直线穿越时间。若初态距壁面仅约一个自由程且方向已经给定，最初一段弹道运动不可由此扩散边界模型准确描述。
\end{solution}

\begin{example}[封闭容器中的浓度均匀化]
一长直容器截面积$A$、长度$L$，两端完全反射，初始密度$n(x,0)=n_0[1+\epsilon\cos(\pi x/L)]$，$|\epsilon|<1$。取恒定$D$，求密度、粒子流及不均匀振幅降为原来的$1/10$所需时间。
\ExpansionDiagram{fig:dep-th1-mix-q}{反射端无粒子流，初态浓度由左向右平滑变化}{
\draw(0,0)rectangle(5,1.6);\draw[very thick](0,0)--(0,1.6);\draw[very thick](5,0)--(5,1.6);\foreach\x/\n in{.5/5,1.5/4,2.5/3,3.5/2,4.5/1}{\foreach\j in{1,...,\n}{\fill(\x,{.2+1.2*\j/(\n+1)})circle(1.2pt);}}\node[below]at(2.5,0){$\partial_xn=0$于两端};}
\end{example}
\begin{solution}
图\ref{fig:dep-th1-mix-q}的无流边界是$\partial_xn=0$。设$n=n_0[1+a(t)\cos(\pi x/L)]$，初值给出$a(0)=\epsilon$。对空间变量逐次求导，得
\[
 \partial_xn=-\frac{\pi n_0a}{L}\sin\frac{\pi x}{L},\qquad
 \partial_x^2n=-\frac{\pi^2n_0a}{L^2}\cos\frac{\pi x}{L}.
\]
第一式在$x=0,L$都为零，满足反射边界。代入扩散方程并比较余弦项的系数，得到$\dot a=-D\pi^2a/L^2$。若$\epsilon\ne0$，分离变量并代入初值，
\[
 \int_\epsilon^{a(t)}\frac{\dd a}{a}
 =-\frac{D\pi^2}{L^2}\int_0^t\dd s,
 \qquad
 \ln\frac{|a(t)|}{|\epsilon|}=-\frac{D\pi^2t}{L^2}.
\]
振幅保持初始符号，故$a(t)=\epsilon\ee^{-\pi^2Dt/L^2}$；$\epsilon=0$时$a(t)=0$也满足方程与初值。再由$j_N=-D\partial_xn$得到
\[
 n(x,t)=n_0\left[1+\epsilon\ee^{-\pi^2Dt/L^2}\cos\frac{\pi x}L\right],
 \quad j_N=\frac{\pi Dn_0\epsilon}{L}\ee^{-\pi^2Dt/L^2}\sin\frac{\pi x}L.
\]
\ExpansionDiagram{fig:dep-th1-mix-a}{常数模保持总粒子数，余弦不均匀模随时间衰减}{
\draw[->](0,0)--(5.5,0)node[right]{$x/L$};\draw[thick,domain=0:5,samples=80]plot(\x,{1.3+.7*cos(180*\x/5)});\draw[dashed,domain=0:5,samples=80]plot(\x,{1.3+.25*cos(180*\x/5)});\draw[dotted](0,1.3)--(5,1.3);\node[right]at(5,1.3){$n_0$};}
图\ref{fig:dep-th1-mix-a}中常数模保持不变。将密度对容积积分，
\[
 N(t)=n_0A\left[x+\frac{\epsilon L}{\pi}
 \ee^{-\pi^2Dt/L^2}\sin\frac{\pi x}{L}\right]_0^L=n_0AL,
\]
可见余弦项不改变总粒子数。当$\epsilon\ne0$时，振幅降为原来的$1/10$要求$\ee^{-\pi^2Dt/L^2}=1/10$，取对数得$-\pi^2Dt/L^2=-\ln10$，所以$t=L^2\ln10/(\pi^2D)$。若$\epsilon=0$，初态已经均匀。若只使容器长度变为两倍而材料和温度不变，均匀化时间增为四倍。反射边界保留常数模，吸收边界则不保留粒子总数；两类问题不能共用同一组末态条件。
\end{solution}

\section{习题与解答}
\begin{exercise}[两种温度下的曲线缩放]
同种气体由 $T$ 变为 $4T$，粒子数不变。求速率分布峰值位置、峰高和区间概率 $P(a<v<b)$ 的变换关系。
\Diagram{fig:scaleq}{不同温度下分布变宽而面积保持为一}{
 \draw[axis](0,0)--(5.6,0)node[right]{$v$};\draw[axis](0,0)--(0,2.4)node[above]{$g$};
 \draw[curve]plot[domain=0:5,samples=80](\x,{4/sqrt(pi)*\x*\x*exp(-\x*\x)});
 \draw[curve,dashed]plot[domain=0:5,samples=80](\x,{1/(2*sqrt(pi))*\x*\x*exp(-\x*\x/4)});\node at(4,1.9){$T\to4T$};
}
\end{exercise}
\begin{solution}
设原密度为 $g_T(v)$，直接代入式\eqref{eq:speeddensity}，有 $g_{4T}(v)=\tfrac12g_T(v/2)$。图\ref{fig:scaleq}中横向伸长两倍、纵向降低一半，面积因而不变。峰值位置 $v_p\to2v_p$，峰高减半。
在变换后的密度中，温度的 $T^{-3/2}$ 因子给出 $1/8$，指数中的 $v^2/T$ 则通过 $v/2$ 保持原形；同时原密度内的速度平方又产生因子4，合起来才是前面的 $1/2$。区间积分另有 $\dd v=2\dd w$，正好与纵向缩放抵消。图中的两段区间代表同一组无量纲端点 $v/\sqrt T$；固定实验室速率端点时，端点相对于热运动尺度的位置已变化，因此其概率不会普遍保持不变。

\Diagram{fig:scalea}{概率区间随无量纲速度一起变换}{
 \draw[axis](0,1.2)--(5,1.2)node[right]{$T$};\draw[axis](0,0)--(5,0)node[right]{$4T$};
 \draw[line width=2pt](1,1.2)--(2,1.2);\draw[line width=2pt](2,0)--(4,0);
 \node[above]at(1,1.2){$a$};\node[above]at(2,1.2){$b$};\node[below]at(2,0){$2a$};\node[below]at(4,0){$2b$};
}
令 $w=v/2$，则 $P_{4T}(a<v<b)=\int_{a/2}^{b/2}g_T(w)\dd w$。图\ref{fig:scalea}表示的是 $P_{4T}(2a<v<2b)=P_T(a<v<b)$；若区间端点不跟随温度缩放，概率通常会改变。
\end{solution}

\begin{exercise}[阈值以上的速率]
求平衡气体中 $v>v_0$ 的比例，并令 $v_0=v_p$ 写成误差函数形式。
\Diagram{fig:speedtailq}{速率阈值右侧的概率}{
 \draw[axis](0,0)--(5,0)node[right]{$v$};\draw[axis](0,0)--(0,2)node[above]{$g$};
 \draw[curve]plot[domain=0:4.5,samples=70](\x,{3*\x*\x*exp(-\x*\x)});\draw[guide](1.5,0)--(1.5,1.7);\node[below]at(1.5,0){$v_0$};
}
\end{exercise}
\begin{solution}
令$t=\sqrt\alpha v$，则$\dd v=\dd t/\sqrt\alpha$，下限$v_0$变为$z=\sqrt\alpha v_0$、上限仍为无穷大。图\ref{fig:speedtailq}右侧概率因而为$(4/\sqrt\pi)\int_z^\infty t^2\ee^{-t^2}\dd t$。分部积分取$t$与$t\ee^{-t^2}\dd t$，得到
\[
 P(v>v_0)=\erfc z+\frac{2z}{\sqrt\pi}\ee^{-z^2}.
\]
分部积分的具体边界项为 $[-t\ee^{-t^2}/2]_z^\infty=z\ee^{-z^2}/2$，剩余积分为原高斯尾的一半。乘上速率密度归一化常数后，便分别形成图中的两项。阈值取最概然速率时 $\alpha v_p^2=1$，所以应把无量纲下限取为1。最概然值只给出密度峰的位置，不把总体分成各占一半的两部分；本题的积分一般也不等于 $1/2$。

\Diagram{fig:speedtaila}{速率尾积分由边界项和高斯尾组成}{
 \node[draw]at(0,0){边界项 $2z\ee^{-z^2}/\sqrt\pi$};\node at(3,0){$+$};\node[draw]at(5,0){$\erfc z$};
}
当 $v_0=v_p$ 时 $z=1$，结果为 $\erfc1+2/(\ee\sqrt\pi)$。图\ref{fig:speedtaila}中的边界项不能省去；速率密度有 $v^2$ 因子，与一维速度分量的尾概率不同。$z=0$ 给出1，$z\to\infty$ 给出0。
\end{solution}

\begin{exercise}[不同温度的相对速度]
两个独立粒子质量为 $m_1,m_2$，各自服从温度 $T_1,T_2$ 的无漂移 Maxwell 分布。求相对速率的平均值和均方值。能否仍把 $T_1$ 或 $T_2$ 直接代入约化质量公式？
\Diagram{fig:twotq}{两个独立速度分布的差}{
 \node[draw,align=center]at(0,0){$m_1,T_1$\\$\vct v_1$};\node[draw,align=center]at(3.4,0){$m_2,T_2$\\$\vct v_2$};\draw[->](.8,-.9)--(2.6,-.9)node[midway,below]{$\vct u=\vct v_1-\vct v_2$};
}
\end{exercise}
\begin{solution}
每个方向的速度分量均为零均值高斯，方差分别为 $\kb T_1/m_1$、$\kb T_2/m_2$。图\ref{fig:twotq}中的两者独立，差的方差相加，记 $s^2=\kb T_1/m_1+\kb T_2/m_2$，有 $u_x,u_y,u_z$ 独立且方差均为 $s^2$。
方差相加之前先展开 $u_x^2=v_{1x}^2+v_{2x}^2-2v_{1x}v_{2x}$。独立且均值均为零时，交叉项平均为 $\langle v_{1x}\rangle\langle v_{2x}\rangle=0$。三个方向具有相同方差，所以均方相对速率是三个分量二阶矩之和。若两组分还存在不同的宏观漂移速度，交叉项的处理要先减去各自均值，相对分布的中心也会偏移；本题给出的零漂移条件正是中心位于原点的依据。

\Diagram{fig:twota}{相对速度各分量的方差相加}{
 \draw[axis](-2.7,0)--(2.7,0)node[right]{$u_x$};\draw[curve]plot[domain=-2.5:2.5,samples=60](\x,{1.7*exp(-\x*\x/2)});\draw[<->](-1,-.4)--(1,-.4)node[midway,below]{$2s$};
}
故三维相对速率仍为 Maxwell 形状，$g_u(u)=\sqrt{2/\pi}\,u^2s^{-3}\ee^{-u^2/(2s^2)}$，于是 $\bar u=\sqrt{8/\pi}\,s$、$\langle u^2\rangle=3s^2$。图\ref{fig:twota}的宽度由两个温度共同决定，计算时应保留 $\kb T_1/m_1+\kb T_2/m_2$ 的完整组合；当 $T_1=T_2=T$ 时，才简化为 $s^2=\kb T/\mu$。
\end{solution}

\begin{exercise}[同种气体的碰撞计数]
体积 $V$ 内有 $N$ 个相同硬球，温度 $T$、有效截面 $\sigma$。求总体每秒碰撞次数，并比较用“每粒子的频率乘粒子数”得到的结果。
\Diagram{fig:paircountq}{同一次碰撞涉及两个粒子}{
 \node[circle,draw,minimum size=9mm](a)at(0,0){$i$};\node[circle,draw,minimum size=9mm](b)at(2,0){$j$};\draw[->](a)--(.7,0);\draw[->](b)--(1.3,0);\node at(1,-.8){一对 $\{i,j\}$};
}
\end{exercise}
\begin{solution}
单个粒子的平均碰撞频率为 $z=n\sigma\sqrt2\bar v$。图\ref{fig:paircountq}中一次碰撞同时出现在 $i$ 和 $j$ 的计数中，因此总体碰撞次数率为
\[
\mathcal R=\frac12Nz\simeq\frac{N^2\sigma}{\sqrt2V}\sqrt{\frac{8\kb T}{\pi m}}.
\]
单粒子频率$z$按参与者计数，$Nz$便把每次碰撞计入两次。图中合并这两条有序记录时，截面和相对运动的平均不变；有限粒子数还须排除与自身配对的项。

\Diagram{fig:paircounta}{有序计数转换为无序粒子对}{
 \node[draw]at(0,0){$i\to j$};\node[draw]at(2,0){$j\to i$};\draw[->](0,-.4)--(1,-1);\draw[->](2,-.4)--(1,-1);\node[below]at(1,-1){同一事件，除以二};
}
图\ref{fig:paircounta}中两条有序记录对应同一次碰撞，因此总计数除以二。有限$N$时还要排除粒子与自身配对，将靶粒子密度改为$(N-1)/V$，同一稀薄无关联模型给
\[\mathcal R=\frac{N(N-1)\sigma}{\sqrt2V}
\sqrt{\frac{8\kb T}{\pi m}}.\]
宏观气体$N\gg1$时，它与前式的相对差为$1/N$。
\end{solution}

\begin{exercise}[小孔两侧的平衡条件]
两容器以很小的薄孔相通，孔的尺度远小于两侧平均自由程。两侧为相同气体，温度 $T_1,T_2$ 由外界维持。求零净粒子通量时 $p_1/p_2$，并说明零粒子通量是否意味着无能量传递。
\Diagram{fig:thermaltransq}{不同温度容器通过分子流小孔相连}{
 \draw(0,0)rectangle(5,1.6);\draw(2.5,0)--(2.5,.65);\draw(2.5,.95)--(2.5,1.6);\node at(1.2,.8){$p_1,T_1$};\node at(3.8,.8){$p_2,T_2$};
}
\end{exercise}
\begin{solution}
图\ref{fig:thermaltransq}两侧各自向孔贡献 $j_i=n_i\bar v_i/4$。零净粒子流要求 $n_1\sqrt{T_1}=n_2\sqrt{T_2}$。用 $p_i=n_i\kb T_i$ 消去数密度，得
\[
 \frac{p_1}{\sqrt{T_1}}=\frac{p_2}{\sqrt{T_2}},\qquad
 \frac{p_1}{p_2}=\sqrt{\frac{T_1}{T_2}}.
\]
热侧在同压下密度较小，而平均速率较大，结合后单向通量正比于$p_i/\sqrt{T_i}$。孔口处是分子流，两侧没有通过密集碰撞建立共同压强；恒温装置仍需补入或取走穿孔传递的净能量。

\Diagram{fig:thermaltransa}{粒子流相抵时能量流仍可不相抵}{
 \draw[axis](0,.5)--(4,.5)node[midway,above]{$j\times2\kb T_1$};\draw[axis](4,-.5)--(0,-.5)node[midway,below]{$j\times2\kb T_2$};
}
对单原子气体，每个逸出粒子带走 $2\kb T_i$，故图\ref{fig:thermaltransa}给出净能流 $2\kb j(T_1-T_2)$。只有温度也相同才无净热传递。这里由外界恒温，不要求两侧构成孤立平衡系统。
\end{solution}

\begin{exercise}[束流的角分布与圆锥筛选]
由薄孔逸出的分子在远场可视为余弦束。收集器只接受与孔法线夹角 $0\le\theta\le\theta_0$ 的分子。求收集到的比例和归一化角密度；若筛选只与方向有关，收集束的平均能量是否改变？
\Diagram{fig:coneq}{小孔法线周围的收集圆锥剖面}{
 \draw[line width=1.2pt](0,-1.2)--(0,-.15);\draw[line width=1.2pt](0,.15)--(0,1.2);
 \draw[axis](0,0)--(4,0)node[right]{法线};\draw[curve](0,0)--(3,1.2);\draw[curve](0,0)--(3,-1.2);
 \draw[->](.9,0)arc(0:21.8:.9);\node at(1.4,.4){$\theta_0$};\node[left]at(0,.3){孔};
}
\end{exercise}
\begin{solution}
图\ref{fig:coneq}是以法线为轴的三维圆锥的剖面。对整个方位角积分后，原束的极角密度为 $q(\theta)=2\sin\theta\cos\theta$，所以
\[
 P_c=\int_0^{\theta_0}q(\theta)\dd\theta
=[\sin^2\theta]_0^{\theta_0}=\sin^2\theta_0,
 \qquad q_c(\theta)=\frac{2\sin\theta\cos\theta}{\sin^2\theta_0}.
\]
收集比例的分母对应整个向外半球。方向密度 $\cos\theta/\pi$ 乘立体角环带 $2\pi\sin\theta\dd\theta$，才成为所用的极角密度。它的积分是 $\sin^2\theta$；不能仅以圆锥立体角占半球的比例 $1-\cos\theta_0$ 代替，因为偏斜方向的通量还受到余弦投影抑制。小角度时收集比例约为 $\theta_0^2$，缩小接受角主要减少粒子数，不额外筛选快慢。

\Diagram{fig:conea}{速度权重与方向权重可分离}{
 \node[draw]at(0,0){$g_e(v)\dd v$};\node at(2,0){$\times$};\node[draw]at(4,0){$q_c(\theta)\dd\theta$};\node at(2,-.8){方向筛选不改变速率密度};
}
原联合密度中速率与方向相乘，图\ref{fig:conea}的筛选只改变角因子的归一化。因此收集束仍有 $\langle E_k\rangle=2\kb T$。若孔道很长，透射概率可能依赖速率、壁面吸附和再发射，便不能仅凭这一角度模型推断。
\end{solution}

\begin{exercise}[重分子与轻分子的泄流分离]
容器内有两种理想气体，质量分别为 $m_a,m_b$，粒子数为 $N_a,N_b$，共同温度 $T$。在温度和容器体积保持不变的小孔泄流过程中，求逸出流的组分比和余气组分随时间的变化。
\Diagram{fig:isotopeq}{同一小孔对不同质量粒子的通量不同}{
 \draw(0,0)rectangle(3,1.6);\fill[white](2.95,.65)rectangle(3.05,.95);\node at(1.5,1.05){$N_a,m_a$};\node at(1.5,.5){$N_b,m_b$};\draw[->](3,.8)--(4.5,.8);
}
\end{exercise}
\begin{solution}
由图\ref{fig:isotopeq}，两种气体共享孔面积 $A_h$ 和体积 $V$。各自数流率为
\[
 J_i=\frac{A_hN_i}{V}\sqrt{\frac{\kb T}{2\pi m_i}}=\kappa_iN_i,
 \qquad \kappa_i=\frac{A_h}{V}\sqrt{\frac{\kb T}{2\pi m_i}}.
\]
所以瞬时逸出流的组分比为
\[
 \frac{J_a}{J_b}=\frac{N_a}{N_b}\sqrt{\frac{m_b}{m_a}}.
\]
微分方程分离为$\dd N_i/N_i=-\kappa_i\dd t$。从$N_i(0)=N_{i0}>0$积分，
\[\int_{N_{i0}}^{N_i(t)}\frac{\dd N}{N}
=-\kappa_i\int_0^t\dd s,
\qquad[\ln N]_{N_{i0}}^{N_i(t)}=-\kappa_i[s]_0^t.\]
所以$\ln[N_i(t)/N_{i0}]=-\kappa_it$，取指数给出下述各组分数目。若$m_a<m_b$，则$\kappa_a>\kappa_b$，轻组分在瞬时逸出流中富集，也在容器内消耗得更快。逸出流比为$J_a/J_b$，余气比为$N_a/N_b$，两者相差一个质量平方根因子。

\Diagram{fig:isotopea}{恒温泄流时两组分分别指数减少}{
 \draw[axis](0,0)--(5,0)node[right]{$t$};\draw[axis](0,0)--(0,2.3)node[above]{$N_i/N_{i0}$};
 \draw[curve]plot[domain=0:4.6](\x,{2*exp(-\x)});\draw[curve,dashed]plot[domain=0:4.6](\x,{2*exp(-\x/2)});\draw(-.06,2)--(.06,2);\node[left]at(0,2){$1$};\node at(3,1.6){轻组分减少较快};
}
两组分数目为$N_i=N_{i0}\ee^{-\kappa_it}$，余气比值为$(N_{a0}/N_{b0})\ee^{-(\kappa_a-\kappa_b)t}$。图\ref{fig:isotopea}实线示意较轻组分。外界持续供热使温度不变，$\kappa_i$才保持常数；若改为绝热容器，泄流系数中的$\sqrt{T(t)}$须与气体能量方程联立求解。
\end{solution}

\begin{exercise}[自由管模型的质量极限]
在例\ref{ex:piston}中，保持初始气体状态不变，分别取 $m_2\to\infty$ 与 $m_1=m_2=M_p$。求活塞实验室速度和管体速度，比较相对速度与可用于喷出活塞的机械能。
\Diagram{fig:recoillimitq}{同一膨胀能量分配给两个可运动物体}{
 \node[draw,minimum width=1.5cm]at(0,0){管 $m_2$};\node[draw]at(3,0){活塞 $m_1$};\draw[axis](-.8,-.6)--(-1.5,-.6);\draw[axis](3,-.6)--(4,-.6);\node at(1.5,.8){总机械能 $E$ 固定};
}
\end{exercise}
\begin{solution}
令 $E=\nu C_{V,m}T_0(1-2^{1-\gamma})$。图\ref{fig:recoillimitq}要求总动量为零，故 $u=\sqrt{2E/\mu_p}$，$v_1=m_2u/(m_1+m_2)$，$v_2=-m_1u/(m_1+m_2)$。
把相对速度代回两物体速度后，动能分配可不依赖具体气体参数写成
\[
 E_1=\frac{m_2}{m_1+m_2}E,\qquad
 E_2=\frac{m_1}{m_1+m_2}E.
\]
两项相加恢复总膨胀能 $E$。较轻的物体获得较大的速度和较多的动能，但两者动量大小始终相等、方向相反。固定管极限和等质量极限都可由这两个分配系数直接读出，避免把相对运动的约化质量误认成活塞在实验室系中的实际质量。

\Diagram{fig:recoillimita}{固定管与等质量反冲的速度分配}{
 \node at(0,1){$m_2\to\infty$};\draw[axis](0,0)--(1.6,0)node[right]{$v_1=\sqrt{2E/m_1}$};
 \node at(0,-1){$m_1=m_2$};\draw[axis](0,-2)--(1.4,-2)node[right]{$v_1=\sqrt{E/M_p}$};\draw[axis](0,-2)--(-1.4,-2)node[left]{$v_2=-v_1$};
}
固定管极限$\mu_p\to m_1$，管速度趋零，全部机械能进入活塞：
\[u\longrightarrow\sqrt{\frac{2E}{m_1}},\qquad
v_1\longrightarrow\sqrt{\frac{2E}{m_1}},\qquad v_2\longrightarrow0.\]
等质量时$\mu_p=M_p/2$，两物体各得$E/2$，且
\[u=2\sqrt{\frac E{M_p}},\qquad
v_1=\frac u2=\sqrt{\frac E{M_p}},\qquad
v_2=-\frac u2=-\sqrt{\frac E{M_p}}.\]
图\ref{fig:recoillimita}中相对速度增大，而活塞分得的能量减少；管体反冲既贡献了相对速度，也带走一半机械能。
\end{solution}

\chapter{泄流、辐射与热力学关系}

\section{小孔泄流中的能量选择}
\label{th:effusion}
第一讲求出的逸出速率密度为 $g_e(v)=2\alpha^2v^3\ee^{-\alpha v^2}$，$\alpha=m/(2\kb T)$。多出的速率因子来自单位时间的穿孔机会。利用 $I_5/I_3$，有
\[
 \langle v^2\rangle_e=\frac{I_5}{I_3}=\frac2\alpha,
 \qquad \langle E_{\rm tr}\rangle_e=2\kb T.
\]
容器中的平均平动能为 $3\kb T/2$，逸出粒子平均多携带 $\kb T/2$。如果分子还有充分激发的转动自由度，而孔不按内部状态筛选，那么每个内部二次项携带的平均能量仍为 $\kb T/2$。设总有效自由度为 $f$，平动占3个，则
\begin{equation}\label{eq:escapeenergy}
 c_v=\frac f2\kb,\qquad
 \varepsilon_e=2\kb T+\frac{f-3}{2}\kb T
 =\frac{f+1}{2}\kb T.
\end{equation}
这里假定内部自由度与平动保持热平衡，且孔的传输不依赖内部能级。常数 $f$ 只计入所讨论温区内已充分激发的自由度；振动逐渐激发的温区应使用随温度变化的热容。

泄流冷却还需要三个尺度条件：孔径远小于平均自由程，孔外压强足够低，容器内重新平衡的时间远短于粒子数显著变化的时间。第一项保证分子流穿孔，第三项使余气仍可用一个温度描述。它们可以同时满足：小孔很小，容器内部仍有足够碰撞。

\subsection{系统边界与粒子携带的能量}
\label{th:openboundary}
对开口容器应用第一定律时，取固定区域内的瞬时存量作为研究对象。设容器中有 $N$ 个粒子，其内部状态可由 $N,T$ 描述，内能为 $U(N,T)$，短时间内逸出粒子数为 $\dd N_{\rm out}>0$，每粒子平均携带能量为 $\varepsilon_{\rm out}$，则容器粒子数变化为 $\dd N=-\dd N_{\rm out}$。若外界还供热 $\delta Q$，边界无宏观机械位移，
\[
 \dd U=\delta Q-\varepsilon_{\rm out}\dd N_{\rm out}
 =\delta Q+\varepsilon_{\rm out}\dd N.
\]
因 $\dd N<0$，后一项为负。对理想独立粒子，$U=Nu(T)$，完整微分为
\[
 \dd U=Nc_v(T)\dd T+u(T)\dd N.
\]
其中 $u\,\dd N$ 表示存量减少本身造成的内能变化，$Nc_v\dd T$ 表示余气温度改变。与边界输运式联立，
\[
 Nc_v(T)\dd T
 =\delta Q+\bigl[\varepsilon_{\rm out}-u(T)\bigr]\dd N.
\]
余气升温或冷却由逸出能量相对于体内平均能量的高低决定。能量零点若平移 $E_\ast$，$u$ 和 $\varepsilon_{\rm out}$ 同时增加 $E_\ast$，差值不变，温度方程也不变。

\Diagram{fig:controlboundary}{固定边界内的存量与穿越边界的能量}{
 \draw[curve](0,0)--(0,2)--(3.4,2)--(3.4,1.15);
 \draw[curve](3.4,.85)--(3.4,0)--(0,0);
 \node at(1.7,1){$U(N,T)$};
 \draw[axis](-1.6,1)--(0,1)node[midway,above]{$\delta Q$};
 \draw[axis](3.4,1)--(5.6,1)node[right]{$\varepsilon_{\rm out}\dd N_{\rm out}$};
 \node at(2,-.55){先选边界，再写储能与输运};
}
图\ref{fig:controlboundary}中，边界能流按穿越的粒子统计。小孔分子流按穿孔事件取样，平动部分为 $2\kb T$，内部自由度若与穿孔概率无关，则仍取体内平均值。这一微观输运模型与有整体流速的连续稳流模型所用能量形式不同。

对一段连续流管，把入口前的流体推入控制体，外界除送入该流体本身的内能外，还需克服入口压强做功。质量 $\dd m$ 占据体积 $\dd m/\rho$，推入功为 $p\,\dd m/\rho$。因此入口携带的内能加流动功合成比焓 $h=u_m+p/\rho$。在稳定流动中，
\[
 \dot Q-\dot W_{\rm shaft}
 =\dot m\left[
 h_2-h_1+\frac{V_2^2-V_1^2}{2}
 +g(z_2-z_1)\right].
\]
这里 $u_m,h$ 均为单位质量的量，$V_i$ 为宏观流速；$\dot W_{\rm shaft}$ 是对外输出的轴功。入口、出口流动功已经计入焓，不能再作为额外项重复加入。分子流小孔则从单粒子动能与内部能量直接计算输运率；若已这样计入穿面能量，再机械地附加一个 $pV$ 项，便改变了模型。

\subsection{流出冷却与稳流节流的不同约束}
刚性容器缓慢泄流时，容器存量随时间改变，能量方程中有 $\dd U/\dd t$。绝热节流则通常选取阀门附近的固定控制体，在稳态下控制体储能不变；无轴功且宏观动能、位能变化可忽略时，得到 $h_1=h_2$。自由膨胀又是封闭气体占据更大可达体积，外边界刚性绝热，给出整个气体的 $U_1=U_2$。三者分别固定存量能量的变化式、入口出口焓差和整体内能差。

这些约束须在状态方程之前确定。理想气体自由膨胀或理想稳流节流的初末温度相同，是因为它的内能或焓只依赖温度；绝热泄流容器却因穿孔样本偏向高能粒子而冷却。对实际气体，等内能与等焓曲线一般不同，两种温变也须分别计算。

\subsection{持续充气中的流动功}
\begin{example}[绝热刚性空罐的充气温度]
\label{ex:tankfilling}
一个初始接近真空、体积为 $V$ 的刚性容器，通过阀门与温度 $T_s$、压强 $p_s$ 恒定的大型气源相连。气体为常热容理想气体，摩尔热容为 $c_{V,m}$、$c_{p,m}$，$\gamma=c_{p,m}/c_{V,m}$。容器与外界绝热，忽略罐壁热容和最终宏观动能。充气后关阀，待罐内达到均匀平衡。求终温，并说明气源做的功进入了哪里。
\Diagram{fig:tankfillingq}{大型气源经阀门向绝热刚性罐充气}{
 \node[draw,minimum width=2.4cm,minimum height=1.6cm](s)at(0,1){$p_s,T_s$};
 \draw[curve](3,0)rectangle(6.4,2);
 \node at(4.7,1){$V,\ \nu(t),T(t)$};
 \draw[axis](1.2,1)--(3,1);
 \draw(2,.75)--(2.4,1.25);\draw(2,1.25)--(2.4,.75);
 \node[below]at(2.2,.5){阀};
 \node at(4.7,-.5){外壁不动，$Q=0$};
}
\end{example}
\begin{solution}
在图\ref{fig:tankfillingq}中，把连接段和容器共同纳入控制体，在气源侧选入口。每进入 $\dd\nu$ 摩尔气体，气源送入摩尔内能 $c_{V,m}T_s$，并通过入口压强做流动功 $RT_s$，合计送入摩尔焓 $c_{p,m}T_s$。若所选入口处流速不可忽略，应计其动能；本题可在大型气源内部流速可忽略处取截面，直接使用源的摩尔焓。
\[
 \dd U=c_{p,m}T_s\,\dd\nu.
\]
从初始空罐积分，$U_f-U_i=c_{p,m}T_s\int_0^{\nu_f}\dd\nu=c_{p,m}T_s[\nu]_0^{\nu_f}=\nu_fc_{p,m}T_s$，其中$U_i=0$。气体最终静止，内能为 $\nu_fc_{V,m}T_f$，于是
\[
 T_f=\gamma T_s.
\]
\Diagram{fig:tankfillinga}{进入罐内的焓分解为原有内能与入口流动功}{
 \node[draw](a)at(0,0){$c_{p,m}T_s$};
 \node[draw](b)at(4,1){$c_{V,m}T_s$};
 \node[draw](c)at(4,-1){$RT_s$};
 \draw[axis](a)--(b);\draw[axis](a)--(c);
 \node[right]at(6,0){均进入最终内能};
}
图\ref{fig:tankfillinga}中，流动功经过射流、黏滞耗散和内部混合，最终转成罐内热运动能。外壁不动只排除了外壁体积功，入口处的压强仍在输送能量。若充气持续到罐压等于 $p_s$，则 $\nu_f=p_sV/(R\gamma T_s)$。

若罐壁初温为 $T_{w0}$、热容为 $C_w$，且最终与气体同温，则同一控制体给出
\[
 (\nu_fc_{V,m}+C_w)T_f
 =\nu_fc_{p,m}T_s+C_wT_{w0}.
\]
这一步只改变储能项；气源输入的摩尔焓不变。真实充气过程不可逆，均匀终态可用状态方程描述，充气途中则无须假定罐内始终处于准静态。
\end{solution}


\begin{example}[刚性绝热容器中的余气冷却]\label{ex:cooling}
刚性容器内最初有 $N_0$ 个单原子理想气体粒子，温度 $T_0$。壁绝热，气体经薄壁小孔向真空缓慢泄流，余气始终充分平衡。求余气温度 $T$ 与粒子数 $N$ 的关系。
\Diagram{fig:coolingq}{粒子穿过系统边界并携带能量}{
 \draw[curve](0,0)rectangle(3,2);\fill[white](2.95,.8)rectangle(3.05,1.2);\node at(1.5,1){$N,T,V$};
 \draw[axis](3,1)--(5,1)node[right]{$-\dd N>0$};\node at(4.1,1.65){每粒子 $2\kb T$};\node at(1.5,-.5){刚性、绝热};
}
\end{example}
\begin{solution}
把图\ref{fig:coolingq}中容器内的余气作为变粒子数系统，取 $\dd N<0$。余气能量变化等于流出能量的负值，故
\[
 \dd U=2\kb T\dd N.
\]
另一方面 $U=3N\kb T/2$，全微分必须同时包含 $N$ 与 $T$ 的改变：
\[
 \frac32\kb(N\dd T+T\dd N)=2\kb T\dd N.
\]
移项得到 $3N\dd T=T\dd N$，分离变量并在初末态之间积分，$3\int_{T_0}^{T}\dd\vartheta/\vartheta=\int_{N_0}^{N}\dd N'/N'$。两边的原函数均为对数，代入上下限得
\[
 3\ln\frac T{T_0}=\ln\frac N{N_0},\qquad
 T=T_0\left(\frac N{N_0}\right)^{1/3}.
\]
这里忽略了二阶小量$\dd N\dd T$。逸出粒子比体内平均粒子多携带$\kb T/2$的能量，故余气降温。粒子离开时的$T$应取瞬时余气温度，不能在整个过程中用初温$T_0$计算流出能量。

\Diagram{fig:coolinga}{粒子数减少时温度和压强的变化}{
 \draw[axis](0,0)--(4.8,0)node[right]{$N/N_0$};\draw[axis](0,0)--(0,2.8)node[above]{归一化量};
 \draw[curve]plot[domain=.01:1,samples=60]({4*\x},{2.3*\x^(1/3)});
 \draw[curve,dashed]plot[domain=0:1,samples=60]({4*\x},{2.3*\x^(4/3)});
 \draw[guide](4,0)--(4,2.3)--(0,2.3);\node[below]at(4,0){$1$};\node[left]at(0,2.3){$1$};
 \node at(1.5,2.5){$T/T_0$};\node at(2.8,.7){$p/p_0$};
}
容器体积固定，$p/p_0=NT/(N_0T_0)=(N/N_0)^{4/3}$，如图\ref{fig:coolinga}。这里没有活塞体积功，却有随粒子跨边界的能量流；把“绝热”理解成余气内能不变，会同时漏掉开放系统的能量交换。
\end{solution}

\begin{example}[单原子、双原子和非线性多原子的比较]\label{ex:degrees}
在上一例条件下，分别讨论 $f=3,5,6$，求 $N(T)$、$p(T)$ 和减少一半粒子数后的温度。内部自由度在全过程保持同一激发状态。
\Diagram{fig:degreesq}{三类气体的平动和转动自由度}{
 \foreach \x/\s in{0/单原子,2.6/双原子,5.2/非线性分子}{\node[draw,minimum height=.7cm]at(\x,0){\s};}
 \node at(0,-.8){$f=3$};\node at(2.6,-.8){$f=5$};\node at(5.2,-.8){$f=6$};
}
\end{example}
\begin{solution}
图\ref{fig:degreesq}三类气体分别满足 $U=fN\kb T/2$。流出能量按式\eqref{eq:escapeenergy}取 $(f+1)\kb T/2$，于是
\[
 f(N\dd T+T\dd N)=(f+1)T\dd N
 \quad\Rightarrow\quad \frac{\dd N}{N}=f\frac{\dd T}{T}.
\]
沿用上一例的对数积分，只将系数$3$换为$f$，有$\ln(N/N_0)=f\ln(T/T_0)$，即$N/N_0=(T/T_0)^f$。体积固定，压强还多含一个温度因子，故$p/p_0=(T/T_0)^{f+1}$。
三种情形的内部能量和逸出能量可以按每粒子逐项列出：平动在体内占 $3\kb T/2$，在逸出样本占 $2\kb T$；双原子的两个转动二次项另加 $\kb T$，非线性分子的三个转动项另加 $3\kb T/2$。因此逸出能量依次为 $2\kb T$、$3\kb T$、$7\kb T/2$。穿孔不改变内部状态的平均分布，所以这些转动项在能量守恒中部分抵消，留下的是不同的储热能力与同一个额外平动能。

\Diagram{fig:degreesa}{相同泄流比例下不同自由度的温度下降}{
 \draw[axis](0,0)--(4.7,0)node[right]{$N/N_0$};\draw[axis](0,0)--(0,2.7)node[above]{$T/T_0$};
 \foreach \f in{3,5,6}{\draw[curve]plot[domain=.02:1,samples=50]({4*\x},{2.3*\x^(1/\f)});}
 \draw[guide](4,0)--(4,2.3)--(0,2.3);\node[below]at(4,0){$1$};\node[left]at(0,2.3){$1$};
 \node at(2,2.8){上至下：$f=6,5,3$};
}
当 $N=N_0/2$，三种温度依次为 $T_02^{-1/3}$、$T_02^{-1/5}$、$T_02^{-1/6}$。图\ref{fig:degreesa}中自由度越多，冷却越弱，因为逸出粒子相对于体内平均粒子多带走的能量始终只有 $\kb T/2$，而每个粒子的热容随 $f$ 增大。
\end{solution}

\begin{example}[温度和粒子数的时间规律]\label{ex:coolingtime}
容器体积 $V$、孔面积 $A_h$，其余条件同前。气体自由度 $f$ 恒定。求 $T(t)$、$N(t)$、$p(t)$，以及温度降至 $T_0/2$ 所需时间。
\Diagram{fig:cooltimeq}{泄流速率由孔面积和余气状态决定}{
 \draw(0,0)rectangle(3,1.8);\fill[white](2.95,.7)rectangle(3.05,1.1);\draw[axis](3,.9)--(5,.9);
 \node at(1.5,.9){$V,N(t),T(t)$};\node at(4,1.6){$A_h$};\node at(4,-.1){$\dot N<0$};
}
\end{example}
\begin{solution}
图\ref{fig:cooltimeq}中穿孔粒子数率是碰壁数率乘孔面积：
\[
 \dot N=-\frac{A_hN}{V}\sqrt{\frac{\kb T}{2\pi m}}.
\]
由 $N\propto T^f$，对时间求导 $\dot N/N=f\dot T/T$，故
\[
 \dot T=-\frac{A_h}{fV}\sqrt{\frac{\kb}{2\pi m}}T^{3/2}.
\]
分离变量并取初末态为积分上下限，
\[
 \left[-2\vartheta^{-1/2}\right]_{T_0}^{T}
 =-\frac{A_h}{fV}\sqrt{\frac{\kb}{2\pi m}}\int_0^t\dd s.
\]
代入端点并约去两边负号，得到
\[
 \frac1{\sqrt T}-\frac1{\sqrt{T_0}}
 =\frac{A_h}{2fV}\sqrt{\frac{\kb}{2\pi m}}t.
\]
定义初始数目衰减率 $\kappa_0=(A_h/V)\sqrt{\kb T_0/(2\pi m)}$ 以及 $t_c=2f/\kappa_0$，则
\[
 \frac{T}{T_0}=(1+t/t_c)^{-2},\quad
 \frac{N}{N_0}=(1+t/t_c)^{-2f},\quad
 \frac{p}{p_0}=(1+t/t_c)^{-2(f+1)}.
\]
定义的$\kappa_0$量纲为时间的倒数，是初始时刻的$-\dot N/N$；后续瞬时衰减率为$\kappa_0\sqrt{T/T_0}$，随降温而减小，所以粒子数遵循幂律而非恒温指数律。

\Diagram{fig:cooltimea}{变换后温度的线性时间关系}{
 \draw[axis](0,0)--(5,0)node[right]{$t/t_c$};\draw[axis](0,0)--(0,3.4)node[above]{$\sqrt{T_0/T}$};
 \draw[curve](0,.7)--(4,2.7);\node[left]at(0,.7){$1$};\node at(3,.8){$1+t/t_c$};
}
图\ref{fig:cooltimea}把$\sqrt{T_0/T}=1+t/t_c$画成直线，斜率为$1/t_c$。取$T/T_0=1/2$，有$1+t/t_c=\sqrt2$，故降至一半初温需$t=(\sqrt2-1)t_c$。
\end{solution}

\begin{example}[温变热容下的泄流]
气体单粒子的平衡内能为 $u(T)$，热容 $c_v(T)=\dd u/\dd T$。内部状态与平动独立，穿孔不选择内部状态。求一般的 $N(T)$ 关系，并取 $c_v(T)=c_0+bT$ 求显式结果。
\Diagram{fig:variablecvq}{逸出只对平动速率加权}{
 \node[draw,align=center]at(0,0){体内\\$u(T)$};\draw[axis](1,0)--(3,0);\node[draw,align=center]at(4,0){逸出\\$u(T)+\kb T/2$};
}
\end{example}
\begin{solution}
图\ref{fig:variablecvq}表示内部平均能量不变，只有平动抽样多出 $\kb T/2$。余气总能量为 $Nu(T)$，所以
\[
 Nc_v(T)\dd T+u(T)\dd N=[u(T)+\kb T/2]\dd N.
\]
消去两侧的 $u\dd N$，有 $\dd N/N=2c_v(T)\dd T/(\kb T)$。因此
\[
 \ln\frac N{N_0}=\int_{T_0}^T\frac{2c_v(\vartheta)}{\kb\vartheta}\dd\vartheta.
\]
若给每个分子加上相同的常数参考能，余气和流出粒子的能量同时改变，温度方程不变。温变热容下，单粒子内能差要积$c_v\dd T$，粒子数关系则要积$c_v\dd T/T$；两种权重须分别保留。

\Diagram{fig:variablecva}{温变热容分别贡献对数项和线性项}{
 \node[draw]at(0,0){$c_0/T$};\draw[->](0,-.35)--(0,-1);\node[below]at(0,-1){$c_0\ln(T/T_0)$};
 \node[draw]at(4,0){$b$};\draw[->](4,-.35)--(4,-1);\node[below]at(4,-1){$b(T-T_0)$};
}
按图\ref{fig:variablecva}将积分分成两项，
\[
 \ln\frac{N}{N_0}
 =\frac{2}{\kb}\left[c_0\ln\vartheta+b\vartheta\right]_{T_0}^{T}
 =\frac{2c_0}{\kb}\ln\frac{T}{T_0}+\frac{2b}{\kb}(T-T_0).
\]
两边取指数，得到
\[
 \frac N{N_0}=\left(\frac T{T_0}\right)^{2c_0/\kb}
 \exp\!\left[\frac{2b}{\kb}(T-T_0)\right].
\]
该式只在给定热容模型的适用温区有效；若内部弛豫比泄流还慢，则内部状态不能由同一瞬时温度描述，需要分别写能量交换方程。
\end{solution}

\section{光子气体的压强与辐射}
\label{th:radiation}
热平衡辐射的传播方向各向同性，光子速率为 $c$。若频率为 $\omega/(2\pi)$ 的单光子能量为 $\varepsilon=\hbar\omega$，其动量大小是 $\varepsilon/c$。为避免把粒子数与能量分布混在一起，直接用频率积分后的能量密度 $u$。每个方向的能量密度为 $u\dd\Omega/(4\pi)$。

单位面积、单位时间穿出小孔的能量为
\begin{equation}\label{eq:radiationflux}
 J_E=uc\int_0^{\pi/2}\frac12\sin\theta\cos\theta\dd\theta=\frac14cu.
\end{equation}
若黑体总辐射通量为 $\sigma_{\rm SB}T^4$，便有 $u=aT^4$，$a=4\sigma_{\rm SB}/c$。这里 $u$ 的单位为能量每体积，$J_E$ 为功率每面积；相差的因子 $c/4$ 有明确的速度和几何意义。只靠经典 Maxwell--Boltzmann 分布不能确定辐射谱和常数 $a$，本讲不把光子当作守恒的经典分子。

\begin{example}[腔内能量与小孔辐射]\label{ex:cavity}
温度 $T$ 的黑体腔体积为 $V$，小孔面积为 $A_h$。已知 $\sigma_{\rm SB}$，求腔内辐射能量与经小孔向冷真空辐射的功率。若外界也是温度 $T_e$ 的黑体辐射场，净功率如何改变？
\Diagram{fig:cavityq}{小孔连接腔内与外界辐射场}{
 \draw(0,0)rectangle(3,1.8);\fill[white](2.95,.7)rectangle(3.05,1.1);\node at(1.5,.9){$T,V$};
 \draw[axis](3,.9)--(5,.9)node[right]{$T_e$};\node at(4,1.5){$A_h$};
}
\end{example}
\begin{solution}
图\ref{fig:cavityq}中腔内能量密度为 $u=4\sigma_{\rm SB}T^4/c$，故 $U=uV=4\sigma_{\rm SB}VT^4/c$。由式\eqref{eq:radiationflux}，出射功率
\[
 P_{\rm out}=A_h\frac{cu}{4}=A_h\sigma_{\rm SB}T^4.
\]
小孔的出射通量由孔口处的各向同性辐射场决定，乘的是孔面积 $A_h$。腔体其余内壁虽然也吸收、发射光，但它们之间的交换留在腔内，不应把整个腔壁面积再乘进外流功率。外界有平衡辐射场时，两个单向功率分别存在，温度相同只意味着净值为零。$U$ 的量纲是焦耳，$P$ 的量纲是瓦特，连接两者的面积、体积和速度因子可由 $P/U=A_hc/(4V)$ 直接辨认。

\Diagram{fig:cavitya}{两侧能流分别计算后相减}{
 \draw[axis](0,.5)--(4,.5)node[midway,above]{$A_h\sigma_{\rm SB}T^4$};
 \draw[axis](4,-.5)--(0,-.5)node[midway,below]{$A_h\sigma_{\rm SB}T_e^4$};
}
图\ref{fig:cavitya}中，外界各向同性辐射的入射功率为 $A_h\sigma_{\rm SB}T_e^4$，相减得净外流功率 $P=A_h\sigma_{\rm SB}(T^4-T_e^4)$。该式要求腔内辐射始终接近平衡，且孔口接收完整的入射半球。孔过大导致明显不均匀，或孔前存在方向筛选时，须根据实际辐射分布重新作角积分。
\end{solution}

\begin{example}[辐射压强]\label{ex:photonpressure}
能量密度为 $u$ 的各向同性辐射照射理想反射平面。由光子动量变化求壁面压强。
\Diagram{fig:photonpressureq}{反射光子的法向动量变化}{
 \draw[line width=1.3pt](3,-1.2)--(3,1.5);\draw[axis](.3,1.2)--(3,0);\draw[axis](3,0)--(.3,-1.2);\draw[guide](1,0)--(4,0);
 \node at(1.4,.9){$\varepsilon/c$};\node at(2,.3){$\theta$};
}
\end{example}
\begin{solution}
图\ref{fig:photonpressureq}的单光子法向冲量为 $2\varepsilon\cos\theta/c$。把方向能量密度 $u\dd\Omega/(4\pi)$ 除以单光子能量，再乘扫掠速度 $c\cos\theta$ 与冲量，能量和速度中的因子分别约去：
\[
 p=\int_{\rm hem}\frac u\varepsilon\frac{\dd\Omega}{4\pi}
 c\cos\theta\frac{2\varepsilon}c\cos\theta
 =u\int_0^{\pi/2}\sin\theta\cos^2\theta\dd\theta
 =u\left[-\frac{\cos^3\theta}{3}\right]_0^{\pi/2}=\frac u3.
\]
频率分布不必预先给定。若用单位频率能量密度 $u_\omega$ 计算，每个频段的动量流同样给出 $p_\omega=u_\omega/3$，再对频率积分即可。因此 $p=u/3$ 依赖的是光子的能量动量关系和角度各向同性，不依赖这里是否知道具体的黑体谱。法向入射的单向光束则没有半球角度平均；对同样的入射能流，吸收与镜面反射的压强还相差两倍。

\Diagram{fig:photonpressurea}{角度平均将动量流分成三个等价方向}{
 \draw[axis](0,0)--(2,0)node[right]{$x$};\draw[axis](0,0)--(0,1.8)node[above]{$z$};\draw[axis](0,0)--(-1.4,-.9)node[left]{$y$};\node at(3,1){$\langle\cos^2\theta\rangle_{4\pi}=1/3$};
}
图\ref{fig:photonpressurea}给出同一结果的对称性解释。对黑色吸收壁，入射光的压强只有上述反射结果的一半；处于辐射平衡的壁还会向腔内发射光，发射反冲贡献另一半，总压强仍为 $u/3$。单向光束则不满足各向同性假设。
\end{solution}

\section{热力学势与偏导关系}
\label{th:potentials}
对成分、物质的量固定的简单可压缩系统，平衡态基本关系为
\begin{equation}\label{eq:fundamental}
 \dd U=T\dd S-p\dd V.
\end{equation}
这是状态空间中的微分恒等式，不要求实际过程可逆；但只有可逆传热才可把 $T\dd S$ 认作过程热量 $\delta Q_{\rm rev}$。不可逆过程中一般 $\delta Q\ne T\dd S$，边界做功也未必是内部平衡压强乘 $\dd V$。

以 $U(S,V)$ 为势，$T=(\partial U/\partial S)_V$、$p=-(\partial U/\partial V)_S$。等温条件下希望把 $T$ 作为自变量，便减去 $TS$；等压条件下希望把 $p$ 作为自变量，便加上 $pV$。于是四种势为
\[
 \begin{array}{c|c|c}
 \text{势}&\text{自然变量}&\text{全微分}\\\hline
 U&S,V&T\dd S-p\dd V\\
 F=U-TS&T,V&-S\dd T-p\dd V\\
 H=U+pV&S,p&T\dd S+V\dd p\\
 G=U+pV-TS&T,p&-S\dd T+V\dd p
 \end{array}
\]
对同一势的混合二阶导数交换次序，依次得到
\begin{align}
 -\left(\frac{\partial p}{\partial S}\right)_V
 &=\left(\frac{\partial T}{\partial V}\right)_S,&
 \left(\frac{\partial S}{\partial V}\right)_T
 &=\left(\frac{\partial p}{\partial T}\right)_V,\\
 \left(\frac{\partial T}{\partial p}\right)_S
 &=\left(\frac{\partial V}{\partial S}\right)_p,&
 \left(\frac{\partial S}{\partial p}\right)_T
 &=-\left(\frac{\partial V}{\partial T}\right)_p.
\end{align}
记忆关系式时保留自然变量，比仅记交换字母可靠。状态方程 $p=p(V,T)$ 可以给出压强的偏导，却不能独自确定任意温度下的内能；还需一个热容函数以及参考值，并满足相容条件。

\subsection{由状态方程求内能}
\label{th:energyidentity}
从 $U(V,S)$ 改用 $V,T$ 作坐标，需要沿等温方向改变 $V$ 时同时计入 $S(V,T)$ 的改变：
\begin{align}
 \left(\frac{\partial U}{\partial V}\right)_T
 &=\left(\frac{\partial U}{\partial V}\right)_S
 +\left(\frac{\partial U}{\partial S}\right)_V
 \left(\frac{\partial S}{\partial V}\right)_T\\
 &=-p+T\left(\frac{\partial p}{\partial T}\right)_V.
\end{align}
又定义 $C_V=(\partial U/\partial T)_V$，所以
\begin{equation}\label{eq:duvt}
 \dd U=C_V\dd T+\left[T\left(\frac{\partial p}{\partial T}\right)_V-p\right]\dd V.
\end{equation}
代入式\eqref{eq:fundamental}，得到计算熵的常用式
\begin{equation}\label{eq:dsvt}
 \dd S=\frac{C_V}{T}\dd T+
 \left(\frac{\partial p}{\partial T}\right)_V\dd V.
\end{equation}
由于$U$是状态函数，$\partial_VU_T=\partial_TU_V$。右侧的乘积求导给出$\partial_T(Tp_T-p)=p_T+Tp_{TT}-p_T=Tp_{TT}$，故
\begin{equation}\label{eq:cvcompat}
 \left(\frac{\partial C_V}{\partial V}\right)_T
 =T\left(\frac{\partial^2p}{\partial T^2}\right)_V.
\end{equation}
因此指定状态方程以后，不能任意再指定一个与它冲突的热容。

\subsection{焓与定压热容}
\label{th:enthalpy}
由$\dd H=T\dd S+V\dd p$，在定压可逆过程中$\delta Q=\dd H$，故$C_p=(\partial H/\partial T)_p$。写$\dd S=S_T|_p\dd T+S_p|_T\dd p$，其中$S_T|_p=C_p/T$；由$G(T,p)$的混合偏导有$S_p|_T=-V_T|_p$。代回焓微分，得到
\begin{equation}\label{eq:dhpt}
 \dd H=C_p\dd T+
 \left[V-T\left(\frac{\partial V}{\partial T}\right)_p\right]\dd p,
 \qquad
 \dd S=\frac{C_p}{T}\dd T-
 \left(\frac{\partial V}{\partial T}\right)_p\dd p.
\end{equation}
在单相稳定区，定义热膨胀系数 $\alpha_p=V^{-1}(\partial V/\partial T)_p$、等温压缩率 $\kappa_T=-V^{-1}(\partial V/\partial p)_T$，可将热容差写为
\begin{equation}\label{eq:cpminuscv}
 C_p-C_V=T\left(\frac{\partial p}{\partial T}\right)_V
 \left(\frac{\partial V}{\partial T}\right)_p
 =\frac{TV\alpha_p^2}{\kappa_T}.
\end{equation}
推导时在式\eqref{eq:dsvt}中沿定压方向取 $\dd V=(\partial V/\partial T)_p\dd T$，再与定压熵微分比较即可。稳定性要求 $\kappa_T>0$，从而 $C_p\ge C_V$。

\subsection{自然变量怎样进入偏导数}
\label{th:naturalderivatives}
势函数的自然变量来自其全微分中独立微分的系数。以 $F=U-TS$ 为例，展开乘积，
\[
 \dd F=\dd U-T\dd S-S\dd T
 =-S\dd T-p\dd V.
\]
$T\dd S$ 的抵消使 $F$ 可直接用 $T,V$ 作为自变量。于是
$S=-F_T|_V$、$p=-F_V|_T$；再求混合导数，
\[
 \left(\frac{\partial S}{\partial V}\right)_T
 =-F_{TV}
 =-F_{VT}
 =\left(\frac{\partial p}{\partial T}\right)_V.
\]
等式中两个负号来自同一个势的两个一阶导数。若从 $U$ 或 $H$ 出发，系数中的正负号不同，所得 Maxwell 关系也随之改变。实际使用时，先写势的全微分，再读出一阶导数，可直接确定符号。

“自然变量”不表示势只能在这些变量下表示。$U(S,V)$ 可改写成 $U(T,V)$，只是后者的偏导数须经过链式法则。沿 $V$ 固定的升温方向，
\[
 \left(\frac{\partial U}{\partial T}\right)_V
 =T\left(\frac{\partial S}{\partial T}\right)_V=C_V.
\]
沿 $T$ 固定的膨胀方向，熵一般同时改变，才有式\eqref{eq:duvt}中的体积系数。若把两条不同方向上的偏导数视为同一个“内能对体积的变化率”，便会丢掉 $T S_V|_T$ 这一项。

\Diagram{fig:derivativepaths}{同一点出发的等温方向与等熵方向}{
 \draw[axis](0,0)--(5.5,0)node[right]{$V$};
 \draw[axis](0,0)--(0,3.2)node[above]{$T$};
 \coordinate(A)at(1,2.4);\fill(A)circle(1.4pt)node[above left]{$A$};
 \draw[axis](A)--(4.6,2.4)node[right]{$T$ 固定};
 \draw[curve]plot[domain=1:4.6](\x,{2.4/\x^.45});
 \draw[->](2.1,1.72)--(2.6,1.56);
 \node at(3.8,.75){$S$ 固定};
}
图\ref{fig:derivativepaths}示意理想气体的两条方向。在等熵方向，内能变化可直接用 $-p\dd V$；沿等温方向则必须同时计算熵变。对于理想气体，两项恰好抵消，使等温内能不变；这一抵消是具体状态方程的结果。

\subsection{全微分、倒数关系与循环关系}
设单相区的状态方程为 $p=p(T,V)$，全微分为
$\dd p=p_T|_V\dd T+p_V|_T\dd V$。沿定压方向令 $\dd p=0$，得到
\[
 \left(\frac{\partial V}{\partial T}\right)_p
 =-\frac{(\partial p/\partial T)_V}
 {(\partial p/\partial V)_T}.
\]
这个比值是状态面上定压切线的斜率。保持同一个第三变量不变时，两个互为反函数的局部斜率满足倒数关系；保持条件不同时，则须从全微分重新推导。将下列三个偏导相乘，得到循环关系
\[
 \left(\frac{\partial p}{\partial T}\right)_V
 \left(\frac{\partial T}{\partial V}\right)_p
 \left(\frac{\partial V}{\partial p}\right)_T=-1.
\]
循环关系中的负号来自 $\dd p=0$ 时两项相抵。若某个分母偏导为零或发散，所用坐标可能失去局部可逆性，应回到全微分或改用另一组状态变量，而不能继续机械使用倒数。

\subsection{热容差、压缩率与绝热斜率}
\label{th:compressibilityratio}
记 $A=(\partial p/\partial T)_V$、$B=(\partial p/\partial V)_T$，其中稳定单相区有 $B<0$。沿定压方向，$\dd V/\dd T=-A/B$，代入熵微分，
\[
 C_p=C_V-\frac{TA^2}{B}.
\]
该式与响应系数形式的式\eqref{eq:cpminuscv}相同。沿等熵方向，先由式\eqref{eq:dsvt}求
\[
 \left(\frac{\partial T}{\partial V}\right)_S
 =-\frac{TA}{C_V},
\]
再代入状态方程全微分，得到
\[
 \left(\frac{\partial p}{\partial V}\right)_S
 =B-\frac{TA^2}{C_V}
 =B\,\frac{C_p}{C_V}.
\]
定义 $\kappa_S=-V^{-1}(\partial V/\partial p)_S$，便有
\[
 \frac{\kappa_T}{\kappa_S}=\frac{C_p}{C_V}.
\]
在 $C_V>0,\kappa_T>0$ 的稳定区，绝热压缩比等温压缩更困难，因为压缩引起的升温会进一步增大压强。理想气体给 $\kappa_T=1/p$、$\kappa_S=1/(\gamma p)$；实际气体也满足上述热容比关系，但热容比一般随状态变化。

小振幅声波中的一小团气体在振动周期内来不及与周围交换足够热量时，其压强密度关系近似等熵。由于 $\rho\propto1/V$，
\[
 c_s^2=\left(\frac{\partial p}{\partial\rho}\right)_S
 =-\frac{V}{\rho}\left(\frac{\partial p}{\partial V}\right)_S
 =\frac1{\rho\kappa_S}.
\]
理想气体因而有 $c_s^2=\gamma RT/M$。这里忽略了黏滞、热传导和内部自由度弛豫对声速的频率依赖；采用哪一热容，应与声波时间尺度内能否建立内部平衡相对应。

\subsection{实际气体中两种压缩斜率的比较}
用摩尔体积 $v$ 表示 van der Waals 状态方程
$p=RT/(v-b)-a/v^2$，摩尔定容热容为常数 $c$。在一摩尔的熵微分中，
$s_T|_v=c/T$、$s_v|_T=R/(v-b)$，于是等熵方向的温度变化为
\[
 \left(\frac{\partial T}{\partial v}\right)_s
 =-\frac{RT}{c(v-b)}.
\]
对状态方程作链式求导，可得
\begin{align*}
 \left(\frac{\partial p}{\partial v}\right)_s
 &=\left(\frac{\partial p}{\partial v}\right)_T+
 \left(\frac{\partial p}{\partial T}\right)_v
 \left(\frac{\partial T}{\partial v}\right)_s\\
 &=-\left(1+\frac Rc\right)\frac{RT}{(v-b)^2}
 +\frac{2a}{v^3}.
\end{align*}
令 $\gamma_0=1+R/c$，便有
\[
 c_s^2=\frac1M\left[
 \frac{\gamma_0RTv^2}{(v-b)^2}-\frac{2a}{v}\right].
\]
这里先完成等熵导数，再乘 $\dd v/\dd\rho=-v^2/M$。$\gamma_0$ 是这个常定容热容模型中出现的组合，实际的 $C_p/C_V$ 应由前文热容差求出，通常与 $\gamma_0$ 不同。

第一项给出热运动及排斥体积对恢复力的贡献，吸引项则削弱这一响应。令 $a,b\to0$，结果恢复为理想气体声速。实际气体的声速由状态方程的等熵斜率确定，既取决于压强，也取决于它随密度和温度变化的方式。

还可在等温失稳边界附近直接比较斜率。若 $p_v|_T\to0^-$，由前面的通式有
$p_v|_s\to-TA^2/C_V$，在 $A$ 与 $C_V$ 有限时仍可为有限负值。等温压缩率发散与等熵响应发散因而不是同一个条件。这一比较限于给定状态方程及常热容近似；接近真实临界点时，热容和弛豫时间本身也可能显著变化。


\begin{example}[理想气体的内能与熵]\label{ex:idealstate}
$\nu$ 摩尔理想气体满足 $pV=\nu RT$，已知 $C_V$ 为常量。相对于参考态 $(V_0,T_0)$，求 $U$、$H$、$S$ 的变化，说明为什么内能与体积无关。
\Diagram{fig:idealstateq}{参考态到任意态的两段积分路径}{
 \draw[axis](0,0)--(4.8,0)node[right]{$V$};\draw[axis](0,0)--(0,3)node[above]{$T$};
 \coordinate(a)at(1,.8);\coordinate(b)at(3.5,.8);\coordinate(c)at(3.5,2.3);
 \draw[axis](a)--(b);\draw[axis](b)--(c);\node[below]at(a){$A(V_0,T_0)$};\node[right]at(c){$C(V,T)$};\node[below]at(b){$B$};
}
\end{example}
\begin{solution}
$p_T|_V=\nu R/V$，所以$U_V|_T=T\nu R/V-p=0$。图\ref{fig:idealstateq}的等温段$A\to B$不改变内能；等容段$B\to C$给出$U-U_0=\int_{T_0}^TC_V\dd\vartheta=[C_V\vartheta]_{T_0}^T=C_V(T-T_0)$。焓增量为
\[
 H-H_0=U-U_0+pV-p_0V_0=(C_V+\nu R)(T-T_0),
\]
故 $C_p=C_V+\nu R$。熵微分为 $\dd S=C_V\dd T/T+\nu R\dd V/V$。
图中第一段固定 $T_0$、第二段固定最终体积 $V$，是为积分选取的可逆参考路径，并不限定系统实际如何由初态到末态。两段内能变化相加时，等温段为零，等容段才贡献热容积分。焓中还含 $pV=\nu RT$，所以定压热容多出 $\nu R$；这里 $C_V,C_p$ 均为整个 $\nu$ 摩尔系统的热容。把它们换成摩尔量时，内能与熵的温度项还要同时乘回物质的量。

\Diagram{fig:idealstatea}{两段路径分别给出体积项与温度项}{
 \node[draw,align=center]at(0,0){$A\to B$ 等温\\$\nu R\ln(V/V_0)$};
 \node[draw,align=center]at(4.4,0){$B\to C$ 等容\\$C_V\ln(T/T_0)$};\draw[->](1.5,-.8)--(3,-.8);
}
按图\ref{fig:idealstatea}，两段熵变分别为$\nu R\int_{V_0}^V\dd w/w=\nu R[\ln w]_{V_0}^V$和$C_V\int_{T_0}^T\dd\vartheta/\vartheta=C_V[\ln\vartheta]_{T_0}^T$，相加得
\[
 S-S_0=C_V\ln\frac T{T_0}+\nu R\ln\frac V{V_0}.
\]
内能仅依赖温度并非“所有气体”的性质，而是理想气体状态方程与热力学相容性的结果。若 $C_V=C_V(T)$，内能温度项改为 $\int C_V\dd T$，熵温度项改为 $\int C_V\dd T/T$。
\end{solution}

\begin{example}[van der Waals 气体的内能与熵]\label{ex:vdwstate}
固定物质的量 $\nu$，气体满足
\[
 p=\frac{\nu RT}{V-B}-\frac A{V^2},\quad A=\nu^2a,\ B=\nu b,
\]
其中 $a,b$ 为摩尔形式的常数。给定常量 $C_V$，求相对于参考态 $(V_0,T_0)$ 的内能和熵变化。只讨论稳定单相区，$V>B$。
\Diagram{fig:vdwstateq}{排除体积改变可用的构型空间}{
 \draw[axis](0,0)--(5,0)node[right]{$V$};\draw[axis](0,0)--(0,2.3)node[above]{$p$};
 \draw[guide](1,0)--(1,2.1)node[above]{$B$};\draw[curve]plot[domain=1.6:4.6,samples=50](\x,{1/(\x-1)-.5/\x^2});\node at(3.6,1.8){稳定区示意};
}
\end{example}
\begin{solution}
图\ref{fig:vdwstateq}的 $B$ 是固定量，求偏导时不随 $T$ 改变。有 $p_T|_V=\nu R/(V-B)$，代入式\eqref{eq:duvt}：
\[
 T p_T-p=\frac A{V^2},\qquad
 \dd U=C_V\dd T+\frac A{V^2}\dd V.
\]
分别积分温度项和体积项。温度积分同上一例；体积积分为$\int_{V_0}^{V}A\dd w/w^2=[-A/w]_{V_0}^{V}=A(1/V_0-1/V)$，所以
\[
 U-U_0=C_V(T-T_0)+A\left(\frac1{V_0}-\frac1V\right).
\]
这里$A,B$均按固定物质的量看作常数。原函数$-A/V$的导数为$+A/V^2$，因此膨胀使负的吸引势能升高。参考常数只改变内能与熵的零点，不改变偏导、热容或两态之间的差。

\Diagram{fig:vdwstatea}{熵微分中的吸引项相消}{
 \node[draw,align=center]at(0,0){$\dd U$ 中\\$+A\dd V/V^2$};\node at(2.7,0){$+$};
 \node[draw,align=center]at(5.1,0){$p\dd V$ 中\\$-A\dd V/V^2$};\node at(2.5,-1){$T\dd S=C_V\dd T+\nu RT\dd V/(V-B)$};
}
按图\ref{fig:vdwstatea}在$\dd U+p\dd V$中消去吸引项，再除以$T$。温度项仍给$C_V\ln(T/T_0)$；体积项令$z=w-B$、$\dd z=\dd w$，上下限变为$V_0-B,V-B$，于是$\nu R\int_{V_0-B}^{V-B}\dd z/z=\nu R[\ln z]_{V_0-B}^{V-B}$。两项相加，
\[
 S-S_0=C_V\ln\frac T{T_0}+\nu R\ln\frac{V-B}{V_0-B}.
\]
吸引参数 $A$ 出现在内能中，但在以 $V,T$ 为自变量的熵差中抵消；排除体积 $B$ 则保留。若把 $A,B$ 当成摩尔常数却仍用总容积 $V$，会漏掉物质的量因子。
\end{solution}

\begin{example}[实际气体的可逆绝热膨胀]\label{ex:vdwadiabat}
上述 van der Waals 气体由 $(V_1,T_1)$ 可逆绝热膨胀至体积 $V_2$，热容为常量。求 $T_2$、$p_2$ 和气体对外做功，并比较同初温、同体积比的理想气体。
\Diagram{fig:vdwadiabatq}{绝热线上的初末两态}{
 \draw[axis](0,0)--(5,0)node[right]{$V$};\draw[axis](0,0)--(0,3)node[above]{$p$};
 \draw[curve]plot[domain=1:4.3](\x,{1.8/(\x-.3)^1.3-.2/\x^2});\fill(1,{1.8/.7^1.3-.2})circle(1.2pt)node[above right]{$1$};\fill(4,{1.8/3.7^1.3-.2/16})circle(1.2pt)node[above right]{$2$};\draw[->]plot[domain=2:2.5](\x,{1.8/(\x-.3)^1.3-.2/\x^2});
}
\end{example}
\begin{solution}
图\ref{fig:vdwadiabatq}的路径可逆绝热，故 $\dd S=0$。由上一例的熵微分
\[
 C_V\frac{\dd T}T+\nu R\frac{\dd V}{V-B}=0.
\]
取初态$1$为积分下限，利用上一例的两个对数积分，得到$C_V\ln(T/T_1)+\nu R\ln[(V-B)/(V_1-B)]=0$，即$T(V-B)^{\nu R/C_V}=T_1(V_1-B)^{\nu R/C_V}$。
所以 $T_2=T_1[(V_1-B)/(V_2-B)]^{\nu R/C_V}$，压强再由状态方程求得
\[
 p_2=\frac{\nu RT_2}{V_2-B}-\frac A{V_2^2}.
\]
令 $s=\nu R/C_V$，将绝热温度关系代入热压强项，可写成 $T(V)=K(V-B)^{-s}$。其功积分中的热运动部分为
\begin{align*}
 \nu RK\int_{V_1}^{V_2}(V-B)^{-s-1}\dd V
 &=\left[-\frac{\nu RK}{s}(V-B)^{-s}\right]_{V_1}^{V_2}\\
 &=\frac{\nu RK}{s}\bigl[(V_1-B)^{-s}-(V_2-B)^{-s}\bigr].
\end{align*}
因$\nu R/s=C_V$，热运动部分为$C_V(T_1-T_2)$。吸引项单独积分为$\int_{V_1}^{V_2}(-A/V^2)\dd V=[A/V]_{V_1}^{V_2}=-A(1/V_1-1/V_2)$，代表克服分子间吸引所需的能量。两项之和即气体对外做的功。

\Diagram{fig:vdwadiabata}{绝热功等于内能减少}{
 \node[draw]at(0,0){热运动能减少};\node[draw]at(4,0){吸引势能增加};\draw[->](0,-.4)--(2,-1);\draw[->](4,-.4)--(2,-1);\node[below]at(2,-1){$W=C_V(T_1-T_2)-A(1/V_1-1/V_2)$};
}
第一定律给 $W=U_1-U_2$，即图\ref{fig:vdwadiabata}的表达式。吸引势能在膨胀时上升，因而减少可供做功的热运动能。对于相同 $V_2/V_1>1$，$B>0$ 使 $(V_1-B)/(V_2-B)<V_1/V_2$，温度比理想气体降得更多；但功还包含吸引能，不能只比较温降。等价的绝热线为 $(p+A/V^2)(V-B)^{1+\nu R/C_V}=\text{常量}$。
\end{solution}

\begin{example}[从辐射压强求能量与熵]\label{ex:photonstate}
平衡光子气体满足 $p=b_rT^4$，$b_r$ 是常量。利用热力学关系求 $U(V,T)$、$S(V,T)$。说明不守恒的光子数为何不妨碍使用 $\dd U=T\dd S-p\dd V$。
\Diagram{fig:photonstateq}{可改变光子数的平衡辐射腔}{
 \draw(0,0)rectangle(3.5,1.8);\draw[decorate,decoration={snake,amplitude=1.3pt,segment length=7pt},->](.4,.5)--(3,.5);
 \draw[decorate,decoration={snake,amplitude=1.3pt,segment length=7pt},->](3,1.3)--(.4,1.3);\node[right]at(3.7,.9){$p=b_rT^4$};
}
\end{example}
\begin{solution}
一般开放系统的基本关系中还含 $\mu\dd N$。图\ref{fig:photonstateq}的热平衡光子可以被壁吸收或发射，光子化学势为零，因此该项消失。用式\eqref{eq:duvt}，
\[
 \left(\frac{\partial U}{\partial V}\right)_T=-b_rT^4+4b_rT^4=3b_rT^4.
\]
先积分得 $U=3b_rVT^4+\phi(T)$。均匀系统在固定温度下能量随体积广延，$U(\lambda V,T)=\lambda U(V,T)$，故不能保留体积无关的 $\phi(T)$。于是 $U=3b_rVT^4$。
广延性也可通过复制腔体理解：同温、同体积的两份辐射合在一起，忽略界面效应，能量应加倍。求出下文的熵后，代入$G=U+pV-TS$得到$G=0$，与平衡光子的零化学势一致。固定粒子数的分子气体仍须保留其粒子数约束。

\Diagram{fig:photonstatea}{在状态平面中积分熵微分}{
 \draw[axis](0,0)--(4.5,0)node[right]{$V$};\draw[axis](0,0)--(0,2.8)node[above]{$T$};
 \draw[axis](1,.7)--(3.2,.7);\draw[axis](3.2,.7)--(3.2,2.1);\node[below]at(1,.7){$(V_0,T_0)$};\node[right]at(3.2,2.1){$(V,T)$};
}
对$U=3b_rVT^4$作乘积微分，$\dd U=3b_rT^4\dd V+12b_rVT^3\dd T$。加$p\dd V$并除以$T$，得$\dd S=4b_rT^3\dd V+12b_rVT^2\dd T$。按图\ref{fig:photonstatea}先沿$T_0$等温段积分，再沿$V$等容段积分，
\begin{align*}
 \Delta S_{\rm iso}&=\int_{V_0}^{V}4b_rT_0^3\dd w=4b_rT_0^3(V-V_0),\\
 \Delta S_{\rm vol}&=\int_{T_0}^{T}12b_rV\vartheta^2\dd\vartheta
 =[4b_rV\vartheta^3]_{T_0}^{T}=4b_rV(T^3-T_0^3).
\end{align*}
相加并消去$VT_0^3$项，得到
\[
 S-S_0=4b_r(VT^3-V_0T_0^3).
\]
取无辐射的零温极限为熵零点，则 $S=4b_rVT^3$。与 $u=aT^4$ 对比有 $a=3b_r$，从而 $S=4aVT^3/3$。
\end{solution}

\begin{example}[状态方程与热容的相容性]
某系统状态方程为 $p=\nu RT/V+qT^2/V^2$，其中 $q>0$。有人同时给定 $C_V=C_0$，这是否相容？若规定在参考体积 $V_*$ 处 $C_V=C_*(T)$，求相容的 $C_V(V,T)$ 与内能。
\Diagram{fig:compatq}{在不同体积处比较定容升温}{
 \draw[axis](0,0)--(4.5,0)node[right]{$V$};\draw[axis](0,0)--(0,2.7)node[above]{$T$};
 \draw[axis](1,.5)--(1,2.1);\draw[axis](3,.5)--(3,2.1);\node[below]at(1,0){$V_*$};\node[below]at(3,0){$V$};
}
\end{example}
\begin{solution}
$p_{TT}|_V=2q/V^2$，故相容式\eqref{eq:cvcompat}要求 $(C_V)_V=2qT/V^2$，不能为零。图\ref{fig:compatq}中两条定容路径的热容存在体积差。积分得
\[
 C_V(V,T)=C_*(T)+2qT\left(\frac1{V_*}-\frac1V\right).
\]
这一结果来自相容条件的体积积分：
\[
 C_V(V,T)-C_V(V_*,T)
 =\int_{V_*}^{V}\frac{2qT}{w^2}\dd w
 =\left[-\frac{2qT}{w}\right]_{V_*}^{V}
 =2qT\left(\frac1{V_*}-\frac1V\right).
\]
再由$U_V=Tp_T-p=qT^2/V^2$积分，得$U=\Phi(T)-qT^2/V$。其温度导数为$\Phi'(T)-2qT/V$；在$V=V_*$匹配热容，确定$\Phi'=C_*(T)+2qT/V_*$。参考热容只规定一条等容线，其余体积处的热容由状态相容性确定。

\Diagram{fig:compata}{同一状态函数的两条积分路径必须闭合}{
 \draw[axis](0,0)--(3,0);\draw[axis](3,0)--(3,1.7);\draw[axis](3,1.7)--(0,1.7);\draw[axis](0,1.7)--(0,0);\node at(1.5,.85){$\oint\dd U=0$};
}
相对于$(V_*,T_*)$，先在$V_*$上从$T_*$升温至$T$，再在终温下从$V_*$改变体积到$V$。第二段的内能增量为$\int_{V_*}^VqT^2w^{-2}\dd w=[-qT^2/w]_{V_*}^V$，所以
\[
 U-U_*=\int_{T_*}^{T}C_*(\vartheta)\dd\vartheta
 +qT^2\left(\frac1{V_*}-\frac1V\right).
\]
图\ref{fig:compata}中，内能沿闭合矩形的增量应为零。若另行指定与状态方程不相容的常热容，两条路径将给出不同内能，模型便失去唯一的状态函数。
\end{solution}

\subsection{由状态关系到过程方程}
\label{th:processclosure}
状态方程规定允许的平衡态，过程约束规定系统怎样在这些状态之间变化。对固定物质的量，式\eqref{eq:duvt}给出内能微分，再用第一定律才能确定热量和功。若过程准静态且只有体积功，$\delta W=p\dd V$，于是
\[
 \delta Q=C_V\dd T+
 T\left(\frac{\partial p}{\partial T}\right)_V\dd V.
\]
这一式与气体的$T\dd S$相同；若外界热源与气体存在有限温差，气体和热源组成的整体仍会产生熵，不能仅凭此式断定整体过程可逆。若外压与内部压强不同，实际功应按边界的外压计算，差值须在不可逆能量转化和熵产生中处理。

给定过程 $V=V(T)$ 后，沿这一路径的热量与温变之比为
\[
 C_{\cal P}=\frac{\delta Q}{\dd T}
 =C_V+T\left(\frac{\partial p}{\partial T}\right)_V
 \frac{\dd V}{\dd T}.
\]
它取决于路径，通常既不是$C_V$也不是$C_p$。路径热容可以为负，即吸热时温度下降。例如对内能只随温度变化且$C_V>0$的理想气体，若吸收热量少于同时对外做的功，差额由内能支付，温度便下降。实际气体的内能还可能含体积项，须联立完整的$\dd U$判断温变。平衡稳定性中的热容条件针对指定约束下的状态响应，应明确保持哪些变量不变。

\begin{example}[多方过程的热量、功与路径热容]
\label{ex:polytropicheat}
$\nu$ 摩尔常热容理想气体经历可逆过程 $pV^a=\text{常量}$，初末体积为 $V_1,V_2$。已知 $C_V=\nu c_{V,m}$、$\gamma=1+\nu R/C_V$。求 $a\ne1$ 时的功、热量和路径热容，并说明 $1<a<\gamma$ 的膨胀过程为什么吸热却降温。
\Diagram{fig:polytropicheatq}{从同一初态出发的等温、多方与绝热线}{
 \draw[axis](0,0)--(5,0)node[right]{$V/V_1$};
 \draw[axis](0,0)--(0,3.2)node[above]{$p/p_1$};
 \foreach \a/\s in{1/solid,1.2/dashed,1.4/densely dotted}
  \draw[curve,\s]plot[domain=1:4.5](\x,{2.5/\x^\a});
 \fill(1,2.5)circle(1.2pt)node[above]{$A$};
 \node at(3.7,2.4){实线 $a=1$};
 \node at(3.7,1.9){虚线 $a=1.2$};
 \node at(3.7,1.4){点线 $\gamma=1.4$};
}
\end{example}
\begin{solution}
图\ref{fig:polytropicheatq}的三条曲线经过同一初态。由理想气体状态方程，
$TV^{a-1}=T_1V_1^{a-1}$，故
$T_2/T_1=(V_1/V_2)^{a-1}$。令 $K=p_1V_1^a$，逐项积分得
\[
 W=\int_{V_1}^{V_2}KV^{-a}\dd V
 =\left[\frac{KV^{1-a}}{1-a}\right]_{V_1}^{V_2}
 =\frac{p_2V_2-p_1V_1}{1-a}
 =\frac{\nu R(T_2-T_1)}{1-a}.
\]
再由 $\Delta U=C_V(T_2-T_1)$ 和 $Q=\Delta U+W$，
\[
 Q=C_V\frac{\gamma-a}{1-a}(T_2-T_1),
 \qquad C_{\cal P}=C_V\frac{\gamma-a}{1-a}.
\]
\Diagram{fig:polytropicheata}{路径热容的符号由两个特征指数分隔}{
 \draw[axis](0,0)--(6.7,0)node[right]{$a$};
 \draw(2.2,-.15)--(2.2,.15)node[above]{$1$};
 \draw(4.4,-.15)--(4.4,.15)node[above]{$\gamma$};
 \node at(1,-.6){$C_{\cal P}>0$};
 \node at(3.3,-.6){$C_{\cal P}<0$};
 \node at(5.5,-.6){$C_{\cal P}>0$};
 \node at(3.3,1){$a=1$ 等温；$a=\gamma$ 可逆绝热};
}
在图\ref{fig:polytropicheata}的中间区间，膨胀给$T_2<T_1$，而$C_{\cal P}<0$，故$Q>0$。功仍为正，且$W>Q$，多出的功来自内能减少。定压过程$a=0$恢复$C_{\cal P}=C_p$；$a=\gamma$给$Q=0$。$a=1$时温变为零，应单独积分：$W=\nu RT\int_{V_1}^{V_2}\dd V/V=\nu RT[\ln V]_{V_1}^{V_2}=\nu RT\ln(V_2/V_1)$，且$\Delta U=0$给$Q=W$，此时热量除以温变不定义有限的路径热容。
\end{solution}

\subsection{内能与熵积分中的参考函数}
由一个偏导积分得到的“常数”可能仍依赖其他变量。例如给出
$U_V|_T=F(T,V)$，积分得到
\[
 U(T,V)=\int_{V_\ast}^{V}F(T,w)\dd w+\Phi(T).
\]
再对温度求导，
\[
 C_V(T,V)=\int_{V_\ast}^{V}
 \left(\frac{\partial F}{\partial T}\right)_w\dd w+\Phi'(T).
\]
指定参考体积上的热容，才确定 $\Phi'(T)$；再给一个内能参考值，才确定最终的加法常数。若热容已经在所有体积上给出，它必须满足式\eqref{eq:cvcompat}，才能与同一内能函数相容。

熵的积分有相同结构。用一条先定容、后等温的路径，
\[
 S(T,V)-S(T_\ast,V_\ast)
 =\int_{T_\ast}^{T}\frac{C_V(t,V_\ast)}{t}\dd t
 +\int_{V_\ast}^{V}\left(\frac{\partial p(T,w)}{\partial T}\right)_w\dd w.
\]
第二项沿终温 $T$ 的等温线积分。对整个结果求温度导数，利用相容关系，可恢复终体积处的 $C_V(T,V)/T$。若选择先等温、后定容的另一条路径，相应的等温段应取初温，定容段也应改在终体积处积分；两段须一起调整。

\subsection{实际气体的绝热线与机械稳定区}
可逆绝热约束由熵微分直接给出，
\[
 \frac{C_V(T,V)}{T}\dd T+
 \left(\frac{\partial p}{\partial T}\right)_V\dd V=0.
\]
只有在系数具有可分离形式时，才可积分成简单幂律。对 $p=\nu RT/(V-B)-A/V^2$ 且总定容热容为常数的模型，
$(\partial p/\partial T)_V=\nu R/(V-B)$，由积分得到温度与有效自由体积的关系。吸引项虽然不直接出现在熵的体积系数中，却仍影响内能、绝热功和机械稳定性。

若路径进入 $(\partial p/\partial V)_T>0$ 的均匀态分支，微小密度扰动会进一步增大，单相均匀模型失稳。此时沿形式上的代数曲线继续积分，并不代表真实平衡过程；应改用两相共存条件及相比例。第四讲讨论的相变平台正是状态描述发生变化的地方。


\section{热容、状态方程与过程约束的综合应用}
状态函数的全微分联系任意相邻平衡态，过程条件则确定其中的一条路径。定压过程取 $\dd p=0$，可逆绝热过程取 $\dd S=0$，稳流节流过程取 $\dd H=0$。应用同一组状态关系，可以得到不同约束下的温度变化与热量交换。

\begin{example}[容器热容不可忽略的泄流冷却]
刚性容器的壁热容为常量 $C_w$，与容器内 $N$ 个自由度为 $f$ 的理想气体粒子保持同温 $T$。整体对外绝热，气体从小孔向真空缓慢泄出。初态 $N_0,T_0$。求 $N(T)$，并讨论极少余气时温度的形式极限。
\Diagram{fig:wallcapacityq}{气体与器壁共同储能，粒子由小孔逸出}{
 \draw[line width=2pt](0,0)rectangle(3.3,1.8);\fill[white](3.22,.73)rectangle(3.38,1.07);
 \node at(1.65,.9){$N,c_v,T$};\node at(1.65,2.3){器壁 $C_w,T$};
 \draw[axis](3.3,.9)--(5,.9);\node at(4.3,.2){$\varepsilon_e$};
}
\end{example}
\begin{solution}
图\ref{fig:wallcapacityq}中系统能量 $U=Nc_vT+C_wT$，$c_v=f\kb/2$。每个流出粒子携能 $(c_v+\kb/2)T$。取 $\dd N<0$，能量守恒为
\[
 (Nc_v+C_w)\dd T+c_vT\dd N
 =(c_v+\kb/2)T\dd N,
\]
即 $(Nc_v+C_w)\dd T=(\kb/2)T\dd N$。把 $N+C_w/c_v$ 作为整体，
\[
 \frac{\dd N}{N+C_w/c_v}=f\frac{\dd T}T.
\]
从初态积分，$[\ln(N'+C_w/c_v)]_{N_0}^{N}=f[\ln\vartheta]_{T_0}^{T}$，即$\ln[(N+C_w/c_v)/(N_0+C_w/c_v)]=f\ln(T/T_0)$。取指数得到
\[
 N+\frac{C_w}{c_v}=\left(N_0+\frac{C_w}{c_v}\right)
 \left(\frac T{T_0}\right)^f.
\]
其中$C_w/c_v$表示与器壁热容相等的气体粒子数。气体流失时，这一等效数不变，器壁所占的热容比例逐渐增大，温度下降也随之放缓。

\Diagram{fig:wallcapacitya}{器壁储能使极少余气的温度保持有限}{
 \draw[axis](0,0)--(4.7,0)node[right]{$N/N_0$};\draw[axis](0,0)--(0,2.9)node[above]{$T/T_0$};\draw[curve]plot[domain=0:1]({4*\x},{2.4*((\x+.3)/1.3)^(1/3)});\draw[guide](0,1.47)--(4,1.47);\node[right]at(4,1.47){$T_{\rm lim}/T_0$};
}
形式上 $N\to0$ 时，图\ref{fig:wallcapacitya}给 $T_{\rm lim}=T_0[C_w/(C_w+N_0c_v)]^{1/f}$。$C_w=0$ 恢复无热容器壁结果；$C_w\gg N_0c_v$ 则温降很小。末期若碰撞不足以使余气平衡，或器壁与余气的热交换太慢，这一共同温度模型会先失效，不能把形式极限当成任意稀薄程度都精确可测的终态。
\end{solution}

\begin{example}[分子气体与辐射共同膨胀]
同一腔内有 $N$ 个单原子理想气体粒子和与之平衡的辐射，温度均为 $T$。气体粒子数固定，辐射能量密度 $aT^4$，腔壁热容忽略。求总压强、总内能、熵差及可逆绝热膨胀的温度体积关系。
\Diagram{fig:gammaradq}{分子和光子共享体积与温度}{
 \draw(0,0)rectangle(4,1.8);\foreach \p in{(.6,.4),(1.2,1.2),(2.8,.5),(3.5,1.3)}\fill\p circle(1.2pt);\draw[decorate,decoration={snake,amplitude=1.2pt,segment length=8pt},->](.5,.8)--(3.4,.8);\node at(2,-.5){$N$ 固定，光子数不固定};
}
\end{example}
\begin{solution}
图\ref{fig:gammaradq}的两部分独立贡献压强和能量，
\[
 p=\frac{N\kb T}V+\frac a3T^4,\qquad
 U=\frac32N\kb T+aVT^4.
\]
总熵差是理想气体与辐射熵差之和：
\begin{align*}
 S-S_0={}&\frac32N\kb\ln\frac T{T_0}+N\kb\ln\frac V{V_0}\\
 &+\frac{4a}3(VT^3-V_0T_0^3).
\end{align*}
可逆绝热令此式为零，便得到有限变化的隐式关系。若求局部温度斜率，更直接从 $\dd U+p\dd V=0$ 出发：
\[
 \left(\frac32N\kb+4aVT^3\right)\dd T
 +\left(\frac{N\kb T}V+\frac{4a}3T^4\right)\dd V=0.
\]
熵的气体项按固定粒子数积分，辐射项按允许光子产生和吸收的平衡场积分，两者虽共享温度，粒子数约束不同。令 $X=aVT^3/(N\kb)$，图中的局部指数为 $\Gamma=(1+4X/3)/(3/2+4X)$。对 $X$ 求导，
\[
 \frac{\dd\Gamma}{\dd X}
 =\frac{\tfrac43(\tfrac32+4X)-4(1+\tfrac43X)}{(3/2+4X)^2}
 =-\frac{2}{(3/2+4X)^2}<0.
\]
因此它连续地从分子气体极限降至辐射极限。把 $X$ 当成常数仅适用于局部斜率的比较；沿实际绝热路径 $V,T$ 都变化，有限变化仍应使用完整的熵方程。

\Diagram{fig:gammarada}{局部绝热指数在两种主导机制间变化}{
 \draw[axis](0,0)--(4.9,0)node[right]{$aVT^3/(N\kb)$};\draw[axis](0,0)--(0,2.8)node[above]{$-\dd\ln T/\dd\ln V$};\draw[curve]plot[domain=0:4.4](\x,{3*(1+4*\x/3)/(1.5+4*\x)});\draw[guide](0,1)--(4.5,1);\node[left]at(0,2){$2/3$};\node[left]at(0,1){$1/3$};
}
因此图\ref{fig:gammarada}的指数为 $(N\kb+4aVT^3/3)/(3N\kb/2+4aVT^3)$。气体主导时趋于 $2/3$，辐射主导时趋于 $1/3$。中间区域不能用一个固定热容比描述整条绝热线。任意两态之间的绝热功仍由 $W=U_1-U_2$ 计算，但要先按熵约束求末温。
\end{solution}

\begin{example}[van der Waals 气体的定压热容]
一摩尔 van der Waals 气体在稳定单相区，$p=RT/(v-b)-a/v^2$，$C_{V,m}=c$。求 $C_{p,m}$，并解释绝热线的指数 $1+R/c$ 是否等于实际的 $C_{p,m}/C_{V,m}$。
\Diagram{fig:vdwcpq}{定压方向与等温线斜率的关系}{
 \draw[axis](0,0)--(4.8,0)node[right]{$v$};\draw[axis](0,0)--(0,3.2)node[above]{$p$};\draw[curve]plot[domain=1.2:4.2](\x,{2/(\x-.3)-.2/\x^2});\draw[curve]plot[domain=1.2:4.2](\x,{2.8/(\x-.3)-.2/\x^2});\draw[axis](1.6,1.460337)--(2.162827,1.460337);\node[above left]at(1.6,1.460337){$T$};\node at(3.4,2.4){$T+\dd T$};
}
\end{example}
\begin{solution}
图\ref{fig:vdwcpq}的定压方向满足 $0=p_T\dd T+p_v\dd v$，所以 $v_T|_p=-p_T/p_v$。热容差为 $C_p-C_V=Tp_Tv_T|_p=-Tp_T^2/p_v$。本题
\[
 p_T=\frac R{v-b},\qquad
 p_v=-\frac{RT}{(v-b)^2}+\frac{2a}{v^3}.
\]
故
\[
 C_{p,m}=c+\frac{TR^2/(v-b)^2}{RT/(v-b)^2-2a/v^3}
 =c+\frac R{1-2a(v-b)^2/(RTv^3)}.
\]
热容差中的分母是 $-p_v|_T$，它同时控制定压升温时所需的体积变化。分母越小，同样的温升引起的膨胀越大，定压供热与定容供热的差也越大。这里 $C_V=c$ 固定不妨碍 $C_p$ 随状态改变，因为附加项还含 $v,T$。前面绝热线中的 $1+R/c$ 由 $\dd S=c\dd T/T+R\dd v/(v-b)$ 积分得到；这一推导从未要求实际气体满足 $C_p-C_V=R$。

\Diagram{fig:vdwcpa}{稳定区趋近机械不稳定边界时热容差增大}{
 \node[draw]at(0,0){$p_v|_T<0$};\draw[axis](1.2,0)--(2.6,0);\node[draw]at(4.6,0){$C_p-C_V=-Tp_T^2/p_v>0$};\node at(2.3,-.9){$p_v\to0^-$ 时可显著增大};
}
图\ref{fig:vdwcpa}中，机械稳定要求分母为正。一般情况下$C_p-C_V$随状态变化，绝热线$(p+a/v^2)(v-b)^{1+R/c}=\text{常量}$中的指数$1+R/c$来自熵积分。当分母中的吸引修正$2a(v-b)^2/(RTv^3)$可忽略时，$C_p-C_V\simeq R$，该指数才近似等于热容比。特别地，$a=0$时这两者精确相等，即使排除体积$b$不为零；不能把这一相等关系当作理想气体的充分条件。
\end{solution}

\begin{example}[实际气体温变热容的绝热线]
一摩尔气体仍满足 van der Waals 方程，但 $C_{V,m}(T)=c_0+qT>0$。给定初态 $(v_1,T_1)$，可逆绝热膨胀至 $v_2$。写出决定 $T_2$ 的方程和做功公式，证明正温度解在适用温区内唯一。
\Diagram{fig:vdwvariableq}{绝热路径的温度斜率由即时热容决定}{
 \draw[axis](0,0)--(4.7,0)node[right]{$v$};\draw[axis](0,0)--(0,3)node[above]{$T$};\draw[curve](1,2.5)to[out=-55,in=170](4,.7);\node[above]at(1,2.5){$1$};\node[right]at(4,.7){$2$};
}
\end{example}
\begin{solution}
由于状态方程对温度仍是线性的，相容条件允许$C_V$为任意温度函数。图\ref{fig:vdwvariableq}的熵微分为$(c_0/T+q)\dd T+R\dd v/(v-b)$。取初末态作积分端点，两个原函数分别是$c_0\ln T+qT$与$R\ln(v-b)$，故
\[
 c_0\ln\frac{T_2}{T_1}+q(T_2-T_1)
 +R\ln\frac{v_2-b}{v_1-b}=0.
\]
左侧对 $T_2$ 的导数为 $c_0/T_2+q=C_V(T_2)/T_2>0$，因而在热容为正的研究温区内至多有一个解。膨胀的体积对数项为正，温度积分必须为负，所以 $T_2<T_1$。
温度积分应按两个不同的被积函数分别处理：熵中是 $(c_0+qT)/T=c_0/T+q$，内能中是 $c_0+qT$。所以前者产生对数加线性项，后者产生线性加二次项。给定 $v_2$ 后，末温方程左侧随 $T_2$ 严格增加，数值求解时只需在模型允许且热容为正的温区内寻找夹住根的两端。单调性保证至多一个解；是否存在适用范围内的解，还要由初末体积和该温区一起判断。

\Diagram{fig:vdwvariablea}{求末温是单调方程的求根问题}{
 \draw[axis](0,0)--(4.8,0)node[right]{$T_2$};\draw[axis](0,-1)--(0,2.4)node[above]{$\Delta S$};\draw[curve](.8,-.7)to[out=30,in=215](4,2);\node[below]at(2,0){根};\node at(3.4,-.6){斜率 $C_V/T_2>0$};
}
内能温度积分为$\int_{T_1}^{T_2}(c_0+qT)\dd T=[c_0T+qT^2/2]_{T_1}^{T_2}$，体积积分为$\int_{v_1}^{v_2}av^{-2}\dd v=[-a/v]_{v_1}^{v_2}$。按图\ref{fig:vdwvariablea}确定末温，绝热功取内能差的负值，得到
\[
 W=c_0(T_1-T_2)+\frac q2(T_1^2-T_2^2)-a(1/v_1-1/v_2).
\]
热容随温度变化时，末温由整个温区的熵积分决定；只取末态热容代入常热容公式会遗漏这部分变化。求出末温后，再由内能差计算功，两次积分各有明确的物理含义。
\end{solution}

\begin{example}[节流反转曲线的完整参数形式]
对一摩尔 van der Waals 气体，利用第二讲节流系数的精确表达，求反转曲线的 $T_i(v)$ 与 $p_i(v)$。讨论正压反转曲线的体积范围和最大压强。
\Diagram{fig:inversionq}{节流温变的符号由反转线分隔}{
 \draw[axis](0,0)--(5,0)node[right]{$T$};\draw[axis](0,0)--(0,3)node[above]{$p$};\draw[curve](.8,0)to[out=80,in=160](2.8,2.4)to[out=-20,in=110](4.4,0);\node at(2.7,1){冷却区示意};
}
\end{example}
\begin{solution}
图\ref{fig:inversionq}的边界满足 $Tv_T|_p=v$。由 $v_T|_p=[R/(v-b)]/[RT/(v-b)^2-2a/v^3]$，交叉相乘得
\[
 \frac{RT}{v-b}=\frac{vRT}{(v-b)^2}-\frac{2a}{v^2}
 \quad\Rightarrow\quad
 T_i(v)=\frac{2a(v-b)^2}{Rb v^2}.
\]
代入状态方程，
\[
 p_i(v)=\frac{2a(v-b)}{bv^2}-\frac a{v^2}
 =\frac{a(2v-3b)}{bv^2}.
\]
对于给定正压，体积参数可由二次方程 $p_ibv^2-2av+3ab=0$ 解出：
\[
 v_\pm=\frac{a\pm\sqrt{a^2-3ab^2p_i}}{bp_i}.
\]
因此低于最大反转压强时通常存在两个参数分支，在最大压强处合并。每个体积再代入 $T_i(v)$ 才得到相应反转温度，不能把单一稀薄反转温度外推到整条曲线。两分支的具体工况还须结合相态判断；机械局部稳定与避开两相共存区是不同的条件。

\Diagram{fig:inversiona}{反转压强作为摩尔体积的函数}{
 \draw[axis](0,0)--(5.3,0)node[right]{$v/b$};\draw[axis](0,0)--(0,2.7)node[above]{$p_i b^2/a$};\draw[curve]plot[domain=1.5:6,samples=70]({.75*\x},{6*(2*\x-3)/(\x*\x)});\draw[guide](2.25,0)--(2.25,2)--(0,2);\node[left]at(0,2){$1/3$};\node[below]at(2.25,0){$3$};
}
正压要求$v>3b/2$。将图\ref{fig:inversiona}的函数写为$p_i=(a/b)(2v^{-1}-3bv^{-2})$，逐项求导，$p_i'=(a/b)(-2v^{-2}+6bv^{-3})=(2a/b)(3b-v)/v^3$。导数在$v=3b$由正变负，故$p_{i,\max}=a/(3b^2)$，对应$T_i=8a/(9Rb)$。当$v\to\infty$，反转温度趋于$2a/(Rb)$，正是稀薄近似的结果；低体积端$v\to3b/2$则趋于$2a/(9Rb)$。

这些参数式属于 van der Waals 单相模型的数学反转线。实际运行还要排除两相区及不稳定状态，图中的几何边界本身不保证每一点都可作为稳定节流工况。
\end{solution}

\begin{example}[由低压体积展开求焓与热容]
一摩尔气体在低压区满足 $v(T,p)=RT/p+B(T)$，忽略更高阶压强项。已知零压极限摩尔焓为 $h_0(T)$。求 $h(T,p)$ 和 $C_{p,m}(T,p)$，并用 $B=b-a/(RT)$ 验证稀薄 van der Waals 节流结果。
\Diagram{fig:virialq}{低压体积相对于理想项的偏移}{
 \draw[axis](0,0)--(4.7,0)node[right]{$1/p$};\draw[axis](0,0)--(0,3)node[above]{$v$};\draw[curve](.3,.8)--(4,2.6);\draw[guide](0,.2)--(4,2.15);\node at(2,2.9){竖直偏差 $B(T)$};
}
\end{example}
\begin{solution}
图\ref{fig:virialq}的定温偏移为 $B(T)$。计算 $v_T|_p=R/p+B'(T)$，焓的等温压强导数为
\[
 h_p|_T=v-Tv_T=B-TB'.
\]
它在保留的近似阶数内与压强无关，于是从零压参考积分，$h-h_0(T)=\int_0^p[B(T)-TB'(T)]\dd p'=[p'(B-TB')]_0^p=p(B-TB')$。再对温度求导，
\[
 C_{p,m}=h_0'(T)-pTB''(T).
\]
对焓修正项求温度导数时，要同时对 $T$ 与 $B'(T)$ 求导：$\dd(TB')/\dd T=B'+TB''$，它与前一个 $B'$ 抵消，才留下 $-TB''$。零压焓 $h_0(T)$ 的导数给零压定压热容，参考内能零点只给焓增加一个常数。这里的体积式是低压展开，故所得焓保留到压强的一阶项；若要继续提高热容的压强精度，需同时提高状态方程的展开阶数。

\Diagram{fig:viriala}{同一修正函数的不同导数组合对应不同性质}{
 \node[draw,align=center]at(0,0){焓修正\\$p(B-TB')$};\node[draw,align=center]at(4.5,0){热容修正\\$-pTB''$};\node at(2.2,-1){节流分子 $Tv_T-v=TB'-B$};
}
取 $B=b-a/(RT)$，有 $B'=a/(RT^2)$，故 $TB'-B=2a/(RT)-b$，与稀薄节流公式一致。图\ref{fig:viriala}列出的各量来自不同的偏导：焓的压强修正含 $B-TB'$，再沿定压方向对温度求导，才得到热容修正中的 $-TB''$。
\end{solution}

\begin{example}[再生器怎样改变循环效率]
$\nu$ 摩尔常热容理想气体经历 Stirling 循环：在 $T_h$ 等温膨胀 $V_a\to V_b$，等容冷却到 $T_c$，低温等温压缩回 $V_a$，再等容加热到 $T_h$。比较没有再生器与理想再生器时的外界吸热和效率。体积比 $r=V_b/V_a>1$。
\Diagram{fig:stirlingq}{两等温段和两等容段组成的循环}{
 \draw[axis](0,0)--(4.7,0)node[right]{$V$};\draw[axis](0,0)--(0,3.5)node[above]{$p$};\draw[curve]plot[domain=1:3.5](\x,{3/\x});\draw[curve]plot[domain=1:3.5](\x,{1.5/\x});\draw[curve](1,1.5)--(1,3);\draw[curve](3.5,.429)--(3.5,.857);\node[below]at(1,0){$V_a$};\node[below]at(3.5,0){$V_b$};\node at(3,2.7){$T_h$};\node at(2.5,.25){$T_c$};
}
\end{example}
\begin{solution}
图\ref{fig:stirlingq}两条等容段不做功。沿等温线积分$p=\nu RT/V$，高温段的功为$W_h=\nu RT_h[\ln V]_{V_a}^{V_b}=\nu RT_h\ln r$，低温段为$W_c=\nu RT_c[\ln V]_{V_b}^{V_a}=-\nu RT_c\ln r$。相加得$W=\nu R(T_h-T_c)\ln r$。没有再生器且等容加热由外部供热时，外界总吸热
\[
 Q_h^{(0)}=\nu RT_h\ln r+C_V(T_h-T_c),
 \quad \eta_0=\frac{\nu R(T_h-T_c)\ln r}{\nu RT_h\ln r+C_V(T_h-T_c)}.
\]
两条等容段的热量分别为 $-C_V(T_h-T_c)$ 与 $+C_V(T_h-T_c)$，沿整个工质循环相加为零，但分别位于放热、吸热阶段。计算从外界取得的总正热量时，不能仅因其代数和为零就把它们事先删去。理想再生器的作用是把前一段释放的热在相应温度处储存，再交还给后一段，使这两段的热量转为装置内部交换。循环功不受再生器改变，效率差别来自外界吸热量的改变。

\Diagram{fig:stirlinga}{理想再生器在两个等容段之间暂存并返还热量}{
 \node[draw]at(0,1.3){等容冷却};\node[draw]at(0,-1.3){等容加热};\node[draw]at(4,0){再生器};\draw[axis](1,1.3)--(3.2,.3)node[midway,above]{$C_V(T_h-T_c)$};\draw[axis](3.2,-.3)--(1,-1.3);
}
图\ref{fig:stirlinga}的理想再生器沿程具有连续温度分布，各微元以无穷小温差存取热量，循环后能量和熵都恢复原值。两等容段于是改为装置内部换热，高温热源只供给$Q_h=\nu RT_h\ln r$，效率为$\eta=1-T_c/T_h$。一块均温金属通常达不到这种再生条件，因为它与各温度的气体接触时仍有有限温差。
\end{solution}

\begin{example}[零粒子流时的能量流与熵产生]
两个单原子气体大容器以分子流小孔相通，温度 $T_1>T_2$ 恒定，数密度满足零净粒子通量条件。以每单位孔面积的单向粒子通量 $j$ 表示，求净能流和两热库的总熵产生率。
\Diagram{fig:effentropyq}{粒子数双向抵消而能量不抵消}{
 \node[draw]at(0,0){$T_1,n_1$};\node[draw]at(5,0){$T_2,n_2$};\draw[axis](1.2,.4)--(3.8,.4)node[midway,above]{$j$};\draw[axis](3.8,-.4)--(1.2,-.4)node[midway,below]{$j$};
}
\end{example}
\begin{solution}
图\ref{fig:effentropyq}两边单向粒子通量相同，要求 $n_1\sqrt{T_1}=n_2\sqrt{T_2}$，并非要求两压强相同。每个从侧1逸出的粒子平均带走 $2\kb T_1$，侧2则带走 $2\kb T_2$。净能流为
\[
 J_E=2\kb j(T_1-T_2).
\]
两容器维持恒温、恒粒子数，稳态时其自身宏观熵不继续积累。热侧恒温装置需补回流出的净能量，冷侧恒温装置需取走同量能量；图中计算的正熵率属于两个维持温度的热库。由于两侧各自的净粒子流为零，状态微分 $\dd S=(\dd U-\mu\dd N)/T$ 中的净 $\mu\dd N$ 项消失。若只说“有粒子双向穿孔”就额外加入两次独立的净粒子熵项，会把流入流出抵消的部分重复计算。

\Diagram{fig:effentropya}{恒温装置从热库1向热库2转移净能量}{
 \node[draw]at(0,0){热库1熵率 $-J_E/T_1$};\node[draw]at(5,0){热库2熵率 $+J_E/T_2$};\node at(2.5,-1){净粒子交换为零};
}
按图\ref{fig:effentropya}相加两个热库的熵变化率，得到
\[
 \dot S_{\rm gen}/A_h=J_E\left(\frac1{T_2}-\frac1{T_1}\right)
 =2\kb j\frac{(T_1-T_2)^2}{T_1T_2}>0.
\]
零净粒子流使两侧粒子数保持不变，熵产生由净能流引起的热库熵变给出。若压强偏离这一条件，粒子数也会变化，还须计入粒子携带的能量及各容器化学势对应的熵变。
\end{solution}

\section{有限换热速度下的气体压缩}
\subsection{等温与绝热之间还存在相位滞后}
一份理想气体在平衡态$(p_0,V_0,T_0)$附近受小幅体积振动，壁面与恒温环境之间的总热导为$H$，总定容热容为$C_V$。忽略气体内部温差，但保留与环境的有限换热速度。由第一定律，线性温度扰动满足
\[
 C_V\dot{\delta T}+H\delta T=-p_0\dot{\delta V}.
\]
右侧使用平衡压强，因为扰动压强与体积变化率的乘积是二阶项。令$\tau=C_V/H$，体积相对扰动为$\delta V/V_0=\Re[\eta\ee^{\mathrm i\omega t}]$。齐次温度方程给$\int\dd(\delta T)/\delta T=-\int\dd t/\tau$，故自由解为$C\ee^{-t/\tau}$。对周期驱动取特解$\delta T=\Re[\widehat T\ee^{\mathrm i\omega t}]$，代入并比较复振幅，
\[
 (H+\mathrm i\omega C_V)\widehat T=-\mathrm i\omega p_0V_0\eta.
\]
总解是该特解加齐次解，常数$C$由初温减去特解在$t=0$的值确定。稳定周期中齐次项已衰减。理想气体有$p_0V_0=(\gamma-1)C_VT_0$，其中$\gamma=C_p/C_V$；以下$\delta T,\delta p$表示周期扰动的复振幅，因此
\[
 \frac{\delta T}{T_0}=-\frac{(\gamma-1)\mathrm i\omega\tau}{1+\mathrm i\omega\tau}\eta,
 \qquad \frac{\delta p}{p_0}=\frac{\delta T}{T_0}-\eta.
\]
把压强写为$\delta p=-K(\omega)\eta$，得到动态体积模量
\[
 K(\omega)=p_0\frac{1+\mathrm i\gamma\omega\tau}{1+\mathrm i\omega\tau}.
\]
以$x=\omega\tau$表示频率参数，分子分母同乘$1-\mathrm ix$，有$K=p_0[1+\gamma x^2+\mathrm i(\gamma-1)x]/(1+x^2)$，故其实部和虚部分别是
\[
 K'=p_0\left[1+\frac{(\gamma-1)x^2}{1+x^2}\right],\quad
 K''=p_0\frac{(\gamma-1)x}{1+x^2}.
\]
低频极限$K'\to p_0$对应等温压缩，高频极限$K'\to\gamma p_0$对应绝热压缩。中间频段还存在由$K''$描述的相位滞后，因而需要复响应系数；只改变实热容的数值会漏掉耗散。

\begin{example}[气体的压强—体积滞回与耗散]
按本节模型，体积作$V(t)=V_0[1+\eta_0\cos\omega t]$，$\eta_0\ll1$。求稳定周期中外界对气体所做的净功，以及频率参数使耗散最大的位置。说明气体每周期内能改变量为何为零。
\ExpansionDiagram{fig:dep-th2-volume-q}{气体由活塞周期压缩，并通过有限热导向恒温环境换热}{
\draw(0,2)--(0,0)--(3,0);\draw(0,2)--(3,2);\draw[very thick](3,0)--(3,2);\draw[<->](2.5,2.4)--(3.5,2.4);\node at(1.5,1){气体};\draw[->](1.5,0)--(1.5,-1)node[below]{$H,T_0$};\node[right]at(4.5,1){$V_0+\delta V$};}
\end{example}
\begin{solution}
图\ref{fig:dep-th2-volume-q}的稳态温度已经排除起始自由弛豫。压强实扰动为
\[
 \delta p=-\eta_0K'\cos\omega t+\eta_0K''\sin\omega t.
\]
外界净功是$W_{\rm on}=-\oint p\dd V$；常数压强项闭合积分为零，所以
\begin{align*}
 W_{\rm on}&=\int_0^{2\pi/\omega}
 \delta p\,V_0\eta_0\omega\sin\omega t\dd t\\
 &=\pi V_0\eta_0^2K''
 =\pi p_0V_0\eta_0^2(\gamma-1)\frac{x}{1+x^2}.
\end{align*}
这里令$\phi=\omega t$，$\dd\phi=\omega\dd t$，上下限为$0,2\pi$。交叉项的积分$\int_0^{2\pi}\sin\phi\cos\phi\dd\phi=[\sin^2\phi/2]_0^{2\pi}=0$；耗散项的积分$\int_0^{2\pi}\sin^2\phi\dd\phi=[\phi/2-\sin(2\phi)/4]_0^{2\pi}=\pi$，分别给出上述两项的系数。
\ExpansionDiagram{fig:dep-th2-volume-a}{压强—体积椭圆的面积对应周期耗散}{
\draw[->](0,0)--(5,0)node[right]{$V$};\draw[->](0,0)--(0,3.5)node[above]{$p$};\draw[thick,domain=0:360,samples=120]plot({2.5+1.5*cos(\x)},{1.8-.8*cos(\x)+.6*sin(\x)});\draw[->,domain=20:40,samples=12]plot({2.5+1.5*cos(\x)},{1.8-.8*cos(\x)+.6*sin(\x)});\node[below]at(2.5,0){$V_0$};\node[right]at(4,2.8){逆时针};}
图\ref{fig:dep-th2-volume-a}按气体坐标画逆时针环路，故$-\oint p\dd V>0$。用商法则，$\dd[x/(1+x^2)]/\dd x=[(1+x^2)-2x^2]/(1+x^2)^2=(1-x^2)/(1+x^2)^2$，导数在$x=1$由正变负，所以每周期耗散在该点最大。这里固定的是体积振幅，若固定驱动力振幅，响应振幅会随频率变化，极值条件需另求。

周期末温度回到同值，理想气体内能也回到同值，$\Delta U=0$。由第一定律，输入功全部变为向环境排出的净热量。环境在恒温$T_0$下的熵增加为$W_{\rm on}/T_0$，与正的耗散相符。低频时换热足够快，高频时几乎来不及换热，两端的每周期耗散均趋零。
\end{solution}

\begin{example}[热弛豫造成的活塞阻尼]
气体推动质量$M$、面积$A$的活塞，活塞还受劲度$k$的弹簧约束，外部常压在平衡点抵消$p_0$。活塞向外位移$x$使$\delta V=Ax$。在气体热修正相对总弹性很小的条件下，求相对等温固有频率的一阶频移和振幅衰减率。
\ExpansionDiagram{fig:dep-th2-piston-q}{弹簧与气体共同提供恢复力，气体还引入热耗散}{
\draw(0,0)rectangle(3,2);\draw[very thick](3,0)--(3,2);\draw(3,1)--(3.2,1)--(3.35,1.2)--(3.55,.8)--(3.75,1.2)--(3.95,.8)--(4.15,1.2)--(4.35,.8)--(4.55,1.2)--(4.75,1)--(5,1);\node[above]at(4,1.3){$k$};\draw[thick](5,0)--(5,2);\node at(1.4,1){$V_0,T_0$};\node[above]at(3,2){$M$};\draw[->](3,-.5)--(4,-.5)node[right]{$x$};}
\end{example}
\begin{solution}
图\ref{fig:dep-th2-piston-q}中压强扰动产生力$A\delta p=-A^2K(\omega)x/V_0$，所以频率方程为
\[
 -M\omega^2+k+\frac{A^2}{V_0}K(\omega)=0.
\]
等温基准频率$\omega_0^2=(k+A^2p_0/V_0)/M$。令$\omega=\omega_0+\delta\omega$、$\delta\omega=\Delta\omega+\mathrm i\Gamma$，由于采用$\ee^{\mathrm i\omega t}$，正$\Gamma$表示衰减。将频率平方展开为$\omega_0^2+2\omega_0\delta\omega$，略去二阶小量；在一阶热修正$K-p_0$内部取$\omega=\omega_0$，得到
\[
 -2M\omega_0\delta\omega+
 \frac{A^2}{V_0}[K(\omega_0)-p_0]=0.
\]
等温零阶项已由$\omega_0$的定义抵消。比较剩余式的实部与虚部，
\[
 \Delta\omega=\frac{A^2[K'(\omega_0)-p_0]}{2MV_0\omega_0},
 \qquad\Gamma=\frac{A^2K''(\omega_0)}{2MV_0\omega_0}.
\]
\ExpansionDiagram{fig:dep-th2-piston-a}{振幅随热耗散衰减，振动频率则由实部弹性决定}{
\draw[->](0,0)--(6,0)node[right]{$t$};\draw[->](0,-1.5)--(0,1.8)node[above]{$x$};\draw[dashed,domain=0:5.5]plot(\x,{1.3*exp(-.2*\x)});\draw[dashed,domain=0:5.5]plot(\x,{-1.3*exp(-.2*\x)});\draw[thick,domain=0:5.5,samples=150]plot(\x,{1.3*exp(-.2*\x)*cos(360*\x)});}
图\ref{fig:dep-th2-piston-a}示意$\ee^{-\Gamma t}$包络。令$x_0=\omega_0\tau$，显式结果为
\[
 \Delta\omega=\frac{A^2p_0(\gamma-1)}{2MV_0\omega_0}
 \frac{x_0^2}{1+x_0^2},\qquad
 \Gamma=\frac{A^2p_0(\gamma-1)}{2MV_0\omega_0}
 \frac{x_0}{1+x_0^2}.
\]
使用这些一阶式需要$A^2p_0(\gamma-1)/(MV_0\omega_0^2)\ll1$；若不满足，应回到热方程与活塞方程联立求完整特征根。仅把压缩指数选为一个实数，会遗漏$\Gamma$，也无法解释振幅逐渐减小。
\end{solution}

\section{习题与解答}
\begin{exercise}[恒温泄流所需热量]
刚性容器内有自由度为 $f$ 的理想气体，经小孔向真空泄流；外界供热使温度保持 $T_0$。求从 $N_0$ 减至 $N_1$ 所需的总热量及 $N(t)$。
\Diagram{fig:isothermalleakq}{恒温装置补热，气体经另一侧小孔泄流}{
 \draw(0,0)rectangle(3,1.7);\fill[white](2.94,.64)rectangle(3.06,.96);
 \draw[axis](-1,.8)--(0,.8)node[midway,above]{$Q$};\draw[axis](3,.8)--(4.6,.8)node[midway,above]{粒子};\node at(1.5,.8){$T_0$ 恒定};
}
\end{exercise}
\begin{solution}
图\ref{fig:isothermalleakq}中余气 $U=fN\kb T_0/2$。开放系统能量式 $\dd U=\delta Q+\varepsilon_e\dd N$ 给出
\[
 \delta Q=\frac f2\kb T_0\dd N-\frac{f+1}2\kb T_0\dd N
 =-\frac12\kb T_0\dd N.
\]
有限变化的热量由 $\int_{N_0}^{N_1}(-\kb T_0/2)\dd N$ 得到，积分下限大于上限时，负号使最终供热为正。每减少一个粒子，余气内能按体内平均能量下降，而流出粒子所带能量更高，外界只补足这个差额。即使有内部转动自由度，它们在“体内减少”和“粒子带走”两项中以相同平均值出现，所以相减后只留下半个平动热能尺度。

\Diagram{fig:isothermalleaka}{每逸出一个粒子所需补热}{
 \node[draw]at(0,0){余气减少 $f\kb T_0/2$};\node[draw]at(4.6,0){逸出带走 $(f+1)\kb T_0/2$};\node at(2.3,-.9){差额由外界供给 $\kb T_0/2$};
}
按图\ref{fig:isothermalleaka}，总热量$Q=(N_0-N_1)\kb T_0/2$，与$f$无关。通量给$\dot N=-\kappa N$，其中$\kappa=(A_h/V)\sqrt{\kb T_0/(2\pi m)}$为常量。分离变量，$\int_{N_0}^{N(t)}\dd N'/N'=-\kappa\int_0^t\dd s$，得到$\ln[N(t)/N_0]=-\kappa t$，故$N=N_0\ee^{-\kappa t}$。保持恒温仍需补热；内能因粒子减少而下降，但不足以支付全部逸出能量。
\end{solution}

\begin{exercise}[辐射腔冷却的两种热容]
体积 $V$ 的腔经小孔 $A_h$ 辐射至冷真空。分别讨论：仅计辐射场能量 $U=aVT^4$；以及腔壁热容为常量 $C_w$ 且主导总能量、腔内辐射始终与壁平衡。求两种近似的 $T(t)$。
\Diagram{fig:radcoolq}{腔壁与辐射场同温，经孔口向外辐射}{
 \draw[line width=2pt](0,0)rectangle(3.4,1.7);\fill[white](3.32,.69)rectangle(3.48,1.01);
 \node at(1.7,.85){$U_r=aVT^4$};\node at(1.7,2.2){壁热容 $C_w$};\draw[axis](3.4,.85)--(5,.85);
}
\end{exercise}
\begin{solution}
净外流功率为$A_hcaT^4/4$。若图\ref{fig:radcoolq}中只计辐射能量，能量方程为
\[
 4aVT^3\dot T=-A_hcaT^4/4.
\]
约去$aT^3$，令$\beta=A_hc/(16V)$，得到$\dot T=-\beta T$。分离变量，$\int_{T_0}^T\dd\vartheta/\vartheta=-\beta\int_0^t\dd s$，故$\ln(T/T_0)=-\beta t$，即
\[
 T=T_0\ee^{-A_hct/(16V)}.
\]

若两种热容都不能忽略，应取总热容$C_w+4aVT^3$。将能量方程除以$T^4$，并从初温积分到$T$，
\[
 \int_{T_0}^{T}\left(C_w\vartheta^{-4}+\frac{4aV}{\vartheta}\right)\dd\vartheta
 =-\frac{A_hca}{4}\int_0^t\dd s.
\]
左边的原函数为$-C_w/(3\vartheta^3)+4aV\ln\vartheta$，代入两端并同时变号，得到
\[
 \frac{C_w}{3}(T^{-3}-T_0^{-3})
 +4aV\ln\frac{T_0}{T}=\frac{A_hca}{4}t.
\]
第一项来自壁热容乘 $T^{-4}$ 的积分，第二项来自辐射热容乘 $T^{-4}$ 后剩下的 $1/T$。温度低于初温时两项都为正。分别忽略其中一个热容贡献，就恢复题设两种极限的显式温度规律。

\Diagram{fig:radcoola}{由能量温度关系决定不同的冷却微分方程}{
 \node[draw,align=center]at(0,0){辐射热容 $4aVT^3$\\$\dot T\propto-T$};
 \node[draw,align=center]at(4.5,0){腔壁热容 $C_w$\\$\dot T\propto-T^4$};
}
若壁主导，忽略混合积分中的辐射热容项，有$C_w(T^{-3}-T_0^{-3})/3=A_hcat/4$。移项并取$-1/3$次幂，得到
\[
 T=T_0\left(1+\frac{3A_hcaT_0^3}{4C_w}t\right)^{-1/3}.
\]
图\ref{fig:radcoola}表明同一个辐射功率律不决定唯一的温度时间律。第一种模型需要有能维持平衡而热容可忽略的壁；真实腔体通常要把壁与辐射热容相加，写 $(C_w+4aVT^3)\dot T=-A_hcaT^4/4$ 后再积分。
\end{solution}

\begin{exercise}[由响应系数求热容差]
某稳定单相物质的体积为 $V$，等温压缩率 $\kappa_T$、热膨胀系数 $\alpha_p$ 已知。推导 $C_p-C_V$，并对理想气体验证。
\Diagram{fig:responsq}{等压升温与等容升温的体积差}{
 \draw[axis](0,0)--(4.7,0)node[right]{$T$};\draw[axis](0,0)--(0,2.8)node[above]{$V$};\draw[axis](1,.8)--(3.5,2);\draw[axis](1,.8)--(3.5,.8);\node[right]at(3.5,2){定压};\node[right]at(3.5,.8){定容};
}
\end{exercise}
\begin{solution}
图\ref{fig:responsq}的定压路径有 $\dd V=V\alpha_p\dd T$。代入式\eqref{eq:dsvt}，再乘 $T/\dd T$，$C_p=C_V+Tp_TV\alpha_p$。状态方程全微分在定容时为 $0=V\alpha_p\dd T-V\kappa_T\dd p$，故 $p_T|_V=\alpha_p/\kappa_T$。
热容差比较的是从同一平衡态出发、同样微小温升所需的可逆热量差。定容路径直接给 $C_V\dd T$；定压路径还改变体积，其熵的额外变化由 $S_V|_T\dd V$ 给出。再用 Maxwell 关系把该偏导改写为 $p_T|_V$，便能完全由可测的响应系数表达。稳定物质 $\kappa_T>0$，故即使 $\alpha_p$ 为负，热容差仍因其平方而非负；负热膨胀本身不会令 $C_p<C_V$。

\Diagram{fig:responsa}{响应系数把状态平面中的斜率相联系}{
 \node[draw]at(0,0){$p_T|_V=\alpha_p/\kappa_T$};\draw[->](2,0)--(3,0);\node[draw]at(5,0){$C_p-C_V=TV\alpha_p^2/\kappa_T$};
}
联立两项偏导关系，得到图\ref{fig:responsa}中的$C_p-C_V=TV\alpha_p^2/\kappa_T$。理想气体$\alpha_p=1/T$、$\kappa_T=1/p$，给出$C_p-C_V=pV/T=\nu R$。若$\alpha_p=0$，两热容相等；若接近临界区响应发散，不能仍使用远离临界点的常数热容近似。
\end{solution}

\begin{exercise}[实际气体的 Joule 自由膨胀]
一摩尔 van der Waals 气体在绝热刚性大容器中从体积 $v_1$ 自由膨胀至 $v_2$，原空区为真空。$C_{V,m}=c$ 恒定。求最终温度和熵增，假设初末态均为稳定单相气体。
\Diagram{fig:jouleq}{撤去隔板后的自由膨胀}{
 \draw(0,0)rectangle(5,1.7);\draw[dashed](2,0)--(2,1.7);\node at(1,.85){$v_1,T_1$};\node at(3.5,.85){真空};\draw[->](2.2,2)--(2.2,1.4);
}
\end{exercise}
\begin{solution}
图\ref{fig:jouleq}的整体边界不移动且绝热，$Q=W=0$，所以 $U_2=U_1$。摩尔内能 $U_m=cT-a/v+\text{常量}$，故
\[
 c(T_2-T_1)+a\left(\frac1{v_1}-\frac1{v_2}\right)=0,
 \qquad T_2=T_1-\frac ac\left(\frac1{v_1}-\frac1{v_2}\right).
\]
熵增的符号还可在状态空间中直接判断。把等内能平衡态连接起来，基本关系给
\[
 0=T\dd S-p\dd V,
 \qquad \left(\frac{\partial S}{\partial V}\right)_U=\frac pT>0.
\]
这条可用于求状态差的等内能曲线，不是实际自由膨胀的准静态路径。只要它处于同一稳定单相、正温正压范围，体积增大就使熵增加。由能量式求得的末温也须先确认在此范围内，再代入熵公式。

\Diagram{fig:joulea}{实际不可逆路径与用于求熵的可逆替代路径}{
 \draw[axis](0,0)--(4.5,0)node[right]{$v$};\draw[axis](0,0)--(0,2.7)node[above]{$T$};
 \draw[axis](1,2)--(3.4,2);\draw[axis](3.4,2)--(3.4,1);\draw[guide](1,2)--(3.4,1);\node[right]at(3.4,1){$2$};\node[left]at(1,2){$1$};
}
利用图\ref{fig:joulea}的替代路径或直接用状态函数，
\[
 \Delta S=c\ln\frac{T_2}{T_1}+R\ln\frac{v_2-b}{v_1-b}.
\]
在满足模型及稳定性的可达状态间，该值非负。温度降低是吸引势能增加的结果；$a=0$ 时温度不变，但熵仍随可用体积增大。若计算给出负温度或跨入相变区，说明单相常热容模型已被用出范围。
\end{solution}

\begin{exercise}[节流系数与反转温度]
稳定流经绝热多孔塞，忽略高度和宏观动能变化。先说明为何焓不变，再对一摩尔 van der Waals 气体求 $\mu_{\rm JT}=(\partial T/\partial p)_H$，给出稀薄极限的反转温度。
\Diagram{fig:jtq}{多孔塞两侧的流动功}{
 \draw(0,0)rectangle(6,1.3);\fill[pattern=north east lines](2.8,0)rectangle(3.2,1.3);\draw[axis](.3,.65)--(1.3,.65);\draw[axis](4.6,.65)--(5.6,.65);\node at(1.5,1.8){$p_1,v_1$};\node at(4.6,1.8){$p_2,v_2$};
}
\end{exercise}
\begin{solution}
取图\ref{fig:jtq}中随一摩尔气体穿塞的能量，入口推动功为 $p_1v_1$，出口排开气体需 $p_2v_2$。绝热稳流能量守恒给 $U_{m,1}+p_1v_1=U_{m,2}+p_2v_2$，故 $H_m$ 不变。由式\eqref{eq:dhpt}，
\[
 \mu_{\rm JT}=\frac{T(\partial v/\partial T)_p-v}{C_{p,m}}.
\]
van der Waals 方程的定压微分满足
\[
 0=\frac R{v-b}\dd T+
 \left[-\frac{RT}{(v-b)^2}+\frac{2a}{v^3}\right]\dd v.
\]
因此$v_T|_p=[R/(v-b)]/[RT/(v-b)^2-2a/v^3]$。令$D_v=RT/(v-b)^2-2a/v^3>0$表示机械稳定分母，则
\[
 Tv_T-v=\frac{RT/(v-b)-vD_v}{D_v}
 =\frac{2a/v^2-RTb/(v-b)^2}{D_v},
\]
从而精确的单相结果为$\mu_{\rm JT}=[2a/v^2-RTb/(v-b)^2]/(C_{p,m}D_v)$。
低压近似体积的来源可从原方程逐项展开。对 $b/v\ll1$，$RT/(v-b)\simeq RT/v+RTb/v^2$，故 $pv\simeq RT+(RTb-a)/v$。在修正项的分母用首项 $v\simeq RT/p$ 代入，再除以 $p$，得到所用的 $v\simeq RT/p+b-a/(RT)$。这样保留的体积修正对压强是零阶，而在 $pv$ 中对应一阶低压修正，其使用范围与稀薄假设一致。

\Diagram{fig:jta}{稀薄节流时吸引与排除体积效应的竞争}{
 \draw[axis](0,0)--(5,0)node[right]{$T$};\draw[axis](0,-1)--(0,2)node[above]{$\mu_{\rm JT}$};\draw[curve]plot[domain=.9:4.5](\x,{2/\x-.8});\node[below]at(2.5,0){$T_i$};
}
稀薄近似 $v\simeq RT/p+b-a/(RT)$，所以 $Tv_T-v\simeq2a/(RT)-b$。于是
\[
 \mu_{\rm JT}\simeq\frac{2a/(RT)-b}{C_{p,m}},\qquad T_i\simeq\frac{2a}{Rb}.
\]
图\ref{fig:jta}中 $\mu_{\rm JT}>0$ 时降压 $\dd p<0$ 才导致冷却 $\dd T<0$。理想气体的焓仅依赖温度，节流不改变温度。
\end{solution}

\begin{exercise}[温度线性状态方程的内能体积项]
某物质在研究范围内满足 $p=p_*(V)+Tq(V)$。证明其内能可写为 $U=\Phi(T)-\int p_*(V)\dd V$；指出 $C_V$ 可以怎样依赖温度和体积。
\Diagram{fig:linearstateq}{在同一体积上压强随温度线性变化}{
 \draw[axis](0,0)--(4.5,0)node[right]{$T$};\draw[axis](0,0)--(0,2.7)node[above]{$p$};\draw[curve](0,.6)--(4,2.2);\node[left]at(0,.6){$p_*$};\node at(3,.7){斜率 $q(V)$};
}
\end{exercise}
\begin{solution}
在图\ref{fig:linearstateq}对应的固定体积上，$p_T=q(V)$。故 $Tp_T-p=-p_*(V)$，沿等温体积方向积分得
\[
 U(V,T)=\Phi(T)-\int^V p_*(w)\dd w.
\]
把得到的内能对温度求导时，体积积分在定容条件下不变，因而热容只有 $\Phi'(T)$。另一方面，$p$ 对温度二阶导为零，相容条件也直接要求 $C_V$ 的体积偏导为零，两种判断一致。若 $p_*(V)$ 为吸引造成的负压强项，负号再乘体积积分后可产生随膨胀上升的势能；若 $p_*$ 为零，则内能只含温度函数，理想气体正属于这一类。

\Diagram{fig:linearstatea}{内能中温度项与体积项分离}{
 \node[draw]at(0,0){$\Phi(T)$};\node at(2,0){$+$};\node[draw]at(4,0){$-\int^V p_*(w)\dd w$};\node at(2,-.8){$C_V=\Phi'(T)$，不依赖 $V$};
}
按图\ref{fig:linearstatea}求温度偏导，$C_V$ 只能依赖 $T$。熵则满足 $\dd S=\Phi'(T)\dd T/T+q(V)\dd V$。状态方程中与温度成正比的项不直接贡献内能体积项，却通过 $q(V)$ 贡献熵的体积变化。题设若再给某个随体积变化的定容热容，便与此线性状态方程不相容。
\end{solution}

\begin{exercise}[光子气体的等温膨胀]
体积由 $V_1$ 变为 $V_2$ 的光子气体与温度 $T$ 的热源接触，过程可逆。求做功、内能变化、吸热和熵增，并比较 $Q$ 与 $W$。
\Diagram{fig:photonisoq}{恒温辐射腔的可动反射壁}{
 \draw(0,0)rectangle(3.2,1.5);\draw[line width=2pt](3.2,0)--(3.2,1.5);\draw[axis](3.2,.75)--(4.5,.75);\node at(1.6,.75){$T$ 恒定};\node at(1.6,-.5){$V_1\to V_2$};
}
\end{exercise}
\begin{solution}
图\ref{fig:photonisoq}中$p=aT^4/3$不随体积变化，所以$W=\int_{V_1}^{V_2}p\dd V=p[V]_{V_1}^{V_2}=p(V_2-V_1)$。光子总数可随腔壁发射而改变，即使温度不变，能量仍随体积增加：$\Delta U=aT^4(V_2-V_1)$。
等温不变的是单位体积辐射能 $u=aT^4$，增加的体积中还要建立同样的辐射场，因此产生非零内能增量。反射壁每移动一小段，所做功为 $(u/3)\dd V$，新增辐射能为 $u\dd V$，两者比例固定。恒温热源需同时承担这两部分；增加的光子由平衡壁吸收、发射调节，其数目不受固定粒子数约束。图中的能流分配据此适用于整个可逆等温过程。

\Diagram{fig:photonisoa}{热量同时增加辐射能与支付膨胀功}{
 \node[draw]at(0,0){$Q$};\draw[->](.5,.1)--(2,1);\draw[->](.5,-.1)--(2,-1);\node[right]at(2,1){$\Delta U=3W$};\node[right]at(2,-1){$W$};
}
由第一定律 $Q=\Delta U+W=4aT^4(V_2-V_1)/3=4W$。图\ref{fig:photonisoa}中温度保持不变，辐射能仍随体积增加；其余吸热用于膨胀功。熵增 $\Delta S=Q/T=4aT^3(V_2-V_1)/3$，与状态函数表达式一致。
\end{solution}

\begingroup
\setlength{\intextsep}{3pt}\setlength{\abovecaptionskip}{4pt}
\begin{exercise}[通过熵函数反推状态方程]
某固定粒子数系统给出的熵差为 $S-S_0=C\ln(T/T_0)+\lambda\ln[(V-B)/(V_0-B)]$，其中 $C,\lambda,B$ 为常量。求与它相容的状态方程的一般形式，并说明为何还不能唯一求出压强。
\Diagram{fig:reverseq}{熵函数的两个独立方向导数}{
 \begin{scope}[yscale=.8]
 \draw[axis](0,0)--(4.4,0)node[right]{$V$};\draw[axis](0,0)--(0,2.6)node[above]{$T$};\draw[axis](1,.8)--(3,.8)node[right]{$S_V|_T$};\draw[axis](1,.8)--(1,2.1)node[above]{$S_T|_V$};
 \end{scope}
}
\end{exercise}
\begin{solution}
图\ref{fig:reverseq}中定容方向给 $C_V=TS_T|_V=C$。Maxwell 关系在等温方向给 $p_T|_V=S_V|_T=\lambda/(V-B)$。对温度积分，
\[
 p(V,T)=\frac{\lambda T}{V-B}+\psi(V).
\]
若进一步给出参考温度 $T_*$ 上的一条压强曲线 $p_*(V)$，则待定函数由 $\psi(V)=p_*(V)-\lambda T_* /(V-B)$ 唯一确定。随后内能的体积积分也随之确定，只剩一个整体参考常数。这个步骤区分了两类信息：熵给出温度变化时的压强斜率，参考等温线给出固定温度时的压强截距。两类信息合用，才足以恢复完整状态方程。

\Diagram{fig:reversea}{熵只确定压强的温度斜率}{
 \begin{scope}[yscale=.65]
 \draw[axis](0,0)--(4,0)node[right]{$T$};\draw[axis](0,0)--(0,2.8)node[above]{$p$};\draw[curve](0,.4)--(3.4,1.9);\draw[guide](0,1)--(3.4,2.5);\node at(3,.5){相同斜率，不同截距};
 \end{scope}
}
图\ref{fig:reversea}的截距 $\psi(V)$ 仍为待定函数，内能相应为 $U=CT-\int^V\psi(w)\dd w+\text{常量}$。取 $\lambda=\nu R$、$\psi=-A/V^2$ 可还原 van der Waals 模型；其他选择给出不同的物质模型。要确定 $\psi(V)$，还需一条参考等温线的压强数据，或等价的能量信息。
\end{solution}
\endgroup

\chapter{第二定律、熵与有限热源}

\section{热机、卡诺循环与效率上界}
\label{th:carnot}
热机以某种工质经历循环，从高温热源吸热 $Q_h>0$，向低温热源放热 $Q_c>0$，对外输出功 $W$。工质回到初态，$\Delta U=0$，第一定律给出
\[
 W=Q_h-Q_c,\qquad \eta=\frac W{Q_h}=1-\frac{Q_c}{Q_h}.
\]
一次循环后工质回到初态，热源和外界则分别发生热量与功的交换。效率按输出功占高温热源供热量的比例定义，分母为 $Q_h$；净吸热 $Q_h-Q_c$ 已等于循环输出功。

卡诺循环由两条可逆等温线和两条可逆绝热线组成。理想气体在高温$T_h$上由$A$膨胀到$B$，再绝热到$C$；在低温$T_c$上由$C$压缩至$D$，再绝热回$A$。等温段内能不变，由$p=\nu RT/V$得$Q=\nu RT\int_{V_i}^{V_f}\dd V/V=\nu RT[\ln V]_{V_i}^{V_f}$。低温段沿压缩方向热量为负，以下取其大小，得到
\[
 Q_h=\nu RT_h\ln\frac{V_B}{V_A},\qquad
 Q_c=\nu RT_c\ln\frac{V_C}{V_D}.
\]
两条绝热线给 $T_hV_B^{\gamma-1}=T_cV_C^{\gamma-1}$、$T_hV_A^{\gamma-1}=T_cV_D^{\gamma-1}$。两式相除，得到 $V_B/V_A=V_C/V_D$，从而
\begin{equation}\label{eq:carnotratio}
 \frac{Q_h}{T_h}=\frac{Q_c}{T_c},\qquad \eta_C=1-\frac{T_c}{T_h}.
\end{equation}
\Diagram{fig:carnottheory}{卡诺循环的两条等温线与两条绝热线}{
 \draw[axis](0,0)--(5.4,0)node[right]{$V$};\draw[axis](0,0)--(0,4.4)node[above]{$p$};
 \pgfmathsetmacro{\vratio}{2^(3/2)}
 \coordinate(A)at(.8,3.75);\coordinate(B)at(1.5,2);\coordinate(C)at({1.5*\vratio},{1.5/(1.5*\vratio)});\coordinate(D)at({.8*\vratio},{1.5/(.8*\vratio)});
 \draw[curve]plot[domain=.8:1.5](\x,{3/\x});\draw[curve]plot[domain=1.5:{1.5*\vratio}](\x,{3*1.5^(2/3)/\x^(5/3)});
 \draw[curve]plot[domain={.8*\vratio}:{1.5*\vratio}](\x,{1.5/\x});\draw[curve]plot[domain=.8:{.8*\vratio}](\x,{3*.8^(2/3)/\x^(5/3)});
 \foreach \p/\pos in{A/above,B/above right,C/right,D/left}{\fill(\p)circle(1.2pt)node[\pos]{$\p$};}
 \node at(2.2,3.2){$T_h$};\node at(3.8,.85){$T_c$};\draw[->](1,3)--(1.2,2.5);\draw[->](3.6,.4167)--(3,.5);
}
图\ref{fig:carnottheory}中的循环面积是 $W$。经过同一点时，绝热线比等温线更陡；高温膨胀与低温压缩之间因而围出正面积。两条等温段的热量由各自的 $p\dd V$ 积分确定。

\begin{example}[用图形伸缩理解卡诺热量比]\label{ex:carnotscale}
理想气体卡诺循环的状态如图\ref{fig:carnottheory}。讨论变换 $(p,V)\mapsto(\alpha p,\beta V)$ 如何改变温度、绝热线和微元热量，并利用它解释两条等温线的热量比。
\Diagram{fig:carnotscaleq}{状态平面中的独立压强伸缩与体积伸缩}{
 \draw[axis](0,0)--(4.3,0)node[right]{$V$};\draw[axis](0,0)--(0,3)node[above]{$p$};
 \fill(1,.9)circle(1.3pt)node[left]{$(p,V)$};\fill(3,2.2)circle(1.3pt)node[right]{$(\alpha p,\beta V)$};\draw[->](1.2,1)--(2.8,2.1);
}
\end{example}
\begin{solution}
图\ref{fig:carnotscaleq}中 $pV$ 乘 $\alpha\beta$，由状态方程知 $T\mapsto\alpha\beta T$。绝热线 $pV^\gamma=C$ 变为 $p'V'^\gamma=\alpha\beta^\gamma C$，仍是同一指数的绝热线。因此可用一次伸缩把绝热段 $AD$ 映到 $BC$；两端分别处于相同高低温，故此时 $\alpha\beta=1$，并有 $V_B/V_A=V_C/V_D$。

再以另一伸缩把高温等温段 $AB$ 映到低温等温段 $DC$，这次 $\alpha\beta=T_c/T_h$。理想气体微元热量为 $\delta Q=C_V\dd T+p\dd V$，两项在映射下都乘同一因子 $\alpha\beta$。
伸缩过程中微元的温度变化也乘 $\alpha\beta$，体积变化则乘 $\beta$，所以 $p'\dd V'=\alpha\beta p\dd V$。对固定物质的量、常热容理想气体，$C_V\dd T'$ 也按同一比例改变，两个能量项相加才保证热量一起缩放。这个几何证明使用了常数 $\gamma$ 的绝热线形式；若热容随温度改变，卡诺热量比仍成立，但应由可逆熵变计算，而不再直接使用同一幂指数的伸缩图形。

\Diagram{fig:carnotscalea}{两条等温段在伸缩下对应的正向积分}{
 \draw[axis](0,1.3)--(4,1.3)node[right]{$A\to B$，$T_h$};\draw[axis](0,0)--(4,0)node[right]{$D\to C$，$T_c$};
 \draw[->](1,1)--(1,.3);\draw[->](3,1)--(3,.3);\node at(2,-.8){$Q_{D\to C}=(T_c/T_h)Q_{A\to B}$};
}
图\ref{fig:carnotscalea}中两段均按膨胀方向积分，故都是吸热；实际循环的低温段方向为 $C\to D$，相应热量为负。取放热量大小后即得 $Q_c/Q_h=T_c/T_h$。伸缩论证保留了体积比和路径对应，不能只凭图上两条曲线“形状相似”就断言热量成比例。
\end{solution}

\subsection{第二定律的两种表述}
\label{th:secondlaw}
开尔文表述：不能制造一种循环工作的装置，其唯一效果是从单一热源吸热并把它全部转为功。克劳修斯表述：不能制造一种装置，其唯一效果是把热量从低温物体传给高温物体。这里“唯一效果”排除了工质未复原、机械储能被消耗等隐藏变化。

两种表述可以通过装置组合互相推出。若有开尔文表述所禁止的装置，就可把它从高温热源吸收的热量全部变成机械功，驱动一台普通制冷机。机械功相抵消后，组合装置只把热从冷源送至热源，违背克劳修斯表述。反过来，若存在不需功的冷到热传热装置，就可把普通热机向冷源排出的热量送回热源。组合装置最终只从高温热源取热并输出同量功，违背开尔文表述。

热量可以全部转成某一过程的功，例如理想气体可逆等温膨胀；这不构成反例，因为工质体积改变，过程不是循环。将其复原所需的代价必须纳入比较。

\subsection{卡诺定理}
\label{th:carnottheorem}
设某热机与可逆卡诺机工作在同一对热源之间，均输出相同的功 $W$。若前者效率更高，它吸收的高温热量更少。把卡诺机反向运行，恰好消耗 $W$，组合后不需外功，却将净热量从低温源转移到高温源，违反第二定律。因此任何热机都满足 $\eta\le1-T_c/T_h$。

若被比较的两台机器均可逆，就可交换它们的角色，同样得到反向不等式，故两者效率相等。于是可逆机效率只依赖两热源温度，与工质无关。不可逆性包括有限温差传热、摩擦、黏滞耗散和未经约束的膨胀；仅把活塞做得很慢，不能消除所有这些因素。

\begin{example}[用装置组合检验宣称的效率]\label{ex:enginepair}
一台热机在 $T_h=600\,\mathrm K$、$T_c=300\,\mathrm K$ 间循环，宣称每输出 $100\,\mathrm J$ 功只吸收 $160\,\mathrm J$ 热。用一台反向卡诺机与它组合，指出违反第二定律的净效果。
\Diagram{fig:enginepairq}{热机和反向可逆机的各项热量与机械交换}{
 \draw(-1,2.4)rectangle(6.3,3);\node at(2.65,2.7){热源 $600\,\mathrm K$};
 \draw(-1,-.8)rectangle(6.3,-.2);\node at(2.65,-.5){冷源 $300\,\mathrm K$};
 \node[draw,circle,minimum size=9mm](a)at(.6,1.1){机};\node[draw,circle,minimum size=9mm](b)at(4.5,1.1){逆};
 \draw[axis](.6,2.4)--(.6,1.6)node[midway,left]{$160\,\mathrm J$};\draw[axis](.6,.6)--(.6,-.2)node[midway,left]{$60\,\mathrm J$};
 \draw[axis](4.5,1.6)--(4.5,2.4)node[midway,right]{$200\,\mathrm J$};\draw[axis](4.5,-.2)--(4.5,.6)node[midway,right]{$100\,\mathrm J$};
 \draw[axis](1.1,1.1)--(4,1.1)node[midway,above]{$W=100\,\mathrm J$};
}
\end{example}
\begin{solution}
宣称热机向冷源放热 $160-100=60\,\mathrm J$。反向卡诺机的正向效率为 $1-300/600=1/2$；反向消耗 $100\,\mathrm J$ 功时，向热源放热 $200\,\mathrm J$，从冷源吸热 $100\,\mathrm J$。
沿每个热源分别作收支表，热源得到的净热为 $200-160$，冷源得到的净热为 $60-100$。两者相加为零，说明组合装置并未制造或损失能量，矛盾来自热量转移的方向及无外功这一条件。两台工质均作循环，机械输出与输入也已逐项抵消，所以没有未计入的储能变化可以承担这一冷到热转移。第二定律正是针对这种所有其他作用都恢复原状后的净效果。

\Diagram{fig:enginepaira}{机械功相抵后只剩冷到热的热量转移}{
 \node[draw]at(0,1.5){热源：净得到 $40\,\mathrm J$};\node[draw]at(0,-.5){冷源：净失去 $40\,\mathrm J$};\draw[axis](0,-.1)--(0,1.1)node[midway,right]{外功为零};
}
图\ref{fig:enginepaira}中的净效果是把 $40\,\mathrm J$ 从冷源移向热源且没有其他变化，故宣称不可能成立。数值上其效率 $100/160=0.625$ 已超过卡诺上界 $0.5$；组合计算进一步明确了矛盾发生在何处。
\end{solution}

\section{熵的定义与不可逆过程}
\label{th:entropy}
对可逆循环，把路径分成许多温差很小的卡诺子循环。相邻循环的内部路径反向重合，其热量贡献相抵；每个子循环满足 $Q_h/T_h-Q_c/T_c=0$。在分割尺度趋零时得到
\[
 \oint\frac{\delta Q_{\rm rev}}T=0.
\]
因此在任意两个平衡态 $A,B$ 之间，沿可逆路径的 $\int_A^B\delta Q_{\rm rev}/T$ 仅由端点决定。定义熵差
\begin{equation}\label{eq:entropydefinition}
 S_B-S_A=\int_A^B\frac{\delta Q_{\rm rev}}T.
\end{equation}
由闭合积分为零定义熵差，与由保守力的环路积分定义势能差具有相同的数学结构。这里成为全微分的是可逆路径上的 $\delta Q_{\rm rev}/T$，热量本身仍依赖路径。经典宏观计算可任选熵的零点，所得熵差相同。

\begin{example}[理想气体熵的不同坐标形式]\label{ex:entropycoordinates}
对 $\nu$ 摩尔常热容理想气体，以 $(p_0,V_0,T_0)$ 为参考态，推导 $S(p,V)$ 与 $S(T,p)$，并证明可逆绝热线是等熵线。
\Diagram{fig:entropycoordq}{先等压、后等容的可逆参考路径}{
 \draw[axis](0,0)--(4.6,0)node[right]{$V$};\draw[axis](0,0)--(0,3)node[above]{$p$};\draw[axis](1,.8)--(3.3,.8);\draw[axis](3.3,.8)--(3.3,2.2);\node[left]at(1,.8){$A$};\node[below]at(3.3,.8){$B$};\node[right]at(3.3,2.2){$C$};
}
\end{example}
\begin{solution}
沿用例~\ref{ex:idealstate}已经求出的可逆熵积分。图\ref{fig:entropycoordq}的等压段$A\to B$有$T_B/T_0=V/V_0$，熵变为$C_p\ln(V/V_0)$；等容段$B\to C$有$T/T_B=p/p_0$，熵变为$C_V\ln(p/p_0)$。于是
\[
 S-S_0=C_V\ln\frac p{p_0}+C_p\ln\frac V{V_0}
 =C_V\ln\frac{pV^\gamma}{p_0V_0^\gamma}.
\]
由 $V/V_0=(T/T_0)(p_0/p)$，又得 $S-S_0=C_p\ln(T/T_0)-\nu R\ln(p/p_0)$。
两段温度比的换算分别使用各自的不变量：等压段有 $T\propto V$，等容段有 $T\propto p$。在 $p,V$ 形式中再提出 $C_V$，利用 $C_p/C_V=\gamma$ 把两项对数合并。若改用 $T,p$，应同时替换体积对数中的温度项和压强项，合并后的温度系数为 $C_V+\nu R=C_p$。三种写法描述同一个熵差，参考态的三个量由状态方程约束，并非三个可独立任取的参数。

\Diagram{fig:entropycoorda}{不同熵值对应一族绝热线}{
 \draw[axis](0,0)--(5,0)node[right]{$V$};\draw[axis](0,0)--(0,3)node[above]{$p$};\foreach \c in{1.2,1.8,2.4}{\draw[curve]plot[domain=1:4.5](\x,{\c/\x^1.4});}
 \node at(3.4,1.8){向上熵增};\draw[axis](3,.6)--(3,1.3);
}
图\ref{fig:entropycoorda}中每条曲线满足 $pV^\gamma$ 恒定。可逆绝热时 $\delta Q_{\rm rev}=0$，故 $\dd S=0$；若绝热但不可逆，系统熵可以增加，不能仍强制初末态在同一条等熵线上。
\end{solution}

\subsection{克劳修斯不等式与熵产生}
\label{th:clausius}
任意循环满足 $\oint\delta Q/T_b\le0$，其中 $T_b$ 是热量通过边界时的热源温度；可逆取等号。可以把每次与外界的热交换通过辅助可逆热机转换成与同一参考热源的交换，再利用单热源循环不能净做正功推出这一结论。

对于非循环过程，系统与环境合起来满足熵不减。若热量 $\delta Q$ 通过温度为 $T_b$ 的边界传入系统，可写
\[
 \dd S=\frac{\delta Q}{T_b}+\dd S_{\rm gen},\qquad \dd S_{\rm gen}\ge0.
\]
孤立系统不与外界交换熵，$\Delta S= S_{\rm gen}\ge0$。有限温差传热时，高温体失去 $\delta Q$、低温体得到 $\delta Q$，整体熵变
\[
 \delta S=\delta Q\left(\frac1{T_c}-\frac1{T_h}\right)>0.
\]
计算不可逆过程的系统熵变，要根据实际初末平衡态另选可逆路径；不能把实际过程的 $\delta Q/T$ 硬积分成熵差。

\subsection{用辅助可逆机推出循环不等式}
\label{th:clausiusconstruction}
把一次循环与外界的热交换分成若干份，装置从温度 $T_k$ 的热库吸收有符号热量 $Q_k$，放热时 $Q_k<0$。另选温度为 $T_0$ 的参考热库。对每个 $T_k$ 配一台辅助可逆机，使它从该热库取得的热量恰为 $-Q_k$；于是原装置与辅助机在 $T_k$ 热库的热交换完全抵消。

设该辅助机从参考热库吸热为 $Q_{0k}$。辅助机可正向也可反向运行，其可逆热量关系用有符号形式写为
\[
 -\frac{Q_k}{T_k}+\frac{Q_{0k}}{T_0}=0,
 \qquad Q_{0k}=T_0\frac{Q_k}{T_k}.
\]
全部工质均作循环，组合装置只与参考热库交换净热
$Q_0=T_0\sum_kQ_k/T_k$。第一定律给组合输出功 $W_{\rm comb}=Q_0$。若 $\sum_kQ_k/T_k>0$，组合的唯一效果便是从单一热库吸热并全部转成功，违反开尔文表述。因此
\[
 \sum_k\frac{Q_k}{T_k}\le0.
\]
把热交换分得足够细便得到循环积分形式。若原循环也可逆，整套装置可以反向运行，同样的不等式施于反循环，便与原式合成等号。

\Diagram{fig:auxiliaryclausius}{辅助可逆机抵消各分温热库的交换}{
 \node[draw](t)at(0,2){热库 $T_k$};
 \node[draw](a)at(-1.6,0){原循环};
 \node[draw](b)at(2,0){辅助可逆机};
 \node[draw](z)at(2,-1.8){参考热库 $T_0$};
 \draw[axis](t)--(a)node[midway,left]{$Q_k$};
 \draw[axis](b)--(t)node[midway,right]{$Q_k$};
 \draw[axis](z)--(b)node[midway,right]{$Q_{0k}$};
 \node at(-1.3,-1.6){箭头按 $Q_k>0$ 表示};
}
图\ref{fig:auxiliaryclausius}画出一份热交换的抵消关系。辅助机向 $T_k$ 送回的热量大小为 $Q_k$，它从参考热库取热；若 $Q_k<0$，相应方向全部反转。比较在全部循环结束后进行，辅助工质已恢复初态，无须再计其能量或熵的存量变化。

\subsection{实际边界温度与可逆替代路径}
熵差只要求初末为可定义的平衡态。实际传热时，系统内部可以存在温度梯度，此时不能随意选一个“平均温度”去除实际热量。熵收支中的温度必须对应所选边界处传递熵的温度；若把接触层一并纳入研究对象，可用外部热库温度计算熵输入，把接触层内的有限温差熵产生也计入 $S_{\rm gen}$。

对从恒温热库 $T_0$ 吸收总热 $Q$ 的封闭系统，
\[
 \Delta S_{\rm sys}-\frac Q{T_0}=S_{\rm gen}\ge0.
\]
左侧的系统熵差可沿任意可逆替代路径计算，热库熵差则按真实热量计算。替代路径只用于求熵差；计算实际热量、功和熵产生时，仍使用真实过程的边界条件。

若小块物体温度为 $T$，通过接触层从 $T_0>T$ 的热库吸收 $\delta Q$，物体在内部近似均温，则物体熵增为 $\delta Q/T$，热库熵变为 $-\delta Q/T_0$，两者之和
\[
 \delta S_{\rm gen}
 =\delta Q\frac{T_0-T}{TT_0}.
\]
温差趋于零时，每传递相同热量的熵产生趋于零；但传递速率通常也随温差降低。可逆极限与有限时间内的最大功率因而属于不同的过程要求。

\subsection{熵温图中的热量与路径方向}
对可逆过程 $\delta Q_{\rm rev}=T\dd S$。在以 $S$ 为横轴、$T$ 为纵轴的图上，沿曲线的有向面积给可逆热量。横向向右对应吸热，向左对应放热；等熵段为竖线，其热量为零。可逆循环满足 $\Delta U=0$，故封闭有向面积同时等于输出功。

\Diagram{fig:entropypatharea}{可逆路径下的有向面积与热量}{
 \draw[axis](0,0)--(5.5,0)node[right]{$S$};
 \draw[axis](0,0)--(0,3.4)node[above]{$T$};
 \fill[gray!15](1,0)--(1,1.3)--(4,2.5)--(4,0)--cycle;
 \draw[curve](1,1.3)--(4,2.5);
 \draw[->](2,1.7)--(3,2.1);
 \draw[guide](1,0)--(1,1.3);\draw[guide](4,0)--(4,2.5);
 \node[below]at(1,0){$S_A$};\node[below]at(4,0){$S_B$};
 \node at(2.5,.65){$\int_A^BT\dd S$};
}
图\ref{fig:entropypatharea}的面积只表示图中可逆路径的热量。不可逆绝热过程可以有 $\Delta S>0$，但实际 $Q=0$；把它的初末点连成一条线再求面积，算到的是另一个人为选择的可逆过程。相变等温段在熵温图上为水平线，面积为 $T\Delta S$，直接对应潜热。


\begin{example}[自由膨胀为何熵增]\label{ex:freeexpansion}
$\nu$ 摩尔理想气体最初占绝热刚性容器的一半体积 $V_0$，另一半为真空。撤去隔板，待其重新平衡。求温度、压强和熵变。
\Diagram{fig:freeexpq}{隔板撤去后气体占据双倍体积}{
 \draw(0,0)rectangle(5,1.7);\draw[dashed](2.5,0)--(2.5,1.7);\node at(1.25,.85){$p_0,T_0,V_0$};\node at(3.75,.85){真空};
}
\end{example}
\begin{solution}
图\ref{fig:freeexpq}的外边界不动，外界又不供热，所以 $\Delta U=0$。理想气体内能仅依赖温度，$T_f=T_0$，由状态方程得 $p_f=p_0/2$。实际过程中不经过一串均匀平衡态，不能对气体内部取一条实际准静态 $p(V)$ 曲线。
初末两态温度相同，只是理想气体内能不随体积改变的结果；在膨胀途中，局部温度和流速可以不均匀。图中的等温线是为了计算熵而构造的另一条可逆路径：这条替代路径有吸热和对外功，二者恰好相等，内能差仍为零。实际自由膨胀则两者均为零。比较两条路径可以看到，内能差和熵差只看端点，热量与功分别依赖过程，因此用替代路径求熵并不等于把其热量当成实际吸热。

\Diagram{fig:freeexpa}{用于计算熵的等温可逆替代路径}{
 \draw[axis](0,0)--(5,0)node[right]{$V$};\draw[axis](0,0)--(0,3)node[above]{$p$};\draw[curve]plot[domain=1:2](\x,{2.4/\x});\draw[->](1.3,1.85)--(1.6,1.5);\node[below]at(1,0){$V_0$};\node[below]at(2,0){$2V_0$};\node at(3.4,2){$T=T_0$};
}
沿图\ref{fig:freeexpa}的可逆等温路径，$\delta Q_{\rm rev}=p\dd V=\nu RT_0\dd V/V$，故 $\Delta S=\nu R\ln2$。实际过程 $Q=0$，不代表 $\Delta S=0$；所增加的熵来自系统内部的不可逆膨胀。
\end{solution}

\subsection{孤立系统的熵最大条件}
\label{th:maxentropy}
在总能量、总体积和总粒子数固定的条件下，允许子系统交换能量、体积或粒子。平衡态是受这些约束的熵极大态。以能量交换为例，$\dd U_2=-\dd U_1$，总熵微分为
\[
 \dd S_{\rm tot}=\left(\frac1{T_1}-\frac1{T_2}\right)\dd U_1.
\]
在允许双向交换的内部平衡点，该式对任意微小交换都应为零，故$T_1=T_2$。若还允许体积和粒子交换，先在单组分基本关系中引入每粒子的化学势$\mu$：$\dd U=T\dd S-p\dd V+\mu\dd N$。两子系统的总量固定，分别取$\dd U_2=-\dd U_1$、$\dd V_2=-\dd V_1$、$\dd N_2=-\dd N_1$，总熵微分为
\[
 \begin{aligned}
 \dd S_{\rm tot}={}&\left(\frac1{T_1}-\frac1{T_2}\right)\dd U_1
 +\left(\frac{p_1}{T_1}-\frac{p_2}{T_2}\right)\dd V_1\\
 &-\left(\frac{\mu_1}{T_1}-\frac{\mu_2}{T_2}\right)\dd N_1.
 \end{aligned}
\]
各交换量若可独立正负变化，其系数分别为零；在已同温的条件下，继续得到$p_1=p_2$及$\mu_1=\mu_2$。固定隔板关闭体积交换，不透粒子的膜关闭粒子交换，对应的平衡条件便不能仅从这类受限变化推出。有外场时还须在能量中计入相应势能。

\section{混合气体与有限物体的熵变}
\label{th:mixing}
理想气体混合物的各组分相互不改变单粒子能量。均匀混合后，组分 $i$ 的分压为 $p_i=\nu_iRT/V$，总压强为 $\sum_i p_i$。熵可按各组分占据的整个可达体积计算后相加。不同组分具有可区分的分子性质，混合可产生构型熵；完全相同的气体撤去隔板，不应再人为贴上“左边粒子、右边粒子”的标签制造混合熵。

\begin{example}[同温同压两种气体混合]\label{ex:equalmix}
绝热刚性容器由隔板分成两半，每半体积 $V_0$，分别装有不同种理想气体，压强 $p_0$、温度 $T_0$ 相同。撤板后求最终状态与总熵增。
\Diagram{fig:equalmixq}{可区分的两种分子分别占据半个容器}{
 \draw(0,0)rectangle(5,1.7);\draw(2.5,0)--(2.5,1.7);\foreach \p in{(.5,.5),(1.1,1.2),(1.8,.6)}\node[circle,draw,inner sep=1.8pt]at\p{};\foreach \p in{(3,.5),(3.7,1.2),(4.5,.6)}\node[rectangle,fill,inner sep=1.8pt]at\p{};
}
\end{example}
\begin{solution}
图\ref{fig:equalmixq}两侧物质的量均为 $\nu_0=p_0V_0/(RT_0)$。总能量守恒，初温相同，故最终仍为 $T_0$。总量 $2\nu_0$ 占体积 $2V_0$，最终压强仍为 $p_0$，每种组分分压为 $p_0/2$。
原先每种分子都只在半个容器中运动，撤板后可达位置加倍，尽管总体温度和总压强都没有改变。求熵时可以设想使用选择膜，把每种组分分别作可逆等温扩展；它们所受的是各自分压，最终各为原压的一半。这条替代路径只用于计算状态熵差。对于同一种气体，原来“来自左边”和“来自右边”并不是可用于选择膜分离的物理组分，按这两个人为标签分开计熵会重复计数。

\Diagram{fig:equalmixa}{混合后每个组分都可占据整个容器}{
 \draw(0,0)rectangle(5,1.6);\foreach \p in{(.5,.5),(2.6,1.1),(4.2,.5)}\node[circle,draw,inner sep=1.8pt]at\p{};\foreach \p in{(1.2,1.1),(2.1,.5),(4.5,1.1)}\node[rectangle,fill,inner sep=1.8pt]at\p{};
 \node at(2.5,-.55){各组分 $V_0\to2V_0$，$T$ 不变};
}
图\ref{fig:equalmixa}中每个组分均经历可达体积加倍的状态变化，所以 $\Delta S_i=\nu_0R\ln2$，总熵增
\[
 \Delta S=2\nu_0R\ln2=\frac{2p_0V_0}{T_0}\ln2.
\]
若两侧是同种气体且初态完全相同，撤板不改变宏观状态，总熵变为零。该变式必须改变计数方式，不能把原式机械套用。
\end{solution}

\begin{example}[不同初温、压强与热容的混合]\label{ex:generalmix}
两种不同理想气体最初分别占体积 $V_1,V_2$，初温 $T_1,T_2$，物质的量为 $\nu_1,\nu_2$，摩尔定容热容分别为常量 $c_1,c_2$。撤去绝热刚性容器内的隔板，求 $T_f,p_f$、分压和总熵变。
\Diagram{fig:generalmixq}{初态可以具有不同温度、密度与热容}{
 \draw(0,0)rectangle(6,1.8);\draw(2.4,0)--(2.4,1.8);\node[align=center]at(1.2,.9){$\nu_1,V_1,T_1$\\$c_1$};\node[align=center]at(4.2,.9){$\nu_2,V_2,T_2$\\$c_2$};
}
\end{example}
\begin{solution}
令 $C_1=\nu_1c_1$、$C_2=\nu_2c_2$。图\ref{fig:generalmixq}的外界不交换热和功，各组分理想气体内能只依赖温度，所以
\[
 C_1T_1+C_2T_2=(C_1+C_2)T_f,
 \qquad T_f=\frac{C_1T_1+C_2T_2}{C_1+C_2}.
\]
最终体积 $V_f=V_1+V_2$，分压 $p_{i,f}=\nu_iRT_f/V_f$，总压强 $p_f=(\nu_1+\nu_2)RT_f/V_f$。
温度项和体积项有不同来源。对每种组分，定容变温给 $\int_{T_i}^{T_f}C_i\dd T/T$，等温改变可达体积给 $\int_{V_i}^{V_f}\nu_iR\dd V/V$；两段分别积分后相加。温度加权用的是系统热容 $C_i$，压强加权用的是粒子数或物质的量，两种权重通常不同。最终总压只反映粒子数之和，求熵时却须保留组分信息，不能把全部气体先合并成一个任意热容的“等效单一气体”。

\Diagram{fig:generalmixa}{每个组分的替代路径：先变温，再扩展可达体积}{
 \node[draw,align=center]at(0,0){$T_i,V_i$};\draw[axis](.9,0)--(2.1,0);\node[draw]at(3,0){$T_f,V_i$};\draw[axis](3.9,0)--(5.1,0);\node[draw]at(6.2,0){$T_f,V_f$};
}
沿图\ref{fig:generalmixa}分别计算两项，
\[
 \Delta S=\sum_{i=1}^2\left[C_i\ln\frac{T_f}{T_i}
 +\nu_iR\ln\frac{V_1+V_2}{V_i}\right].
\]
温度项由热平衡贡献，体积项由不同组分的混合贡献。若题目给 $p_i,V_i,T_i$ 而未给物质的量，应先以 $\nu_i=p_iV_i/(RT_i)$ 换算。最终组分的分压通常不等于其初压。
\end{solution}

\begin{example}[两个物体直接接触的熵增]\label{ex:directcontact}
热容分别为常量 $C_1,C_2$ 的两个固体，初温 $T_1>T_2>0$。整体绝热，体积变化可忽略，直接接触到平衡。求终温、传热量和总熵增，并用传热量作积分变量重新推导熵变。
\Diagram{fig:directq}{热量从高温物体传向低温物体}{
 \node[draw,minimum height=1cm]at(0,0){$C_1,T_1$};\node[draw,minimum height=1cm]at(4,0){$C_2,T_2$};\draw[axis](1,0)--(3,0)node[midway,above]{$Q$};
}
\end{example}
\begin{solution}
图\ref{fig:directq}的整体能量守恒给出算术加权终温
\[
 T_d=\frac{C_1T_1+C_2T_2}{C_1+C_2},\qquad
 Q_f=C_1(T_1-T_d)=\frac{C_1C_2}{C_1+C_2}(T_1-T_2).
\]
实际传递累计热量 $Q$ 时，两体温度分别为 $t_1=T_1-Q/C_1$、$t_2=T_2+Q/C_2$。固体各自可近似均温，取微小状态变化的可逆替代路径，
\[
 \dd S_{\rm tot}=\left(-\frac1{t_1}+\frac1{t_2}\right)\dd Q.
\]
以累计传热量为变量时，两项积分的原函数分别为 $C_1\ln(T_1-Q/C_1)$ 和 $C_2\ln(T_2+Q/C_2)$。两者对 $Q$ 求导，恰好给 $-1/t_1$ 与 $+1/t_2$。积分到 $Q_f$ 时，两体温度才同时变为 $T_d$。在此前始终 $t_1>t_2$，因而瞬时总熵变化为正；到末态温差消失，其变化率才趋零。用共同终温替代整个过程中的瞬时温度，会把这段有限温差熵产生漏掉。

\Diagram{fig:directa}{累计热量决定瞬时温差和熵产生率}{
 \draw[axis](0,0)--(5,0)node[right]{$Q$};\draw[axis](0,0)--(0,3)node[above]{$t$};\draw[curve](0,2.5)--(4,1.3);\draw[curve](0,.5)--(4,1.3);\node[right]at(4,1.3){$T_d$};\node[left]at(0,2.5){$T_1$};\node[left]at(0,.5){$T_2$};
}
图\ref{fig:directa}两条直线在 $Q_f$ 处相交。分别积分，
\[
 \Delta S=C_1\ln\frac{T_d}{T_1}+C_2\ln\frac{T_d}{T_2}\ge0.
\]
由对数的凹性，$\ln T_d\ge[C_1\ln T_1+C_2\ln T_2]/(C_1+C_2)$，等号仅在初温相同取到，所以温差非零时总熵严格增加。计算不曾假设实际传热可逆；可逆路径只是求各固体状态熵差的工具。
\end{solution}

\subsection{熵极值的一阶条件与稳定性}
\label{th:entropystability}
把两个固定体积的子系统组成孤立整体，总内能 $U_\Sigma$ 固定。以 $U_1$ 为独立变量，
$S_\Sigma(U_1)=S_1(U_1)+S_2(U_\Sigma-U_1)$。
一阶导数为$1/T_1-1/T_2$，其零点给共同温度。还须检查二阶项。由$\dd T_1/\dd U_1=1/C_1$、$\dd T_2/\dd U_1=-1/C_2$，链式法则给$\dd(T_1^{-1})/\dd U_1=-1/(T_1^2C_1)$、$\dd(-T_2^{-1})/\dd U_1=-1/(T_2^2C_2)$。在共同温度$T$处，
\[
 \frac{\dd^2S_\Sigma}{\dd U_1^2}
 =-\frac1{T^2C_1}-\frac1{T^2C_2}<0,
\]
其中 $C_1,C_2>0$ 是相应固定体积下的热容。偏离平衡一个微小能量 $\delta U$ 时，
\[
 S_\Sigma-S_{\Sigma,\rm eq}
 =-\frac{(\delta U)^2}{2T^2}
 \left(\frac1{C_1}+\frac1{C_2}\right)+O((\delta U)^3).
\]
一次项消失，二次项为负，所以微小自发能量交换会把系统带回熵更高的附近状态。这一二阶条件使熵的驻点成为稳定最大值。

\Diagram{fig:entropymaximum}{能量分配方向上的熵最大值}{
 \draw[axis](0,0)--(5.5,0)node[right]{$U_1$};
 \draw[axis](0,0)--(0,3.2)node[above]{$S_\Sigma$};
 \draw[curve]plot[domain=.7:4.5](\x,{2.7-.4*(\x-2.6)^2});
 \draw[guide](2.6,0)--(2.6,2.7);
 \node[below]at(2.6,0){$U_{1,\rm eq}$};
 \draw[->](1.3,1.7)--(1.9,2.3);
 \draw[->](4,1.65)--(3.4,2.3);
}
图\ref{fig:entropymaximum}的横向变化只重新分配整体已有能量，未改变总能量。若允许体积、粒子数或外场参数变化，应连同新的约束重新建立总熵函数。隔板固定、可移动、可导热或选择透过粒子，分别开放不同的交换方向，所需平衡条件也不同。

\subsection{可达状态与最大功的几何关系}
能量守恒与熵不减共同限制可达终态，但并不替代对装置操作条件的说明。例如两个有限物体之间可输出机械功时，物体总内能可以下降；若禁止对外功且最终机械储能恢复，物体总内能必须等于初值。前者在不同能量水平上寻找尽可能低的终态，后者只能沿固定能量曲线移动。

对常热容物体，用温度坐标表示的能量曲线是直线
$C_1T_1+C_2T_2=\text{常量}$，熵曲线则为
$C_1\ln(T_1/T_\ast)+C_2\ln(T_2/T_\ast)=\text{常量}$。
这里 $T_\ast>0$ 是共同参考温度，所有对数的自变量无量纲。等温线 $T_1=T_2$ 与能量直线的交点给直接接触终态；与初始等熵曲线的交点给可逆做功的共同终态。这两个交点一般不重合。

暂存机械功可以使物体总内能在中途变化，最后若把机械储能完全归还，两物体又回到原能量直线上。此时所求“某物体最高温”是这条能量直线与熵允许区间的端点，而非两物体共同终温。前文温度转移题的两个根分别对应初态和另一个等熵端点，必须根据所求物体及允许操作选择。

\subsection{热容随温度变化时的统一端点方程}
若各物体热容为 $C_i(T)>0$，先分别定义相对于任意固定参考态的能量和熵，
\[
 U_i(T)-U_i(T_\ast)=\int_{T_\ast}^{T}C_i(t)\dd t,
 \qquad
 S_i(T)-S_i(T_\ast)=\int_{T_\ast}^{T}\frac{C_i(t)}t\dd t.
\]
直接接触的共同终温由
\[
 \sum_i\int_{T_i}^{T_d}C_i(t)\dd t=0
\]
确定；可逆热机取最大功的共同终温由
\[
 \sum_i\int_{T_i}^{T_r}\frac{C_i(t)}t\dd t=0
\]
确定。两式的被积函数不同，分别代表能量和熵约束。前一式左侧对终温的导数为 $\sum_iC_i(T)>0$，后一式为 $\sum_iC_i(T)/T>0$，故在正热容温区内各有至多一个共同终温根。

最大功再由
$W_{\max}=\sum_i[U_i(T_i)-U_i(T_r)]$ 求得。参考常数在初末相减时全部抵消，既无需给每个物体规定绝对内能，也无需把温变热容强行替换成某个平均常数。数值求根时，初始最低温与最高温可作为这类共同终温的自然夹逼区间。

\subsection{不同热容流率的并流与逆流换热}
\label{th:unequalexchanger}
稳态换热应使用每单位时间通过装置的热容，即热容流率
$\mathcal C=\dot m c_p$，其中 $c_p$ 为单位质量的定压热容，$\mathcal C$ 的量纲为功率除温度。设冷热流的热容流率 $\mathcal C_h,\mathcal C_c$ 沿程为常量，不发生相变，传热系数与单位长度面积的乘积 $k_\ell$ 也为常量，总长为 $L$。管外绝热、无其他做功时，总传热率满足
\[
 \dot Q=\mathcal C_h(T_{h,\rm in}-T_{h,\rm out})
 =\mathcal C_c(T_{c,\rm out}-T_{c,\rm in}).
\]
这与间歇地把热水份移走的协议不同：稳态两流在沿程持续交换能量，出口物质不断离开装置。

并流时，坐标 $x$ 沿两流共同方向，局部温差
$\Delta=T_h-T_c$ 满足
\[
 \frac{\dd\Delta}{\dd x}
 =-k_\ell\left(\frac1{\mathcal C_h}+\frac1{\mathcal C_c}\right)\Delta.
\]
分离变量，从入口到出口积分，$[\ln\Delta]_{\Delta_{\rm in}}^{\Delta_{\rm out}}=-k_\ell(1/\mathcal C_h+1/\mathcal C_c)[x]_0^L$。取指数给出$\Delta_{\rm out}=\Delta_{\rm in}\ee^{-\mathcal N(1+r)}$。再用出口温差等于
$\Delta_{\rm in}-\dot Q(1/\mathcal C_h+1/\mathcal C_c)$，可求传热率。
记 $\mathcal C_{\min}$ 为两者中较小值，
$r=\mathcal C_{\min}/\mathcal C_{\max}$、
$\mathcal N=k_\ell L/\mathcal C_{\min}$，并定义
$\varepsilon=\dot Q/[\mathcal C_{\min}(T_{h,\rm in}-T_{c,\rm in})]$，得到
\[
 \varepsilon_{\parallel}
 =\frac{1-\ee^{-\mathcal N(1+r)}}{1+r}.
\]
并流无限加长时，两股出口流趋于共同温度，换热效能趋于 $1/(1+r)$。热量守恒决定这个共同温度，两股流均无法完全达到对方的入口温度。

\Diagram{fig:unequalexchanger}{逆流换热的两端边界位于不同位置}{
 \draw(0,0)rectangle(6,1.6);
 \draw[axis](.3,1.2)--(5.7,1.2);
 \draw[axis](5.7,.4)--(.3,.4);
 \node[left]at(0,1.2){$T_{h,\rm in}$};
 \node[right]at(6,1.2){$T_{h,\rm out}$};
 \node[left]at(0,.4){$T_{c,\rm out}$};
 \node[right]at(6,.4){$T_{c,\rm in}$};
 \draw[axis](0,-.5)--(6,-.5)node[right]{$x$};
}
逆流时按图\ref{fig:unequalexchanger}仍取 $x$ 沿热流方向，冷流实际流向与坐标相反。因此
$\dd T_h/\dd x=-k_\ell\Delta/\mathcal C_h$，
$\dd T_c/\dd x=-k_\ell\Delta/\mathcal C_c$。
为写出边界消元的步骤，先取 $\mathcal C_h\ge\mathcal C_c$，令
$D=T_{h,\rm in}-T_{c,\rm in}$，则
\[
 \Delta(0)=D-\frac{\dot Q}{\mathcal C_c},\qquad
 \Delta(L)=D-\frac{\dot Q}{\mathcal C_h}.
\]
相减两条温度方程，得$\Delta'=k_\ell(1/\mathcal C_c-1/\mathcal C_h)\Delta$。分离变量后，$[\ln\Delta]_{\Delta(0)}^{\Delta(L)}=k_\ell(1/\mathcal C_c-1/\mathcal C_h)[x]_0^L$，故
\[
 \frac{\Delta(L)}{\Delta(0)}
 =\exp\!\left[k_\ell L\left(
 \frac1{\mathcal C_c}-\frac1{\mathcal C_h}\right)\right]
 =\ee^{\mathcal N(1-r)}.
\]
记$E=\ee^{\mathcal N(1-r)}$。代入两端温差，$D-\dot Q/\mathcal C_h=E(D-\dot Q/\mathcal C_c)$，移项得$\dot Q(E/\mathcal C_c-1/\mathcal C_h)=D(E-1)$。因$\mathcal C_c=\mathcal C_{\min}$，除以$D\mathcal C_c$并约去$E$，得到
\[
 \varepsilon_{\rm counter}
 =\frac{1-\ee^{-\mathcal N(1-r)}}
 {1-r\ee^{-\mathcal N(1-r)}}\quad(r<1),\qquad
 \varepsilon_{\rm counter}\big|_{r=1}
 =\frac{\mathcal N}{1+\mathcal N}.
\]
取$r\to1$时，令$\delta=1-r$，则$e^{-\mathcal N\delta}=1-\mathcal N\delta+O(\delta^2)$。分子为$\mathcal N\delta+O(\delta^2)$，分母为$1-(1-\delta)[1-\mathcal N\delta+O(\delta^2)]=(1+\mathcal N)\delta+O(\delta^2)$，约去$\delta$便得到$\mathcal N/(1+\mathcal N)$；例~\ref{ex:heatexchanger}另从恒定沿程温差直接求解这一情形。

另一种热容流率大小关系可交换冷热流的命名后得到同一形式。无限长时，逆流可使较小热容流率的一侧趋近对方入口温度。冷流出口可高于热流出口，因为它们位于不同位置；任意同一截面仍须保持热流温度不低于冷流，局部传热方向才与第二定律一致。


\begin{example}[有限热源间的最大功]\label{ex:finitework}
上述两物体之间接一台可逆热机，整体绝热，热机循环工作，直到两物体温度相同。求可获得的最大功和最终温度，并与直接接触比较。
\Diagram{fig:finiteworkq}{有限热源的温度随做功不断变化}{
 \node[draw]at(0,1.5){$C_1,t_1$};\node[draw,circle]at(0,0){机};\node[draw]at(0,-1.5){$C_2,t_2$};
 \draw[axis](0,1.1)--(0,.45);\draw[axis](0,-.45)--(0,-1.1);\draw[axis](.5,0)--(2,0)node[right]{$\dd W$};
}
\end{example}
\begin{solution}
图\ref{fig:finiteworkq}热机每个微循环从热体吸收 $\delta Q_h=-C_1\dd t_1$，向冷体排出 $\delta Q_c=C_2\dd t_2$。可逆条件 $\delta Q_h/t_1=\delta Q_c/t_2$ 给
\[
 C_1\frac{\dd t_1}{t_1}+C_2\frac{\dd t_2}{t_2}=0.
\]
从初态积分至共同终温$T_r$，即$C_1[\ln t_1]_{T_1}^{T_r}+C_2[\ln t_2]_{T_2}^{T_r}=0$。展开端点并合并$\ln T_r$，得到$(C_1+C_2)\ln T_r=C_1\ln T_1+C_2\ln T_2$，所以
\[
 T_r=T_1^{C_1/(C_1+C_2)}T_2^{C_2/(C_1+C_2)}.
\]
最大功的上界可从任意终态直接证明。若热机最终回到初态、两物体达到共同温度 $T_f$，孤立整体的熵不减给出
\[
 C_1\ln\frac{T_f}{T_1}+C_2\ln\frac{T_f}{T_2}\ge0.
\]
因此 $T_f\ge T_r$。而输出功是初始能量减去 $(C_1+C_2)T_f$，随终温升高而减少。可逆机使熵不等式取等号，所以其结果达到上界。这个证明同时排除了靠摩擦或有限温差传热获得更多功的可能，而不依赖某一台热机的具体结构。

\Diagram{fig:finiteworka}{可逆过程在温度平面中的熵约束}{
 \draw[axis](0,0)--(4.5,0)node[right]{$t_1$};\draw[axis](0,0)--(0,3)node[above]{$t_2$};\draw[guide](0,0)--(3,3);\draw[curve]plot[domain=1.1:3.8](\x,{2.8/\x});\node at(3.2,2.4){$t_1^{C_1}t_2^{C_2}=\text{常量}$};\node at(2.7,.5){等热容示意};
}
由整体能量守恒，最大功为
\[
 W_{\max}=C_1T_1+C_2T_2-(C_1+C_2)T_r.
\]
图\ref{fig:finiteworka}的可逆曲线与等温直线交点为终态。直接接触的终温为加权算术平均，可逆做功的终温为加权几何平均，故 $T_d\ge T_r$。两种终态的内能差正好对应直接接触所损失的输出功。
\end{solution}

\begin{example}[热容随温度改变的有限热源]
两个物体热容分别为 $C_1(T)=a_1T$、$C_2(T)=a_2T$，其中 $a_i>0$。初温 $T_1>T_2$。分别求直接接触终温和可逆热机做功终温及最大功。
\Diagram{fig:linearcfiniteq}{热容随温度线性增长的两个物体}{
 \draw[axis](0,0)--(4.2,0)node[right]{$T$};\draw[axis](0,0)--(0,2.8)node[above]{$C$};\draw[curve](0,0)--(3,2.4)node[right]{$a_1T$};\draw[curve](0,0)--(3.7,1.5)node[right]{$a_2T$};
}
\end{example}
\begin{solution}
图\ref{fig:linearcfiniteq}中能量变化为$\Delta U_i=\int_{T_i}^{T_f}a_iT\dd T=[a_iT^2/2]_{T_i}^{T_f}=a_i(T_f^2-T_i^2)/2$。直接接触不做功，故
\[
 T_d=\sqrt{\frac{a_1T_1^2+a_2T_2^2}{a_1+a_2}}.
\]
可逆机则保持总熵，$\Delta S_i=\int_{T_i}^{T_f}a_i\dd T=[a_iT]_{T_i}^{T_f}=a_i(T_f-T_i)$，从而
\[
 T_r=\frac{a_1T_1+a_2T_2}{a_1+a_2}.
\]
本题中的 $a_i$ 量纲为热容除以温度。定容能量积分给温度平方，而熵积分中热容的一个温度因子被分母抵消，所以只剩温度一次项。这是直接接触终温为平方平均、可逆终温为算术平均的根源。把 $T_r$ 代回最大功式时，可先通分，利用
\[
 (a_1+a_2)(a_1T_1^2+a_2T_2^2)
 -(a_1T_1+a_2T_2)^2=a_1a_2(T_1-T_2)^2,
\]
便得到非负的温差平方形式。

\Diagram{fig:linearcfinitea}{求能量与熵时使用不同的积分权重}{
 \node[draw]at(0,0){$\Delta U=\int C(T)\dd T$};\node[draw]at(5,0){$\Delta S=\int C(T)\dd T/T$};\node at(2.5,-.9){不再使用常热容的对数公式};
}
按图\ref{fig:linearcfinitea}计算初末内能之差，
\[
 W_{\max}=\frac12[a_1T_1^2+a_2T_2^2-(a_1+a_2)T_r^2]
 =\frac{a_1a_2}{2(a_1+a_2)}(T_1-T_2)^2.
\]
本题的可逆终温是算术加权平均，不是前一例的几何平均；终温的形式由热容函数决定，不能脱离积分条件背诵。
\end{solution}

\section{分批接触、换热与温度转移}
\begin{example}[分批热水加热冷水的极限]\label{ex:batchwater}
各 $100\,\mathrm g$ 的热水和冷水，初温分别为 $100\,^\circ\mathrm C$、$0\,^\circ\mathrm C$，水比热视为常量且无散热。把热水等分为 $n$ 份，每次只让一份与全部冷水隔壁接触到热平衡，随后移走该热水份，再换下一份。求冷水最终温度；与直接接触、分两份比较，并求 $n\to\infty$ 极限。
\Diagram{fig:batchq}{每份热水只使用一次，冷水主体始终保留}{
 \draw(0,0)rectangle(2,1.4);\node at(1,.7){冷水 $C$};\foreach \x in{3,3.8,4.6,5.4}{\draw(\x,0)rectangle++(.55,1.4);}\node at(4.5,1.9){热水各 $C/n$};\draw[axis](3,-.4)--(1,-.4);
}
\end{example}
\begin{solution}
用 $\vartheta_h,\vartheta_0$ 表示摄氏温度。图\ref{fig:batchq}每次接触只用能量守恒，温差和加权平均允许用摄氏温标计算。第 $j$ 次接触后冷水温度满足
\[
 C\vartheta_{j-1}+\frac Cn\vartheta_h
 =\left(C+\frac Cn\right)\vartheta_j,
 \quad\vartheta_j-\vartheta_h=\frac n{n+1}(\vartheta_{j-1}-\vartheta_h).
\]
迭代 $n$ 次，
\[
 \vartheta_n=\vartheta_h+
 \left(\frac n{n+1}\right)^n(\vartheta_0-\vartheta_h).
\]
第一次接触后的主体仍只有冷水的热容 $C$，刚使用的一份热水已被移走；第二次并不把它的热容继续加入主体。这是递推分母始终为 $C+C/n$ 的原因。两份方案中，冷水先升到 $100/3\,^\circ\mathrm C$，再与另一份初温 $100\,^\circ\mathrm C$ 的热水接触。若把已经用过的热水也一直混留在容器内，最终是全部冷热水直接混合的能量平均，操作协议与本题不同，不能沿用这里的递推。

\Diagram{fig:batcha}{递推温差按相同因子缩小}{
 \node[draw]at(0,0){$\vartheta_{j-1}-\vartheta_h$};\draw[axis](1.6,0)--(3.1,0)node[midway,above]{$n/(n+1)$};\node[draw]at(4.8,0){$\vartheta_j-\vartheta_h$};
}
图\ref{fig:batcha}中 $n=1$ 得 $50\,^\circ\mathrm C$，$n=2$ 得 $100(1-4/9)=55.56\,^\circ\mathrm C$；$n\to\infty$，因 $(1+1/n)^{-n}\to\ee^{-1}$，最终为 $100(1-\ee^{-1})\,^\circ\mathrm C\approx63.21\,^\circ\mathrm C$。

这个极限对应题设规定的操作：全部冷水保留，热水分份使用后移走。足够长的逆流换热器，或配有储功装置的可逆热机，允许不同的操作方式，其最高可达温度须另行求解。这里的能量式可用摄氏温差，计算熵时则须把温度换成开尔文温度。
\end{solution}

\begin{example}[等热容流率的并流和逆流换热]\label{ex:heatexchanger}
两股不混合液流的热容流率相等，均为 $\dot C$，入口温度为 $T_h>T_c$。总传热系数乘面积为 $K$，无对外散热。求并流和逆流时的换热量率及出口温度；讨论 $K/\dot C\to\infty$ 的极限。
\Diagram{fig:exchangerq}{并流与逆流的流向}{
 \draw(0,.2)rectangle(4,1.6);\draw[axis](.3,1.2)--(3.7,1.2);\draw[axis](.3,.6)--(3.7,.6);\node at(2,-.3){并流};
 \begin{scope}[xshift=5.4cm]\draw(0,.2)rectangle(4,1.6);\draw[axis](.3,1.2)--(3.7,1.2);\draw[axis](3.7,.6)--(.3,.6);\node at(2,-.3){逆流};\end{scope}
}
\end{example}
\begin{solution}
以累计传热能力$k$作长度坐标，$\dd\dot Q=(t_h-t_c)\dd k$。并流两股都沿$k$增大流动，故$\dot C\dd t_h=-\dd\dot Q$、$\dot C\dd t_c=\dd\dot Q$。温差$D=t_h-t_c$满足$\dd D/\dd k=-2D/\dot C$。以$D(0)=T_h-T_c$积分，$\int_{T_h-T_c}^{D_o}\dd D/D=-(2/\dot C)\int_0^K\dd k$，得到$\ln[D_o/(T_h-T_c)]=-2K/\dot C$，故$D_o=(T_h-T_c)\ee^{-2K/\dot C}$。两股温度和保持$T_h+T_c$，所以
\[
 \dot Q_{\parallel}=\frac{\dot C}{2}(T_h-T_c)(1-\ee^{-2K/\dot C}).
\]
逆流的冷端边界位于热流坐标的末端，而非起点，因此两条温度线沿同一坐标可以同时下降。每一局部的熵产生率为
\[
 \dd\dot S_{\rm gen}=
 (t_h-t_c)\dd k\left(\frac1{t_c}-\frac1{t_h}\right)
 =\frac{(t_h-t_c)^2}{t_ht_c}\dd k>0.
\]
有限 $K$ 下总有正温差。只有传热能力趋于无穷大时，逆流局部温差才可趋零，达到入口温度近乎互换的极限。冷流出口比热流出口热并不表示局部逆向传热，因为这两个出口位于装置的不同端。

\Diagram{fig:exchangera}{沿热流方向观察的温度分布}{
 \draw[axis](0,0)--(4.4,0)node[right]{$k$};\draw[axis](0,0)--(0,2.8)node[above]{$t$};\draw[curve]plot[domain=0:4](\x,{1.3+1.1*exp(-\x)});\draw[curve]plot[domain=0:4](\x,{1.3-1.1*exp(-\x)});\node at(2.5,2.6){并流};
 \begin{scope}[xshift=5.8cm]\draw[axis](0,0)--(4.2,0)node[right]{$k$};\draw[axis](0,0)--(0,2.8)node[above]{$t$};\draw[curve](0,2.5)--(3.8,1.3);\draw[curve](0,1.6)--(3.8,.4);\node at(2.2,2.8){逆流};\end{scope}
}
逆流时仍以热流方向增大的 $k$ 为坐标，冷流沿反方向运动，所以 $\dd t_c/\dd k=-D/\dot C$，与热流温度斜率相同，图\ref{fig:exchangera}两温度线的差 $D$ 恒定。设换热量率为 $\dot Q$，冷流在热流入口处的温度为 $T_c+\dot Q/\dot C$，于是 $D=T_h-T_c-\dot Q/\dot C$；又 $\dot Q=KD$，解得
\[
 \dot Q_{\rm counter}=\dot C(T_h-T_c)\frac{K/\dot C}{1+K/\dot C}.
\]
两种情形的热、冷出口分别为 $T_h-\dot Q/\dot C$、$T_c+\dot Q/\dot C$。无限传热能力时，并流两出口趋于算术平均；逆流则趋于交换入口温度。逆流每一局部仍由高温传向低温，并未违反第二定律。
\end{solution}

\begin{example}[用可逆装置转移温度的可达端点]\label{ex:temperaturetransfer}
两个不混合物体热容为常量 $C$、$\xi C$，初温 $T_1>T_2$。允许在两者之间运行可逆热机或热泵，并暂存机械功；全过程与外界绝热，结束时储功装置恢复原状，净外功为零。求第二物体能够达到的最高温度的方程，并解 $\xi=1$ 的情形。
\Diagram{fig:transferq}{可逆装置和临时储功装置都在整体系统内}{
 \node[draw]at(0,1.3){$C,T_1$};\node[draw]at(4,1.3){$\xi C,T_2$};\node[draw]at(2,0){可逆机／热泵};\node[draw]at(2,-1.3){储功装置};\draw[<->](0,.9)--(1.4,.35);\draw[<->](4,.9)--(2.6,.35);\draw[<->](2,-.35)--(2,-.95);
}
\end{example}
\begin{solution}
图\ref{fig:transferq}的净外功为零，故两物体总内能不变。设最终 $T_{1f}=xT_1$、$T_{2f}=yT_1$，并记 $\eta=T_2/T_1<1$，有
\[
 x+\xi y=1+\xi\eta.
\]
对于给定总能量，最终熵随第二物体温度先增加后减少。最高可达 $y$ 位于高温分支与初始总熵相等处；若允许更高温度，总熵将低于初始值，违反第二定律。可逆过程满足
\[
 \ln x+\xi\ln\frac y\eta=0,
 \quad x\left(\frac y\eta\right)^\xi=1.
\]
消去 $x$ 得到端点方程
\begin{equation}\label{eq:transferendpoint}
 1+\xi\eta-\xi y=\left(\frac\eta y\right)^\xi.
\end{equation}
取 $x=1+\xi\eta-\xi y$，在允许的正温度范围内，终态相对于初态的熵除以 $C$ 为 $F(y)=\ln x+\xi\ln(y/\eta)$。它的导数为
\[
 F'(y)=\xi\left(\frac1y-\frac1x\right),\qquad
 F''(y)=-\frac\xi{y^2}-\frac{\xi^2}{x^2}<0.
\]
所以等温点是唯一最大点；初态根位于其左侧，右侧另有唯一的等初熵根。再向右超过该根，总熵会低于初始值，故该根才给最高可达温度。图中的定能量约束针对储功已归零的终态，过程途中两物体还可把一部分能量暂存为机械功。

\Diagram{fig:transfera}{定能量直线上总熵的两个等初熵端点}{
 \draw[axis](0,0)--(5.3,0)node[right]{$y$};\draw[axis](0,-.8)--(0,2.5)node[above]{$(S-S_i)/C$};
 \draw[curve]plot[domain=.45:1.05,samples=80]({4*\x},{10*ln(2*\x*(1.5-\x))});\fill(2,0)circle(1.2pt)node[below]{$\eta$};\fill(4,0)circle(1.2pt)node[below]{$y_{\max}$};\draw(-.06,1)--(.06,1);\node[left]at(0,1){$0.1$};\node at(2.8,2.1){示例：$\xi=1,\ \eta=1/2$};
}
式\eqref{eq:transferendpoint}已有根 $y=\eta$，对应初态，所求是图\ref{fig:transfera}右侧另一根。$\xi=1$ 时，$x+y=1+\eta$、$xy=\eta$，两根为 $y=\eta,1$，故最高可达温度为 $T_1$，两物体恰好交换温度。

可达过程可分两段：先用热机取得功，直到两体同温；再用储功驱动热泵，把热量向原第二物体集中，直至储功归零。若不允许热泵、临时储功或温度越过相等点，只允许热机一直输出功，则终态应是上一节的共同几何平均温度，不能求本题的另一根。
\end{solution}

\subsection{环境约束下的可用功与势函数}
\label{th:availabilityprinciple}
设封闭系统与温度 $T_0$ 的大热库交换热量 $Q$，从初态变化到终态，对外输出总机械功 $W$。第一定律与整体熵不减分别给
\[
 W=Q-\Delta U,\qquad
 \Delta S-\frac Q{T_0}=S_{\rm gen}\ge0.
\]
消去 $Q$，得到
\[
 W=-\Delta(U-T_0S)-T_0S_{\rm gen}.
\]
在端点给定的条件下，可逆过程的输出功最大；实际功少去的量正是 $T_0S_{\rm gen}$。这里 $T_0$ 是外界恒温热库的温度，即使系统初态与热库不同温，该式也成立。只有系统端点本身处于 $T_0$ 时，$U-T_0S$ 才直接等于这些端点的 Helmholtz 自由能。

若系统还需推动压强为 $p_0$ 的环境，体积变化占用的功为 $p_0\Delta V$。定义可供负载使用的功为 $W_{\rm use}=W-p_0\Delta V$，则
\[
 W_{\rm use}
 =-\Delta(U+p_0V-T_0S)-T_0S_{\rm gen}.
\]
因此在同一 $T_0,p_0$ 下比较初末平衡态，最大有用功为 $-\Delta G$。没有其他功输出的自发过程满足 $\Delta G\le0$。固定体积、同一温度下则得到 $\Delta F\le0$。使用势函数判据时，应同时保留其恒温、定容或定压的环境条件。

\Diagram{fig:availabilitybalance}{能量变化、环境排热与有用功的分配}{
 \node[draw](s)at(0,0){系统初态 $\to$ 终态};
 \node[draw](e)at(4,1.4){环境 $T_0,p_0$};
 \node[draw](w)at(4,-1.3){机械负载};
 \draw[<->](s)--(e)node[midway,above left]{$Q,\ p_0\Delta V$};
 \draw[axis](s)--(w)node[midway,below left]{$W_{\rm use}$};
}
图\ref{fig:availabilitybalance}中，推动环境的功与负载取得的有用功分开计数。负载获得的功须扣除推动环境所占用的部分。指定端点和环境后，势函数差已经确定；实际过程中的熵产生再使有用输出减少。

\begin{example}[给定输入功对两个有限物体制冷]
\label{ex:finitefridge}
两个相同物体的热容均为常数 $C$，初温均为 $T_0$。用可逆装置在两者之间传热，同时向整体输入机械功 $W_{\rm in}>0$，外界不供热。求冷物体能达到的最低温度及热物体终温。
\Diagram{fig:finitefridgeq}{输入功使原先同温的两物体分离为冷热端}{
 \node[draw](c)at(0,0){$C,T_c$};
 \node[draw](m)at(3,0){可逆装置};
 \node[draw](h)at(6,0){$C,T_h$};
 \draw[axis](c)--(m);\draw[axis](m)--(h);
 \draw[axis](3,-1.6)--(m)node[midway,right]{$W_{\rm in}$};
 \node at(3,1.1){初始均为 $T_0$};
}
\end{example}
\begin{solution}
图\ref{fig:finitefridgeq}中机械功最终全部增加两物体总内能。装置循环复原，故
$C(T_h+T_c-2T_0)=W_{\rm in}$。令 $w=W_{\rm in}/C$，能量约束给 $T_h+T_c=2T_0+w$。可逆过程使两物体总熵不变，
\[
 C\ln\frac{T_h}{T_0}+C\ln\frac{T_c}{T_0}=0,
 \qquad T_hT_c=T_0^2.
\]
把 $T_h=2T_0+w-T_c$ 代入第二式，得到
$T_c^2-(2T_0+w)T_c+T_0^2=0$，较小根为冷端，
\[
 T_{c,\min}=T_0+\frac w2-\sqrt{T_0w+\frac{w^2}{4}},
 \qquad
 T_h=T_0+\frac w2+\sqrt{T_0w+\frac{w^2}{4}}.
\]
\Diagram{fig:finitefridgea}{能量直线与等熵曲线确定最冷端点}{
 \draw[axis](0,0)--(5,0)node[right]{$T_c$};
 \draw[axis](0,0)--(0,4.5)node[above]{$T_h$};
 \draw[curve]plot[domain=.6:4.4](\x,{2.25/\x});
 \draw[curve](.55,3.45)--(3.45,.55);
 \fill(.677,3.323)circle(1.4pt)node[right]{$A$};
 \fill(1.5,1.5)circle(1.3pt)node[below left]{$(T_0,T_0)$};
 \node at(3.25,3.5){$T_cT_h=T_0^2$};
}
图\ref{fig:finitefridgea}用一个正输入功的能量直线示意端点 $A$。若存在不可逆性，则同一输入功下仍固定温度和，而熵不减要求 $T_hT_c\ge T_0^2$。在冷支 $T_c<(T_h+T_c)/2$ 上，乘积随 $T_c$ 增加而增加，故任何不可逆终态的冷端温度都不低于上述小根。$W_{\rm in}\to0$ 时两根合并为 $T_0$；有限正功给正的冷端温度。在模型中要使 $T_c\to0$，所需输入功趋于无穷，且常热容近似在极低温前通常已失效。
\end{solution}

\begin{example}[有限供热速率下的最大功率]
\label{ex:finiteconductance}
恒温热源为 $T_h$，冷源为 $T_c<T_h$。热源通过热导 $K$ 向一台内部可逆热机供热，热机吸热端温度为 $T_x$，排热端与冷源可逆接触。假定供热率为 $\dot Q_h=K(T_h-T_x)$，忽略其他损失。求稳态最大输出功率及相应效率，说明它与卡诺效率上界的关系。
\Diagram{fig:finiteconductanceq}{有限温差供热与内部可逆热机}{
 \node[draw](h)at(0,0){$T_h$};
 \node[draw](x)at(3,0){热机 $T_x,T_c$};
 \node[draw](c)at(6,0){$T_c$};
 \draw[axis](h)--(x);\node at(1.5,.8){$K(T_h-T_x)$};
 \draw[axis](x)--(c);
 \draw[axis](3,-.4)--(3,-1.5)node[below]{$P$};
}
\end{example}
\begin{solution}
图\ref{fig:finiteconductanceq}中，热机本体在 $T_x,T_c$ 间可逆循环，效率为 $1-T_c/T_x$。供热接触层却跨越有限温差，故整个装置不可逆。功率为
\[
 P(T_x)=K(T_h-T_x)\left(1-\frac{T_c}{T_x}\right)
 =K\left(T_h+T_c-T_x-\frac{T_hT_c}{T_x}\right).
\]
允许区间为 $T_c<T_x<T_h$。求导，
\[
 P'=K\left(-1+\frac{T_hT_c}{T_x^2}\right),\qquad
 P''=-\frac{2KT_hT_c}{T_x^3}<0.
\]
唯一极大点是 $T_x=\sqrt{T_hT_c}$，故
\[
 P_{\max}=K(\sqrt{T_h}-\sqrt{T_c})^2,\qquad
 \eta_{\rm mp}=1-\sqrt{\frac{T_c}{T_h}}.
\]
\Diagram{fig:finiteconductancea}{最大功率位于两个零功率边界之间}{
 \draw[axis](0,0)--(5.5,0)node[right]{$T_x$};
 \draw[axis](0,0)--(0,2.5)node[above]{$P$};
 \draw[curve]plot[domain=1:4](\x,{2*(5-\x-4/\x)});
 \draw[guide](2,0)--(2,2);
 \node[below]at(1,0){$T_c$};
 \node[below]at(2,0){$\sqrt{T_hT_c}$};
 \node[below]at(4,0){$T_h$};
}
图\ref{fig:finiteconductancea}示意 $T_h/T_c=4$。当 $T_x\to T_h$，效率趋近卡诺上界，但供热率趋于零；当 $T_x\to T_c$，效率趋零。最大功率介于两者之间。所得 $\eta_{\rm mp}$ 只属于题设的单侧线性传热模型，若冷端也有有限热导或传热律改变，最优温度与效率均须重新求解。
\end{solution}


\section{过程设计与可用功}
比较传热与做功方案时，应先约定允许的操作。外界热源、机械储能和选择膜都会改变可达末态。以下各题分别给出这些约束，再用能量与熵确定过程的限度。

\begin{example}[分批接触为何仍有熵产生]
热容均为 $C$ 的热水和冷水初温为 $T_h>T_c$。仍采用热水等分 $n$ 份、逐份与全部冷水接触后移走的协议。求有限 $n$ 的总熵增表达式，并证明 $n\to\infty$ 时过程仍不可逆。
\Diagram{fig:batchentropyq}{各份热水最终温度不同，必须逐份计熵}{
 \node[draw]at(0,0){冷水 $C,T_n$};\foreach \x/\s in{3.1/1,5.5/2,7.9/n}{\node[draw,align=center]at(\x,0){第\s 份热水\\$C/n,T_{\s}$};}\node at(3.9,-1){末态包含冷水与全部移走的热水};
}
\end{example}
\begin{solution}
设 $r=n/(n+1)$，前文递推给 $T_j=T_h+(T_c-T_h)r^j$。图\ref{fig:batchentropyq}中，第 $j$ 份热水移走时的温度为 $T_j$，以后不再变化；只有一直保留的冷水最终达到 $T_n$。因此
\[
 \Delta S_n=C\ln\frac{T_n}{T_c}
 +\frac Cn\sum_{j=1}^n\ln\frac{T_j}{T_h}.
\]
有限 $n$ 时，每次有有限温差接触，故每段熵增都非负。考察连续分份极限，令累计使用份额 $s=j/n\in[0,1]$，则 $T(s)=T_h-(T_h-T_c)\ee^{-s}$，满足 $\dd T/\dd s=T_h-T$。
全部冷水和移走的热水共同构成绝热整体。连续极限下，热水各份的熵变之和成为$C\int_0^1\ln[T(s)/T_h]\dd s$，积分变量$s$是已使用的热容份额。

\Diagram{fig:batchentropya}{每个无穷小热水份仍从同一个高初温冷却}{
 \draw[axis](0,0)--(4.8,0)node[right]{$s$};\draw[axis](0,0)--(0,3)node[above]{$T$};\draw[curve]plot[domain=0:4](\x,{2.6-1.7*exp(-\x/4)});\draw[guide](0,2.6)--(4.2,2.6)node[right]{$T_h$};\node[left]at(0,.9){$T_c$};
}
图\ref{fig:batchentropya}中的一小份热水热容为 $C\dd s$。它冷却到当时冷水的温度，熵变为 $C\ln(T/T_h)\dd s$；冷水升温的熵变为 $C\dd T/T$。相加得
\[
 \frac{\dd S}{C\dd s}=\frac{T_h-T}{T}+\ln\frac T{T_h}
 =z-1-\ln z,\qquad z=T_h/T>1.
\]
因 $\ln z<z-1$，每一阶段仍有正熵产，所以极限总熵
\[
 \Delta S_\infty=C\ln\frac{T(1)}{T_c}
 +C\int_0^1\ln\frac{T(s)}{T_h}\dd s>0.
\]
每份热水虽趋于无穷小，仍须从$T_h$冷却到当前冷水温度；这段有限温差传热保留了熵产生。代入$100^\circ\mathrm C$与$0^\circ\mathrm C$的数据时，熵公式中的温度须取$373.15\,\mathrm K$和$273.15\,\mathrm K$。
\end{solution}

\begin{example}[不同热容物体的温度转移端点]
一个物体热容为 $C$、初温 $600\,\mathrm K$，另一物体热容为 $2C$、初温 $300\,\mathrm K$。允许前述可逆机、热泵与归零的储功装置，整体绝热且净外功为零。求第二物体最高温度、第一物体末温，以及两者同温时暂存的最大机械功。
\Diagram{fig:transfernumericq}{先取得机械功，再将其用于热泵}{
 \node[draw]at(0,0){$C,600\,\mathrm K$};\node[draw]at(5,0){$2C,300\,\mathrm K$};\node[draw]at(2.5,-1.3){可逆装置＋储功};\draw[<->](.8,-.3)--(1.6,-1);\draw[<->](4.2,-.3)--(3.4,-1);
}
\end{example}
\begin{solution}
图\ref{fig:transfernumericq}对应 $\xi=2$、$\eta=1/2$。令两末温分别为 $600x$、$600y$，能量与熵约束为
\[
 x+2y=2,\qquad x(2y)^2=1.
\]
消去 $x$，$2y^3-2y^2+1/4=0$。已知初态根为 $y=1/2$，多项式因式分解为 $2(y-1/2)(y^2-y/2-1/4)=0$。另一正根
\[
 y_{\max}=\frac{1+\sqrt5}{4},\quad x=\frac{3-\sqrt5}{2}.
\]
因此第二物体最高温度 $150(1+\sqrt5)\,\mathrm K\approx485.41\,\mathrm K$，第一物体末温 $300(3-\sqrt5)\,\mathrm K\approx229.18\,\mathrm K$。
因式分解后的二次因子还有一负根 $(1-\sqrt5)/4$，它对应负温度，不在本题模型范围内；初态根则只代表未作温度转移。所求正根使原热物体低于原冷物体初温，表明第二阶段必须反向作为热泵运行。两者同温时储功达到最大，因为沿总熵不变的轨迹，其总内能在这一点最小；继续把能量向第二物体集中，需要从储功中支付机械能，直至另一端点把它恰好用完。

\Diagram{fig:transfernumerica}{同熵曲线上暂存功先增加后减少到零}{
 \draw[axis](0,0)--(5,0)node[right]{过程进程};\draw[axis](0,0)--(0,2.7)node[above]{$W_{\rm store}$};\draw[curve](.4,0)to[out=70,in=180](2.5,2)to[out=0,in=110](4.6,0);\node[above]at(2.5,2){两体同温};\node[below]at(.4,0){初态};\node[below]at(4.6,0){转移端点};
}
两体同温时 $T_r=600^{1/3}300^{2/3}=600\,2^{-2/3}\,\mathrm K$。图\ref{fig:transfernumerica}顶点的储功为
\[
 W_{\rm store,max}=C[1200-1800\,2^{-2/3}]\,\mathrm K
 \approx66.07\,C\,\mathrm K.
\]
此后把这些功全部用于热泵，直到储功归零，才能达到另一等初熵端点。第一物体温度低于原冷物体初温并不矛盾：整个过程同时改变了两者，且并未把热量从低温送往高温而不付任何中间代价。
\end{solution}

\begin{example}[两级非理想热机的中间温度]
两台热机串联工作在恒温热源 $T_h>T_c$ 之间，中间热源温度为 $T_m$。每台实际效率均为同温区卡诺效率的 $q$ 倍，$0<q<1$。第一级排出的热量全部成为第二级的吸热。求总效率及使之最大的 $T_m$，再讨论 $q=1$。
\Diagram{fig:cascadeq}{两级热机通过中间热源传递热量}{
 \node[draw]at(0,2.2){$T_h$};\node[draw,circle]at(0,1){1};\node[draw]at(0,-.2){$T_m$};\node[draw,circle]at(0,-1.4){2};\node[draw]at(0,-2.6){$T_c$};
 \foreach \y in{1.8,.6,-.6,-1.8}{\draw[axis](0,\y)--(0,\y-.4);}\draw[axis](.5,1)--(2,1)node[right]{$W_1$};\draw[axis](.5,-1.4)--(2,-1.4)node[right]{$W_2$};
}
\end{example}
\begin{solution}
图\ref{fig:cascadeq}第一级吸热为 $Q_h$，其排热为 $Q_m=(1-\eta_1)Q_h$，第二级排热为 $Q_c=(1-\eta_2)Q_m$。因此总效率
\[
 \eta=1-(1-\eta_1)(1-\eta_2),
 \quad\eta_1=q(1-T_m/T_h),\quad\eta_2=q(1-T_c/T_m).
\]
展开剩余热量比例，
\[
 1-\eta=(1-q)^2+q^2\frac{T_c}{T_h}
 +q(1-q)\left(\frac{T_m}{T_h}+\frac{T_c}{T_m}\right).
\]
当$0<q<1$时，只有括号内的两项随$T_m$变化，且其系数$q(1-q)>0$。提高总效率等价于使这两项的和最小。

\Diagram{fig:cascadea}{需要最小化的中间温度函数}{
 \draw[axis](0,0)--(5,0)node[right]{$T_m$};\draw[axis](0,0)--(0,2.8)node[above]{$1-\eta$};\draw[curve](.7,2.3)to[out=-65,in=180](2.5,.8)to[out=0,in=230](4.5,2.1);\node[below]at(2.5,0){$\sqrt{T_hT_c}$};
}
图\ref{fig:cascadea}中唯一变项的导数为 $1/T_h-T_c/T_m^2$，二阶导为 $2T_c/T_m^3>0$，故最优 $T_m=\sqrt{T_hT_c}$。总效率最大值
\[
 \eta_{\max}=1-\left(1-q+q\sqrt{T_c/T_h}\right)^2.
\]
若 $q=1$，总效率对任何允许的中间温度均为 $1-T_c/T_h$。因此中间温度取几何平均是本题非理想效率模型的最优选择；串联可逆机的总效率则与这一选择无关。
\end{solution}

\begin{example}[导热活动隔板后的平衡状态]
绝热刚性容器内有一块可导热、无摩擦的活动隔板，两侧气体不能混合。隔板初始静止，热容忽略。两侧初态为$N_1,V_1,T_1$和$N_2,V_2,T_2$，均是常热容理想气体，系统定容热容为$C_1,C_2$。允许气体内部耗散使其最终静止平衡，求终温、终压、两侧终体积与熵增。
\Diagram{fig:conductpistonq}{隔板允许传热和位移，禁止粒子穿过}{
 \draw(0,0)rectangle(6,1.8);\draw[line width=2pt](2.5,0)--(2.5,1.8);\node[align=center]at(1.2,.9){$N_1,V_1,T_1$\\$C_1$};\node[align=center]at(4.3,.9){$N_2,V_2,T_2$\\$C_2$};\draw[<->](2,2.2)--(3,2.2);
}
\end{example}
\begin{solution}
图\ref{fig:conductpistonq}整体绝热且外壁不动。隔板内部做功和传热会改变各侧能量，却在整体能量式中相抵消。最终两侧同温 $T_f$，所以
\[
 T_f=\frac{C_1T_1+C_2T_2}{C_1+C_2}.
\]
无摩擦静止隔板要求两侧同压，$N_1\kb T_f/V_{1f}=N_2\kb T_f/V_{2f}$，再用总体积 $V=V_1+V_2$ 固定，得
\[
 V_{if}=\frac{N_i}{N_1+N_2}V,\qquad
 p_f=\frac{(N_1+N_2)\kb T_f}{V}.
\]
也可从熵极值推导终态。允许微小能量交换，总熵的相应导数为 $1/T_1-1/T_2$，平衡条件给出同温；再允许隔板微移，导数为 $(p_1-p_2)/T$，给出同压。结合整体能量守恒确定终温，再由状态方程确定隔板位置。两侧粒子数各自固定，同温、同压使最终体积比等于粒子数之比。

\Diagram{fig:conductpistona}{最终隔板位置按粒子数分配体积}{
 \draw(0,0)rectangle(6,1.6);\draw[line width=2pt](2,0)--(2,1.6);\node[align=center]at(1,.8){$T_f,p_f$\\$V_{1f}$};\node[align=center]at(4,.8){$T_f,p_f$\\$V_{2f}$};\node at(3,-.6){$V_{1f}/V_{2f}=N_1/N_2$};
}
图\ref{fig:conductpistona}中各组分仍被隔板隔开，熵变应按自己的终体积计算：
\[
 \Delta S=\sum_{i=1}^2\left[C_i\ln\frac{T_f}{T_i}
 +N_i\kb\ln\frac{V_{if}}{V_i}\right].
\]
活动隔板保留了两侧各自的粒子边界，所以熵公式使用 $V_{if}$；撤去隔板的混合题才让各组分占据总容积 $V$。此外，题设的静止终态依赖内部弛豫。若完全没有阻尼，隔板可能持续振荡，不能按已经达到平衡处理。
\end{solution}

\begin{example}[等温压强不均匀气体的最大功]
两侧理想气体用不透粒子的活动隔板分开，均由温度 $T_0$ 的同一热源维持恒温，粒子数 $N_1,N_2$ 固定。初体积为 $V_1,V_2$，总体积不变。通过可逆外部负载控制隔板移动至机械平衡，求最大输出功。
\Diagram{fig:pressureworkq}{隔板对外负载输出压差功}{
 \draw(0,0)rectangle(6,1.7);\draw[line width=2pt](2.4,0)--(2.4,1.7);\node at(1.2,.85){$p_1,T_0$};\node at(4.2,.85){$p_2,T_0$};\draw[axis](2.4,2)--(4,2)node[right]{负载};
}
\end{example}
\begin{solution}
图\ref{fig:pressureworkq}中左侧体积增加 $\dd V_1$，右侧减少同量。两侧对隔板的净功为 $(p_1-p_2)\dd V_1$，可逆负载将其全部输出。令总体积 $V$、变化变量为 $v=V_1$，则
\[
 W=\kb T_0\int_{V_1}^{V_{1f}}
 \left(\frac{N_1}v-\frac{N_2}{V-v}\right)\dd v.
\]
第二项令$w=V-v$，$\dd w=-\dd v$，于是原函数为$N_2\ln(V-v)$。两项合起来的端点差为$[N_1\ln v+N_2\ln(V-v)]_{V_1}^{V_{1f}}$，所以
\[
 W_{\max}=\kb T_0\left[N_1\ln\frac{V_{1f}}{V_1}
 +N_2\ln\frac{V_{2f}}{V_2}\right],
 \quad V_{if}=\frac{N_i}{N_1+N_2}V.
\]
功积分的上限先作为未知末位置处理，对其求导得到 $\dd W/\dd V_{1f}=p_{1f}-p_{2f}$。只有压差为零时，继续移动隔板不再能取得一阶正功。若初态恰已等压，上限就等于原位置，所有对数项相消，最大功为零；若初态不等压，可逆负载使隔板逐步释放压差功。

\Diagram{fig:pressureworka}{功作为隔板位置的函数在等压点最大}{
 \draw[axis](0,0)--(4.8,0)node[right]{$V_{1f}$};\draw[axis](0,0)--(0,2.8)node[above]{$W$};\draw[curve](.6,.4)to[out=60,in=180](2.5,2.2)to[out=0,in=115](4.3,.5);\node at(2.6,2.7){$p_1=p_2$};
}
图\ref{fig:pressureworka}的二阶导为 $-\kb T_0[N_1/v^2+N_2/(V-v)^2]<0$，所以等压点为唯一最大值。气体总内能因恒温不变，输出功全部来自热源供热 $Q=W_{\max}$；气体熵增为 $W_{\max}/T_0$，热源熵减恰好抵消。隔板不透气，两侧粒子数$N_1,N_2$始终不变；改变的是两侧体积和压强。只要取得了正功，气体终态便不同于初态，因而这不是单热源循环做功。
\end{solution}

\begin{example}[理想混合气体的可逆分离功]
温度 $T_0$、总压强 $p$ 的二元理想气体混合物，共 $\nu$ 摩尔，摩尔分数为 $x$、$1-x$。利用对各组分选择透过的理想膜，在同温同压条件下将其分离成纯组分。求最小输入功、放热量与总熵变化。
\Diagram{fig:separationq}{理想选择膜可分别限制两种分子的可达区域}{
 \draw(0,0)rectangle(5.5,1.7);\draw[dashed](1.8,0)--(1.8,1.7);\draw[dotted,line width=1pt](3.7,0)--(3.7,1.7);\node at(.9,.85){组分1};\node at(2.75,.85){混合};\node at(4.6,.85){组分2};
}
\end{example}
\begin{solution}
混合物总体积 $V=\nu RT_0/p$。图\ref{fig:separationq}中的理想膜允许分别控制两种组分，各组分起初都可达 $V$，分离后纯组分体积分别为 $xV$、$(1-x)V$。对每组分作可逆等温压缩，输入功
\[
 W_{\rm on,1}=\nu xRT_0\ln\frac1x,\quad
 W_{\rm on,2}=\nu(1-x)RT_0\ln\frac1{1-x}.
\]
每个组分的压缩功来自 $W_{{\rm on},i}=-\int_V^{V_i}p_i\dd V_i$，其中过程中的 $p_i=\nu_iRT_0/V_i$。积分上限小于下限，因此外界输入功为正。末态纯组分体积按其物质的量分配，使各自压强仍等于 $p$；总体积并未凭空减少。恒温压缩的内能不变，输入机械功须以热量排给恒温热源。对可逆选择膜取到的是最小功，真实膜阻力或混杂损失会增加所需功。

\Diagram{fig:separationa}{分离后的体积之和保持不变}{
 \draw(0,0)rectangle(5.5,1.6);\draw[line width=1.4pt](2,0)--(2,1.6);\node at(1,.8){$xV$};\node at(3.75,.8){$(1-x)V$};\node at(2.75,-.6){两纯组分均为 $T_0,p$};
}
图\ref{fig:separationa}末态的总内能与初态相同，因此最小功
\[
 W_{\min}=-\nu RT_0[x\ln x+(1-x)\ln(1-x)]
 =T_0\Delta S_{\rm mix}.
\]
系统向热源放热同量，系统熵减少 $\Delta S_{\rm mix}$，热源熵增加相同量，总熵变为零。自由混合能量不变，却使此后可逆分离需要功；混合不可逆性的代价体现在熵和自由能，而不一定体现在内能上。
\end{solution}

\begin{example}[直接传热损失的可用功]
两个常热容物体初温 $T_1,T_2$ 均高于环境温度 $T_0$。比较两种方案：分别通过可逆热机冷却到环境；先直接接触到 $T_d$，再共同通过可逆热机冷却到环境。求第二方案比第一方案少得到的功。
\Diagram{fig:lostworkq}{先直接接触会改变后续可用功}{
 \node[draw]at(0,1){$C_1,T_1$};\node[draw]at(0,-1){$C_2,T_2$};\draw[axis](1.2,1)--(3,0);\draw[axis](1.2,-1)--(3,0);\node[draw]at(4.3,0){$C_1+C_2,T_d$};
}
\end{example}
\begin{solution}
单个物体相对环境的最大功由习题中“有限热物体对无限冷源的最大功”的完整积分给出，$A_i=C_i[(T_i-T_0)-T_0\ln(T_i/T_0)]$。方案一得到$W_1=A_1+A_2$。图\ref{fig:lostworkq}的直接接触保持总内能，所以$(C_1+C_2)T_d=C_1T_1+C_2T_2$；方案二得到
\[
 W_2=(C_1+C_2)[(T_d-T_0)-T_0\ln(T_d/T_0)].
\]
两方案最终都降至$T_0$，初末总能量差相同。直接接触保持两物体总内能，却产生熵，改变了后续热机可取得的功。

\Diagram{fig:lostworka}{能量项相消后损失由熵产决定}{
 \node[draw]at(0,0){$W_1-W_2$};\node at(2.2,0){$=$};\node[draw]at(4.7,0){$T_0\Delta S_{\rm contact}$};
}
两式相减，线性能量项消去，得
\[
 W_1-W_2=T_0\left[(C_1+C_2)\ln T_d-C_1\ln T_1-C_2\ln T_2\right]
 =T_0\Delta S_{\rm contact}\ge0.
\]
图\ref{fig:lostworka}描述的是给定环境下的可用功减少。式中的$T_0$为环境温度；同样的熵产生在不同环境温度下，对应不同的输出功损失。
\end{solution}

\begin{example}[有限物体间的辐射换热时间]
两个常热容均为 $C$ 的均温物体构成绝热整体，彼此净辐射换热率为 $P=K(t_1^4-t_2^4)$，$K>0$，初温 $T_1>T_2$。求温差的时间关系与总熵变化，忽略传播时延和辐射场储能。
\Diagram{fig:radcontactq}{两有限物体仅通过净辐射交换能量}{
 \node[draw,minimum height=1cm]at(0,0){$C,t_1$};\node[draw,minimum height=1cm]at(5,0){$C,t_2$};\draw[decorate,decoration={snake,amplitude=1.2pt,segment length=8pt},->](1.2,0)--(3.8,0);\node at(2.5,.8){$K(t_1^4-t_2^4)$};
}
\end{example}
\begin{solution}
图\ref{fig:radcontactq}的能量式给 $C\dot t_1=-P$、$C\dot t_2=P$，所以 $t_1+t_2=2\bar T$ 恒定，$\bar T=(T_1+T_2)/2$。令 $u=(t_1-t_2)/(2\bar T)$，即 $t_{1,2}=\bar T(1\pm u)$。展开四次方差，$t_1^4-t_2^4=8\bar T^4u(1+u^2)$，从而
\[
 \dot u=-\frac{8K\bar T^3}C u(1+u^2).
\]
分离变量，用$1/[u(1+u^2)]=1/u-u/(1+u^2)$。后一项令$z=1+u^2$、$\dd z=2u\dd u$，故原函数为$\ln u-\tfrac12\ln(1+u^2)$。取$u_0\to u(t)$和$0\to t$作两侧上下限，得到
\[
 \ln\frac{u}{\sqrt{1+u^2}}-
 \ln\frac{u_0}{\sqrt{1+u_0^2}}
 =-\frac{8K\bar T^3}C t.
\]
为得到显式温差，可令 $Y(t)=u_0\ee^{-\gamma t}/\sqrt{1+u_0^2}$，其中 $\gamma=8K\bar T^3/C$。积分结果成为 $u/\sqrt{1+u^2}=Y$，两边平方并解出正根，得到 $u=Y/\sqrt{1-Y^2}$。正根由始终 $t_1\ge t_2$ 决定。温差很小时 $u\ll1$，分母近于1，才恢复通常的指数衰减；大温差时仍须保留这一非线性反解。

\Diagram{fig:radcontacta}{变换后的温差按指数规律衰减}{
 \draw[axis](0,0)--(4.8,0)node[right]{$t$};\draw[axis](0,0)--(0,2.8)node[above]{$u/\sqrt{1+u^2}$};\draw[curve]plot[domain=0:4.4](\x,{2.4*exp(-.7*\x)});
}
图\ref{fig:radcontacta}给出的是变换后的量，并非 $u$ 在大温差下严格指数衰减。总熵相对于初态为
\[
 \Delta S(t)=C\ln\frac{t_1t_2}{T_1T_2}
 =C\ln\frac{1-u^2}{1-u_0^2}.
\]
随 $u$ 下降，熵单调增加，最终熵增为 $C\ln[\bar T^2/(T_1T_2)]$。改变传热定律会改变到达过程的时间，却不改变在同一能量约束下的直接接触终态和总熵增。
\end{solution}

\section{连续导热中的温度场与熵产生}
\subsection{能量流守恒而可用功能沿途损失}
在稳态一维导热中，若无体热源，单位时间通过各截面的热量$J_Q$相同。截面积为$A$、导热系数为$\kappa(T)$时，傅里叶定律给
\[
 J_Q=-A\kappa(T)\frac{\dd T}{\dd x}.
\]
热量传递方向取$+x$，温度沿该方向下降。随热流携带的熵流为$J_Q/T$，其空间导数为
\[
 \frac{\dd}{\dd x}\frac{J_Q}{T}
 =J_Q\frac{\dd(1/T)}{\dd x}
 =\frac{A\kappa(T)}{T^2}\left(\frac{\dd T}{\dd x}\right)^2\geq0.
\]
右侧是每单位长度、每单位时间的熵产生。它不是额外的能量产生：热流在各处相同，但低温端同一热量对应的熵更多。整体积分给$\dot S_{\rm gen}=J_Q(1/T_c-1/T_h)$，与两热源熵变相同。

若环境温度为$T_e$，可逆热机从温度$T$处取得一小份热量$\delta Q$时，最多输出$(1-T_e/T)\delta Q$功。导热把热量从$T_h$降到$T_c$而没有输出这部分功，损失的最大做功机会为$T_e\dot S_{\rm gen}$。因此稳恒能量流的守恒与可用功的减少可以同时成立。
\ExpansionDiagram{fig:dep-th3-entropyflow}{同一热流在低温处携带更大的熵流}{
\draw(0,0)rectangle(5,1.2);\node[left]at(0,.6){$T_h$};\node[right]at(5,.6){$T_c$};\draw[->,thick](.5,.6)--(4.5,.6)node[midway,above]{$J_Q$};\draw[->](.5,-.5)--(1.7,-.5)node[midway,below]{$J_Q/T_h$};\draw[->](3,-.5)--(5,-.5)node[midway,below]{$J_Q/T_c$};}
图\ref{fig:dep-th3-entropyflow}的两条熵流箭头长度只示意大小变化。局部熵产生的位置取决于真实温度场和材料性质，不能仅按杆长平均分配。

\subsection{温变导热系数应在温度上积分}
当$\kappa$不是常数，稳态方程分离为$J_Q\dd x=-A\kappa(T)\dd T$。对全长积分，
\[
 J_Q=\frac A L\int_{T_c}^{T_h}\kappa(T)\dd T.
\]
所以可以定义区间平均导热系数，但它由积分决定，不一般等于中点温度处的值。若材料分层，各层热流相同、界面温度待求，应按各层各自的积分关系接续；总温差不能先按几何长度平分。

\begin{example}[两层温变导体的界面温度]
两段导体串联，总长$L$、截面积$A$，各长$L/2$。第一段位于热端，$\kappa_1(T)=k_1T/T_r$；第二段位于冷端，$\kappa_2(T)=k_2T/T_r$，其中$T_r$为固定参考温度。两端维持$T_h>T_c>0$，界面理想接触。求界面温度、稳态热流、各段熵产生及温度分布。
\ExpansionDiagram{fig:dep-th3-layers-q}{两层热流相同，界面温度由材料参数共同决定}{
\draw(0,0)rectangle(3,1.2);\fill[gray!15](3,0)rectangle(6,1.2);\draw(3,0)rectangle(6,1.2);\node at(1.5,.6){$k_1T/T_r$};\node at(4.5,.6){$k_2T/T_r$};\node[left]at(0,.6){$T_h$};\node[right]at(6,.6){$T_c$};\node[above]at(3,1.2){$T_m$};\draw[->](1,-.6)--(5,-.6)node[midway,below]{$J_Q$};}
\end{example}
\begin{solution}
图\ref{fig:dep-th3-layers-q}中第一段从$T_h$降至$T_m$。将傅里叶定律分离为
\[
 J_Q\dd x=-(Ak_1/T_r)T\dd T,
\]
从$x=0$到$L/2$积分，得$J_QL/2=-(Ak_1/T_r)[T^2/2]_{T_h}^{T_m}$，整理为
\[
 J_Q=\frac{Ak_1}{LT_r}(T_h^2-T_m^2).
\]
第二段同理为$J_Q=Ak_2(T_m^2-T_c^2)/(LT_r)$。令两式相等，得
\[
 T_m^2=\frac{k_1T_h^2+k_2T_c^2}{k_1+k_2},\qquad
 J_Q=\frac{Ak_1k_2}{LT_r(k_1+k_2)}(T_h^2-T_c^2).
\]
注意是$T^2$按系数加权。第一段从热端积分至任意$x$，有$J_Qx=-(Ak_1/2T_r)[T^2(x)-T_h^2]$，解出
\[
 T^2(x)=T_h^2-\frac{2T_rJ_Q}{Ak_1}x,
 \quad 0<x<L/2;
\]
第二段以界面为起点，
\[
 T^2(x)=T_m^2-\frac{2T_rJ_Q}{Ak_2}(x-L/2),
 \quad L/2<x<L.
\]
\ExpansionDiagram{fig:dep-th3-layers-a}{温度平方在各层中为直线，界面处斜率可以改变}{
\draw[->](0,0)--(6,0)node[right]{$x$};\draw[->](0,0)--(0,3.5)node[above]{$T^2$};\draw[thick](0,3)--(2.5,2.3)--(5,.6);\draw[dashed](2.5,0)--(2.5,2.3);\node[below]at(2.5,0){$L/2$};\node[below]at(5,0){$L$};\node[above right]at(2.5,2.3){$T_m^2$};}
图\ref{fig:dep-th3-layers-a}画$T^2$而非$T$，以显示分段线性结构；界面温度连续，热流连续，而温度梯度一般不连续。

两段的熵产生分别为
\[
 \dot S_1=J_Q(1/T_m-1/T_h),\quad
 \dot S_2=J_Q(1/T_c-1/T_m).
\]
相加后界面项抵消，得到$J_Q(1/T_c-1/T_h)$。若$k_1=k_2,T_h=400\,\mathrm K,T_c=300\,\mathrm K$，则$T_m=\sqrt{125000}\,\mathrm K\simeq353.55\,\mathrm K$，冷半段约贡献总熵产生的60.6\%。即使两层材料相同、长度相同，局部温度不同仍使熵产生不均匀。
\end{solution}

\subsection{有限接触热阻的熵产生不能漏计}
若两段接触处有温度跃变，热侧界面温度为$T_{m+}$、冷侧为$T_{m-}$，且热流为$J_Q$，界面熵产生是
\[
 \dot S_{\rm contact}=J_Q(1/T_{m-}-1/T_{m+})>0.
\]
此时应分别对两段内部积分，再加上这个界面项。若只积分两段光滑温度梯度，会遗漏温度跳跃处的贡献。若把接触层视为很薄但导热系数也很小的连续层，积分式在薄层极限恰给出同一个结果。

热源的熵变、导体内部及接触界面的熵产生，应分别在相应位置计算。导体处于稳态时，流出的熵比流入的熵多，恰好带走内部产生的部分，所以导体总熵保持不变。整体熵平衡须同时计入流入、流出和产生项。

\section{习题与解答}
\begin{exercise}[可逆制冷机的功与制热量]
\enlargethispage{\baselineskip}
制冷机从温度 $T_c$ 的冷库吸热 $Q_c$，向温度 $T_h>T_c$ 的环境放热。求最小耗功和放热量；若同一装置作为热泵，其制热系数是多少？
\Diagram{fig:fridgeq}{制冷与制热使用同一能量流}{
 \node[draw]at(0,1.5){$T_h$};\node[draw,circle]at(0,0){逆};\node[draw]at(0,-1.5){$T_c$};\draw[axis](0,-1.1)--(0,-.45);\draw[axis](0,.45)--(0,1.1);\draw[axis](-2,0)--(-.5,0)node[midway,above]{$W$};
}
\end{exercise}
\begin{solution}
图\ref{fig:fridgeq}的可逆条件为 $Q_h/T_h=Q_c/T_c$，因此 $Q_h=Q_cT_h/T_c$。由 $Q_h=Q_c+W$，最小耗功 $W_{\min}=Q_c(T_h-T_c)/T_c$。不可逆过程产生熵，需付出更多功才能转移相同冷量。
最小功的依据也可由熵不等式直接写出：一次循环后热源总熵变 $Q_h/T_h-Q_c/T_c\ge0$，因此给定冷量时至少要向高温源排出 $Q_cT_h/T_c$。第一定律再把多排出的部分与输入功对应。制冷系数用的是搬出的冷量，制热系数用的是送入热端的总热量，后者还包含机械输入最终转化的热，所以两种性能系数恰差1，而不是两种不同的能量守恒规律。

\Diagram{fig:fridgea}{同一耗功对应的两种性能定义}{
 \node[draw]at(0,0){制冷系数 $Q_c/W$};\node[draw]at(5,0){制热系数 $Q_h/W$};\node at(2.5,-.9){两者相差1};
}
图\ref{fig:fridgea}中$\mathrm{COP}_{\rm cool}=T_c/(T_h-T_c)$，$\mathrm{COP}_{\rm heat}=T_h/(T_h-T_c)$。两者比较的是搬运热量与输入功，可以大于1。
\end{solution}

\begin{exercise}[等热容有限热源与直接传热的损失]
两个热容均为 $C$ 的物体初温 $T_h,T_c$。比较直接接触与可逆做功的终温，求最大功，并在 $T_h=T_0+\Delta T/2$、$T_c=T_0-\Delta T/2$ 时给出小温差近似。
\Diagram{fig:equalfiniteq}{初温关于平均温度对称}{
 \draw[axis](0,0)--(5,0)node[right]{$T$};\foreach \x/\s in{1/T_c,2.5/T_0,4/T_h}{\draw(\x,-.12)--(\x,.12);\node[above]at(\x,.15){$\s$};}\draw[<->](1,-.5)--(4,-.5)node[midway,below]{$\Delta T$};
}
\end{exercise}
\begin{solution}
图\ref{fig:equalfiniteq}中直接接触终温 $T_d=(T_h+T_c)/2=T_0$，可逆机终温 $T_r=\sqrt{T_hT_c}$。最大功
\[
 W_{\max}=C(T_h+T_c-2\sqrt{T_hT_c})
 =C(\sqrt{T_h}-\sqrt{T_c})^2.
\]
最大功中的平方根可以先合成 $T_hT_c=T_0^2-\Delta T^2/4$。当 $|\Delta T|\ll T_0$ 时，展开 $\sqrt{1-z}=1-z/2+O(z^2)$，这里 $z=\Delta T^2/(4T_0^2)$ 已是二阶小量，所以保留它的一次项就足以求出温差二阶的功。初温差为零时，两体本就等温，最大功也为零；两者温度对调则功不变，与结果只含偶次温差相符。

\Diagram{fig:equalfinitea}{温差较小时可用功是二阶小量}{
 \draw[axis](0,0)--(4.8,0)node[right]{$\Delta T$};\draw[axis](0,0)--(0,2.7)node[above]{$W_{\max}$};\draw[curve]plot[domain=0:4](\x,{.12*\x*\x});\node at(2.7,2.5){$C\Delta T^2/(4T_0)$};
}
用 $T_r=T_0\sqrt{1-\Delta T^2/(4T_0^2)}\simeq T_0-\Delta T^2/(8T_0)$，有 $W_{\max}\simeq C\Delta T^2/(4T_0)$。图\ref{fig:equalfinitea}中可用功从温差的二阶项开始增长。热机运行时两个有限物体都在变温，能输出的功远小于直接以 $C\Delta T$ 估算的传热尺度。
\end{solution}

\begin{exercise}[有限热物体对无限冷源的最大功]
热容为常量 $C$ 的物体初温 $T_i>T_0$，允许通过可逆热机向温度恒为 $T_0$ 的无限热源放热。求最大功与冷源得到的热量。
\Diagram{fig:availabilityq}{有限热物体逐渐降至环境温度}{
 \node[draw]at(0,1.4){$C,T$};\node[draw,circle]at(0,0){机};\node[draw]at(0,-1.4){$T_0$ 恒定};\draw[axis](.5,0)--(2,0)node[right]{$\dd W$};\draw[axis](0,1)--(0,.45);\draw[axis](0,-.45)--(0,-1);
}
\end{exercise}
\begin{solution}
图\ref{fig:availabilityq}中物体降温 $\dd T<0$ 时释放热量 $-C\dd T$，瞬时可逆效率 $1-T_0/T$，所以
\[
 W_{\max}=\int_{T_0}^{T_i}C\left(1-\frac{T_0}T\right)\dd T
 =C[T-T_0\ln T]_{T_0}^{T_i}
 =C[(T_i-T_0)-T_0\ln(T_i/T_0)].
\]
积分变量是热物体当时的温度，因其放热而下降，先将吸热量写成 $-C\dd T$，再反转积分上下限。冷源收到的热量可逐次写为 $(-C\dd T)T_0/T$，从初温积分到环境温度便产生对数项。热物体的温度越接近 $T_0$，瞬时卡诺效率越小；虽然它仍在放热，越来越多的热量必须排给环境，可用于做功的比例趋于零。

\Diagram{fig:availabilitya}{总释放能量分为输出功与向环境放热}{
 \node[draw]at(0,0){$C(T_i-T_0)$};\draw[->](1.5,.2)--(3,1);\draw[->](1.5,-.2)--(3,-1);\node[right]at(3,1){$W_{\max}$};\node[right]at(3,-1){$CT_0\ln(T_i/T_0)$};
}
由图\ref{fig:availabilitya}的能量分配，环境吸热 $Q_0=CT_0\ln(T_i/T_0)$。环境熵增加 $Q_0/T_0$，恰好抵消物体熵减少 $C\ln(T_0/T_i)$，总熵为零。直接把所有 $C(T_i-T_0)$ 视作可用功，忽略了必须排给环境的热量。
\end{solution}

\begin{exercise}[混合物组成与混合熵]
温度和压强相同的 $s$ 种不同理想气体混合，总物质的量为 $\nu$，各组分摩尔分数 $x_i$。在总温度、总压强不变的条件下推导混合熵，并求二元混合物何时最大。
\Diagram{fig:mixfractionq}{各纯组分原来占据与粒子数成正比的体积}{
 \draw(0,0)rectangle(6,1.6);\draw(1.5,0)--(1.5,1.6);\draw(4,0)--(4,1.6);\node at(.75,.8){$x_1V$};\node at(2.75,.8){$x_2V$};\node at(5,.8){$\cdots$};
}
\end{exercise}
\begin{solution}
同温同压时图\ref{fig:mixfractionq}的 $V_i=x_iV$。混合后每种组分可达体积由 $x_iV$ 增至 $V$，温度不变，所以
\[
 \Delta S_{\rm mix}=\sum_i\nu x_iR\ln\frac1{x_i}
 =-\nu R\sum_i x_i\ln x_i.
\]
体积比$V_i/V=x_i$来自同温同压下的$V_i=\nu_iRT_0/p$与$V=\nu RT_0/p$。组分$i$的体积熵差为$\nu_iR\ln(V/V_i)$，逐组相加便得到上式。交换二元体系的两个组分相当于$x\mapsto1-x$，熵曲线因而关于$x=1/2$对称。

\Diagram{fig:mixfractiona}{二元混合熵随组分比例的变化}{
 \draw[axis](0,0)--(5,0)node[right]{$x$};\draw[axis](0,0)--(0,2.6)node[above]{$\Delta S/(\nu R)$};\draw[curve]plot[domain=.001:.999,samples=70]({4.5*\x},{3*(-\x*ln(\x)-(1-\x)*ln(1-\x))});\node[below]at(2.25,0){$1/2$};
}
二元情形 $\Delta S/(\nu R)=-x\ln x-(1-x)\ln(1-x)$，一阶导为 $\ln[(1-x)/x]$，二阶导为 $-1/x-1/(1-x)<0$。故图\ref{fig:mixfractiona}在 $x=1/2$ 取得唯一最大值 $\nu R\ln2$。端点取 $x\ln x\to0$，纯物质没有不同组分的混合熵。
\end{solution}

\begin{exercise}[三物体的最大做功]
三个常热容物体 $C_i$、初温 $T_i$ 可通过可逆热机互相交换能量，整体绝热。求最终共同温度和最大输出功。说明先后连接次序是否影响可逆终态。
\Diagram{fig:threefiniteq}{三有限物体之间允许可逆交换}{
 \node[draw](a)at(0,1){$C_1,T_1$};\node[draw](b)at(4,1){$C_2,T_2$};\node[draw](c)at(2,-1){$C_3,T_3$};\draw[<->](a)--(b);\draw[<->](b)--(c);\draw[<->](c)--(a);
}
\end{exercise}
\begin{solution}
图\ref{fig:threefiniteq}中的各微循环可逆，总熵保持不变，故 $\sum_iC_i\ln(T_f/T_i)=0$。整理得到
\[
 T_f=\prod_{i=1}^3T_i^{C_i/C_\Sigma},\qquad C_\Sigma=\sum_iC_i.
\]
各热机作为辅助装置均循环复原，因而整体熵只需统计三个物体。把熵守恒式中的 $\ln T_f$ 合并，其系数是总热容 $C_\Sigma$；除以它后再取指数，得到权重和为1的几何平均。功的能量零点可任选，因为每个物体的参考常数在初末态相减时抵消。若只让其中两个物体先达到共同温度，再与第三个可逆交换，只要没有直接接触损失，所得权重仍会合并为同一组 $C_i/C_\Sigma$。

\Diagram{fig:threefinitea}{最终温度由全部初始熵确定}{
 \node[draw]at(0,0){$\sum_iC_i\ln T_i$};\draw[axis](2,0)--(3.2,0);\node[draw]at(5,0){$C_\Sigma\ln T_f$};
}
初末能量差给 $W_{\max}=\sum_iC_iT_i-C_\Sigma T_f$。图\ref{fig:threefinitea}的熵约束只含端点；全部交换可逆且最终同温时，连接次序不影响 $T_f$ 和总功。若中途直接接触，产生的熵会使最终共同温度升高，相应地减少输出功。
\end{solution}

\begin{exercise}[接触过程的熵产生随温差的展开]
两个热容 $C_1,C_2$ 的物体初温均接近 $T_0$，记差值为 $\Delta T$，并令能量加权平均温度恰为 $T_0$。求直接接触的熵增到 $\Delta T^2$ 阶。
\Diagram{fig:smallentropyq}{偏离平均温度的两端由热容加权约束}{
 \draw[axis](0,0)--(5.5,0)node[right]{$T$};\node[above]at(1,0){$T_2$};\node[above]at(3,0){$T_0$};\node[above]at(4.5,0){$T_1$};\draw[<->](1,-.5)--(4.5,-.5)node[midway,below]{$\Delta T$};
}
\end{exercise}
\begin{solution}
按图\ref{fig:smallentropyq}设 $T_1=T_0+\delta_1$、$T_2=T_0+\delta_2$，则 $C_1\delta_1+C_2\delta_2=0$、$\delta_1-\delta_2=\Delta T$。解得 $\delta_1=C_2\Delta T/(C_1+C_2)$、$\delta_2=-C_1\Delta T/(C_1+C_2)$。
二阶项的系数可先计算分子
\[
 C_1\delta_1^2+C_2\delta_2^2
 =\frac{C_1C_2^2+C_2C_1^2}{(C_1+C_2)^2}\Delta T^2
 =\frac{C_1C_2}{C_1+C_2}\Delta T^2.
\]
一阶项因为加权能量平均取在 $T_0$ 而精确消失，二阶项则保留两热容的组合。对等热容情形，它给 $C\Delta T^2/(4T_0^2)$；乘环境温度 $T_0$ 后，与相应小温差最大功的二阶量级一致。

\Diagram{fig:smallentropya}{对数展开的一阶项因能量约束相消}{
 \node[draw]at(0,0){$-\sum_iC_i\delta_i/T_0=0$};\node[draw]at(5,0){$\sum_iC_i\delta_i^2/(2T_0^2)>0$};
}
熵增 $\Delta S=-\sum_iC_i\ln(1+\delta_i/T_0)$。利用 $\ln(1+x)=x-x^2/2+O(x^3)$，图\ref{fig:smallentropya}的一阶项消去，二阶项为
\[
 \Delta S=\frac{C_1C_2}{2(C_1+C_2)T_0^2}\Delta T^2+O(\Delta T^3).
\]
熵产随小温差的平方增长，与微元传热公式中的倒温差和累计传热量各贡献一个温差相吻合。
\end{solution}

\begin{exercise}[有限热源与相变冷源]
热容为常量 $C$ 的热物体初温 $T_i>T_m$。冷物体处在温度 $T_m$ 的固液共存态，尚可吸收潜热 $L_{
\rm avail}$ 而不升温。二者之间可逆做功。求热物体降到 $T_m$ 所需冷源潜热，及潜热不足时停止在何温度。
\Diagram{fig:latentengineq}{相变冷源在潜热耗尽前保持恒温}{
 \node[draw]at(0,1.4){$C,T$};\node[draw,circle]at(0,0){机};\node[draw]at(0,-1.4){固液共存 $T_m$};\draw[axis](0,1)--(0,.4);\draw[axis](0,-.4)--(0,-1);\draw[axis](.5,0)--(2,0)node[right]{$W$};
}
\end{exercise}
\begin{solution}
图\ref{fig:latentengineq}中每次从热体取热$-C\dd T$，可逆机向冷源排热$\delta Q_c=(-C\dd T)T_m/T$。降至$T_m$时，总排热为$-CT_m\int_{T_i}^{T_m}\dd T/T=-CT_m[\ln T]_{T_i}^{T_m}$，即
\[
 Q_{c,\rm req}=CT_m\ln\frac{T_i}{T_m}.
\]
可用潜热不足时，由 $L_{\rm avail}<CT_m\ln(T_i/T_m)$ 立刻得到 $T_*>T_m$，说明热物体尚未降到冷源温度，相变库存就已耗尽。若可用潜热恰为零，公式回到 $T_*=T_i$、这一阶段做功为零；若恰等于所需潜热，则停止温度正好为 $T_m$。这些边界反映的是同一个能量和可逆熵约束，而不是把冷物体在相变后仍无限维持在熔点。

\Diagram{fig:latentenginea}{可逆排热累计到可用潜热上限}{
 \draw[axis](0,0)--(5,0)node[right]{$\ln(T_i/T)$};\draw[axis](0,0)--(0,2.8)node[above]{$Q_c$};\draw[curve](0,0)--(4.2,2.2);\draw[guide](0,1.5)--(4.4,1.5)node[right]{$L_{\rm avail}$};
}
若 $L_{\rm avail}\ge Q_{c,\rm req}$，最大功为 $C(T_i-T_m)-CT_m\ln(T_i/T_m)$。若不足，图\ref{fig:latentenginea}先触及潜热上限，停止温度满足 $CT_m\ln(T_i/T_*)=L_{\rm avail}$，即 $T_*=T_i\ee^{-L_{\rm avail}/(CT_m)}$。在相变冷源模型有效阶段的功为 $C(T_i-T_*)-L_{\rm avail}$；其后是否还能做功，需给出冷源相变结束后的热容，不能继续假定它恒温。
\end{solution}

\chapter{相变、界面与外场平衡}

\section{选择状态变量与守恒的能量}
\label{th:coordinates}
本讲将前面的状态函数方法用于界面和外场。简单可压缩物质的体积功为$-p\dd V$；界面变化产生表面功，外场还会引入势能，须先写清系统所含的能量与全部广义功。第~\ref{th:energyidentity}节已用链式法则证明$U_V|_T=-p+Tp_T|_V$，其中等温与等熵偏导取沿不同的状态方向；改用新的坐标时，仍须分别固定其余变量。

体积偏导只确定$U$的体积依赖，积分后留下的函数$\Phi(T)$还要由定容热容确定。下面先处理总量与摩尔量，再把同一方法用于表面面积及外场参数。

\subsection{广义坐标与相应的功}
\label{th:generalizedwork}
体积、界面面积和外场参数改变时，能量中出现的功项各有不同。把一个可控参数记为 $a$，其共轭广义力记为 $Y$，对固定粒子数的平衡态可写
\[
 \dd U=T\dd S-p\dd V+Y\dd a.
\]
这里 $Y\dd a$ 按外界对系统做功为正。界面以 $a=\Sigma$ 为坐标时 $Y=\sigma$；改变均匀重力场强度时，若系统能量计入重力势能，$a=g$，则 $Y$ 是全部粒子的质量乘高度之和的平均值。共轭量的符号由所选参数及功的约定确定。

定义 $F=U-TS$，相同的参数项保留为
\[
 \dd F=-S\dd T-p\dd V+Y\dd a.
\]
固定其他变量后，有 $Y=F_a|_T$、$S_a|_T=-Y_T|_a$。进一步改用 $T,a$ 作变量，
\[
 \left(\frac{\partial U}{\partial a}\right)_T
 =Y-T\left(\frac{\partial Y}{\partial T}\right)_a.
\]
它把第二讲的体积公式与表面公式统一起来：取 $a=V,Y=-p$，体积系数为 $-p+Tp_T$；取 $a=\Sigma,Y=\sigma$，面积系数为 $\sigma-T\sigma_T$。两个符号差别源于两种广义功的定义。

\Diagram{fig:generalizedcoordinates}{不同约束方向对应不同的能量系数}{
 \node[draw](u)at(0,0){$U(S,V,a)$};
 \node[draw](s)at(4,1.6){熵方向：$T\dd S$};
 \node[draw](v)at(4,0){体积方向：$-p\dd V$};
 \node[draw](a)at(4,-1.6){参数方向：$Y\dd a$};
 \draw[axis](u)--(s);\draw[axis](u)--(v);\draw[axis](u)--(a);
}
图\ref{fig:generalizedcoordinates}中，每一方向都规定其他独立坐标保持不变。若实际装置同时改变体积和外场参数，两个功项应同时进入第一定律；在实验室能量与旋转参照系的有效能量之间切换时，也要保留变换所引入的角动量项。

\subsection{从微观能量求广义力}
设可分辨的微观状态能量为 $E_i(a)$，没有显式温度依赖。固定 $T,a$ 的配分和为 $Z=\sum_i\ee^{-\beta E_i(a)}$。对参数求导，
\[
 \left(\frac{\partial\ln Z}{\partial a}\right)_T
 =-\beta\left\langle\frac{\partial E_i}{\partial a}\right\rangle,
 \qquad
 \left(\frac{\partial F}{\partial a}\right)_T
 =\left\langle\frac{\partial E_i}{\partial a}\right\rangle.
\]
逐项求导时，$Z_a=\sum_i(-\beta E_{i,a})\ee^{-\beta E_i}$；除以$Z$就得到第一式。再用$F=-\kb T\ln Z$，有$F_a=-\kb T(-\beta\langle E_{i,a}\rangle)=\langle E_{i,a}\rangle$。因此微观参数变化所引起的能量增量，平均后正是宏观广义功。在经典连续相空间中，把求和改成积分也得到同样关系，前提是积分区域和测度的参数依赖已正确处理。

同一配分和还可给出熵。能量$E_i$不显含温度时，固定外参有
\[
 \left(\frac{\partial\ln Z}{\partial T}\right)_a
 =\frac1Z\sum_i\frac{E_i}{\kb T^2}\ee^{-E_i/(\kb T)}
 =\frac{\langle E\rangle}{\kb T^2}.
\]
对$F=-\kb T\ln Z$用乘积法则，得到
\begin{equation}\label{eq:canonicalentropy}
 \begin{aligned}
 S=-\left(\frac{\partial F}{\partial T}\right)_a
 &=\kb\ln Z+\kb T\left(\frac{\partial\ln Z}{\partial T}\right)_a\\
 &=\kb\ln Z+\frac{\langle E\rangle}{T}.
 \end{aligned}
\end{equation}
对$N$个独立经典粒子，$Z_N=Z_1^N/N!$；固定$N$比较状态时，$\ln N!$为常数，在熵差中相消。因而可由单粒子配分积分求出每粒子的熵差，再乘$N$。

若能量零点只是与$a$无关的常数，占有概率、热容与广义力都不变，平均能量只平移同一常数；若减去一个随$a$变化的函数，则$Y$也随之改变。外场参数可变时，不能在不同参数状态各自任选势能零点后再把差值丢掉。例如重力势能应始终相对于同一固定几何高度定义，才能正确计算改变$g$所做的功。

\subsection{广延性、摩尔量与附加界面能}
均匀体相系统在短程相互作用、边界效应可忽略时满足
$U(\lambda S,\lambda V,\lambda N)=\lambda U(S,V,N)$。
对 $\lambda$ 在1处求导，得到 Euler 关系
\[
 U=TS-pV+\mu N.
\]
对右侧三个乘积分别微分，得
\[
 \dd U=T\dd S+S\dd T-p\dd V-V\dd p+\mu\dd N+N\dd\mu.
\]
与$\dd U=T\dd S-p\dd V+\mu\dd N$比较，相同项相消后便有$S\dd T-V\dd p+N\dd\mu=0$。单组分系统的Gibbs自由能为$G=\mu N$；用摩尔量时，化学势按每摩尔定义，便等于摩尔Gibbs自由能。

含显著界面能时，按相同形状把线尺度放大 $\ell$ 倍，体积乘 $\ell^3$，面积却只乘 $\ell^2$。因而把小液滴当作只有体积的均匀广延系统，会漏掉界面项；液滴内外压差与蒸气压的曲率修正都由此产生。远离界面的液相和气相仍按各自的体相摩尔关系计算，界面自由能另行加入，并与面积、体积的几何约束联立。


\begin{example}[总量与摩尔量：实际气体的完整状态函数]\label{ex:vdwmolar}
$\nu$ 摩尔 van der Waals 气体的摩尔体积为 $v=V/\nu$，满足 $(p+a/v^2)(v-b)=RT$。给定常量 $C_{V,m}=c$，推导摩尔内能、摩尔熵和可逆绝热线，再写回总量形式。说明总容积增大与物质的量增大有什么区别。
\Diagram{fig:molarq}{固定温度下分别改变体积和系统尺度}{
 \draw(0,0)rectangle(2,1.4);\node at(1,.7){$\nu,V$};\draw[axis](2.5,.7)--(3.5,.7);\draw(4,0)rectangle(7,1.4);\node at(5.5,.7){$\lambda\nu,\lambda V$};\node at(3.5,-.7){尺度扩大时 $v=V/\nu$ 不变};
}
\end{example}
\begin{solution}
在固定 $\nu$ 下对摩尔体积求偏导，$\partial U_m/\partial v|_T=Tp_T-p=a/v^2$，且 $\partial U_m/\partial T|_v=c$。因此
\[
 U_m-U_{m0}=c(T-T_0)+a(1/v_0-1/v),
\]
\[
 S_m-S_{m0}=c\ln(T/T_0)+R\ln[(v-b)/(v_0-b)].
\]
等熵时 $(v-b)T^{c/R}$ 恒定，等价于 $T(v-b)^{R/c}$ 恒定。图\ref{fig:molarq}保持 $v$ 的尺度变化下，$U$、$S$ 随 $\nu$ 成比例。
从摩尔式写回总量时，先用 $U=\nu U_m$、$S=\nu S_m$，再把所有摩尔体积都换成 $V/\nu$。内能吸引项因而由 $-\nu a/v$ 变成 $-\nu^2a/V$；熵中的体积比则让分子、分母共同含有的 $1/\nu$ 约去。这里的初末态比较始终固定物质的量。随后讨论的 $\lambda$ 倍尺度扩大，是把参考态也复制相同倍数的另一种操作，不能在同一个固定粒子数体积偏导中把 $\nu$ 一起改变。

\Diagram{fig:molara}{摩尔式写回总量后吸引项与排除体积的变化}{
 \node[draw]at(0,0){$a/v^2$};\draw[axis](.8,0)--(2.1,0);\node[draw]at(3.5,0){$\nu^2a/V^2$};
 \node[draw]at(0,-1){$v-b$};\draw[axis](.8,-1)--(2.1,-1);\node[draw]at(3.5,-1){$(V-\nu b)/\nu$};
}
以相同物质的量的参考态比较，
\[
 \Delta U=\nu c\Delta T+\nu^2a(1/V_0-1/V),
\]
\[
 \Delta S=\nu c\ln(T/T_0)+\nu R\ln[(V-\nu b)/(V_0-\nu b)].
\]
图\ref{fig:molara}说明 $a$ 的总容积形式含 $\nu^2$，而排除体积为 $\nu b$。固定 $\nu$ 的等温膨胀会改变摩尔体积和势能；同时放大 $V,\nu$ 则只复制系统，不能把两种变化混用在同一偏导中。
\end{solution}

\subsection{从辐射压强独立推出四次方律}
\label{th:photonlaw}
第二讲由辐射通量的实验温度律写出 $u\propto T^4$。现在反过来，只用各向同性光子的 $p=u/3$、均匀辐射能的广延性 $U=Vu(T)$ 和热力学恒等式。固定温度改变体积时 $U_V=u$，而
\[
 u=-p+Tp_T=-\frac u3+\frac T3\frac{\dd u}{\dd T}.
\]
于是$T\dd u/\dd T=4u$。对$u>0,T>0$分离变量，$\int_{u_0}^{u}\dd u'/u'=4\int_{T_0}^{T}\dd\vartheta/\vartheta$，得到$\ln(u/u_0)=4\ln(T/T_0)$，即$u=aT^4$、$a=u_0/T_0^4$。热力学只能确定幂指数，常数$a$需要辐射谱或实验给出。再由$J_E=cu/4$，得到Stefan--Boltzmann通量律及$\sigma_{\rm SB}=ca/4$。

\begin{example}[辐射腔的可逆绝热关系]\label{ex:radiationadiabat}
腔内只有平衡光子气体，体积由 $V_1$ 缓慢可逆地增至 $V_2$，不与外界换热。求温度、压强、总能量及做功的变化，并与单原子理想气体比较。
\Diagram{fig:radadiabatq}{反射活塞使辐射腔缓慢膨胀}{
 \draw(0,0)rectangle(3.5,1.6);\draw[line width=2pt](3.5,0)--(3.5,1.6);\draw[axis](3.5,.8)--(5,.8);\node at(1.75,.8){$U=aVT^4$};\node at(1.75,-.5){$Q=0$，过程可逆};
}
\end{example}
\begin{solution}
图\ref{fig:radadiabatq}的辐射熵为 $S=4aVT^3/3$。可逆绝热时熵不变，所以
\[
 TV^{1/3}=\text{常量},\quad pV^{4/3}=\text{常量},\quad
 U=aVT^4\propto V^{-1/3}.
\]
因此 $T_2/T_1=(V_1/V_2)^{1/3}$，$p_2/p_1=(V_1/V_2)^{4/3}$，$U_2/U_1=(V_1/V_2)^{1/3}$。
功积分也可逐步与内能差对应。令 $K=p_1V_1^{4/3}$，则
\[
 W=K\int_{V_1}^{V_2}V^{-4/3}\dd V
 =[-3KV^{-1/3}]_{V_1}^{V_2}
 =3K(V_1^{-1/3}-V_2^{-1/3}).
\]
用$U_1=3p_1V_1$改写首项，所得功等于上面的内能减少。能量密度按$T^4$变化，总能量还须乘以当前体积，二者具有不同的体积指数。平衡光子数不受独立的守恒约束；但黑体光子数密度正比于$T^3$，在本题这个保持$VT^3$不变的理想可逆绝热过程中，平均总光子数恰好不变，不能由“光子可产生、湮灭”推断它在每种过程中都会改变。

\Diagram{fig:radadiabata}{等熵膨胀时辐射能降低并转化为功}{
 \draw[axis](0,0)--(5,0)node[right]{$V/V_1$};\draw[axis](0,0)--(0,2.6)node[above]{$U/U_1$};\draw[curve]plot[domain=1:4.6](\x,{2/\x^(1/3)});\node[left]at(0,2){$1$};\node[below]at(1,0){$1$};
}
图\ref{fig:radadiabata}中的能量差为$W=U_1-U_2=U_1[1-(V_1/V_2)^{1/3}]$，与上述压强积分一致。单原子理想气体的绝热指数为$5/3$，平衡辐射则为$4/3$；后者应由辐射状态方程和熵约束推导。
\end{solution}

\section{表面张力系统}
\label{th:surface}
液面面积增大时需要外界做功。以 $\Sigma$ 表示实际界面总面积、$\sigma$ 表示表面张力，在忽略体积功和其他功的模型中
\begin{equation}\label{eq:surfacefund}
 \dd U=T\dd S+\sigma\dd\Sigma.
\end{equation}
这与体积系统的 $-p\dd V$ 符号不同：增大面积需要对系统做正功。表面 Helmholtz 自由能满足 $\dd F=-S\dd T+\sigma\dd\Sigma$，故
\begin{equation}\label{eq:surfacemaxwell}
 \left(\frac{\partial S}{\partial\Sigma}\right)_T
 =-\left(\frac{\partial\sigma}{\partial T}\right)_\Sigma.
\end{equation}
改用 $\Sigma,T$ 作变量，链式法则给
\begin{equation}\label{eq:surfaceenergy}
 \left(\frac{\partial U}{\partial\Sigma}\right)_T
 =\sigma-T\left(\frac{\partial\sigma}{\partial T}\right)_\Sigma,
 \quad
 \dd S=\frac{C_\Sigma}{T}\dd T-\sigma_T\dd\Sigma.
\end{equation}
因此表面内能不一般等于 $\sigma\Sigma$；表面张力是自由能的面积导数。薄膜有两面时 $\Sigma=2A$，单液面才是 $\Sigma=A$。若还计入液体体积热容，必须说明 $C_\Sigma$ 是整个研究系统的定面积热容还是仅表面过剩热容。

\subsection{由几何变化求毛细压强}
\label{th:laplacegeometry}
表面张力给定的是增加实际界面面积所需的等温可逆功。对半径为 $r$ 的球形单界面，
$\Sigma=4\pi r^2$、$V=4\pi r^3/3$，所以
\[
 \frac{\dd\Sigma}{\dd V}
 =\frac{8\pi r}{4\pi r^2}=\frac2r.
\]
虚拟地把界面向外移动 $\dd r$，内外压差提供功
$(p_{\rm in}-p_{\rm out})\dd V$，界面自由能增加
$\sigma\dd\Sigma$。机械平衡要求两者相等，得
\[
 p_{\rm in}-p_{\rm out}=\frac{2\sigma}{r}.
\]
肥皂泡有两层半径近似相同的界面，实际总面积是 $8\pi r^2$，因此压差为 $4\sigma/r$。液滴的液气界面只有一层，不能使用肥皂泡的因子4。

\Diagram{fig:laplacegeometry}{虚位移把球面面积增量与体积增量联系起来}{
 \draw[curve](0,0)circle(1.15);\draw[guide](0,0)circle(1.45);
 \draw[axis](0,0)--(1.15,0)node[midway,below]{$r$};
 \draw[<->](.813,.813)--(1.025,1.025);
 \node[above right]at(1,1){$\dd r$};
 \node[right,align=left]at(2,.4){$\dd\Sigma=8\pi r\,\dd r$\\$\dd V=4\pi r^2\dd r$};
}
图\ref{fig:laplacegeometry}中，虚位移比较相邻几何状态，内外压差与表面张力在这一步作平衡功。圆柱单界面侧面的面积为 $2\pi rL$，体积为 $\pi r^2L$，固定 $L$ 时相同方法给压差 $\sigma/r$。两种几何的差别来自两个主方向中有几个发生弯曲。

\subsection{表面自由能、表面内能与吸热}
对同一面积增量，等温外功为 $\sigma\dd\Sigma$，内能增量为
$(\sigma-T\sigma_T)\dd\Sigma$，可逆吸热为
$-T\sigma_T\dd\Sigma$。三者满足第一定律
$\dd U=\delta Q+\delta W_{\rm on}$。温度升高时表面张力若下降，$\sigma_T<0$，等温拉伸便需要从外界吸热；这些热量进入表面内能，使内能面积系数高于自由能面积系数。

温度改变而面积不变时，热量由定面积热容决定。若把界面和一定量的液体作为整体，$C_\Sigma$ 包含该整体的储热能力；若只讨论表面过剩量，则须把体相的热容另加。选择研究边界后，同一式子中的 $U,S,C_\Sigma$ 均须按这个边界定义。

\begin{example}[液滴合并时的温升与熵变]
\label{ex:dropcoalescence}
$n_d\ge2$ 个相同球形小液滴初始半径为 $r$、温度为 $T_i$，最终合并成一个静止球形液滴。液体不可压缩，忽略蒸发、外界做功和热交换，合并过程产生的宏观动能最终全部耗散。整个液体的定面积热容为常数 $C$，表面张力在所涉温区内为 $\sigma(T)=\sigma_0+\alpha T>0$，其中 $\sigma_0>0$。求终半径、终温和熵变。
\Diagram{fig:dropcoalescenceq}{总体积不变的球形液滴合并}{
 \foreach \x in{0,2.2}\draw[curve](\x,0)circle(.35);
 \node at(1.1,0){$\cdots$};
 \node at(1.1,-.75){$n_d$ 个，半径 $r$};
 \draw[axis](3,0)--(4.1,0);
 \draw[curve](5.2,0)circle(.75);
 \draw[axis](5.2,0)--(5.95,0)node[midway,below]{$R$};
 \node at(5.2,-1.1){$Q=W_{\rm ext}=0$};
}
\end{example}
\begin{solution}
图\ref{fig:dropcoalescenceq}中液体总体积不变，故
$n_dr^3=R^3$，即 $R=n_d^{1/3}r$。初末实际界面面积分别为
\[
 \Sigma_i=4\pi n_dr^2,\qquad
 \Sigma_f=4\pi n_d^{2/3}r^2.
\]
记面积减少量 $D=\Sigma_i-\Sigma_f>0$。由表面内能公式，
$U(T,\Sigma)=CT+\sigma_0\Sigma+\text{常数}$。
整体绝热且无外界功，初末动能均为零，因此
\[
 C(T_f-T_i)=\sigma_0D,\qquad
 T_f=T_i+\frac{\sigma_0D}{C}.
\]
\Diagram{fig:dropcoalescencea}{减少的表面内能转为液体的热运动能}{
 \node[draw](a)at(0,0){$\sigma_0D$};
 \node[draw](b)at(4,0){$C(T_f-T_i)$};
 \draw[axis](a)--(b)node[midway,above]{能量转化};
 \node at(2,-1){外界吸热为零，系统温度仍可升高};
}
图\ref{fig:dropcoalescencea}中的表面能系数是 $\sigma_0=\sigma-T\sigma_T$，而不是初温下的 $\sigma(T_i)$。熵是状态函数，用
$\dd S=C\dd T/T-\alpha\dd\Sigma$ 积分，
\[
 \Delta S=C\ln\frac{T_f}{T_i}+\alpha D.
\]
尽管 $\alpha$ 可为负，总熵在题设有效温区内仍增加。沿相同总能量的连续状态曲线，以已减少面积 $D'$ 为参数，
$T=T_i+\sigma_0D'/C$，于是
\[
 \frac{\dd S}{\dd D'}=\frac{\sigma_0}{T}+\alpha
 =\frac{\sigma(T)}T>0.
\]
积分到 $D$ 即证明 $\Delta S>0$。这个等能状态曲线用于计算和判断端点，真实合并过程中的流动与耗散不必准静态。若算出的终温超出线性表面张力或常热容的适用温区，应重新使用实际函数积分。
\end{solution}


\begin{example}[线性表面张力的内能与熵]\label{ex:linearsurface}
某界面系统在研究温区内满足 $\sigma=\sigma_0+\alpha T$，$\sigma_0,\alpha$ 为常量；定面积热容 $C_\Sigma=C$ 为常量。求相对于 $(\Sigma_0,T_0)$ 的 $U,S$，并检验可积条件。
\Diagram{fig:surfacelinearq}{温度改变时表面张力线性变化}{
 \draw[axis](0,0)--(4.5,0)node[right]{$T$};\draw[axis](0,0)--(0,2.8)node[above]{$\sigma$};\draw[curve](0,2.4)--(4,.8);\node[left]at(0,2.4){$\sigma_0$};\node at(3.3,2.3){示意 $\alpha<0$};
}
\end{example}
\begin{solution}
由图\ref{fig:surfacelinearq}的斜率 $\sigma_T=\alpha$，式\eqref{eq:surfaceenergy}给出 $U_\Sigma|_T=\sigma_0+\alpha T-T\alpha=\sigma_0$。再用 $U_T|_\Sigma=C$，积分得
\[
 U-U_0=\sigma_0(\Sigma-\Sigma_0)+C(T-T_0).
\]
Maxwell 关系给 $S_\Sigma|_T=-\alpha$，定面积方向给 $S_T|_\Sigma=C/T$，所以
\[
 S-S_0=-\alpha(\Sigma-\Sigma_0)+C\ln\frac T{T_0}.
\]
温度项的积分沿定面积路径进行：$\Delta S_T=\int_{T_0}^T C\dd\vartheta/\vartheta$。面积项沿等温路径进行：$\Delta S_\Sigma=\int_{\Sigma_0}^{\Sigma}(-\alpha)\dd\Sigma'$。两者相加时，温度积分产生的对数只含无量纲比值，而面积项的量纲为 $(\mathrm{J\,m^{-2}K^{-1}})\mathrm{m^2}$，也是熵。表面张力虽包含温度项，内能的面积导数中该项被 $T\sigma_T$ 抵消，这正是表面自由能和表面内能需要区分的地方。

\Diagram{fig:surfacelineara}{两条积分方向的熵贡献}{
 \draw[axis](0,0)--(4.6,0)node[right]{$\Sigma$};\draw[axis](0,0)--(0,2.8)node[above]{$T$};\draw[axis](1,.7)--(3.2,.7);\draw[axis](3.2,.7)--(3.2,2.1);\node at(2.1,.3){$-\alpha\Delta\Sigma$};\node[right]at(3.2,1.6){$C\ln(T/T_0)$};
}
图\ref{fig:surfacelineara}的温度项必须为对数，不能写成 $C(T-T_0)$：后者量纲为能量，不是熵。一般可积条件为 $(\partial C_\Sigma/\partial\Sigma)_T=-T\sigma_{TT}$，本题两侧均为零，因此线性表面张力与常热容相容。所有式子只在 $\sigma>0$ 且界面模型成立的温区使用。
\end{solution}

\begin{example}[液膜等温拉伸的热与功]
长 $\ell$ 的活动杆拉伸肥皂膜，膜有两面。温度恒为 $T$，表面张力为 $\sigma_0+\alpha T$。活动杆沿垂直杆方向移动距离 $x$，忽略重力和膜厚变化的体积功。求外力、外界做功、膜吸热和内能变化。
\Diagram{fig:filmq}{两面液膜的活动边界}{
 \draw[curve](0,0)--(0,2)--(4,2);\draw[curve](0,0)--(4,0);\draw[line width=2pt](2.5,0)--(2.5,2);\draw[guide](3.7,0)--(3.7,2);\draw[<->](4.3,0)--(4.3,2)node[midway,right]{$\ell$};\draw[axis](2.5,1)--(3.7,1)node[midway,above]{$x$};
}
\end{example}
\begin{solution}
图\ref{fig:filmq}中单面面积增加 $\ell x$，实际总界面增加 $\Delta\Sigma=2\ell x$。两面各对杆施加 $\sigma\ell$，故缓慢拉伸外力 $F=2\sigma\ell$，对膜做功 $W_{\rm on}=2\sigma\ell x$。
图中杆长 $\ell$ 不变，位移 $\dd x$ 使每一面面积增加 $\ell\dd x$，所以外界微元功为 $F\dd x=\sigma\dd\Sigma$。这同时给出力的二倍系数与功的面积因子，两者应使用同一个总界面面积。等温时膜的熵变化只来自面积项；若液面张力随温度下降，即 $\alpha<0$，新生成界面具有正的熵增，需要从恒温热源吸热。内能式中的 $C(T-T_0)$ 在本题等温条件下为零，但定面积热容本身并不因此为零。

\Diagram{fig:filma}{等温拉伸时的自由能、内能与热量}{
 \node[draw]at(0,0){$\Delta F=\sigma\Delta\Sigma$};\node[draw]at(5,0){$\Delta U=\sigma_0\Delta\Sigma$};\node at(2.5,-.9){$Q=T\Delta S=-\alpha T\Delta\Sigma$};
}
由图\ref{fig:filma}，等温熵增$\Delta S=-\alpha\Delta\Sigma$，吸热$Q=-2\alpha T\ell x$，内能增加$\Delta U=2\sigma_0\ell x$。两种能量输入相加，$Q+W_{\rm on}=2(\sigma_0+\alpha T)\ell x-2\alpha T\ell x=\Delta U$，满足第一定律。
\end{solution}

\begin{example}[界面的绝热温度变化]
线性表面张力、常定面积热容的界面系统，从 $(\Sigma_0,T_0)$ 可逆绝热拉伸到面积 $\Sigma$。求 $T(\Sigma)$、表面张力和外界做功。假设温区内模型始终成立。
\Diagram{fig:surfaceadq}{与外界不换热的可动液膜}{
 \draw(0,0)rectangle(4,1.6);\draw[line width=2pt](2.5,0)--(2.5,1.6);\draw[axis](2.5,.8)--(3.8,.8);\node at(1.2,.8){$Q=0$};\node at(2,-.5){$C_\Sigma=C$};
}
\end{example}
\begin{solution}
图\ref{fig:surfaceadq}的过程可逆，故$\dd S=0$。用已求得的熵微分，$C\dd T/T=\alpha\dd\Sigma$；从$(\Sigma_0,T_0)$积分，有$C[\ln T']_{T_0}^{T}=\alpha[\Sigma']_{\Sigma_0}^{\Sigma}$，即$C\ln(T/T_0)=\alpha(\Sigma-\Sigma_0)$，取指数得
\[
 T=T_0\exp[\alpha(\Sigma-\Sigma_0)/C],\quad
 \sigma=\sigma_0+\alpha T_0\exp[\alpha(\Sigma-\Sigma_0)/C].
\]
外界做功积分中的变温部分以指数函数作原函数，得到
\[
 \int_{\Sigma_0}^{\Sigma}\alpha T_0
 \ee^{\alpha(\Sigma'-\Sigma_0)/C}\dd\Sigma'
 =\left[CT_0\ee^{\alpha(\Sigma'-\Sigma_0)/C}\right]_{\Sigma_0}^{\Sigma}
 =CT_0\bigl[\ee^{\alpha\Delta\Sigma/C}-1\bigr].
\]
它对应内能中的温度变化项，符号在 $\alpha<0$ 的拉伸过程中为负；只要全过程张力为正，总输入功仍为正。膜的降温减小了一部分所需输入，但不取消创建界面的能量代价。

\Diagram{fig:surfaceada}{可逆绝热拉伸时的温度变化}{
 \draw[axis](0,0)--(4.7,0)node[right]{$\Sigma-\Sigma_0$};\draw[axis](0,0)--(0,2.7)node[above]{$T$};\draw[curve]plot[domain=0:4.3](\x,{2.3*exp(-.25*\x)});\node at(2.7,2.4){$\alpha<0$};
}
图\ref{fig:surfaceada}中若 $\alpha<0$，拉伸使温度降低。外界做功等于内能增量，
\[
 W_{\rm on}=\sigma_0\Delta\Sigma+C(T-T_0).
\]
常数张力项给出$\sigma_0\Delta\Sigma$，前面求出的指数项给出$C(T-T_0)$。张力随温度变化，故整个过程的功须沿路径积分。
\end{solution}

\section{相平衡与汽化热}
\label{th:clapeyron}
同一物质的液、气两相共存时，温度和压强相等，两相摩尔 Gibbs 自由能也相等：$g_g=g_l$。沿共存线的相邻状态仍满足此等式。利用 $\dd g=-s\dd T+v\dd p$，
\[
 -(s_g-s_l)\dd T+(v_g-v_l)\dd p=0.
\]
定温定压的可逆汽化每摩尔吸热为 $L_m=T(s_g-s_l)$，从而
\begin{equation}\label{eq:clapeyron}
 \frac{\dd p_s}{\dd T}=\frac{L_m}{T(v_g-v_l)}.
\end{equation}
这就是 Clapeyron 方程。汽化时 $L_m>0$、$v_g>v_l$，共存线斜率为正。熔化时应代入液相减固相的摩尔体积差；冰熔化的体积差为负，熔化潜热为正，故其共存线斜率为负。

还可从跨越两条相邻等温线的微小两相卡诺循环推导。上等温线汽化一摩尔吸热 $L_m$，两条绝热线温差为 $\dd T$；循环面积近似为 $(v_g-v_l)\dd p_s$。可逆效率为 $\dd T/T$，故 $(v_g-v_l)\dd p_s/L_m=\dd T/T$，得到同一方程。用循环法时，面积与吸热必须对应同一物质的量。

\subsection{相平衡条件与稳定相的选择}
\label{th:phasestability}
同温同压的两相处于机械、热学平衡后，仍可能发生物质转移。设从相1向相2可逆转移 $\dd\nu$ 摩尔，总 Gibbs 自由能变化为
\[
 \dd G=(g_2-g_1)\dd\nu.
\]
若 $g_2<g_1$，向相2转移会降低总 $G$，因此相2更稳定；只有 $g_2=g_1$ 时，两相才可在不产生净转移趋势的条件下共存。因而相平衡除了同温、同压，还要求转移物质时两相之间没有自由能降低的方向。

沿共存线改变温度、压强时，两相的 $g$ 继续相等。前文 Clapeyron 方程就是这一相等条件的全微分。斜率的分子是相变熵差，分母是体积差：熵差规定升温倾向哪一相，体积差规定加压倾向哪一相。对通常的液气相变，升温有利于气相，加压有利于液相，两种影响共同形成向右上升的饱和曲线。

\Diagram{fig:phasegibbs}{两相摩尔自由能在共存压强处相等}{
 \draw[axis](0,0)--(5.5,0)node[right]{$p$};
 \draw[axis](0,0)--(0,3.5)node[above]{$g$};
 \draw[curve](.6,1.5)--(4.8,2.2)node[right]{$g_l$};
 \draw[curve](.6,.65)--(4.8,2.95)node[right]{$g_g$};
 \draw[guide](2.831,0)--(2.831,1.872);
 \node[below]at(2.831,0){$p_s(T)$};
 \node at(2.2,3.1){固定温度，斜率为摩尔体积};
}
图\ref{fig:phasegibbs}示意共存点附近的变化。因 $g_p|_T=v$，气相的曲线斜率远大于液相；低于饱和压强时气相自由能较低，高于饱和压强时液相较低。图中的直线只表示局部线性化，跨越较大压强区间时应按各相体积积分。

\subsection{两相平台、杠杆规则与可压缩性}
固定温度下，若每摩尔系统的平均体积 $v$ 位于 $v_l$ 与 $v_g$ 之间，稳定平衡可由液气两相组成。令气相摩尔分数为 $x$，则
\[
 v=(1-x)v_l+xv_g,\qquad
 x=\frac{v-v_l}{v_g-v_l}.
\]
压强保持 $p_s(T)$，体积增大通过增加气相比例实现。此时形式上的等温压缩率很大，反映相比例调整，不能再用单一均匀相的微小密度响应解释。若相变成核受阻，实际系统可能暂时停留在亚稳态，但这需要额外的动力学条件。

\Diagram{fig:leverrule}{共存平台上的总体积由两端相体积加权}{
 \draw[axis](0,0)--(6,0)node[right]{$v$};
 \foreach \x/\t in{1/v_l,3.3/v,5/v_g}{
  \draw(\x,-.13)--(\x,.13);\node[below]at(\x,-.15){$\t$};
 }
 \draw[<->](1,.7)--(3.3,.7)node[midway,above]{$x(v_g-v_l)$};
 \draw[<->](3.3,1.5)--(5,1.5)node[midway,above]{$(1-x)(v_g-v_l)$};
}
图\ref{fig:leverrule}中，靠近液相端点意味着气相比例小。本式使用摩尔体积，权重相应为摩尔分数；采用质量分数或体积分数时，应连同被加权的量一起变换。计算两相热容时，显热项与 $x(T)$ 的变化项都从这一约束出发。

\subsection{等面积条件与自由能相等}
在可解析延拓的状态方程中，沿一条等温线有 $\dd g=v\dd p$。假设两端 $v_l,v_g$ 对应同一压强 $p_s$，积分并分部积分，
\[
 g_g-g_l=\int_{v_l}^{v_g}v\,p'(v)\dd v
 =p_s(v_g-v_l)-\int_{v_l}^{v_g}p(v)\dd v.
\]
要求两端自由能相等，便得到
\[
 \int_{v_l}^{v_g}[p(v)-p_s]\dd v=0.
\]
这正是等面积条件。积分中使用的中间均匀态分支是状态方程的数学延拓，不表示系统在平衡相变时真实经过它们。真实等温压缩沿共存平台改变相比例；曲线在平台上下的有向面积相抵，是两端化学势相等的另一种表达。

\subsection{饱和蒸气压近似的使用层次}
从 Clapeyron 方程到常见的指数蒸气压式，需要依次采用三个近似：液相比气相体积可忽略；气相满足理想气体状态方程；潜热在所讨论温区内近似不变。前两项把斜率化为 $L_m p_s/(RT^2)$，最后一项才允许把 $L_m$ 提到积分号外。

温度接近临界点时，两相体积差和潜热同时趋小，忽略液相体积的近似失效；在高压区，气相也可能明显偏离理想气体。对较宽温区，即使气相足够稀薄，也应保留 $L_m(T)$ 的变化。后面的变潜热模型只改进第三层近似，仍保留前两层，所得幂律乘指数式也应在这一范围内理解。


\begin{example}[饱和蒸气压的指数近似]\label{ex:vaporpressure}
在一段温区内忽略液体摩尔体积，气相视为理想气体，摩尔汽化热 $L_m=L$ 为常量。已知 $p_s(T_0)=p_0$，求 $p_s(T)$ 和压强为 $p$ 时的沸点。
\Diagram{fig:vaporq}{液气共存线上给定一个参考点}{
 \draw[axis](0,0)--(4.8,0)node[right]{$T$};\draw[axis](0,0)--(0,3)node[above]{$p_s$};\draw[curve]plot[domain=.8:4.2](\x,{.12*exp(.65*\x)});\fill(2.3,.535)circle(1.3pt)node[above left]{$T_0,p_0$};
}
\end{example}
\begin{solution}
由图\ref{fig:vaporq}沿共存线的体积差 $v_g-v_l\simeq v_g=RT/p_s$，式\eqref{eq:clapeyron}变为 $\dd p_s/p_s=L\dd T/(RT^2)$。积分时 $\int T^{-2}\dd T=-1/T$，所以
\[
 \ln\frac{p_s(T)}{p_0}=\frac LR\left(\frac1{T_0}-\frac1T\right),\qquad
 p_s(T)=p_0\exp\!\left[\frac LR\left(\frac1{T_0}-\frac1T\right)\right].
\]
求沸点时，先把对数式中的 $1/T_b$ 单独移项，再取倒数，不能直接把 $\ln p$ 看作与温度线性。共存线的局部斜率始终正：增压要求更高的沸点。参考点规定了积分常数，因此指数中的每一项都按相同温标和单位计算；$L$ 是摩尔汽化热，必须与摩尔气体常数 $R$ 配套。若使用单位质量潜热，应同时把 $R$ 换成单位质量气体常数。

\Diagram{fig:vapora}{对数蒸气压与倒温度的线性关系}{
 \draw[axis](0,0)--(4.8,0)node[right]{$1/T$};\draw[axis](0,0)--(0,3)node[above]{$\ln p_s$};\draw[curve](.6,2.5)--(4.2,.5);\node at(3.1,2.5){斜率 $-L/R$};
}
图\ref{fig:vapora}给出实验上常用的线性化。沸点满足 $p_s(T_b)=p$，因此 $T_b=[T_0^{-1}-(R/L)\ln(p/p_0)]^{-1}$。高海拔外压较低，$\ln(p/p_0)<0$，沸点随之降低。靠近临界点时潜热和体积差都变化显著，本近似不能外推。
\end{solution}

\begin{example}[汽化热随温度变化的修正]\label{ex:variablelatent}
仍忽略液体体积，气相理想。气相摩尔定容热容 $c_{v,g}$、液相比热 $c_l$ 均视为常量，已知参考温度 $T_0$ 的汽化热为 $L_0$、蒸气压为 $p_0$。推导 $L_m(T)$ 和修正后的 $p_s(T)$。
\Diagram{fig:latentvariableq}{汽化焓差由两相能量和气相膨胀功组成}{
 \node[draw]at(0,0){液体 $U_{m,l}$};\draw[axis](1.3,0)--(2.8,0)node[midway,above]{$L_m$};\node[draw]at(4.5,0){气体 $U_{m,g}+RT$};
}
\end{example}
\begin{solution}
图\ref{fig:latentvariableq}的定温定压汽化热为焓差 $L_m=U_{m,g}-U_{m,l}+p_s(v_g-v_l)$。在题设近似下 $p_sv_g=RT$，且 $\dd U_{m,g}/\dd T=c_{v,g}$、$\dd U_{m,l}/\dd T\simeq c_l$。定义 $q=c_{v,g}+R-c_l$，则
\[
 L_m(T)=L_0+q(T-T_0)=\Lambda+qT,\quad\Lambda=L_0-qT_0.
\]
代入 $\dd\ln p_s=L_m(T)\dd T/(RT^2)$，分成两项积分，
\[
 \ln\frac{p_s(T)}{p_0}=
 \frac qR\ln\frac T{T_0}+\frac\Lambda R\left(\frac1{T_0}-\frac1T\right).
\]
常数 $\Lambda$ 与参考点汽化热 $L_0$ 不同，它是把线性 $L_m(T)$ 外推到零温时的截距，通常不具备真实零温汽化热的含义。积分时分别有 $\int q\dd T/(RT)=(q/R)\ln T$、$\int\Lambda\dd T/(RT^2)=-\Lambda/(RT)$，所以两个因子中的符号由原函数确定。将最终结果对温度求导，会恢复 $\dd\ln p_s/\dd T=L_m(T)/(RT^2)$，也说明在参考温度附近它与给定的 $L_0$ 连续匹配。

\Diagram{fig:latentvariablea}{修正项分别产生幂律因子和指数因子}{
 \node[draw,align=center]at(0,0){$q/T$\\$(T/T_0)^{q/R}$};\node[draw,align=center]at(5,0){$\Lambda/T^2$\\$\exp[\Lambda(1/T_0-1/T)/R]$};
}
所以图\ref{fig:latentvariablea}的两个因子相乘即 $p_s/p_0$。$q$ 可正可负，取决于两相热容；不能预先把幂指数指定为正。$q=0$ 时回到常潜热公式。这里还忽略了液体焓的压强依赖和气相非理想性，修正温变潜热并不自动消除其他近似误差。
\end{solution}

\begin{example}[熔化曲线的负斜率]
某物质熔化潜热为正，固相摩尔体积 $v_s$ 大于液相 $v_l$。求熔点对压强的局部变化率，并以给定常量近似求压强变化 $\Delta p$ 引起的熔点变化。
\Diagram{fig:meltingq}{固、液体积差决定共存线的斜率}{
 \draw[axis](0,0)--(4.8,0)node[right]{$T$};\draw[axis](0,0)--(0,3)node[above]{$p$};\draw[curve](1,2.6)--(2.8,.4);\node at(1,1.1){固};\node at(3.5,1.8){液};
}
\end{example}
\begin{solution}
图\ref{fig:meltingq}从固到液的 $\Delta v=v_l-v_s<0$，熔化熵增为 $L_m/T>0$。Clapeyron 方程给 $\dd p/\dd T=L_m/[T(v_l-v_s)]<0$。局部取倒数，
\[
 \frac{\dd T_m}{\dd p}=\frac{T_m(v_l-v_s)}{L_m}.
\]
取倒数后，潜热分母为正，斜率的符号由 $v_l-v_s$ 决定。在参考态附近用 $T_{m0}$ 计算斜率，保留的是压强变化的一阶项。压强变化较大时，须沿共存线积分 $\dd T/T=(v_l-v_s)\dd p/L_m$，同时计入各量随状态的变化。图中“加压有利于液相”表示同温下液相的 Gibbs 自由能相对降低，因而共存熔点下降。

\Diagram{fig:meltinga}{加压有利于摩尔体积较小的液相}{
 \node[draw]at(0,0){固体积较大};\draw[axis](1.3,0)--(2.7,0)node[midway,above]{加压};\node[draw]at(4.4,0){液体更有利};\node at(2.2,-.8){共存熔点降低};
}
若变化足够小，可把右侧参数取参考值，$\Delta T_m\simeq T_{m0}(v_l-v_s)\Delta p/L_m$。图\ref{fig:meltinga}从自由能角度说明符号：压强增加时，大体积相的 Gibbs 自由能增加更快。不能仅凭“压强抑制相变”一类措辞判断方向。
\end{solution}

\section{相空间分布与外场中的粒子}
\label{th:phaseprobability}
经典独立粒子的单粒子能量为
\[
 E(\vct r,\vct p)=\frac{p^2}{2m}+E_p(\vct r).
\]
与温度 $T$ 的热源平衡时，相空间密度为 $f(\vct r,\vct p)=Z^{-1}\ee^{-E/(\kb T)}$，概率测度为 $\dd^3r\dd^3p$。若用 $\dd^3r\dd^3p/h^3$ 计数微观状态，则归一化因子也相应改变；固定粒子数的经典位置概率和速度概率不受该常数尺度影响。

由于动能与势能可分离，配分积分也分离：
\[
 Z=\int_{\cal V}\ee^{-E_p/(\kb T)}\dd^3r
 \int_{\mathbb R^3}\ee^{-p^2/(2m\kb T)}\dd^3p.
\]
对位置积分得到 Maxwell 动量分布；对动量积分得到 Boltzmann 位置分布，
\begin{equation}\label{eq:boltzmannposition}
 \rho_P(\vct r)=\frac{\ee^{-E_p(\vct r)/(\kb T)}}
 {\int_{\cal V}\ee^{-E_p(\vct r')/(\kb T)}\dd^3r'}.
\end{equation}
温度相同不意味着粒子密度均匀；势能较低处的粒子更多。反过来，密度随位置改变也不意味着不同高度的平均平动能不同。在等温平衡且动能可分离的条件下，每个位置处的条件速度分布相同。

\subsection{局部密度、位置概率与几何测度}
\label{th:spatialmeasure}
式\eqref{eq:boltzmannposition}给出的是每单位真实空间体积的概率密度。若装置截面随坐标改变，单独观测某一坐标的分布还须计入这一几何因素。例如轴对称圆筒内，体密度为 $n(r)$ 时，半径在 $r$ 与 $r+\dd r$ 间的粒子数为
\[
 \dd N=2\pi rL\,n(r)\dd r.
\]
因而径向概率密度正比于 $r n(r)$，即使 $n(r)$ 为常数，径向样本也更多出现在较大的 $r$。采用 $s=r^2/R_c^2$ 时，$2r\dd r=R_c^2\dd s$，几何因子才被吸收到新变量的微分中。这一几何计数与第一讲速度空间的球壳因子相同：等坐标间隔内所包含的体积会随坐标变化。

\Diagram{fig:spatialannulus}{圆筒横截面的环带测度}{
 \draw[curve](0,0)circle(1.5);
 \fill[gray!20,even odd rule](0,0)circle(1.15)(0,0)circle(.95);
 \draw[curve](0,0)circle(.95);\draw[curve](0,0)circle(1.15);
 \draw[axis](0,0)--(35:.95)node[midway,above]{$r$};
 \draw[<->](315:.95)--(315:1.15);
 \node[right]at(1.05,-.9){$\dd r$};
 \node[right,align=left]at(2.2,.2){体积元 $2\pi rL\dd r$\\径向计数含因子 $r$};
}
图\ref{fig:spatialannulus}中的环带是位置计数区域。重力气柱若截面恒定，其高度概率密度与 $n(h)$ 成正比；变截面容器则应乘截面积 $A(h)$。比较不同高度或半径的“粒子多寡”时，须先说明比较的是等体积、等高度厚度还是等半径厚度。

\subsection{位置涨落怎样决定外场热容}
设外场势能 $\varepsilon(\vct r)$ 不显含温度，容器边界固定。位置配分积分为 $Z_r=\int\ee^{-\beta\varepsilon}\dd^3r$。对 $\beta$ 求导，
\[
 \langle\varepsilon\rangle=-\frac{\partial\ln Z_r}{\partial\beta},
 \qquad
 \frac{\partial\langle\varepsilon\rangle}{\partial\beta}
 =-\bigl(\langle\varepsilon^2\rangle-\langle\varepsilon\rangle^2\bigr).
\]
这里$Z_r'=-\int\varepsilon\ee^{-\beta\varepsilon}\dd^3r$、$Z_r''=\int\varepsilon^2\ee^{-\beta\varepsilon}\dd^3r$，故$\partial_\beta(-Z_r'/Z_r)=-Z_r''/Z_r+(Z_r'/Z_r)^2=-\langle\varepsilon^2\rangle+\langle\varepsilon\rangle^2$。再用$\dd\beta/\dd T=-1/(\kb T^2)$，得单粒子势能对热容的贡献
\[
 \frac{\dd\langle\varepsilon\rangle}{\dd T}
 =\frac{\operatorname{Var}(\varepsilon)}{\kb T^2}\ge0.
\]
因此固定外场中的升温会增加平均势能，增加多少由能量涨落确定。对独立单原子粒子，完整热容为
$C=3N\kb/2+N\operatorname{Var}(\varepsilon)/(\kb T^2)$。
若能量函数本身显含温度，或控制参数在过程中改变，则还需计入相应导数项，不能直接沿用这个固定参数公式。

均匀重力场给 $\varepsilon=mgh$，故气柱的附加热容为
$N(mg)^2\operatorname{Var}(h)/(\kb T^2)$。在有限高度且高温接近均匀分布时，$\operatorname{Var}(h)\to H^2/12$，附加热容趋零；在尺度高度远小于 $H$ 时，高度近似指数分布，其方差为 $(\kb T/(mg))^2$，附加热容趋于 $N\kb$。两个极限也将由后文的有限气柱能量积分得到；方差形式还直接保证附加热容非负。

\subsection{参数变化与可逆绝热降温}
\begin{example}[谐势阱缓慢减弱时的温度与功]
\label{ex:trapadiabat}
$N$ 个经典独立单原子粒子在三维谐势
$\varepsilon=(\kappa_xx^2+\kappa_yy^2+\kappa_zz^2)/2$
中保持热平衡，初温为 $T_1$。在与外界不交换热量的条件下，缓慢改变 $\kappa_x$，使其由 $\kappa_{x1}$ 变为 $\kappa_{x2}>0$，其他两个刚度不变。假定内部碰撞能及时建立平衡，粒子不逃逸，求终温、空间宽度和外界对气体所做的功。
\Diagram{fig:trapadiabatq}{同一谐势方向减弱后的势能曲线}{
 \draw[axis](-2.4,0)--(2.6,0)node[right]{$x$};
 \draw[axis](0,0)--(0,3)node[above]{$\varepsilon_x$};
 \draw[curve]plot[domain=-1.6:1.6](\x,{\x*\x});
 \draw[curve,dashed]plot[domain=-2.2:2.2](\x,{.35*\x*\x});
 \node[right]at(2.2,2.7){初始 $\kappa_{x1}$};
 \node[right]at(2.8,1.2){减弱后 $\kappa_{x2}$};
}
\end{example}
\begin{solution}
图\ref{fig:trapadiabatq}中，势阱减弱改变了可用的位置空间。对动量和位置分别作三个高斯积分，单粒子配分积分的参数依赖为
$Z_1\propto T^3/(\kappa_x\kappa_y\kappa_z)^{1/2}$。其中三个动量高斯积分各给$T^{1/2}$，三个位置高斯积分各给$(T/\kappa_i)^{1/2}$，归一化使用第一讲已证明的高斯积分。按式\eqref{eq:canonicalentropy}，每粒子的熵为$\kb[\ln Z_1+T\partial_T\ln Z_1]=\kb(\ln Z_1+3)$，再加与本题状态无关的固定粒子数常数。比较两态时常数相消，熵差为
\[
 \Delta S=N\kb\left[
 3\ln\frac{T_2}{T_1}
 -\frac12\ln\frac{\kappa_{x2}\kappa_{y2}\kappa_{z2}}
 {\kappa_{x1}\kappa_{y1}\kappa_{z1}}\right].
\]
准静态且内部平衡、无耗散的绝热过程满足 $\Delta S=0$，故记 $a=\kappa_{x2}/\kappa_{x1}$，有
\[
 T_2=T_1a^{1/6}.
\]
每个动能项和势能项各贡献 $\kb T/2$，总能量为 $U=3N\kb T$。因此外界对系统做功
\[
 W_{\rm on}=U_2-U_1=3N\kb T_1(a^{1/6}-1).
\]
\Diagram{fig:trapadiabata}{等熵条件把势阱刚度与温度联系起来}{
 \draw[axis](0,0)--(5,0)node[right]{$\ln(\kappa_x/\kappa_{x1})$};
 \draw[axis](2,-1.1)--(2,1.8)node[above]{$\ln(T/T_1)$};
 \draw[curve](.3,-.283)--(4.6,.433);
 \node[right]at(3.3,1.05){斜率 $1/6$};
 \draw[->](1.4,-.1)--(.7,-.217);
 \node[below]at(1,-.55){减弱势阱};
}
图\ref{fig:trapadiabata}中向左下移动对应 $a<1$：气体降温，并向控制势阱的装置输出功。还可直接计算广义功验证符号，
$\delta W_{\rm on}=\langle\sum_i x_i^2/2\rangle\dd\kappa_x
 =N\kb T\,\dd\kappa_x/(2\kappa_x)$。
令$z=\kappa_x/\kappa_{x1}$，则$\dd\kappa_x/\kappa_x=\dd z/z$，上下限为$1,a$。代入$T=T_1z^{1/6}$，功为$(N\kb T_1/2)\int_1^a z^{-5/6}\dd z=(N\kb T_1/2)[6z^{1/6}]_1^a=3N\kb(T_2-T_1)$，与内能差相同。

位置二阶矩为 $\langle x^2\rangle=\kb T/\kappa_x$，所以
$\sqrt{\langle x^2\rangle_2/\langle x^2\rangle_1}=a^{-5/12}$；
另两方向因刚度固定，宽度之比均为 $a^{1/12}$。
减弱一个方向的束缚会使该方向显著展宽，同时通过碰撞降温使另两方向收窄。若三个刚度一同乘 $a$，等熵式改给 $T_2/T_1=a^{1/2}$，可见不同参数操作不能共用同一指数。
\end{solution}


\subsection{等温大气的两种推导}
\label{th:barometric}
设重力加速度$g$恒定，气体摩尔质量为$M$、温度为常量$T$。高度$h$到$h+\dd h$的薄层受力平衡，$p(h)S-p(h+\dd h)S-\rho Sg\dd h=0$，所以$\dd p/\dd h=-\rho g$。理想气体$\rho=Mp/(RT)$，分离变量并从底面积分至高度$h$，有
\[
 \int_{p(0)}^{p(h)}\frac{\dd p'}{p'}
 =-\frac{Mg}{RT}\int_0^h\dd h',\qquad
 \ln\frac{p(h)}{p(0)}=-\frac{Mgh}{RT}.
\]
取指数，再用$n=p/(\kb T)$，得到
\begin{equation}\label{eq:barometric}
 p(h)=p(0)\ee^{-Mgh/(RT)},\qquad
 n(h)=n(0)\ee^{-mgh/(\kb T)}.
\end{equation}
用式\eqref{eq:boltzmannposition}，单粒子势能为 $mgh$，立即得到同一密度指数。两条路线分别从宏观受力和微观概率进入，指数尺度一致，因为 $M/R=m/\kb$。尺度高度为 $H_T=RT/(Mg)$。地球实际大气温度随高度变化，不能在任意高度直接套用等温指数律。

\begin{example}[气柱压强与总重量]\label{ex:atmosphere}
截面积为 $S$、高 $H$ 的竖直刚性柱内，封闭 $N$ 个质量为 $m$ 的理想气体粒子，处于温度 $T$ 的平衡。求底面与顶面压强，并证明其差产生的合力等于气体重量。
\Diagram{fig:columnpressureq}{上下两端的压强不同}{
 \draw(0,0)rectangle(2,3);\draw[axis](2.6,0)--(2.6,3.5)node[above]{$h$};\node[left]at(0,0){$p_b$};\node[left]at(0,3){$p_t$};\node at(1,1.5){$N,T$};\draw[<->](3.3,0)--(3.3,3)node[midway,right]{$H$};
}
\end{example}
\begin{solution}
取 $a_g=mg/(\kb T)$、$x=a_gH$，密度为 $n(h)=n_b\ee^{-a_gh}$。图\ref{fig:columnpressureq}的固定粒子数要求
\[
 N=S\int_0^Hn(h)\dd h=Sn_b\frac{1-\ee^{-x}}{a_g},
 \quad n_b=\frac{Na_g}{S(1-\ee^{-x})}.
\]
因此
\[
 p_b=\frac{Nmg}{S(1-\ee^{-x})},\qquad
 p_t=p_b\ee^{-x}.
\]
这里高度积分为$\int_0^H\ee^{-a_gh}\dd h=[-\ee^{-a_gh}/a_g]_0^H=(1-\ee^{-x})/a_g$。底部密度由总粒子数决定，并非外界气源预先固定的密度。乘上$\kb T$才得到压强，$a_g\kb T=mg$使显式分子中的温度因子约去，但温度仍保留在指数里。

\Diagram{fig:columnpressurea}{气体整体的竖直受力平衡}{
 \draw(0,0)rectangle(2,2.5);\draw[axis](1,0)--(1,1)node[right]{$Sp_b$};\draw[axis](1,2.5)--(1,1.8)node[right]{$Sp_t$};\draw[axis](2.6,1.6)--(2.6,.5)node[right]{$Nmg$};
}
图\ref{fig:columnpressurea}中的气体所受底面向上力减去顶面向下力为 $S(p_b-p_t)=Nmg$，恰好平衡总重力。升温时底面密度归一化常数也会改变，不能把 $n_b$ 当与温度无关的固定量。
\end{solution}

\begin{example}[有限高度气柱的总热容]\label{ex:columnheat}
上述刚性封闭气柱中装有 $\nu$ 摩尔单原子理想气体，摩尔质量 $M$。缓慢供热使每一状态都处于同一均匀温度的热平衡。求重力势能 $E_p(T)$ 和容器的热容 $C=\delta Q/\dd T$，忽略容器自身热容。
\Diagram{fig:columnheatq}{升温除了增大热运动，还改变高度分布}{
 \draw(0,0)rectangle(2,3);\foreach \y in{.25,.55,.85,1.25,1.8,2.5}{\fill(.8,\y)circle(1.2pt);}\draw[axis](-1,1.3)--(0,1.3)node[midway,above]{$Q$};\node[right]at(2.2,1.5){$E_p=\int_0^H Mgh\,\dd\nu$};
}
\end{example}
\begin{solution}
图\ref{fig:columnheatq}的气体密度从底部向上衰减，令 $x=MgH/(RT)$、$y=h/H$。按上一例归一化，单粒子高度的无量纲密度为 $w(y)=x\ee^{-xy}/(1-\ee^{-x})$。于是
\[
 E_p=\nu MgH\frac{x}{1-\ee^{-x}}\int_0^1y\ee^{-xy}\dd y.
\]
分部积分取 $u=y$、$\dd v=\ee^{-xy}\dd y$，
\[
 \int_0^1y\ee^{-xy}\dd y
 =\left[-\frac{y\ee^{-xy}}x\right]_0^1
   +\frac1x\left[-\frac{\ee^{-xy}}x\right]_0^1
 =-\frac{\ee^{-x}}x+\frac{1-\ee^{-x}}{x^2}.
\]
所以
\begin{equation}\label{eq:columnpotential}
 E_p=\nu MgH\left[\frac1x-\frac1{\ee^x-1}\right]
 =\nu RT-\frac{\nu MgH}{\ee^x-1}.
\end{equation}
对势能式求温度导数时，第二项的全部温度依赖都藏在 $x=MgH/(RT)$ 中。链式法则给
\[
 \frac{\dd}{\dd T}\frac1{\ee^x-1}
 =-\frac{\ee^x}{(\ee^x-1)^2}\left(-\frac xT\right)
 =\frac{x\ee^x}{T(\ee^x-1)^2}.
\]
再用 $MgH=xRT$，把这一项的系数写为无量纲 $x^2$，便得到热容中的修正分式。系统能量包含外场势能，所以固定几何体积下的供热量不只是平动内能变化；若把势能留在系统之外，则还应等价地计入重力对粒子所做的功。

\Diagram{fig:columnheata}{刚性容器的供热分给平动能与重力势能}{
 \node[draw]at(0,0){$\delta Q$};\draw[axis](.8,.2)--(2.3,1);\draw[axis](.8,-.2)--(2.3,-1);\node[right]at(2.3,1){$\dd U_{\rm tr}=3\nu R\dd T/2$};\node[right]at(2.3,-1){$\dd E_p$};
}
外壁刚性，无边界体积功。按图\ref{fig:columnheata}将重力势能计入系统，有$\delta Q=\dd(3\nu RT/2+E_p)$。用前面求出的势能导数，再加上平动热容，得到
\begin{equation}\label{eq:columncapacity}
 C=\nu R\left[\frac52-\frac{x^2\ee^x}{(\ee^x-1)^2}\right].
\end{equation}
式中第一项 $3\nu R/2$ 是局部平动热容，额外部分来自高度分布重排。气体内部没有活塞运动，仍可有超过定容平动热容的供热需求，因为外场势能随温度改变。
\end{solution}

\begin{example}[气柱热容的两端极限]\label{ex:columnlimits}
分析式\eqref{eq:columncapacity}在 $x\ll1$、$x\gg1$ 时的行为；求平均高度的相应极限，并说明先令 $H\to\infty$ 与先令 $T\to\infty$ 为什么不同。
\Diagram{fig:columnlimitq}{尺度高度相对于容器高度的两种情形}{
 \draw(0,0)rectangle(1.8,3);\foreach \y in{.3,.8,1.3,1.8,2.3,2.8}\fill(.9,\y)circle(1.3pt);\node at(.9,-.5){$x\ll1$};
 \draw(3.4,0)rectangle(5.2,3);\foreach \y in{.15,.3,.5,.75,1.1,1.7}\fill(4.3,\y)circle(1.3pt);\node at(4.3,-.5){$x\gg1$};
}
\end{example}
\begin{solution}
图\ref{fig:columnlimitq}左图$H\ll H_T$，密度近乎均匀。先从指数的Taylor展开求倒数。记$F(x)=1/(\ee^x-1)$，写成
\[
 \begin{aligned}
 \ee^x-1&=x\left(1+\frac x2+\frac{x^2}6+\frac{x^3}{24}
 +\frac{x^4}{120}+\cdots\right),\\
 F(x)&=\frac1x(1+a_1x+a_2x^2+a_3x^3+a_4x^4+\cdots).
 \end{aligned}
\]
两式相乘须为1，各次幂系数依次满足
\[
 \begin{aligned}
 a_1+\tfrac12&=0,\\
 a_2+\tfrac12a_1+\tfrac16&=0,\\
 a_3+\tfrac12a_2+\tfrac16a_1+\tfrac1{24}&=0,\\
 a_4+\tfrac12a_3+\tfrac16a_2+\tfrac1{24}a_1+\tfrac1{120}&=0.
 \end{aligned}
\]
逐式回代得$a_1=-1/2$、$a_2=1/12$、$a_3=0$、$a_4=-1/720$。又有$F(-x)=-1-F(x)$，故$F(x)+1/2$的Laurent展开只有奇次幂，得到
\[
 F(x)=\frac1x-\frac12+\frac x{12}-\frac{x^3}{720}+O(x^5).
\]
式\eqref{eq:columnpotential}中的$1/x$先相消，平均高度为
\[
 \langle h\rangle=H\left(\frac1x-F(x)\right)
 =H\left(\frac12-\frac x{12}+\frac{x^3}{720}+O(x^5)\right).
\]
因此首两项为$H/2-MgH^2/(12RT)$。热容须对能量求导，不能把高度展开直接当成热容展开。对$F$逐项求导，
\[
 F'(x)=-\frac1{x^2}+\frac1{12}-\frac{x^2}{240}+O(x^4).
\]
由$E_p=\nu RT-\nu MgHF(x)$及$\dd x/\dd T=-x/T$，有
\[
 \frac{\dd E_p}{\dd T}=\nu R+\frac{\nu MgHx}{T}F'(x)
 =\nu R[1+x^2F'(x)].
\]
加上平动热容，得到
\[
 C=\nu R\left[\frac52+x^2F'(x)\right]
 =\nu R\left[\frac32+\frac{x^2}{12}-\frac{x^4}{240}+O(x^6)\right].
\]
有限顶壁使平均势能在高温下趋于$\nu MgH/2$，附加热容从正值趋零。右图$H\gg H_T$时，$F(x)=\ee^{-x}[1+O(\ee^{-x})]$，顶壁影响指数小，$E_p\simeq\nu RT$、$\langle h\rangle\simeq RT/(Mg)$，故$C\simeq5\nu R/2$。

\Diagram{fig:columnlimita}{气柱热容随无量纲高度单调升高}{
 \draw[axis](0,0)--(5.5,0)node[right]{$x$};\draw[axis](0,0)--(0,3.1)node[above]{$C/(\nu R)$};
 \draw[curve]plot[domain=.05:6,samples=80]({.8*\x},{2.5-\x*\x*exp(-\x)/(1-exp(-\x))^2});\draw[guide](0,2.5)--(5,2.5);\node[left]at(0,1.5){$3/2$};\node[left]at(0,2.5){$5/2$};
}
图\ref{fig:columnlimita}的单调性也可由公式确定。将修正分式写成$h(x)^2$，其中$h(x)=x/[2\sinh(x/2)]$。令$y=x/2>0$，商法则给
\[
 h'(x)=\frac{\sinh y-y\cosh y}{2\sinh^2y}.
\]
分子反号后的函数$f(y)=y\cosh y-\sinh y$满足$f(0)=0$、$f'(y)=y\sinh y>0$，故$h'<0$，于是$\dd[C/(\nu R)]/\dd x=-2hh'>0$，曲线从$3/2$单调升至$5/2$。

固定有限$H$先升到高温，粒子高度受顶壁限制，额外热容趋零；先把顶壁移到无穷远，平均高度可以随温度不断增长，始终保留$\nu R$的重力热容。两种操作改变了不同的几何约束，极限不交换。
\end{solution}

\section{界面、相变与外场的综合模型}
液膜拉伸时，外界做功先与表面自由能相联系；计算供热还需表面内能。外场中的气体则同时涉及热运动、势能与外界机械功。以下各题按图中的系统边界分别列出这些能量项。

\begin{example}[非线性表面张力与相容热容]
界面满足 $\sigma=\sigma_0-a_sT^2$，$a_s>0$，热容在零面积参考处为常量 $C_0$。忽略其他广义功。求相容的 $C_\Sigma(\Sigma,T)$、$U$ 与 $S$；说明能否同时规定定面积热容处处等于 $C_0$。
\Diagram{fig:nonlinearsurfaceq}{表面张力曲率影响面积方向的热容变化}{
 \draw[axis](0,0)--(4.6,0)node[right]{$T$};\draw[axis](0,0)--(0,2.8)node[above]{$\sigma$};\draw[curve]plot[domain=0:4](\x,{2.5-.1*\x*\x});\node at(3,2.6){$\sigma_{TT}=-2a_s$};
}
\end{example}
\begin{solution}
图\ref{fig:nonlinearsurfaceq}的 $U_\Sigma=\sigma-T\sigma_T=\sigma_0+a_sT^2$，先对面积积分，$U=(\sigma_0+a_sT^2)\Sigma+\Phi(T)$。在零面积参考处 $U_T=C_0$，故 $\Phi=C_0T+\text{常量}$。因此
\[
 C_\Sigma=C_0+2a_sT\Sigma,
 \qquad U=\sigma_0\Sigma+a_s\Sigma T^2+C_0T+\text{常量}.
\]
Maxwell 关系给 $S_\Sigma=-\sigma_T=2a_sT$，积分并匹配 $S_T=C_\Sigma/T$，得到
\[
 S=2a_s\Sigma T+C_0\ln T+\text{常量}.
\]
给定的 $C_0$ 用来确定温度积分函数 $\Phi(T)$，并非规定所有面积上的热容都相同。对 $U$ 中的 $a_s\Sigma T^2$ 求温度导数，产生不可略去的 $2a_s\Sigma T$。再用所得熵求回基本关系，温度方向有 $T S_T=C_0+2a_s\Sigma T$，面积方向有 $T S_\Sigma+\sigma=\sigma_0+a_sT^2$，分别与内能的两个偏导匹配。两个方向同时匹配，才说明 $U,S$ 属于同一个相容模型。

\Diagram{fig:nonlinearsurfacea}{两种顺序求混合偏导必须一致}{
 \node[draw]at(0,0){$\partial_\Sigma C_\Sigma$};\node at(2.3,0){$=$};\node[draw]at(4.6,0){$-T\sigma_{TT}=2a_sT$};
}
图\ref{fig:nonlinearsurfacea}的等式要求热容随面积变化。若处处规定 $C_\Sigma=C_0$，左侧为零而右侧非零，模型便不相容。线性 $\sigma(T)$ 的二阶温度导数为零，才允许采用面积无关的热容。零面积参考值用于指定积分函数，物理应用范围仍须满足界面可作连续描述且 $\sigma>0$。
\end{solution}

\begin{example}[封闭肥皂泡的等温膨胀]
肥皂泡中有 $\nu$ 摩尔理想气体，温度保持 $T$。薄膜张力 $\sigma=\sigma_0+\alpha T$，两面面积近似相同，泡半径由 $r_1$ 缓慢变为 $r_2$。外压受控制，忽略膜厚、重力和膜体积功。求平衡外压、气体做功、对外负载做功及气体和膜整体吸热。
\Diagram{fig:bubbleq}{泡内外压强差由两层界面提供}{
 \draw[curve](0,0)circle(1.5);\draw[curve](0,0)circle(1.57);\draw[axis](0,0)--(1.5,0)node[midway,above]{$r$};\node at(0,-.65){$p_{\rm in}$};\node at(2.8,.7){$p_{\rm ext}$};\node at(3,-.2){$p_{\rm in}-p_{\rm ext}=4\sigma/r$};
}
\end{example}
\begin{solution}
图\ref{fig:bubbleq}中 $V=4\pi r^3/3$、实际界面面积 $\Sigma=8\pi r^2$，理想气体压强 $p_{\rm in}=\nu RT/V$。机械平衡要求
\[
 p_{\rm ext}(r)=\frac{3\nu RT}{4\pi r^3}-\frac{4\sigma}{r}.
\]
路径只能在外压条件可实现的半径区间运行。气体等温做功为
\[
 W_g=\nu RT\int_{V_1}^{V_2}\frac{\dd V}{V}
 =\nu RT[\ln V]_{V_1}^{V_2}=3\nu RT\ln\frac{r_2}{r_1}.
\]
其中一部分增加膜的自由能，另一部分付给外压负载：
\[
 W_{\rm ext}=\int p_{\rm ext}\dd V
 =W_g-\sigma\Delta\Sigma,
 \quad \Delta\Sigma=8\pi(r_2^2-r_1^2).
\]
膜所需的功也可从气体内外压差直接积分。球体体积微元为 $\dd V=4\pi r^2\dd r$，因此
\[
 \int_{r_1}^{r_2}(p_{\rm in}-p_{\rm ext})\dd V
 =\int_{r_1}^{r_2}16\pi\sigma r\dd r
 =[8\pi\sigma r^2]_{r_1}^{r_2}=8\pi\sigma(r_2^2-r_1^2).
\]
这与两面总面积的自由能增量相同。气体把这部分能量交给膜，膜又与恒温热源交换热；对气体和膜的整体列式时，这笔内部功抵消，外界负载只收到由外压积分计算的那一部分。

\Diagram{fig:bubblea}{以气体和膜为整体时，膜功属于内部能量交换}{
 \node[draw]at(0,0){气体 $W_g$};\draw[axis](1.2,.2)--(2.8,1);\draw[axis](1.2,-.2)--(2.8,-1);\node[right]at(2.8,1){膜自由能增加 $\sigma\Delta\Sigma$};\node[right]at(2.8,-1){外界负载 $W_{\rm ext}$};
}
气体内能不变，膜内能增量为 $\sigma_0\Delta\Sigma$。按图\ref{fig:bubblea}选择整体边界，只有 $W_{\rm ext}$ 是对外功，故
\[
 Q=\Delta U_{\rm total}+W_{\rm ext}
 =W_g-\alpha T\Delta\Sigma.
\]
也可由整体熵变 $\Delta S=\nu R\ln(V_2/V_1)-\alpha\Delta\Sigma$ 得 $Q=T\Delta S$。若把气体做功和膜功都再次计为整体对外功，会重复计算内部交换。
\end{solution}

\begin{example}[定容两相系统的表观热容]
$\nu$ 摩尔物质在刚性容器中处于液气共存，摩尔平均体积为 $\bar v=V/\nu$。气相理想，液体体积可忽略，气相摩尔定容热容为 $c_{v,g}$，液相比热为 $c_l$。沿共存线缓慢加热，摩尔汽化热为 $L_m(T)$。求气相摩尔分数和系统定容供热率 $C=\delta Q/\dd T$。
\Diagram{fig:twophasecapacityq}{容器定容，液气比例随温度改变}{
 \draw(0,0)rectangle(3,2.6);\fill[gray!20](0,0)rectangle(3,.8);\draw[curve](0,.8)--(3,.8);\node at(1.5,.4){液相};\node at(1.5,1.7){气相 $x\nu$};\draw[axis](-1,1)--(0,1)node[midway,above]{$Q$};
}
\end{example}
\begin{solution}
图\ref{fig:twophasecapacityq}的气体占据几乎全部体积，故 $V\simeq x\nu RT/p_s(T)$，于是
\[
 x(T)=\frac{\bar v p_s(T)}{RT},\qquad
 \frac{x'}x=\frac{L_m(T)}{RT^2}-\frac1T.
\]
只在 $0<x<1$ 且两相近似成立的区间使用。总内能为 $U=\nu[(1-x)u_l+xu_g]$，沿共存线求导要保留组分比例的变化：
\[
 \frac{\dd U}{\dd T}=\nu[(1-x)c_l+xc_{v,g}+(u_g-u_l)x'].
\]
由汽化焓差 $L_m=u_g-u_l+RT$，得到
\[
 C=\nu\left[(1-x)c_l+xc_{v,g}
 +(L_m-RT)x\left(\frac{L_m}{RT^2}-\frac1T\right)\right].
\]
由 $x'/x=(L_m-RT)/(RT^2)$，相比例项还可化成
\[
 (L_m-RT)x'=\frac{x(L_m-RT)^2}{RT^2}\ge0.
\]
因此在本题近似下，这一额外储热能力非负。它与气相比例改变的方向无关：升温可使 $x$ 增大，也可在某些参数范围使其减小，相变对应的内能变化会同时换号。容器一旦完全气化或液化，$x$ 不能再按共存线公式继续变化，需改用对应单相的状态方程和热容。

\Diagram{fig:twophasecapacitya}{供热同时改变各相温度和相的数量}{
 \node[draw,align=center]at(0,0){显热项\\$(1-x)c_l+xc_{v,g}$};\node at(2.8,0){$+$};\node[draw,align=center]at(5.1,0){相比例项\\$(L_m-RT)x'$};
}
图\ref{fig:twophasecapacitya}中的相比例项可能很大。容器定容只使外壁体积功为零，内部相变仍会贡献热容。潜热 $L_m$ 是两相焓差，扣除气相形成所需的 $RT$ 后，才得到内能差 $L_m-RT$，故应由它乘以 $x'$。更精确的模型还须保留液体体积及各相沿共存线的压强效应。
\end{solution}

\begin{example}[临界点与等面积条件]
对一摩尔 van der Waals 气体求临界参数。低于临界温度时，说明为什么不能把等温线中 $p_v|_T>0$ 的部分作为稳定平衡路径，并从两相 Gibbs 自由能相等导出等面积条件。
\Diagram{fig:criticalvdwq}{低于临界温度的等温线与共存平台}{
 \draw[axis](0,0)--(5.5,0)node[right]{$v/v_c$};\draw[axis](0,0)--(0,3)node[above]{$p/p_c$};
 \draw[curve]plot[domain=.55:4,samples=120]({1.2*\x},{2*(7.2/(3*\x-1)-3/(\x*\x))});
 \draw[guide](.7241,1.294)--(2.8186,1.294);\node at(3.7,2.5){$T=0.9T_c$};\node[below]at(.7241,0){$v_l/v_c$};\node[below]at(2.8186,0){$v_g/v_c$};
}
\end{example}
\begin{solution}
临界等温线在拐点有 $p_v=0$、$p_{vv}=0$：
\[
 -\frac{RT_c}{(v_c-b)^2}+\frac{2a}{v_c^3}=0,\qquad
 \frac{2RT_c}{(v_c-b)^3}-\frac{6a}{v_c^4}=0.
\]
两式相除消去 $T_c$，得 $2/(v_c-b)=3/v_c$，所以
\[
 v_c=3b,\quad T_c=\frac{8a}{27Rb},\quad p_c=\frac a{27b^2}.
\]
图\ref{fig:criticalvdwq}中向右上升的曲线段满足等温压缩率为负，不能构成稳定均匀相。平衡时由液气两相共存平台替代。
等面积条件不是对实际两相过程使用不稳定均匀状态，而是借数学状态方程比较两端纯相的自由能。实际共存状态由相比例改变连接两端，压强保持 $p_s$。对于平均摩尔体积 $\bar v$ 介于两端之间，气相摩尔分数由 $\bar v=(1-x)v_l+xv_g$ 给出，故 $x=(\bar v-v_l)/(v_g-v_l)$。平台内部的不同横坐标代表不同相比例，而不是均匀气体沿回折线逐点移动。

\Diagram{fig:criticalvdwa}{平台上下的有向面积必须相抵}{
 \draw[axis](0,0)--(5.4,0)node[right]{$v/v_c$};\draw[axis](0,0)--(0,2.6)node[above]{$p/p_c$};
 \draw[curve]plot[domain=.6034:2.34884,samples=100]({1.8*\x},{2.5*(7.2/(3*\x-1)-3/(\x*\x))});
 \draw[curve](1.0861,1.6175)--(4.2279,1.6175);\node at(1.7,.8){负面积};\node at(3.2,2){正面积};
}
在定温下 $\dd g=v\dd p$，沿数学等温线从 $v_l$ 积分到 $v_g$，分部积分得到
\[
 g_g-g_l=[vp]_{v_l}^{v_g}-\int_{v_l}^{v_g}p(v,T)\dd v.
\]
两端共存压强均为 $p_s$，且 $g_g=g_l$，所以
\[
 \int_{v_l}^{v_g}[p(v,T)-p_s]\dd v=0.
\]
这就是图\ref{fig:criticalvdwa}的等面积条件。与两端方程 $p(v_l,T)=p(v_g,T)=p_s$ 联立，可确定液、气摩尔体积与饱和压强。图中取 $T/T_c=0.9$，对应 $v_l/v_c\approx0.60340$、$v_g/v_c\approx2.34884$、$p_s/p_c\approx0.64700$，平台按这些数值绘制。
\end{solution}

\begin{example}[有限气柱的等温升顶做功]
气柱截面积 $S$ 固定，内有 $N$ 个理想气体粒子，处于均匀温度 $T$ 的重力平衡。顶壁可缓慢从高度 $H_1$ 升至 $H_2$，粒子数固定，过程可逆等温。求顶壁做功和吸热；讨论 $H_2\to\infty$ 时的功是否发散。
\Diagram{fig:raisetopq}{升高顶壁时只受顶端局部压强做功}{
 \draw(0,0)--(0,3.4);\draw(2.2,0)--(2.2,3.4);\draw(0,0)--(2.2,0);\draw[line width=1.4pt](0,2)--(2.2,2);\draw[guide](0,3.1)--(2.2,3.1);\draw[axis](1.1,2.1)--(1.1,3);\node[right]at(2.4,2){$H_1$};\node[right]at(2.4,3.1){$H_2$};
}
\end{example}
\begin{solution}
由有限气柱归一化，图\ref{fig:raisetopq}的顶端压强为$p_t=Nmg/[S(\ee^{a_gH}-1)]$，$a_g=mg/(\kb T)$。气体向外做功只在顶壁。令$x=a_gH$，则$\dd H=\dd x/a_g$，上下限由$H_1,H_2$变为$a_gH_1,a_gH_2$，所以
\begin{align*}
 W&=\int_{H_1}^{H_2}Sp_t\dd H
 =N\kb T\int_{a_gH_1}^{a_gH_2}\frac{\dd x}{\ee^x-1}\\
 &=N\kb T\ln\frac{1-\ee^{-a_gH_2}}{1-\ee^{-a_gH_1}}.
\end{align*}
最后一步利用 $\dd\ln(1-\ee^{-x})/\dd x=1/(\ee^x-1)$。温度不变，平动内能不变，重力势能却改变，
\[
 \Delta E_p=Nmg\left[\frac{H_1}{\ee^{a_gH_1}-1}
 -\frac{H_2}{\ee^{a_gH_2}-1}\right].
\]
吸热中的重力势能项取决于平均高度，顶壁功则由顶端的局部压强决定。也可用配分积分求熵：固定温度下，位置积分 $Z_h=(1-\ee^{-a_gH})/a_g$ 随顶壁高度改变，熵差为 $N\kb\ln(Z_{h,2}/Z_{h,1})+\Delta E_p/T$。乘以温度，第一项给出顶壁功，第二项给出势能增量，与第一定律一致。位置积分在无限高时收敛，功也趋于有限值。

\Diagram{fig:raisetopa}{等温供热同时支付重力势能和顶壁功}{
 \node[draw]at(0,0){$Q$};\draw[axis](.7,.2)--(2.3,1);\draw[axis](.7,-.2)--(2.3,-1);\node[right]at(2.3,1){$\Delta E_p$};\node[right]at(2.3,-1){$W$};
}
图\ref{fig:raisetopa}给 $Q=\Delta E_p+W$。$H_2\to\infty$ 时功趋于有限值 $-N\kb T\ln(1-\ee^{-a_gH_1})$，因为高处密度指数小，后续顶壁压强很快趋零。忽略重力时的 $N\kb T\ln(H_2/H_1)$ 发散结论不能用于这个极限；应先比较尺度高度与容器高度。
\end{solution}

\begin{example}[恒定角速度下的旋转气体热容]
圆柱容器在电机维持的恒定角速度 $\Omega$ 下转动，半径 $R_c$，内有 $N$ 个质量为 $m$ 的单原子理想气体粒子。忽略重力，气体相对容器始终处于温度 $T$ 的平衡。求平均转动惯量，正确计入电机功后求供热热容。
\Diagram{fig:rotcapacityq}{固定角速度下的供热与有符号机械功}{
 \draw[curve](0,0)circle(1.6);\draw[->](77:1.9)arc(77:170:1.9);\node at(-1.5,2.2){$\Omega$ 固定};
 \draw[<->](2.2,0)--(3.6,0)node[right]{电机 $\Omega\dd L_z$};\draw[axis](-3,0)--(-1.7,0)node[midway,above]{$Q$};
}
\end{example}
\begin{solution}
令 $a=m\Omega^2/(2\kb T)$、$b_\Omega=aR_c^2$。径向密度 $\propto\ee^{ar^2}$，体积元含 $r\dd r$。对归一化积分求参数导数，
\[
 \langle r^2\rangle=\frac{\dd}{\dd a}\ln\frac{\ee^{aR_c^2}-1}{a}
 =\frac{R_c^2}{1-\ee^{-b_\Omega}}-\frac1a.
\]
转动惯量 $I=Nm\langle r^2\rangle$，角动量 $L_z=I\Omega$。实验室系总能量为平动热能加整体旋转能，$E=3N\kb T/2+I\Omega^2/2$。图\ref{fig:rotcapacityq}中电机为保持 $\Omega$ 还要做功 $\delta W_{\rm on}=\Omega\dd L_z=\Omega^2\dd I$，故
\[
 \delta Q=\dd E-\Omega\dd L_z
 =\frac32N\kb\dd T-\frac12\Omega^2\dd I.
\]
将惯量写为$I=NmR_c^2/(1-\ee^{-b_\Omega})-2N\kb T/\Omega^2$。因$\dd b_\Omega/\dd T=-b_\Omega/T$，第一项求导为$NmR_c^2b_\Omega\ee^{-b_\Omega}/[T(1-\ee^{-b_\Omega})^2]$，第二项为$-2N\kb/\Omega^2$。再用$mR_c^2=2\kb Tb_\Omega/\Omega^2$，合并得
\[
 \frac{\dd I}{\dd T}=-\frac{2N\kb}{\Omega^2}
 \left[1-\frac{b_\Omega^2\ee^{-b_\Omega}}
 {(1-\ee^{-b_\Omega})^2}\right].
\]
括号非负，故升温时径向密度趋于均匀，平均惯量减小。角速度固定，电机须取走一部分角动量，电机对气体做的功因而为负。把这一有符号的机械功从实验室系能量增量中减去，才得到供热量。

\Diagram{fig:rotcapacitya}{升温使径向位置分布变均匀，同时与电机交换功}{
 \draw[axis](0,0)--(4.8,0)node[right]{$s=r^2/R_c^2$};\draw[axis](0,0)--(0,4.7)node[above]{$q(s)$};
 \draw[curve]plot[domain=0:1,samples=70]({4*\x},{exp(\x-1)/(1-exp(-1))});
 \draw[curve,dashed]plot[domain=0:1,samples=70]({4*\x},{4*exp(4*(\x-1))/(1-exp(-4))});
 \node[below]at(4,0){$1$};\node at(1.3,1.6){高温};\node at(2.8,3.7){低温};
 \node at(2.4,-.85){$\delta Q=\dd E-\Omega\dd L_z$};
}
将已经求出的$\dd I/\dd T$代入图\ref{fig:rotcapacitya}对应的供热式，得到
\[
 C_\Omega=N\kb\left[\frac52-
 \frac{b_\Omega^2\ee^{-b_\Omega}}{(1-\ee^{-b_\Omega})^2}\right].
\]
弱旋转极限回到 $3N\kb/2$，强旋转时边壁附近的分布重排再贡献约 $N\kb$。该供热热容已经扣除了电机的机械交换。若装置改为与外界隔绝、角动量固定，升温时 $\Omega$ 也会改变，需要按新的约束重新求导。
\end{solution}

\begin{example}[重力气柱从低处泄流可以升温]
高 $H$、截面积 $S$ 的绝热刚性气柱中有 $N$ 个单原子理想气体粒子，温度 $T$ 均匀。高度 $h_0$ 处开小孔向真空缓慢泄流，气柱始终重新平衡，重力势能计入系统能量。写出 $T(N)$ 的微分关系；再求无限高气柱从底部泄流的温度粒子数关系。
\Diagram{fig:gravityleakq}{出口高度决定粒子携带的外场势能}{
 \draw(0,0)rectangle(2,3.2);\fill[white](1.94,1.05)rectangle(2.06,1.35);
 \draw[axis](2,1.2)--(4,1.2);\draw[guide](-.3,1.2)--(2,1.2);\node[left]at(-.3,1.2){$h_0$};
 \node at(3.2,1.9){$2\kb T+mgh_0$};\node at(1,-.5){势能零点在底部};
}
\end{example}
\begin{solution}
图\ref{fig:gravityleakq}每个逸出粒子平均携带 $\varepsilon_e=2\kb T+mgh_0$。余气单粒子平均总能量 $u(T)=3\kb T/2+mg\langle h\rangle$，固定 $H$ 的单粒子热容 $c(T)=\dd u/\dd T$。开放系统能量守恒为
\[
 u\dd N+Nc\dd T=\varepsilon_e\dd N,
 \qquad Nc\dd T=
 \left[\frac12\kb T+mg(h_0-\langle h\rangle)\right]\dd N.
\]
因此泄流是否冷却取决于出口能量相对于体内平均总能量的高低，不能只比较平动能。若 $h_0<\langle h\rangle-\kb T/(2mg)$，右侧方括号为负，因 $\dd N<0$，温度反而上升。
固定的重力场不显含时间，所以粒子穿孔时只把其所在高度的势能带出，不再额外计一次外场参数功。余气能量的全微分同时含 $u\dd N$ 和 $Nc\dd T$；前者是按体内平均能量减少粒子的贡献，后者才反映温度改变。底部逸出的粒子虽然平动能高于体内平动平均，却缺少体内平均具有的重力势能。比较总能量后，差值反而为负，这才使无限高气柱的底部泄流导致升温。

\Diagram{fig:gravityleaka}{无限高气柱底部逸出的总能量低于体内平均}{
 \node[draw,align=center]at(0,0){体内平均\\$5\kb T/2$};\node[draw,align=center]at(4.6,0){底部逸出\\$2\kb T$};\node at(2.3,-1){留下的粒子平均总能量升高};
}
无限高气柱有$mg\langle h\rangle=\kb T$、$c=5\kb/2$。从底部$h_0=0$泄流，图\ref{fig:gravityleaka}给
\[
 \frac52N\kb\dd T=-\frac12\kb T\dd N,
 \qquad T=T_0(N/N_0)^{-1/5}.
\]
分离变量后的初值积分为$5\int_{T_0}^{T}\dd\vartheta/\vartheta=-\int_{N_0}^{N}\dd N'/N'$，即$5\ln(T/T_0)=-\ln(N/N_0)$，给出上式。余气升温的原因是底部逸出粒子的总能量低于体内平均值。无限高度近似只适用于尺度高度远小于容器高度的阶段；温度升高使尺度高度增大后，应改用有限$H$的微分式描述。
\end{solution}

\begin{example}[光子工质的卡诺循环]
以平衡光子气体为工质，工作在 $T_h,T_c$ 之间。高温等温膨胀体积从 $V_A$ 到 $V_B$，随后可逆绝热到低温，再等温压缩和绝热返回。求其余两状态体积、两热源热量与循环功，验证卡诺效率。
\Diagram{fig:photoncarnotq}{光子气体的等温线是水平线}{
 \draw[axis](0,0)--(5.2,0)node[right]{$V$};\draw[axis](0,0)--(0,3.6)node[above]{$p$};
 \draw[curve](1,3)--(2,3);\draw[curve]plot[domain=2:3.90625](\x,{3*(2/\x)^(4/3)});
 \draw[curve](1.953125,1.2288)--(3.90625,1.2288);\draw[curve]plot[domain=1:1.953125](\x,{3/\x^(4/3)});
 \node[above]at(1,3){$A$};\node[above]at(2,3){$B$};\node[right]at(3.90625,1.2288){$C$};\node[left]at(1.953125,1.2288){$D$};\node at(3.9,3){$T_c/T_h=0.8$ 示意};
}
\end{example}
\begin{solution}
图\ref{fig:photoncarnotq}的绝热关系是 $VT^3=\text{常量}$，所以
\[
 V_C=V_B(T_h/T_c)^3,\qquad V_D=V_A(T_h/T_c)^3.
\]
高温等温段压强为 $aT_h^4/3$，内能也随体积增加。由 $Q=\Delta U+W$，
\[
 Q_h=\frac{4a}3T_h^4(V_B-V_A).
\]
低温段放热量大小 $Q_c=(4a/3)T_c^4(V_C-V_D)$，代入体积关系得到 $Q_c/Q_h=T_c/T_h$。
高温等温膨胀中，光子工质的熵增为$\Delta S_h=(4a/3)T_h^3(V_B-V_A)$。两条绝热段保持各端点的熵，因而低温压缩的熵减大小相同。低温体积差是高温体积差的$(T_h/T_c)^3$倍，而能量密度缩小为$(T_c/T_h)^4$，相乘后热量比只剩一次温度比。转到图\ref{fig:photoncarnota}的熵温平面，矩形面积为$(T_h-T_c)\Delta S_h$，等于两热源热量之差；原压强—体积图中的循环并非矩形。

\Diagram{fig:photoncarnota}{在熵温平面中不同工质具有相同的可逆循环结构}{
 \draw[axis](0,0)--(4.8,0)node[right]{$S$};\draw[axis](0,0)--(0,3.2)node[above]{$T$};\draw[curve](1,1)rectangle(3.6,2.5);\node[left]at(0,2.5){$T_h$};\node[left]at(0,1){$T_c$};\draw[axis](1.5,2.5)--(2.5,2.5);\draw[axis](3,1)--(2,1);\node at(2.3,1.8){循环功};
}
因此$W=Q_h(1-T_c/T_h)$，效率仍为卡诺值。图\ref{fig:photoncarnota}的矩形由两条等温线和两条等熵线组成；光子状态方程改变了各状态的体积及热量，不改变这一可逆循环的温度比效率。
\end{solution}

\section{磁性工质的熵与绝热降温}
\subsection{由两态分布接到广义功}
考虑$N$个相互独立的局域磁矩，磁矩大小为$\mu$，在外加磁感应强度$B>0$中有两个能级$\mp\mu B$。相互作用及内部磁场暂忽略，外场作为受控参数。令$x=\mu B/(\kb T)$，单个配分和为$2\cosh x$，于是
\[
 F_s=-N\kb T\ln(2\cosh x),\quad
 M=-\left(\frac{\partial F_s}{\partial B}\right)_T=N\mu\tanh x,
\]
\[
 U_s=-N\mu B\tanh x,\qquad
 S_s=N\kb[\ln(2\cosh x)-x\tanh x].
\]
这里$M$是样品沿场方向的总磁矩。所选系统能量包含磁矩与外加场的耦合能，磁场发生器及其场能在系统之外，相应热力学微分为$\dd U_s=T\dd S_s-M\dd B$。若改用单位体积磁化强度或把场源纳入系统，须同时调整能量与功的定义。

先有$\partial_Bx|_T=\mu/(\kb T)=x/B$，所以$M=-F_B=N\kb T\tanh x\,(x/B)=N\mu\tanh x$。从能级平均得到$U_s=-N\mu B\tanh x$后，用$S_s=(U_s-F_s)/T$得到上述熵。固定$B$时$\partial_Tx=-x/T$，因此$\partial_TU_s=-N\mu B\operatorname{sech}^2x(-x/T)$，再用$\mu B=\kb Tx$，得到
\[
 C_{B,s}=N\kb x^2\operatorname{sech}^2x.
\]
熵对$x$的导数为$N\kb[\tanh x-(\tanh x+x\operatorname{sech}^2x)]=-N\kb x\operatorname{sech}^2x$。再分别乘$x_T=-x/T$、$x_B=x/B$，得到
\[
 \dd S_s=\frac{C_{B,s}}T\dd T-\frac{C_{B,s}}B\dd B.
\]
对上述独立磁矩，可逆绝热过程满足$\dd T/T=\dd B/B$，即$B/T$保持不变，两个能级的占有比例也保持不变。推广到其他磁性物质或温度范围前，须重新计入相互作用及其他自由度的熵。
\ExpansionDiagram{fig:dep-th4-spins}{外场改变两能级间隔，热平衡占有比由$B/T$决定}{
\draw[->](0,0)--(0,3)node[above]{能量};\draw[thick](.6,.5)--(2.6,.5)node[right]{$-\mu B$};\draw[thick](.6,2.3)--(2.6,2.3)node[right]{$+\mu B$};\draw[<->](1.6,.6)--(1.6,2.2)node[midway,right]{$2\mu B$};\node[right]at(4,1.4){$N_+/N_-=\ee^{-2x}$};}
图\ref{fig:dep-th4-spins}中的能级标记以能量为准：$N_+$指上能级占有数，$N_-$指下能级。磁矩沿场方向对应低能级，不能把磁矩方向的正号与能级正号混用。

\subsection{晶格热容限制实际温降}
若样品还有与磁矩保持内部热平衡的晶格，晶格热容为$C_l(T)$且暂不依赖$B$，总熵微分为
\[
 \dd S=\frac{C_l+C_{B,s}}T\dd T-\frac{C_{B,s}}B\dd B.
\]
可逆绝热条件给
\[
 \left(\frac{\dd T}{\dd B}\right)_S
 =\frac TB\frac{C_{B,s}}{C_l+C_{B,s}}.
\]
晶格必须随磁矩一起降温，温降幅度因而小于独立自旋系统的$T\propto B$结果。若晶格与磁矩尚未来得及热平衡，则单一温度$T$已不足以描述过程，应另列两子系统的能量交换方程。

\begin{example}[有限晶格热容的绝热去磁]
样品的磁矩服从上述两态模型，晶格热容近似常数$C_l>0$。初态为$(B_i,T_i)$，其中$x_i=\mu B_i/(\kb T_i)$。样品与环境绝热并足够缓慢地减小外场到零，求理想模型给出的最终温度。特别取$x_i=1$、$C_l=N\kb$。
\ExpansionDiagram{fig:dep-th4-demag-q}{磁矩与晶格内部换热，外部绝热，仅改变受控磁场}{
\draw[dashed](-.4,-.4)rectangle(6,2);\draw(0,0)rectangle(2.2,1.5)node[midway]{磁矩};\draw(3,0)rectangle(5.4,1.5)node[midway]{晶格};\draw[<->](2.2,.7)--(3,.7);\node[above]at(2.7,2.1){总熵守恒};\draw[->](6.8,1.3)--(6.8,.1)node[below]{$B\downarrow$};}
\end{example}
\begin{solution}
图\ref{fig:dep-th4-demag-q}要求保存磁矩与晶格的总熵。晶格熵差是$C_l\ln(T_f/T_i)$。只要$T_f>0$，末场$B_f=0$时$x_f=0$，两个能级等概率，磁矩熵为$N\kb\ln2$。初末熵相等给
\[
 C_l\ln\frac{T_f}{T_i}+N\kb\ln2
 =N\kb[\ln(2\cosh x_i)-x_i\tanh x_i].
\]
所以
\[
 \frac{T_f}{T_i}=
 \exp\!\left\{\frac{N\kb}{C_l}
 [\ln\cosh x_i-x_i\tanh x_i]\right\}.
\]
括号在$x_i>0$时为负：其导数是$\tanh x_i-[\tanh x_i+x_i\operatorname{sech}^2x_i]=-x_i\operatorname{sech}^2x_i<0$，且原点为零。因此$0<T_f<T_i$，但有限正$C_l$使这个近似模型中的$T_f$不为零。
\ExpansionDiagram{fig:dep-th4-demag-a}{磁矩熵增加由晶格熵减少补偿；总熵保持不变}{
\draw[->](0,0)--(6,0)node[right]{过程};\draw[->](0,0)--(0,3.9)node[above]{$S$};\draw[thick](.5,1)--(5,2.2)node[right]{$S_s$};\draw[dashed](.5,2.4)--(5,1.2)node[right]{$S_l$};\draw[dotted](.5,3.4)--(5,3.4)node[right]{$S_{\rm total}$};}
图\ref{fig:dep-th4-demag-a}中磁矩熵的增加由晶格熵的减少补偿，线段仅表示变化趋势；温度由总熵方程求得。取$x_i=1,C_l=N\kb$，有$\ln\cosh1\simeq0.433781$、$\tanh1\simeq0.761594$，故$T_f/T_i\simeq\ee^{-0.327813}\simeq0.7205$。

若错误地仅保留磁矩子系统并令$T/B$不变，会预测$B\to0$时温度也趋零，遗漏了冷却晶格所需的熵转移。更低温时晶格热容并非常数，磁矩相互作用和内部场也可能不可忽略；届时应把实际$C_l(T)$积分写入总熵方程，而不是无限外推本题的常热容近似。
\end{solution}

\subsection{降温与系统能量降低不是同一个条件}
对只有磁矩的可逆绝热分支，$x$固定，$U_s=-N\mu B\tanh x$。减小$B$使$T$下降，却使负的$U_s$向零增加；由$\dd U_s=-M\dd B$，增加的能量来自受控外场参数的做功。温度反映占有分布相对能级间隔的尺度，不能脱离能量零点和外参变化，直接把降温解释为内能必定减少。

若要用这一工质进行制冷，应先等温磁化向环境排热，再绝热减场降温，随后与较冷负载接触吸热。各段固定的量不同，热量分别按$T\dd S$积分，磁场控制功按本节约定计入。只有把样品、负载与场源的边界交代清楚，才能比较一个完整制冷循环的能量与熵。

\section{习题与解答}
\begin{exercise}[等温混合大气的组分分层]
两种不发生反应的理想气体处于温度 $T$ 的重力平衡，摩尔质量为 $M_1,M_2$，地面分压为 $p_{10},p_{20}$。求各分压、摩尔分数与高度的关系，并讨论较重组分的比例。
\Diagram{fig:atmixq}{两种质量的粒子共享温度而尺度高度不同}{
 \draw(0,0)rectangle(2.6,3);\node at(1.3,.4){$p_{10},p_{20}$};\draw[guide](0,2)--(2.6,2);\node[right]at(2.8,2){$h$};\draw[axis](-.6,0)--(-.6,3.5);
}
\end{exercise}
\begin{solution}
各组分分别满足静力平衡，图\ref{fig:atmixq}中 $p_i(h)=p_{i0}\ee^{-M_igh/(RT)}$。因此 $x_i(h)=p_i(h)/[p_1(h)+p_2(h)]$，组分比为
\[
 \frac{x_2(h)}{x_1(h)}=\frac{p_{20}}{p_{10}}
 \exp\!\left[-\frac{(M_2-M_1)gh}{RT}\right].
\]
同一种外场中，不同组分的势能增量按各自摩尔质量计算；温度相同只意味着它们共享热能尺度 $RT$，并不使势能指数相同。总压 $p_1+p_2$ 通常不是用一个固定平均摩尔质量写成的单一指数。若 $M_2>M_1$，组分2的尺度高度更小，其相对比例随高度下降；若两摩尔质量恰相同，则指数相消，摩尔分数保持地面比例，恢复未发生质量分层的特例。

\Diagram{fig:atmixa}{较重组分的相对浓度随高度下降}{
 \draw[axis](0,0)--(4.7,0)node[right]{$h$};\draw[axis](0,0)--(0,2.8)node[above]{$x_2/x_1$};\draw[curve]plot[domain=0:4.3](\x,{2.3*exp(-.4*\x)});\node at(3,2.4){$M_2>M_1$};
}
图\ref{fig:atmixa}对应静态的分子扩散平衡，各组分按自己的摩尔质量分层。强湍流可以使较大高度范围内的组分比例近乎均匀，但这需要持续的混合作用，超出了此处静态模型的条件。
\end{solution}

\begin{exercise}[温度随高度线性下降的大气]
气体处于静力平衡，给定 $T(h)=T_0-\beta h>0$，$\beta>0$。摩尔质量 $M$、重力 $g$ 恒定。求 $p(h)$、$n(h)$。若局部满足理想气体的干绝热梯度，$\beta$ 应是多少？
\Diagram{fig:lapseq}{温度剖面由外部条件给定}{
 \draw[axis](0,0)--(4.6,0)node[right]{$T$};\draw[axis](0,0)--(0,3.2)node[above]{$h$};\draw[curve](4,0)--(1,2.7);\node[below]at(4,0){$T_0$};\node at(3,2.5){$T_0-\beta h$};
}
\end{exercise}
\begin{solution}
图\ref{fig:lapseq}不是全局等温平衡，因此不能直接写 $n\propto\ee^{-Mgh/(RT(h))}$。应从 $\dd p/p=-Mg\dd h/[R(T_0-\beta h)]$ 积分，得到
\[
 p=p_0\left(1-\frac{\beta h}{T_0}\right)^{Mg/(R\beta)},\qquad
 n=n_0\left(1-\frac{\beta h}{T_0}\right)^{Mg/(R\beta)-1}.
\]
积分中令$z=T_0-\beta h$，有$\dd h=-\dd z/\beta$，下限$h=0$变为$T_0$，上限$h$变为$T_0-\beta h$。所以
\[
 \ln\frac p{p_0}=\frac{Mg}{R\beta}
 \int_{T_0}^{T_0-\beta h}\frac{\dd z}{z}
 =\frac{Mg}{R\beta}[\ln z]_{T_0}^{T_0-\beta h}.
\]
数密度还要除以局部温度$T(h)$，因此幂指数比压强少1。若使$\beta\to0$，利用$\ln(1-\beta h/T_0)\simeq-\beta h/T_0$，所得压强恢复等温指数律；不能直接把变温度塞进等温指数的分母。

\Diagram{fig:lapsea}{微小气团绝热位移时焓变化与重力势能平衡}{
 \node[draw]at(0,0){$C_{p,m}\dd T$};\node at(2.3,0){$+$};\node[draw]at(4.2,0){$Mg\dd h$};\node at(2.1,-.8){$\dd H_m+Mg\dd h=0$};
}
按图\ref{fig:lapsea}，气团绝热位移时满足$\dd s=C_{p,m}\dd T/T-R\dd p/p=0$，所以$\dd p/p=(C_{p,m}/R)\dd T/T$。与静力式$\dd p/p=-Mg\dd h/(RT)$联立，约去$RT$后有
\[
 C_{p,m}\dd T=-Mg\dd h,\qquad
 \frac{\dd T}{\dd h}=-\frac{Mg}{C_{p,m}},\qquad
 \beta=\frac{Mg}{C_{p,m}}.
\]
这也等价于气团的摩尔焓增加与重力势能增加相抵。常热容时梯度为常数；若$C_{p,m}$随温度变化，应积分$\int C_{p,m}(T)\dd T=-Mg\int\dd h$。这种绝热分层或气团位移条件与全局等温平衡不同。
\end{solution}

\begin{exercise}[旋转容器内的气体分布]
圆柱容器半径 $R_c$、高 $L$，装有 $N$ 个质量 $m$ 的理想气体粒子，以角速度 $\Omega$ 匀速转动。温度 $T$ 均匀，忽略重力，气体相对容器达到平衡。求径向数密度和轴心、壁面压强之比。
\Diagram{fig:rotateq}{圆柱横截面中的半径与旋转方向}{
 \draw[curve](0,0)circle(1.8);\fill(0,0)circle(1.2pt);\draw[axis](0,0)--(1.5,.6)node[midway,above]{$r$};
 \draw[->](70:2.05)arc(70:160:2.05);\node at(-1.6,2.2){$\Omega$};\node at(2.8,0){$R_c$};
}
\end{exercise}
\begin{solution}
图\ref{fig:rotateq}在旋转系中有效势为 $E_p=-m\Omega^2r^2/2$，故 $n(r)=n_c\ee^{ar^2}$，$a=m\Omega^2/(2\kb T)$。归一化要用柱坐标体积元 $2\pi rL\dd r$，
\[
 N=2\pi Ln_c\int_0^{R_c}r\ee^{ar^2}\dd r
 =\frac{\pi Ln_c}{a}(\ee^{aR_c^2}-1).
\]
径向积分可令 $u=r^2$，使 $2r\dd r=\dd u$，原函数变为 $\ee^{ar^2}/(2a)$。若漏掉柱坐标的 $r$ 因子，就无法得到正确的归一化系数。气体向外受离心有效力，在旋转系的平衡式中由向内的压强合力抵消，故压强必须向外增大。轴心密度随总粒子数约束一起变化，升高角速度会使粒子向外缘集中，并不意味着可在保持原轴心密度的同时仍保持总量不变。

\Diagram{fig:rotatea}{密度向外增大，压强梯度提供向心加速度}{
 \draw[axis](0,0)--(4.7,0)node[right]{$r$};\draw[axis](0,0)--(0,2.8)node[above]{$n$};\draw[curve]plot[domain=0:4](\x,{.4*exp(.1*\x*\x)});\node at(2.7,2.5){$\dd p/\dd r=nm\Omega^2r$};
}
所以 $n_c=Na/[\pi L(\ee^{aR_c^2}-1)]$，而 $p(R_c)/p(0)=\ee^{aR_c^2}$。图\ref{fig:rotatea}的压强向外增大，其合力指向内侧，提供实验室系的向心加速度。$\Omega\to0$ 时用 $\ee^{aR_c^2}-1\simeq aR_c^2$，恢复均匀密度 $N/(\pi R_c^2L)$。
\end{solution}

\begin{exercise}[谐势阱的分布与热容]
$N$ 个经典独立粒子处于三维谐势 $E_p=(\kappa_xx^2+\kappa_yy^2+\kappa_zz^2)/2$ 中，温度为 $T$，空间无限，$\kappa_i>0$。求位置概率密度、各坐标均方值、总能量和热容。
\Diagram{fig:harmonicq}{各方向均有二次型束缚}{
 \draw[axis](-2.3,0)--(2.3,0)node[right]{$x$};\draw[axis](0,0)--(0,2.8)node[above]{$E_p$};\draw[curve]plot[domain=-1.8:1.8](\x,{.6*\x*\x});\node at(2,2.6){$\kappa_xx^2/2$};
}
\end{exercise}
\begin{solution}
图\ref{fig:harmonicq}每一方向的势能为正二次型，位置积分分别归一化，
\[
 \rho_P(x,y,z)=\frac{\sqrt{\kappa_x\kappa_y\kappa_z}}{(2\pi\kb T)^{3/2}}
 \exp\!\left[-\frac{\kappa_xx^2+\kappa_yy^2+\kappa_zz^2}{2\kb T}\right].
\]
每个坐标方向令$\alpha_i=\kappa_i/(2\kb T)$，已有高斯积分给
\[
 Z_i=\int_{-\infty}^{\infty}\ee^{-\alpha_ix_i^2}\dd x_i
 =\sqrt\pi\alpha_i^{-1/2},
\]
三个$Z_i$相乘就是上式的归一化分母。对参数求导，
\[
 -Z_i'=\int x_i^2\ee^{-\alpha_ix_i^2}\dd x_i
 =(\sqrt\pi/2)\alpha_i^{-3/2},
\]
因此
\[
 \langle x_i^2\rangle=-\frac{Z_i'}{Z_i}=\frac1{2\alpha_i}
 =\frac{\kb T}{\kappa_i}.
\]
每个势能项的平均值为$\kappa_i\langle x_i^2\rangle/2=\kb T/2$，总势能为$3N\kb T/2$。增大刚度缩小了空间宽度，经典平衡时每个二次项的平均能量仍只由温度决定。

\Diagram{fig:harmonica}{三个动能项和三个势能项各贡献半个热能尺度}{
 \node[draw,align=center]at(0,0){动能3项\\$3\kb T/2$};\node at(2.5,0){$+$};\node[draw,align=center]at(4.8,0){势能3项\\$3\kb T/2$};\node at(2.4,-.9){单粒子总能量 $3\kb T$};
}
图\ref{fig:harmonica}再加三个动能二次项，得到总能量 $U_{\rm tot}=3N\kb T$，热容 $3N\kb$。若只束缚 $d$ 个位置方向，其他方向由固定硬壁限制，则势能贡献为 $d\kb T/2$，总热容 $(3+d)N\kb/2$。量子低温下经典等分不再适用。
\end{solution}

\begin{exercise}[有限气柱高度的涨落]
单粒子在高 $H$ 的等温重力气柱中，令 $a_g=mg/(\kb T)$、$x=a_gH$。求平均高度与高度方差，证明 $\dd\langle mgh\rangle/\dd T=(mg)^2\operatorname{Var}(h)/(\kb T^2)$，并联系气柱的额外热容。
\Diagram{fig:heightvarq}{归一化的截断指数高度分布}{
 \draw[axis](0,0)--(4.8,0)node[right]{$h$};\draw[axis](0,0)--(0,3)node[above]{$w(h)$};\draw[curve]plot[domain=0:4](\x,{2.5*exp(-.5*\x)});\draw[guide](4,0)--(4,2.6);\node[below]at(4,0){$H$};
}
\end{exercise}
\begin{solution}
图\ref{fig:heightvarq}的归一化积分 $Z_h=\int_0^H\ee^{-a_gh}\dd h=(1-\ee^{-a_gH})/a_g$。对参数求导，$\langle h\rangle=-\partial\ln Z_h/\partial a_g$，再求一次导数得到方差：
\[
 \langle h\rangle=\frac1{a_g}-\frac H{\ee^x-1},\quad
 \operatorname{Var}(h)=\frac1{a_g^2}-\frac{H^2\ee^x}{(\ee^x-1)^2}.
\]
对显式的$\ln Z_h=\ln(1-\ee^{-a_gH})-\ln a_g$求导，得$\partial_{a_g}\ln Z_h=H/(\ee^{a_gH}-1)-1/a_g$。再求导时，$\partial_{a_g}[H/(\ee^{a_gH}-1)]=-H^2\ee^{a_gH}/(\ee^{a_gH}-1)^2$，$\partial_{a_g}(-1/a_g)=1/a_g^2$，因此$\partial_{a_g}^2\ln Z_h$等于上式方差。其概率意义已由位置配分积分的商法则证明：$-Z_h''/Z_h+(Z_h'/Z_h)^2=-\operatorname{Var}(h)$。这里高度积分的边界必须保持固定。

\Diagram{fig:heightvara}{参数导数把平均位置与位置涨落相联系}{
 \node[draw]at(0,0){$-\partial_{a_g}\ln Z_h$};\draw[axis](1.6,0)--(2.8,0);\node[draw]at(4.6,0){$\partial_{a_g}^2\ln Z_h$};\node at(2.3,-.9){平均高度的响应由方差决定};
}
由图\ref{fig:heightvara}，$\partial\langle h\rangle/\partial a_g=-\operatorname{Var}(h)$，且 $\dd a_g/\dd T=-mg/(\kb T^2)$，链式相乘即得题给关系。乘 $N$ 后，额外热容为 $N(mg)^2\operatorname{Var}(h)/(\kb T^2)\ge0$，与式\eqref{eq:columncapacity}一致。热容非负在此也由概率方差直接保证。
\end{solution}

\begin{exercise}[有限气柱的绝热外场变化（选做）]
单原子气体位于高 $H$ 的刚性气柱中，粒子数固定，初温 $T$。外场加速度 $g$ 缓慢改变，过程可逆绝热。用单粒子高度分布的配分积分推导 $T$ 对 $g$ 的局部变化率，并讨论无限高气柱极限。
\Diagram{fig:changegq}{外场改变会做功，即使容器不变形}{
 \draw(0,0)rectangle(2,3);\node at(1,1.5){$N,H$};\draw[axis](3,2.5)--(3,.7)node[midway,right]{$g\to g+\dd g$};\node at(1,-.5){$Q=0$};
}
\end{exercise}
\begin{solution}
图\ref{fig:changegq}中参数$g$改变，外界对粒子做广义功$Nm\langle h\rangle\dd g$，不能再写总能量不变。单粒子配分积分中动量项$\propto T^{3/2}$，位置项为$Z_h$。由式\eqref{eq:canonicalentropy}，再用$\langle E\rangle/T=3\kb/2+mg\langle h\rangle/T$，忽略与温度和$g$无关的常数，单粒子熵为
\[
 s=\kb\left[\frac32\ln T+\ln Z_h+\frac32+a_g\langle h\rangle\right].
\]
取 $T$ 不变求 $g$ 导数，利用 $\partial_{a_g}\ln Z_h=-\langle h\rangle$ 与 $\partial_{a_g}\langle h\rangle=-\operatorname{Var}(h)$，得到
\[
 \left(\frac{\partial S}{\partial g}\right)_T
 =-\frac{Nm^2g}{\kb T^2}\operatorname{Var}(h).
\]
具体地，固定$T$时$a_{g,g}=m/(\kb T)$，对位置熵两项求导有$\partial_g[\ln Z_h+a_g\langle h\rangle]=[-\langle h\rangle+\langle h\rangle-a_g\operatorname{Var}(h)]m/(\kb T)$。前两项抵消，乘$N\kb$即得上述结果。等熵微分还使用$S_T|_g=C/T$，这里$C$包含外场势能响应；只用平动热容会漏掉位置分布的储能变化。

\Diagram{fig:changega}{等熵约束联系温度响应与外场响应}{
 \node[draw]at(0,0){$(C/T)\dd T$};\node at(2.4,0){$+$};\node[draw]at(4.7,0){$S_g|_T\dd g$};\node at(2.4,-.8){总和为零};
}
按图\ref{fig:changega}，
\[
 \left(\frac{\dd T}{\dd g}\right)_S
 =\frac{Nm^2g\operatorname{Var}(h)}{\kb TC}.
\]
无限高气柱中方差$1/a_g^2$、$C=5N\kb/2$，故$\dd T/T=(2/5)\dd g/g$。从初值$(g_i,T_i)$积分，$[\ln\vartheta]_{T_i}^{T}=(2/5)[\ln g']_{g_i}^{g}$，即$T=T_i(g/g_i)^{2/5}$。外场增强使气柱压低，同时因广义做功而升温。
\end{solution}

\begin{exercise}[球面液滴与蒸气压的曲率修正（选做）]
不可压缩液体摩尔体积为 $v_l$，温度 $T$、表面张力 $\sigma$ 恒定。平面液面饱和蒸气压为 $p_\infty$。半径 $r$ 的液滴与理想蒸气平衡，忽略蒸气压改变量对液体压强的小贡献。推导液滴周围的平衡蒸气压。
\Diagram{fig:kelvinq}{球形界面的内外压强差}{
 \draw[curve](0,0)circle(1.5);\draw[axis](0,0)--(1.5,0)node[midway,above]{$r$};\node at(0,-.7){$p_l$};\node at(2.7,1){$p_v$};\node at(2.8,0){$p_l-p_v=2\sigma/r$};
}
\end{exercise}
\begin{solution}
图\ref{fig:kelvinq}的球面机械平衡给 $p_l-p_v=2\sigma/r$。相对于平面共存，液体化学势增量近似 $\Delta\mu_l=v_l\Delta p_l\simeq2\sigma v_l/r$。理想蒸气定温化学势变化为 $\Delta\mu_v=RT\ln(p_v/p_\infty)$。
液体侧定温积分为$\Delta\mu_l=\int_{p_\infty}^{p_l}v_l\dd p'=v_l(p_l-p_\infty)$，再忽略$p_v-p_\infty$，保留曲率压差。蒸气侧为$\Delta\mu_v=\int_{p_\infty}^{p_v}RT\dd p'/p'=RT[\ln p']_{p_\infty}^{p_v}$。两相都从平面共存参考态出发，化学势初值相等，故末态相平衡要求两增量相等。指数修正仅在不可压缩液体、理想蒸气及上述压差近似同时成立时适用。

\Diagram{fig:kelvina}{相平衡要求两相化学势增量相等}{
 \node[draw]at(0,0){$2\sigma v_l/r$};\node at(2.4,0){$=$};\node[draw]at(5,0){$RT\ln(p_v/p_\infty)$};
}
按图\ref{fig:kelvina}令两者相等，得
\[
 p_v=p_\infty\exp\!\left(\frac{2\sigma v_l}{rRT}\right).
\]
表面曲率提高了液体化学势，因而小液滴需要更高的蒸气压才能维持相平衡。液滴只有一个液气界面，压差为 $2\sigma/r$；肥皂泡的两层界面才给出 $4\sigma/r$。半径接近分子尺度时，连续介质表面张力模型失效。
\end{solution}

\begin{exercise}[共存线上潜热的局部变化]
不忽略两相摩尔体积，已知各相 $C_{p,m}$ 与 $v_T|_p$。从摩尔焓差 $L_m=h_g-h_l$ 出发，求沿共存线的 $\dd L_m/\dd T$；说明它何时化为前面的热容差近似。
\Diagram{fig:latentderivq}{沿共存曲线同时改变温度和压强}{
 \draw[axis](0,0)--(4.7,0)node[right]{$T$};\draw[axis](0,0)--(0,3)node[above]{$p$};\draw[curve](1,.4)to[out=30,in=240](3.7,2.5);\draw[axis](2,1)--(2.6,1.6);\node[right]at(3,1){$\dd p=p_s'(T)\dd T$};
}
\end{exercise}
\begin{solution}
图\ref{fig:latentderivq}的方向不是定压方向。对每一相，摩尔焓微分为 $\dd h=C_{p,m}\dd T+[v-Tv_T|_p]\dd p$，两相相减后除以 $\dd T$：
\[
 \frac{\dd L_m}{\dd T}=\Delta C_{p,m}
 +\left[\Delta v-T\Delta(v_T|_p)\right]\frac{\dd p_s}{\dd T}.
\]
沿共存线求导时，每个相的 $\dd p$ 都由同一个 $p_s'(T)\dd T$ 给出，所以压强项在两相相减后留下的是 $[v_g-v_l-T((v_g)_T-(v_l)_T)]p_s'$。其中的温度偏导按定压取，而整个 $L_m$ 的导数按沿共存线取，这两种导数有不同的约束。把它们都写成未注明条件的“对温度求导”，容易把第二项误删。最后再用 Clapeyron 式替换 $p_s'$，即可把修正完全用共存两相的性质表示。

\Diagram{fig:latentderiva}{潜热变化包含定压项与沿共存线的压强修正}{
 \node[draw]at(0,0){$\Delta C_{p,m}$};\node at(2.2,0){$+$};\node[draw]at(4.8,0){$[\Delta v-T\Delta v_T]p_s'$};
}
图\ref{fig:latentderiva}中第二项是沿共存线变压带来的修正。气相理想时 $v_g-T(v_g)_T|_p=0$，再忽略液体体积及热膨胀项，便得到 $\dd L_m/\dd T\simeq C_{p,m,g}-C_{p,m,l}=c_{v,g}+R-c_l$。前面的线性汽化热公式正是在这些近似下成立。
\end{solution}

\end{document}
